% Generated mechanically from the frozen stacks-morphisms.tex target.
% Do not edit this generated file; edit the bound locale source instead.

% OK, start here.
%

\setcounter{chapter}{100}
\chapter{代数叠的态射}



\phantomsection
\label{stacks-morphisms-section-phantom}


\section{引言}
\label{stacks-morphisms-section-introduction}

\noindent
本章介绍若干类代数叠态射。对于对角态射可表示的拟分离代数叠，
可参考 \cite{LM-B}。

\medskip\noindent
我们的目标是把各类代数空间态射的定义推广到代数叠态射。
各情形略有不同，因此最好分别处理。

\medskip\noindent
对于可由代数空间表示的代数叠态射，我们已经定义了许多类态射；
见《叠的性质》第
\ref{stacks-properties-section-properties-morphisms} 节。
本章处理每个相应情形时，都必须确保一般情形下的定义与该处的定义相容。




\section{约定与术语滥用}
\label{stacks-morphisms-section-conventions}

\noindent
我们继续采用《叠的性质》第
\ref{stacks-properties-section-conventions} 节引入的约定与术语滥用。




\section{对角态射的性质}
\label{stacks-morphisms-section-diagonals}

\noindent
代数叠的对角态射与 $\mathit{Isom}$-层密切相关；见《代数叠》引理
\ref{algebraic-lemma-representable-diagonal}。
由代数叠定义中的第二个性质，这些 $\mathit{Isom}$-层总是代数空间。

\begin{lemma}
\label{stacks-morphisms-lemma-isom-locally-finite-type}
设 $\mathcal{X}$ 是代数叠。设 $T$ 是概形，且 $x,y$ 是
$\mathcal{X}$ 在 $T$ 上的纤维范畴中的对象。则态射
$\mathit{Isom}_\mathcal{X}(x, y) \to T$ 局部有限型。
\end{lemma}

\begin{proof}
由《代数叠》引理 \ref{algebraic-lemma-stack-presentation}，
可设对某个代数空间中的光滑群胚有
$\mathcal{X} = [U/R]$。由《空间的下降》引理
\ref{spaces-descent-lemma-descending-property-locally-finite-type}，
只需在 $T$ 上 fppf 局部地检验该性质。因此可设 $x,y$ 来自态射
$x', y' : T \to U$。由《空间中的群胚》引理
\ref{spaces-groupoids-lemma-quotient-stack-morphisms}，此时
$\mathit{Isom}_\mathcal{X}(x, y) = T \times_{(y', x'), U \times_S U} R$.
故只需证明 $R \to U \times_S U$ 局部有限型。这由复合
$s : R \to U \times_S U \to U$ 光滑而得（因而局部有限型；见
《空间的态射》引理
\ref{spaces-morphisms-lemma-smooth-locally-finite-presentation} 和
\ref{spaces-morphisms-lemma-finite-presentation-finite-type})
以及《空间的态射》引理
\ref{spaces-morphisms-lemma-permanence-finite-type}。
\end{proof}

\begin{lemma}
\label{stacks-morphisms-lemma-isom-pseudo-torsor-aut}
设 $\mathcal{X}$ 是代数叠。设 $T$ 是概形，且 $x,y$ 是
$\mathcal{X}$ 在 $T$ 上的纤维范畴中的对象。则
\begin{enumerate}
\item $\mathit{Isom}_\mathcal{X}(y, y)$ 是 $T$ 上的群代数空间，并且
\item $\mathit{Isom}_\mathcal{X}(x, y)$ 是 $T$ 上
$\mathit{Isom}_\mathcal{X}(y, y)$ 的伪挠子。
\end{enumerate}
\end{lemma}

\begin{proof}
见《空间中的群胚》定义
\ref{spaces-groupoids-definition-group-space} 和
\ref{spaces-groupoids-definition-pseudo-torsor}。
该引理由 $\mathcal{X}$ 是群胚中的叠立即得到。
\end{proof}

\noindent
设 $f : \mathcal{X} \to \mathcal{Y}$ 是代数叠的态射。
{\it $f$ 的对角态射}是态射
$$
\Delta_f :
\mathcal{X}
\longrightarrow
\mathcal{X} \times_\mathcal{Y} \mathcal{X}
$$
每个对角态射都具有以下两个性质。

\begin{lemma}
\label{stacks-morphisms-lemma-properties-diagonal}
\begin{slogan}
代数叠态射的对角态射可由代数空间表示，且局部有限型。
\end{slogan}
设 $f : \mathcal{X} \to \mathcal{Y}$ 是代数叠的态射。则
\begin{enumerate}
\item $\Delta_f$ 可由代数空间表示，并且
\item $\Delta_f$ 局部有限型。
\end{enumerate}
\end{lemma}

\begin{proof}
设 $T$ 是概形，且
$a : T \to \mathcal{X} \times_\mathcal{Y} \mathcal{X}$
是态射。由纤维积的定义和 $2$-Yoneda 引理，态射 $a$ 由三元组
$a = (x, x', \alpha)$ 给出，其中 $x,x'$ 是 $\mathcal{X}$
在 $T$ 上的对象，而 $\alpha : f(x) \to f(x')$ 是
$\mathcal{Y}$ 在 $T$ 上的纤维范畴中的态射。由代数叠的定义，层
$\mathit{Isom}_\mathcal{X}(x, x')$ 和
$\mathit{Isom}_\mathcal{Y}(f(x), f(x'))$ 是 $T$ 上的代数空间。
在此表述下，$\alpha$ 定义态射
$\mathit{Isom}_\mathcal{Y}(f(x), f(x')) \to T$ 的一个截面。对于
$T$ 上的概形 $T' \to T$，
$\mathcal{X} \times_{\mathcal{X} \times_\mathcal{Y} \mathcal{X}, a} T$
的一个 $T'$-值点等同于一个同构
$x|_{T'} \to x'|_{T'}$，其在 $f$ 下的像为 $\alpha|_{T'}$。
因此可见
\begin{equation}
\label{stacks-morphisms-equation-diagonal}
\vcenter{
\xymatrix{
\mathcal{X} \times_{\mathcal{X} \times_\mathcal{Y} \mathcal{X}, a} T
\ar[d] \ar[r] &
\mathit{Isom}_\mathcal{X}(x, x') \ar[d] \\
T\ar[r]^-\alpha &
\mathit{Isom}_\mathcal{Y}(f(x), f(x'))
}
}
\end{equation}
是 $T$ 上层的纤维方块。特别地，
$\mathcal{X} \times_{\mathcal{X} \times_\mathcal{Y} \mathcal{X}, a} T$
是代数空间，这证明了引理的 (1)。

\medskip\noindent
为证明第二个断言，需证明图 (\ref{stacks-morphisms-equation-diagonal}) 的左竖直箭头
局部有限型。由引理 \ref{stacks-morphisms-lemma-isom-locally-finite-type}，
代数空间 $\mathit{Isom}_\mathcal{X}(x, x')$ 在 $T$ 上局部有限型。
因此图 (\ref{stacks-morphisms-equation-diagonal}) 的右竖直箭头局部有限型；见
《空间的态射》引理 \ref{spaces-morphisms-lemma-permanence-finite-type}。
再由《空间的态射》引理
\ref{spaces-morphisms-lemma-base-change-finite-type} 即得结论。
\end{proof}

\begin{lemma}
\label{stacks-morphisms-lemma-properties-diagonal-representable}
设 $f : \mathcal{X} \to \mathcal{Y}$ 是可由代数空间表示的代数叠态射。则
\begin{enumerate}
\item $\Delta_f$ 可表示（由概形表示），
\item $\Delta_f$ 局部有限型，
\item $\Delta_f$ 是单态射，
\item $\Delta_f$ 分离，并且
\item $\Delta_f$ 局部拟有限。
\end{enumerate}
\end{lemma}

\begin{proof}
引理 \ref{stacks-morphisms-lemma-properties-diagonal} 已说明 $\Delta_f$
可由代数空间表示。因此断言 (2) -- (5) 有意义；见《叠的性质》第
\ref{stacks-properties-section-properties-morphisms} 节。
引理 \ref{stacks-morphisms-lemma-properties-diagonal} 也保证 (2) 成立。设
$T \to \mathcal{X} \times_\mathcal{Y} \mathcal{X}$ 是态射，并考察
图 (\ref{stacks-morphisms-equation-diagonal})。由《代数叠》引理
\ref{algebraic-lemma-criterion-map-representable-spaces-fibred-in-groupoids}
可知，右竖直箭头作为层态射是单射，即为代数空间的单态射。因此态射
$T \times_{\mathcal{X} \times_\mathcal{Y} \mathcal{X}} \mathcal{X} \to T$
也是单态射，故 (3) 成立。我们已经知道
$T \times_{\mathcal{X} \times_\mathcal{Y} \mathcal{X}} \mathcal{X} \to T$
局部有限型。因此《空间的态射》引理
\ref{spaces-morphisms-lemma-monomorphism-loc-finite-type-loc-quasi-finite}
说明
$T \times_{\mathcal{X} \times_\mathcal{Y} \mathcal{X}} \mathcal{X} \to T$
局部拟有限且分离。这证明了 (4) 和 (5)。最后，《空间的态射》命题
\ref{spaces-morphisms-proposition-locally-quasi-finite-separated-over-scheme}
说明
$T \times_{\mathcal{X} \times_\mathcal{Y} \mathcal{X}} \mathcal{X}$
是概形，从而证明 (1)。
\end{proof}

\begin{lemma}
\label{stacks-morphisms-lemma-representable-separated-diagonal-closed}
设 $f : \mathcal{X} \to \mathcal{Y}$ 是可由代数空间表示的代数叠态射。
则下列条件等价：
\begin{enumerate}
\item $f$ 分离；
\item $\Delta_f$ 是闭浸入；
\item $\Delta_f$ 是真态射；
\item $\Delta_f$ 泛闭。
\end{enumerate}
\end{lemma}

\begin{proof}
“$f$ 分离”“$\Delta_f$ 是闭浸入”“$\Delta_f$ 泛闭”以及
“$\Delta_f$ 是真态射”均指《叠的性质》第
\ref{stacks-properties-section-properties-morphisms} 节定义的概念。
选取概形 $V$ 以及满光滑态射 $V \to \mathcal{Y}$。
令 $U = \mathcal{X} \times_\mathcal{Y} V$；由假设它是代数空间，
而态射 $U \to \mathcal{X}$ 满且光滑。由《范畴》引理
\ref{categories-lemma-base-change-diagonal} 和《叠的性质》引理
\ref{stacks-properties-lemma-check-property-covering} 可知，对任意性质
$P$（如该引理中），$\Delta_f$ 具有 $P$ 当且仅当
$\Delta_{U/V} : U \to  U \times_V U$ 具有 $P$。
把《空间的态射》引理
\ref{spaces-morphisms-lemma-separated-diagonal-proper}
应用于 $U \to V$，即得 (2)、(3)、(4) 等价。
此外，若 (1) 成立，则 $U \to V$ 分离，故
$\Delta_{U/V}$ 是闭浸入，即 (2) 成立。
最后，设 (2) 成立。设 $T$ 是概形，且 $a : T \to \mathcal{Y}$
是态射。令 $T' = \mathcal{X} \times_\mathcal{Y} T$。
为证明 (1)，需证明代数空间的态射 $T' \to T$ 分离。
再次使用《范畴》引理 \ref{categories-lemma-base-change-diagonal}，
可见 $\Delta_{T'/T}$ 是 $\Delta_f$ 的基变换。
因此假设 (2) 蕴含 $\Delta_{T'/T}$ 是闭浸入，从而
$T' \to T$ 分离，正合所需。
\end{proof}

\begin{lemma}
\label{stacks-morphisms-lemma-representable-quasi-separated-diagonal-quasi-compact}
设 $f : \mathcal{X} \to \mathcal{Y}$ 是可由代数空间表示的代数叠态射。
则下列条件等价：
\begin{enumerate}
\item $f$ 拟分离；
\item $\Delta_f$ 拟紧；
\item $\Delta_f$ 有限型。
\end{enumerate}
\end{lemma}

\begin{proof}
“$f$ 拟分离”“$\Delta_f$ 拟紧”以及“$\Delta_f$ 有限型”
均指《叠的性质》第
\ref{stacks-properties-section-properties-morphisms} 节定义的概念。
注意，由引理 \ref{stacks-morphisms-lemma-properties-diagonal-representable}
（以及《代数叠》引理
\ref{algebraic-lemma-representable-transformations-property-implication}），
$\Delta_f$ 局部有限型，故 (2) 与 (3) 等价。
选取概形 $V$ 以及满光滑态射 $V \to \mathcal{Y}$。
令 $U = \mathcal{X} \times_\mathcal{Y} V$；由假设它是代数空间，
而态射 $U \to \mathcal{X}$ 满且光滑。由《范畴》引理
\ref{categories-lemma-base-change-diagonal} 和《叠的性质》引理
\ref{stacks-properties-lemma-check-property-covering} 可知：
$\Delta_f$ 拟紧当且仅当
$\Delta_{U/V} : U \to  U \times_V U$ 拟紧。
若 (1) 成立，则 $U \to V$ 拟分离，故
$\Delta_{U/V}$ 拟紧，即 (2) 成立。
设 (2) 成立。设 $T$ 是概形，且 $a : T \to \mathcal{Y}$
是态射。令 $T' = \mathcal{X} \times_\mathcal{Y} T$。
为证明 (1)，需证明代数空间的态射 $T' \to T$ 拟分离。
再次使用《范畴》引理 \ref{categories-lemma-base-change-diagonal}，
可见 $\Delta_{T'/T}$ 是 $\Delta_f$ 的基变换。
因此假设 (2) 蕴含 $\Delta_{T'/T}$ 拟紧，从而
$T' \to T$ 拟分离，正合所需。
\end{proof}

\begin{lemma}
\label{stacks-morphisms-lemma-representable-locally-separated-diagonal-immersion}
设 $f : \mathcal{X} \to \mathcal{Y}$ 是可由代数空间表示的代数叠态射。
则下列条件等价：
\begin{enumerate}
\item $f$ 局部分离；
\item $\Delta_f$ 是浸入。
\end{enumerate}
\end{lemma}

\begin{proof}
“$f$ 局部分离”和“$\Delta_f$ 是浸入”均指《叠的性质》第
\ref{stacks-properties-section-properties-morphisms} 节定义的概念。
证明略。提示：仿照引理
\ref{stacks-morphisms-lemma-representable-separated-diagonal-closed} 和
\ref{stacks-morphisms-lemma-representable-quasi-separated-diagonal-quasi-compact}
的证明论证。
\end{proof}






\section{分离公理}
\label{stacks-morphisms-section-separated}

\noindent
设 $\mathcal{X} = [U/R]$ 是代数叠的一个表现。则
$\mathcal{X}$ 在 $S$ 上的对角态射所具有的性质，就是态射
$j : R \to U \times_S U$ 所具有的性质。例如，若
$\mathcal{X} = [S/G]$，其中 $G$ 是 $S$ 上代数空间中的某个光滑群，
则 $j$ 是结构态射 $G \to S$。因此，对角态射本身并不自动分离
（这与概形和代数空间的情形不同）。说 $[S/G]$ 在 $S$ 上拟分离，
显然应蕴含 $G \to S$ 拟紧；但若不同时要求态射 $G \to S$ 拟分离，
我们不愿把 $[S/G]$ 称为在 $S$ 上拟分离。换言之，仅要求对角态射
拟紧并不真正符合我们对“拟分离代数叠”的直觉；还应要求对角态射
本身拟分离。

\medskip\noindent
那么“分离代数叠”应如何定义？《空间的态射》引理
\ref{spaces-morphisms-lemma-separated-diagonal-proper}
说明，代数空间分离当且仅当其对角态射为真态射。通常也用这一条件
定义分离代数叠。在上述例子 $[S/G] \to S$ 中，这意味着
$G \to S$ 是真群概形。因此，若 $E$ 是 $k$ 上的椭圆曲线，
则形如 $[\Spec(k)/E]$ 的代数叠在 $k$ 上为真
（此处以后补充参考文献）。在某些情形下，假设对角态射有限可能更自然。

\begin{definition}
\label{stacks-morphisms-definition-separated}
设 $f : \mathcal{X} \to \mathcal{Y}$ 是代数叠的态射。
\begin{enumerate}
\item 若 $\Delta_f$ 非分歧，则称 $f$ 为 {\it DM}\footnote{字母 DM
代表 Deligne--Mumford。若 $f$ 为 DM，则对任意概形 $T$ 和任意态射
$T \to \mathcal{Y}$，纤维积
$\mathcal{X}_T = \mathcal{X} \times_\mathcal{Y} T$
是 $T$ 上对角态射非分歧的代数叠，即 $\mathcal{X}_T$ 为 DM。
这蕴含 $\mathcal{X}_T$ 是 Deligne--Mumford 叠；见定理
\ref{stacks-morphisms-theorem-DM}。换言之，DM 态射的“纤维”是 Deligne--Mumford
叠。这至少为该术语提供了一些动机。}。
\item 若 $\Delta_f$ 局部拟有限，则称 $f$ 为 {\it 拟 DM}
\footnote{若 $f$ 为拟 DM，则 $\mathcal{X} \to \mathcal{Y}$ 的“纤维”
$\mathcal{X}_T$ 为拟 DM。代数叠 $\mathcal{X}$ 为拟 DM，当且仅当
存在概形 $U$ 以及有限表示、满、平坦且局部拟有限的态射
$U \to \mathcal{X}$；见定理 \ref{stacks-morphisms-theorem-quasi-DM}。
注意，这与 Deligne--Mumford 性质相似；后者以存在概形给出的
\'etale 覆盖来定义。}。
\item 若 $\Delta_f$ 为真态射，则称 $f$ {\it 分离}。
\item 若 $\Delta_f$ 拟紧且拟分离，则称 $f$ {\it 拟分离}。
\end{enumerate}
\end{definition}

\noindent
这里用到了 $\Delta_f$ 可由代数空间表示，并借助《叠的性质》第
\ref{stacks-properties-section-properties-morphisms} 节，
使对 $\Delta_f$ 施加这些条件有意义。注意，当 $f$ 可由代数空间表示时，
这些定义与已有概念并不冲突；见引理
\ref{stacks-morphisms-lemma-representable-quasi-separated-diagonal-quasi-compact} 和
\ref{stacks-morphisms-lemma-representable-separated-diagonal-closed}。
还可通过考察高阶对角态射来刻画这些条件；见引理
\ref{stacks-morphisms-lemma-definition-separated}。

\begin{definition}
\label{stacks-morphisms-definition-absolute-separated}
设 $\mathcal{X}$ 是基概形 $S$ 上的代数叠。记
$p : \mathcal{X} \to S$ 为结构态射。
\begin{enumerate}
\item 若 $p : \mathcal{X} \to S$ 为 DM，则称 $\mathcal{X}$
在 $S$ 上为 {\it DM}。
\item 若 $p : \mathcal{X} \to S$ 为拟 DM，则称 $\mathcal{X}$
在 $S$ 上{\it 拟 DM}。
\item 若 $p : \mathcal{X} \to S$ 分离，则称 $\mathcal{X}$
在 $S$ 上{\it 分离}。
\item 若 $p : \mathcal{X} \to S$ 拟分离，则称 $\mathcal{X}$
在 $S$ 上{\it 拟分离}。
\item 若 $\mathcal{X}$ 在 $\Spec(\mathbf{Z})$ 上为
DM\footnote{定理 \ref{stacks-morphisms-theorem-DM} 说明，这等价于
$\mathcal{X}$ 是 Deligne--Mumford 叠。}，则称 $\mathcal{X}$ 为 {\it DM}。
\item 若 $\mathcal{X}$ 在 $\Spec(\mathbf{Z})$ 上拟 DM，
则称 $\mathcal{X}$ {\it 拟 DM}。
\item 若 $\mathcal{X}$ 在 $\Spec(\mathbf{Z})$ 上分离，
则称 $\mathcal{X}$ {\it 分离}。
\item 若 $\mathcal{X}$ 在 $\Spec(\mathbf{Z})$ 上拟分离，
则称 $\mathcal{X}$ {\it 拟分离}。
\end{enumerate}
在最后四个定义中，我们通过《代数叠》定义
\ref{algebraic-definition-viewed-as}，
把 $\mathcal{X}$ 视为 $\Spec(\mathbf{Z})$ 上的代数叠。
\end{definition}

\noindent
因此，每个情形都有一个绝对概念和一个相对于给定基概形的概念
（按照《叠的性质》第 \ref{stacks-properties-section-conventions}
节引入的术语滥用，我们通常省略对基概形的说明）。
引理 \ref{stacks-morphisms-lemma-separated-implies-morphism-separated} 将说明
(1) $\Leftrightarrow$ (5) 且 (2) $\Leftrightarrow$ (6)。
下面花一些篇幅证明这些概念的若干标准结果。

\begin{lemma}
\label{stacks-morphisms-lemma-trivial-implications}
设 $f : \mathcal{X} \to \mathcal{Y}$ 是代数叠的态射。
\begin{enumerate}
\item 若 $f$ 分离，则 $f$ 拟分离。
\item 若 $f$ 为 DM，则 $f$ 拟 DM。
\item 若 $f$ 可由代数空间表示，则 $f$ 为 DM。
\end{enumerate}
\end{lemma}

\begin{proof}
为证 (1)，注意代数空间的真态射拟紧且拟分离；见《空间的态射》定义
\ref{spaces-morphisms-definition-proper}。为证 (2)，注意代数空间的
非分歧态射局部拟有限；见《空间的态射》引理
\ref{spaces-morphisms-lemma-unramified-quasi-finite}。
最后，(3) 由引理 \ref{stacks-morphisms-lemma-properties-diagonal-representable} 得到。
\end{proof}

\begin{lemma}
\label{stacks-morphisms-lemma-base-change-separated}
定义 \ref{stacks-morphisms-definition-separated} 中列出的所有分离公理都在基变换下稳定。
\end{lemma}

\begin{proof}
设 $f : \mathcal{X} \to \mathcal{Y}$ 和
$\mathcal{Y}' \to \mathcal{Y}$ 是代数叠的态射。设
$f' : \mathcal{Y}' \times_\mathcal{Y} \mathcal{X} \to \mathcal{Y}'$
是 $f$ 沿 $\mathcal{Y}' \to \mathcal{Y}$ 的基变换。
则 $\Delta_{f'}$ 是 $\Delta_f$ 沿态射
$\mathcal{X}' \times_{\mathcal{Y}'} \mathcal{X}' \to
\mathcal{X} \times_\mathcal{Y} \mathcal{X}$ 的基变换；见《范畴》引理
\ref{categories-lemma-base-change-diagonal}。
由《叠的性质》第
\ref{stacks-properties-section-properties-morphisms} 节的结果，
定义 \ref{stacks-morphisms-definition-separated} 中对角态射所用的每个性质都在基变换下稳定。
故引理成立。
\end{proof}

\begin{lemma}
\label{stacks-morphisms-lemma-check-separated-covering}
设 $f : \mathcal{X} \to \mathcal{Y}$ 是代数叠的态射。设 $W$ 是代数空间，
且 $W \to \mathcal{Y}$ 满、平坦并局部有限表示。若基变换
$W \times_\mathcal{Y} \mathcal{X} \to W$ 具有定义
\ref{stacks-morphisms-definition-separated} 中的某个分离性质，则 $f$ 也具有该性质。
\end{lemma}

\begin{proof}
记 $g : W \times_\mathcal{Y} \mathcal{X} \to W$ 为该基变换。
则 $\Delta_g$ 是 $\Delta_f$ 沿态射
$q : W \times_\mathcal{Y} (\mathcal{X} \times_\mathcal{Y} \mathcal{X})
\to \mathcal{X} \times_\mathcal{Y} \mathcal{X}$ 的基变换。由于 $q$
是 $W \to \mathcal{Y}$ 的基变换，$q$ 可由代数空间表示，并且满、
平坦且局部有限表示。因此结论由《叠的性质》引理
\ref{stacks-properties-lemma-check-property-weak-covering} 得到。
\end{proof}

\begin{lemma}
\label{stacks-morphisms-lemma-change-of-base-separated}
设 $S$ 是概形。在 $S$ 上拟 DM、在 $S$ 上拟分离或在 $S$ 上分离的性质
（见定义 \ref{stacks-morphisms-definition-absolute-separated}）在更换基概形下稳定；
见《代数叠》定义 \ref{algebraic-definition-change-of-base}。
\end{lemma}

\begin{proof}
由引理 \ref{stacks-morphisms-lemma-base-change-separated} 立即得到。
\end{proof}

\begin{lemma}
\label{stacks-morphisms-lemma-fibre-product-after-map}
设 $f : \mathcal{X} \to \mathcal{Z}$、$g : \mathcal{Y} \to \mathcal{Z}$
和 $\mathcal{Z} \to \mathcal{T}$ 是代数叠的态射。考察诱导态射
$i : \mathcal{X} \times_\mathcal{Z} \mathcal{Y} \to
\mathcal{X} \times_\mathcal{T} \mathcal{Y}$.
则
\begin{enumerate}
\item $i$ 可由代数空间表示，且局部有限型；
\item 若 $\Delta_{\mathcal{Z}/\mathcal{T}}$ 拟分离，则 $i$ 拟分离；
\item 若 $\Delta_{\mathcal{Z}/\mathcal{T}}$ 分离，则 $i$ 分离；
\item 若 $\mathcal{Z} \to \mathcal{T}$ 为 DM，则 $i$ 非分歧；
\item 若 $\mathcal{Z} \to \mathcal{T}$ 拟 DM，则 $i$ 局部拟有限；
\item 若 $\mathcal{Z} \to \mathcal{T}$ 分离，则 $i$ 为真态射；
\item 若 $\mathcal{Z} \to \mathcal{T}$ 拟分离，则 $i$ 拟紧且拟分离。
\end{enumerate}
\end{lemma}

\begin{proof}
下图
$$
\xymatrix{
\mathcal{X} \times_\mathcal{Z} \mathcal{Y} \ar[r]_i \ar[d] &
\mathcal{X} \times_\mathcal{T} \mathcal{Y} \ar[d] \\
\mathcal{Z} \ar[r]^-{\Delta_{\mathcal{Z}/\mathcal{T}}} \ar[r] &
\mathcal{Z} \times_\mathcal{T} \mathcal{Z}
}
$$
是 $2$-纤维积图；见《范畴》引理
\ref{categories-lemma-fibre-product-after-map}。
因此 $i$ 是对角态射 $\Delta_{\mathcal{Z}/\mathcal{T}}$ 的基变换。
故引理由引理 \ref{stacks-morphisms-lemma-properties-diagonal} 以及《叠的性质》第
\ref{stacks-properties-section-properties-morphisms} 节的内容得到。
\end{proof}

\begin{lemma}
\label{stacks-morphisms-lemma-semi-diagonal}
设 $\mathcal{T}$ 是代数叠。设 $g : \mathcal{X} \to \mathcal{Y}$
是 $\mathcal{T}$ 上代数叠的态射。考察 $g$ 的图
$i : \mathcal{X} \to \mathcal{X} \times_\mathcal{T} \mathcal{Y}$。则
\begin{enumerate}
\item $i$ 可由代数空间表示，且局部有限型；
\item 若 $\mathcal{Y} \to \mathcal{T}$ 为 DM，则 $i$ 非分歧；
\item 若 $\mathcal{Y} \to \mathcal{T}$ 拟 DM，则 $i$ 局部拟有限；
\item 若 $\mathcal{Y} \to \mathcal{T}$ 分离，则 $i$ 为真态射；
\item 若 $\mathcal{Y} \to \mathcal{T}$ 拟分离，则 $i$ 拟紧且拟分离。
\end{enumerate}
\end{lemma}

\begin{proof}
这是把引理 \ref{stacks-morphisms-lemma-fibre-product-after-map} 应用于态射
$\mathcal{X} = \mathcal{X} \times_\mathcal{Y} \mathcal{Y} \to
\mathcal{X} \times_\mathcal{T} \mathcal{Y}$ 所得的特殊情形。
\end{proof}

\begin{lemma}
\label{stacks-morphisms-lemma-section-immersion}
设 $f : \mathcal{X} \to \mathcal{T}$ 是代数叠的态射。
设 $s : \mathcal{T} \to \mathcal{X}$ 是态射，使得
$f \circ s$ 与 $\text{id}_\mathcal{T}$ $2$-同构。则
\begin{enumerate}
\item $s$ 可由代数空间表示，且局部有限型；
\item 若 $f$ 为 DM，则 $s$ 非分歧；
\item 若 $f$ 拟 DM，则 $s$ 局部拟有限；
\item 若 $f$ 分离，则 $s$ 为真态射；
\item 若 $f$ 拟分离，则 $s$ 拟紧且拟分离。
\end{enumerate}
\end{lemma}

\begin{proof}
这是把引理 \ref{stacks-morphisms-lemma-semi-diagonal} 应用于
$g = s$ 和 $\mathcal{Y} = \mathcal{T}$ 所得的特殊情形；此时
$i : \mathcal{T} \to \mathcal{T} \times_\mathcal{T} \mathcal{X}$
与 $s$ $2$-同构。
\end{proof}

\begin{lemma}
\label{stacks-morphisms-lemma-composition-separated}
定义 \ref{stacks-morphisms-definition-separated} 中列出的所有分离公理都在态射复合下稳定。
\end{lemma}

\begin{proof}
设 $f : \mathcal{X} \to \mathcal{Y}$ 和
$g : \mathcal{Y} \to \mathcal{Z}$ 是满足所论公理的代数叠态射。
对角态射 $\Delta_{\mathcal{X}/\mathcal{Z}}$ 是复合
$$
\mathcal{X} \longrightarrow
\mathcal{X} \times_\mathcal{Y} \mathcal{X} \longrightarrow
\mathcal{X} \times_\mathcal{Z} \mathcal{X}.
$$
所论分离公理通过要求对角态射具有某个性质 $\mathcal{P}$ 来定义。
由上面的引理 \ref{stacks-morphisms-lemma-fibre-product-after-map}，
第二个箭头也具有该性质。由于具有性质 $\mathcal{P}$ 且可由代数空间
表示的态射之复合仍具有性质 $\mathcal{P}$，引理即得；见
《叠的性质》第 \ref{stacks-properties-section-properties-morphisms} 节
的一般讨论以及《空间的态射》引理
\ref{spaces-morphisms-lemma-composition-unramified},
\ref{spaces-morphisms-lemma-composition-quasi-finite},
\ref{spaces-morphisms-lemma-composition-proper},
\ref{spaces-morphisms-lemma-composition-quasi-compact} 和
\ref{spaces-morphisms-lemma-composition-separated}.
\end{proof}

\begin{lemma}
\label{stacks-morphisms-lemma-separated-over-separated}
设 $f : \mathcal{X} \to \mathcal{Y}$ 是基概形 $S$ 上代数叠的态射。
\begin{enumerate}
\item 若 $\mathcal{Y}$ 在 $S$ 上为 DM 且 $f$ 为 DM，
则 $\mathcal{X}$ 在 $S$ 上为 DM。
\item 若 $\mathcal{Y}$ 在 $S$ 上拟 DM 且 $f$ 拟 DM，
则 $\mathcal{X}$ 在 $S$ 上拟 DM。
\item 若 $\mathcal{Y}$ 在 $S$ 上分离且 $f$ 分离，
则 $\mathcal{X}$ 在 $S$ 上分离。
\item 若 $\mathcal{Y}$ 在 $S$ 上拟分离且 $f$ 拟分离，
则 $\mathcal{X}$ 在 $S$ 上拟分离。
\item 若 $\mathcal{Y}$ 为 DM 且 $f$ 为 DM，则 $\mathcal{X}$ 为 DM。
\item 若 $\mathcal{Y}$ 拟 DM 且 $f$ 拟 DM，则 $\mathcal{X}$ 拟 DM。
\item 若 $\mathcal{Y}$ 分离且 $f$ 分离，则 $\mathcal{X}$ 分离。
\item 若 $\mathcal{Y}$ 拟分离且 $f$ 拟分离，则 $\mathcal{X}$ 拟分离。
\end{enumerate}
\end{lemma}

\begin{proof}
(1)、(2)、(3)、(4) 由引理 \ref{stacks-morphisms-lemma-composition-separated}
和定义 \ref{stacks-morphisms-definition-absolute-separated} 立即得到。
对于 (5)、(6)、(7)、(8)，把 $\mathcal{X}$ 和 $\mathcal{Y}$
视为 $\Spec(\mathbf{Z})$ 上的代数叠，并应用引理
\ref{stacks-morphisms-lemma-composition-separated}。细节略。
\end{proof}

\noindent
下一个引理与代数空间的相应结果略有不同。比较可见《空间的态射》引理
\ref{spaces-morphisms-lemma-compose-after-separated}。

\begin{lemma}
\label{stacks-morphisms-lemma-compose-after-separated}
设 $f : \mathcal{X} \to \mathcal{Y}$ 和
$g : \mathcal{Y} \to \mathcal{Z}$ 是代数叠的态射。
\begin{enumerate}
\item 若 $g \circ f$ 为 DM，则 $f$ 也为 DM。
\item 若 $g \circ f$ 拟 DM，则 $f$ 也拟 DM。
\item 若 $g \circ f$ 分离且 $\Delta_g$ 分离，则 $f$ 分离。
\item 若 $g \circ f$ 拟分离且 $\Delta_g$ 拟分离，则 $f$ 拟分离。
\end{enumerate}
\end{lemma}

\begin{proof}
考察 $g \circ f$ 的对角态射的分解
$$
\mathcal{X} \to
\mathcal{X} \times_\mathcal{Y} \mathcal{X} \to
\mathcal{X} \times_\mathcal{Z} \mathcal{X}
$$
两个态射都可由代数空间表示；见引理
\ref{stacks-morphisms-lemma-properties-diagonal} 和
\ref{stacks-morphisms-lemma-fibre-product-after-map}。
因此，对任意概形 $T$ 和态射
$T \to \mathcal{X} \times_\mathcal{Y} \mathcal{X}$
得到代数空间的态射
$$
A = \mathcal{X} \times_{(\mathcal{X} \times_\mathcal{Z} \mathcal{X})} T
\longrightarrow
B = (\mathcal{X} \times_\mathcal{Y} \mathcal{X})
\times_{(\mathcal{X} \times_\mathcal{Z} \mathcal{X})} T
\longrightarrow
T
$$
若 $g \circ f$ 为 DM（相应地，拟 DM），则复合 $A \to T$
非分歧（相应地，局部拟有限）。因此，由《空间的态射》引理
\ref{spaces-morphisms-lemma-permanence-unramified}
（相应地，《空间的态射》引理
\ref{spaces-morphisms-lemma-permanence-quasi-finite}），
$A \to B$ 非分歧（相应地，局部拟有限）。现考察图
$$
\xymatrix{
A \times_B T \ar[r] \ar[d] & T \ar[d] \\
A \ar[r] \ar[d] & B \ar[r] \ar[d] & T \ar[d] \\
\mathcal{X} \ar[r] &
\mathcal{X} \times_\mathcal{Y} \mathcal{X} \ar[r] &
\mathcal{X} \times_\mathcal{Z} \mathcal{X}
}
$$
其中所有方块均为 $2$-笛卡儿方块，而箭头 $T \to B$ 来自一开始的箭头
$T \to \mathcal{X} \times_\mathcal{Y} \mathcal{X}$。
由《空间的态射》引理
\ref{spaces-morphisms-lemma-base-change-unramified}
（相应地，《空间的态射》引理
\ref{spaces-morphisms-lemma-base-change-quasi-finite}），
箭头 $A \times_B T \to T$ 非分歧（相应地，局部拟有限）。
注意到 $A \times_B T$ 等于
$\mathcal{X} \times_{\mathcal{X} \times_\mathcal{Y} \mathcal{X}} T$
可知 $f$ 为 DM（相应地，拟 DM）。这证明了 (1) 和 (2)。

\medskip\noindent
(4) 的证明。设 $g \circ f$ 拟分离且 $\Delta_g$ 拟分离。考察
$g \circ f$ 的对角态射的分解
$$
\mathcal{X} \to
\mathcal{X} \times_\mathcal{Y} \mathcal{X} \to
\mathcal{X} \times_\mathcal{Z} \mathcal{X}
$$
两个态射都可由代数空间表示，且第二个态射拟分离；见引理
\ref{stacks-morphisms-lemma-properties-diagonal} 和
\ref{stacks-morphisms-lemma-fibre-product-after-map}。
因此，对任意概形 $T$ 和态射
$T \to \mathcal{X} \times_\mathcal{Y} \mathcal{X}$
得到代数空间的态射
$$
A = \mathcal{X} \times_{(\mathcal{X} \times_\mathcal{Z} \mathcal{X})} T
\longrightarrow
B = (\mathcal{X} \times_\mathcal{Y} \mathcal{X})
\times_{(\mathcal{X} \times_\mathcal{Z} \mathcal{X})} T
\longrightarrow
T
$$
其中 $B \to T$ 拟分离。由于假设 $g \circ f$ 拟分离，
复合 $A \to T$ 拟紧且拟分离。因此，由《空间的态射》引理
\ref{spaces-morphisms-lemma-compose-after-separated}，
$A \to B$ 拟分离；由《空间的态射》引理
\ref{spaces-morphisms-lemma-quasi-compact-permanence}，
$A \to B$ 拟紧。故 $f$ 拟分离。

\medskip\noindent
(3) 的证明。设 $g \circ f$ 分离且 $\Delta_g$ 分离。考察
$g \circ f$ 的对角态射的分解
$$
\mathcal{X} \to
\mathcal{X} \times_\mathcal{Y} \mathcal{X} \to
\mathcal{X} \times_\mathcal{Z} \mathcal{X}
$$
两个态射都可由代数空间表示，且第二个态射分离；见引理
\ref{stacks-morphisms-lemma-properties-diagonal} 和
\ref{stacks-morphisms-lemma-fibre-product-after-map}。
因此，对任意概形 $T$ 和态射
$T \to \mathcal{X} \times_\mathcal{Y} \mathcal{X}$
得到代数空间的态射
$$
A = \mathcal{X} \times_{(\mathcal{X} \times_\mathcal{Z} \mathcal{X})} T
\longrightarrow
B = (\mathcal{X} \times_\mathcal{Y} \mathcal{X})
\times_{(\mathcal{X} \times_\mathcal{Z} \mathcal{X})} T
\longrightarrow
T
$$
其中 $B \to T$ 分离。由于我们假设 $g \circ f$ 拟分离，
复合 $A \to T$ 为真态射。因此，由《空间的态射》引理
\ref{spaces-morphisms-lemma-universally-closed-permanence}，
$A \to B$ 为真态射；这意味着 $f$ 分离。
\end{proof}

\begin{lemma}
\label{stacks-morphisms-lemma-separated-implies-morphism-separated}
设 $\mathcal{X}$ 是基概形 $S$ 上的代数叠。
\begin{enumerate}
\item
$\mathcal{X}$ 为 DM $\Leftrightarrow$
$\mathcal{X}$ 在 $S$ 上为 DM。
\item
$\mathcal{X}$ 拟 DM $\Leftrightarrow$
$\mathcal{X}$ 在 $S$ 上拟 DM。
\item 若 $\mathcal{X}$ 分离，则 $\mathcal{X}$ 在 $S$ 上分离。
\item 若 $\mathcal{X}$ 拟分离，则 $\mathcal{X}$ 在 $S$ 上拟分离。
\end{enumerate}
设 $f : \mathcal{X} \to \mathcal{Y}$ 是基概形 $S$ 上代数叠的态射。
\begin{enumerate}
\item[(5)] 若 $\mathcal{X}$ 在 $S$ 上为 DM，则 $f$ 为 DM。
\item[(6)] 若 $\mathcal{X}$ 在 $S$ 上拟 DM，则 $f$ 拟 DM。
\item[(7)] 若 $\mathcal{X}$ 在 $S$ 上分离且
$\Delta_{\mathcal{Y}/S}$ 分离，则 $f$ 分离。
\item[(8)] 若 $\mathcal{X}$ 在 $S$ 上拟分离且
$\Delta_{\mathcal{Y}/S}$ 拟分离，则 $f$ 拟分离。
\end{enumerate}
\end{lemma}

\begin{proof}
(5)、(6)、(7)、(8) 由引理 \ref{stacks-morphisms-lemma-compose-after-separated}
和《空间》定义 \ref{spaces-definition-separated} 立即得到。
为证明 (3) 和 (4)，把 $X$ 和 $Y$ 视为
$\Spec(\mathbf{Z})$ 上的代数叠，并应用引理
\ref{stacks-morphisms-lemma-compose-after-separated}。
类似地，为证明 (1) 和 (2)，把 $\mathcal{X}$ 视为
$\Spec(\mathbf{Z})$ 上的代数叠，并考察态射
$$
\mathcal{X} \longrightarrow
\mathcal{X} \times_S \mathcal{X} \longrightarrow
\mathcal{X} \times_{\Spec(\mathbf{Z})} \mathcal{X}
$$
两个箭头都可由代数空间表示。第二个箭头是浸入
$\Delta_{S/\mathbf{Z}}$ 的基变换，因而非分歧且局部拟有限。
所以该复合非分歧（相应地，局部拟有限）当且仅当第一个箭头
非分歧（相应地，局部拟有限）；见《空间的态射》引理
\ref{spaces-morphisms-lemma-composition-unramified} 和
\ref{spaces-morphisms-lemma-permanence-unramified}
（相应地，《空间的态射》引理
\ref{spaces-morphisms-lemma-composition-quasi-finite} 和
\ref{spaces-morphisms-lemma-permanence-quasi-finite}).
\end{proof}

\begin{lemma}
\label{stacks-morphisms-lemma-properties-covering-imply-diagonal}
设 $\mathcal{X}$ 是代数叠。设 $W$ 是代数空间，且
$f : W \to \mathcal{X}$ 是满、平坦、局部有限表示的态射。
\begin{enumerate}
\item 若 $f$ 非分歧（即为 \'etale；也即 $\mathcal{X}$ 为
Deligne--Mumford），则 $\mathcal{X}$ 为 DM。
\item 若 $f$ 局部拟有限，则 $\mathcal{X}$ 拟 DM。
\end{enumerate}
\end{lemma}

\begin{proof}
注意，若 $f$ 非分歧，则由《空间的态射》引理
\ref{spaces-morphisms-lemma-unramified-flat-lfp-etale}，
该态射为 \'etale。这解释了 (1) 中的括注。
设 $f$ 非分歧（相应地，局部拟有限）。需证明
$\Delta_\mathcal{X} : \mathcal{X} \to \mathcal{X} \times \mathcal{X}$
非分歧（相应地，局部拟有限）。注意
$W \times W \to \mathcal{X} \times \mathcal{X}$ 也
满、平坦且局部有限表示。因此只需证明
$$
W \times_{\mathcal{X} \times \mathcal{X}, \Delta_\mathcal{X}} \mathcal{X}
=
W \times_\mathcal{X} W
\longrightarrow
W \times W
$$
非分歧（相应地，局部拟有限）；见《叠的性质》引理
\ref{stacks-properties-lemma-check-property-covering}。
由假设，态射 $\text{pr}_i : W \times_\mathcal{X} W \to W$
非分歧（相应地，局部拟有限）。因此，由《空间的态射》引理
\ref{spaces-morphisms-lemma-permanence-unramified}
（相应地，《空间的态射》引理
\ref{spaces-morphisms-lemma-permanence-quasi-finite}），
所显示的箭头非分歧（相应地，局部拟有限）。
\end{proof}

\begin{lemma}
\label{stacks-morphisms-lemma-monomorphism-separated}
代数叠的单态射分离且为 DM。代数叠的浸入亦如此。
\end{lemma}

\begin{proof}
若 $f : \mathcal{X} \to \mathcal{Y}$ 是代数叠的单态射，
则 $\Delta_f$ 是同构；见《叠的性质》引理
\ref{stacks-properties-lemma-monomorphism}。
代数空间的同构为真且非分歧，故 $f$ 分离且为 DM。
由于浸入是单态射，第二个断言由第一个断言得到；见《叠的性质》引理
\ref{stacks-properties-lemma-immersion-monomorphism}。
\end{proof}

\begin{lemma}
\label{stacks-morphisms-lemma-separation-properties-residual-gerbe}
设 $\mathcal{X}$ 是代数叠，且 $x \in |\mathcal{X}|$。
假设 $\mathcal{X}$ 在 $x$ 处的剩余胚 $\mathcal{Z}_x$ 存在。
若 $\mathcal{X}$ 为 DM、相应地拟 DM、相应地分离、相应地拟分离，
则 $\mathcal{Z}_x$ 也如此。
\end{lemma}

\begin{proof}
这是因为 $\mathcal{Z}_x \to \mathcal{X}$ 是单态射，故由引理
\ref{stacks-morphisms-lemma-monomorphism-separated}，它为 DM 且分离。
应用引理 \ref{stacks-morphisms-lemma-separated-over-separated} 即得结论。
\end{proof}


















\section{惯性叠}
\label{stacks-morphisms-section-inertia}

\noindent
群胚中的叠的（相对）惯性叠定义于《叠》第
\ref{stacks-section-the-inertia-stack} 节。
纤维范畴框架下的实际构造及其若干性质见《范畴》第
\ref{categories-section-inertia} 节。

\begin{lemma}
\label{stacks-morphisms-lemma-inertia}
设 $\mathcal{X}$ 是代数叠。则惯性叠
$\mathcal{I}_\mathcal{X}$ 也是代数叠。态射
$$
\mathcal{I}_\mathcal{X} \longrightarrow \mathcal{X}
$$
可由代数空间表示，且局部有限型。更一般地，设
$f : \mathcal{X} \to \mathcal{Y}$ 是代数叠的态射。
则相对惯性叠 $\mathcal{I}_{\mathcal{X}/\mathcal{Y}}$ 是代数叠，且态射
$$
\mathcal{I}_{\mathcal{X}/\mathcal{Y}} \longrightarrow \mathcal{X}
$$
可由代数空间表示，且局部有限型。
\end{lemma}

\begin{proof}
由《范畴》引理 \ref{categories-lemma-inertia-fibred-category}，有等价
$$
\mathcal{I}_\mathcal{X} \to
\mathcal{X} \times_{\Delta, \mathcal{X} \times_S \mathcal{X}, \Delta}
\mathcal{X}
\quad\text{且}\quad
\mathcal{I}_{\mathcal{X}/\mathcal{Y}} \to
\mathcal{X}
\times_{\Delta, \mathcal{X} \times_\mathcal{Y} \mathcal{X}, \Delta}
\mathcal{X}
$$
这说明这些惯性叠是代数叠。设态射 $T \to \mathcal{X}$
由 $\mathcal{X}$ 在 $T$ 上的纤维范畴中的对象 $x$ 给出。
于是得到 $2$-纤维积方块
$$
\xymatrix{
\mathit{Isom}_\mathcal{X}(x, x) \ar[d] \ar[r] &
\mathcal{I}_\mathcal{X} \ar[d] \\
T \ar[r]^x & \mathcal{X}
}
$$
这由 $\mathcal{I}_\mathcal{X}$ 的定义立即得到。
由于 $\mathit{Isom}_\mathcal{X}(x, x)$ 总是 $T$ 上局部有限型的
代数空间（见引理 \ref{stacks-morphisms-lemma-isom-locally-finite-type}），
可知 $\mathcal{I}_\mathcal{X} \to \mathcal{X}$
可由代数空间表示，且局部有限型。最后，对相对惯性叠有
$$
\vcenter{
\xymatrix{
\mathit{Isom}_\mathcal{X}(x, x) \ar[d] &
K \ar[l] \ar[d] \ar[r] &
\mathcal{I}_{\mathcal{X}/\mathcal{Y}} \ar[d] \\
\mathit{Isom}_\mathcal{Y}(f(x), f(x)) &
T \ar[l]_-e \ar[r]^x & \mathcal{X}
}
}
$$
其中两个方块都是 $2$-纤维积。这由《范畴》引理
\ref{categories-lemma-relative-inertia-as-fibre-product} 得到。
左竖直箭头是 $T$ 上局部有限型的代数空间之间的态射，因而局部有限型；
见《空间的态射》引理
\ref{spaces-morphisms-lemma-permanence-finite-type}。
故 $K$ 是代数空间，且 $K \to T$ 局部有限型。
这证明了关于相对惯性叠的断言。
\end{proof}

\begin{remark}
\label{stacks-morphisms-remark-inertia-is-group-in-spaces}
设 $f : \mathcal{X} \to \mathcal{Y}$ 是代数叠的态射。
《叠的性质》注
\ref{stacks-properties-remark-representable-over} 已说明，以
$\mathcal{X}$ 为靶、可由代数空间表示的态射
$\mathcal{Z} \to \mathcal{X}$ 所成的 $2$-范畴构成一个范畴。
在此范畴中，$\mathcal{X}/\mathcal{Y}$ 的惯性叠是一个{\it 群对象}。
回忆
$\mathcal{I}_{\mathcal{X}/\mathcal{Y}}$
的对象就是一对 $(x, \alpha)$，其中 $x$ 是 $\mathcal{X}$ 的对象，
而 $\alpha$ 是 $x$ 所在的 $\mathcal{X}$ 纤维范畴中 $x$ 的自同构，
并满足 $f(\alpha) = \text{id}$。复合
$$
c :
\mathcal{I}_{\mathcal{X}/\mathcal{Y}}
\times_\mathcal{X} \mathcal{I}_{\mathcal{X}/\mathcal{Y}}
\longrightarrow
\mathcal{I}_{\mathcal{X}/\mathcal{Y}}
$$
在对象上由规则
$$
((x, \alpha), (x', \alpha'), \beta) \mapsto
(x, \alpha \circ \beta^{-1} \circ \alpha' \circ \beta)
$$
给出。按照我们的纤维积定义，$\beta : x \to x'$ 是纤维范畴中的同构，
故该规则有意义。中性元
$e : \mathcal{X} \to \mathcal{I}_{\mathcal{X}/\mathcal{Y}}$
由函子 $x \mapsto (x, \text{id}_x)$ 给出。
群对象公理成立的证明从略。
\end{remark}

\noindent
设 $f : \mathcal{X} \to \mathcal{Y}$ 是代数叠的态射，并设
$\mathcal{I}_{\mathcal{X}/\mathcal{Y}}$ 为其惯性叠。
设 $T$ 是概形，且 $x$ 是 $\mathcal{X}$ 在 $T$ 上的对象。
令 $y = f(x)$。引理 \ref{stacks-morphisms-lemma-inertia} 的证明已说明：
对任意概形 $T$ 和 $\mathcal{X}$ 在 $T$ 上的对象 $x$，
有群层的正合列
\begin{equation}
\label{stacks-morphisms-equation-exact-sequence-isom}
0 \to
\mathit{Isom}_{\mathcal{X}/\mathcal{Y}}(x, x) \to
\mathit{Isom}_\mathcal{X}(x, x) \to
\mathit{Isom}_\mathcal{Y}(y, y)
\end{equation}
第二项和第三项上的群结构由引理
\ref{stacks-morphisms-lemma-isom-pseudo-torsor-aut} 定义，而该序列赋予第一项意义。
此外，有典范笛卡儿方块
$$
\xymatrix{
\mathit{Isom}_{\mathcal{X}/\mathcal{Y}}(x, x) \ar[d] \ar[r] &
\mathcal{I}_{\mathcal{X}/\mathcal{Y}} \ar[d] \\
T \ar[r]^x & \mathcal{X}
}
$$
事实上，注 \ref{stacks-morphisms-remark-inertia-is-group-in-spaces} 中讨论的
$\mathcal{I}_{\mathcal{X}/\mathcal{Y}}$ 上的群结构诱导出
$\mathit{Isom}_{\mathcal{X}/\mathcal{Y}}(x, x)$ 上的群结构。
这使我们也能对从代数空间到 $\mathcal{X}$ 的态射定义层
$\mathit{Isom}_{\mathcal{X}/\mathcal{Y}}(x, x)$
。下述定义将其形式化。

\begin{definition}
\label{stacks-morphisms-definition-isom}
设 $f : \mathcal{X} \to \mathcal{Y}$ 是代数叠的态射。
设 $Z$ 是代数空间。
\begin{enumerate}
\item 设 $x : Z \to \mathcal{X}$ 是态射。令
$$
\mathit{Isom}_{\mathcal{X}/\mathcal{Y}}(x, x) =
Z \times_{x, \mathcal{X}} \mathcal{I}_{\mathcal{X}/\mathcal{Y}}
$$
通过拉回注 \ref{stacks-morphisms-remark-inertia-is-group-in-spaces} 中讨论的复合法则，
赋予它 $Z$ 上群代数空间的结构。我们有时把
$\mathit{Isom}_{\mathcal{X}/\mathcal{Y}}(x, x)$
称为{\it $x$ 的相对自同构层}。
\item 设 $x_1, x_2 : Z \to \mathcal{X}$ 是态射。令
$y_i = f \circ x_i$。设 $\alpha : y_1 \to y_2$ 是 $2$-态射。
则 $\alpha$ 确定态射
$\Delta^\alpha : Z \to Z \times_{y_1, \mathcal{Y}, y_2} Z$，并令
$$
\mathit{Isom}_{\mathcal{X}/\mathcal{Y}}^\alpha(x_1, x_2) =
(Z \times_{x_1, \mathcal{X}, x_2} Z)
\times_{Z \times_{y_1, \mathcal{Y}, y_2} Z, \Delta^\alpha} Z.
$$
我们有时把
$\mathit{Isom}_{\mathcal{X}/\mathcal{Y}}^\alpha(x_1, x_2)$
称为{\it 从 $x_1$ 到 $x_2$ 的相对同构层}。
\end{enumerate}
若 $\mathcal{Y} = \Spec(\mathbf{Z})$，或更一般地，若 $\mathcal{Y}$
是代数空间，则使用记号
$\mathit{Isom}_\mathcal{X}(x, x)$ 和 $\mathit{Isom}_\mathcal{X}(x_1, x_2)$
，并使用术语{\it $x$ 的自同构层}和{\it 从 $x_1$ 到 $x_2$ 的同构层}。
\end{definition}

\begin{lemma}
\label{stacks-morphisms-lemma-isom-pseudo-torsor-aut-over-space}
设 $f : \mathcal{X} \to \mathcal{Y}$ 是代数叠的态射。
设 $Z$ 是代数空间，且 $x_i : Z \to \mathcal{X}$（$i = 1, 2$）
是态射。则
\begin{enumerate}
\item $\mathit{Isom}_{\mathcal{X}/\mathcal{Y}}(x_2, x_2)$
是 $Z$ 上的群代数空间；
\item 有群的正合列
$$
0 \to \mathit{Isom}_{\mathcal{X}/\mathcal{Y}}(x_2, x_2)
\to \mathit{Isom}_\mathcal{X}(x_2, x_2)
\to \mathit{Isom}_\mathcal{Y}(f \circ x_2, f \circ x_2)
$$
\item 有代数空间的态射
$
\mathit{Isom}_\mathcal{X}(x_1, x_2)
\to \mathit{Isom}_\mathcal{Y}(f \circ x_1, f \circ x_2)
$
使得对任意 $2$-态射 $\alpha : f \circ x_1 \to f \circ x_2$，
得到笛卡儿图
$$
\xymatrix{
\mathit{Isom}_{\mathcal{X}/\mathcal{Y}}^\alpha(x_1, x_2) \ar[d] \ar[r] &
Z \ar[d]^\alpha \\
\mathit{Isom}_\mathcal{X}(x_1, x_2) \ar[r] &
\mathit{Isom}_\mathcal{Y}(f \circ x_1, f \circ x_2)
}
$$
\item 对任意 $2$-态射 $\alpha : f \circ x_1 \to f \circ x_2$，
代数空间 $\mathit{Isom}_{\mathcal{X}/\mathcal{Y}}^\alpha(x_1, x_2)$
是 $Z$ 上 $\mathit{Isom}_{\mathcal{X}/\mathcal{Y}}(x_2, x_2)$ 的伪挠子。
\end{enumerate}
\end{lemma}

\begin{proof}
(1) 由定义 \ref{stacks-morphisms-definition-isom} 得到。
(2) 在 $Z$ 上 \'etale 局部地由正合列
(\ref{stacks-morphisms-equation-exact-sequence-isom}) 得到。展开定义即可看出 (3)。
在 $Z$ 的 \'etale 拓扑下局部地，(4) 归结为引理
\ref{stacks-morphisms-lemma-isom-pseudo-torsor-aut} 的 (2)。
\end{proof}

\begin{lemma}
\label{stacks-morphisms-lemma-cartesian-square-inertia}
设 $\pi : \mathcal{X} \to \mathcal{Y}$ 和
$f : \mathcal{Y}' \to \mathcal{Y}$ 是代数叠的态射。令
$\mathcal{X}' = \mathcal{X} \times_\mathcal{Y} \mathcal{Y}'$。
则下图中的两个方块
$$
\xymatrix{
\mathcal{I}_{\mathcal{X}'/\mathcal{Y}'} \ar[r]
\ar[d]_{
\text{《范畴》，公式}\ (\ref{categories-equation-functorial})
} &
\mathcal{X}' \ar[r]_{\pi'} \ar[d] & \mathcal{Y}' \ar[d]^f \\
\mathcal{I}_{\mathcal{X}/\mathcal{Y}} \ar[r] &
\mathcal{X} \ar[r]^\pi & \mathcal{Y}
}
$$
均为纤维积方块。
\end{lemma}

\begin{proof}
惯性叠 $\mathcal{I}_{\mathcal{X}'/\mathcal{Y}'}$ 定义为诸对
$(x', \alpha')$ 所成的范畴，其中 $x'$ 是 $\mathcal{X}'$ 的对象，
而 $\alpha'$ 是满足 $\pi'(\alpha') = \text{id}$ 的 $x'$ 的自同构；
见《范畴》第 \ref{categories-section-inertia} 节。
设 $x'$ 位于概形 $U$ 上，并映到 $\mathcal{X}$ 的对象 $x$。
由《范畴》引理 \ref{categories-lemma-2-product-categories-over-C}
中 $2$-纤维积的构造，
$x' = (U, x, y', \beta)$，其中 $y'$ 是 $\mathcal{Y}'$ 在 $U$ 上的对象，
而 $\beta$ 是 $\mathcal{Y}$ 在 $U$ 上的纤维范畴中的同构
$\beta : \pi(x) \to f(y')$。
由 $2$-纤维积的构造本身，自同构 $\alpha'$ 是一对
$(\alpha, \gamma)$，其中 $\alpha$ 是 $x$ 在 $U$ 上的自同构，
$\gamma$ 是 $y'$ 在 $U$ 上的自同构，并且
$\alpha$ 与 $\gamma$ 通过 $\beta$ 相容。条件
$\pi'(\alpha') = \text{id}$ 意味着 $\gamma = \text{id}$；
此时 $\alpha,\beta,\gamma$ 相容恰好等价于
$\pi(\alpha) = \text{id}$，即恰好意味着 $(x, \alpha)$ 是
$\mathcal{I}_{\mathcal{X}/\mathcal{Y}}$ 的对象。
由此可见左方块是纤维积方块（略去若干细节）。
\end{proof}

\begin{lemma}
\label{stacks-morphisms-lemma-monomorphism-cartesian-square-inertia}
设 $f : \mathcal{X} \to \mathcal{Y}$ 是代数叠的单态射。则图
$$
\xymatrix{
\mathcal{I}_\mathcal{X} \ar[r] \ar[d] &
\mathcal{X} \ar[d] \\
\mathcal{I}_\mathcal{Y} \ar[r] &
\mathcal{Y}
}
$$
是纤维积方块。
\end{lemma}

\begin{proof}
这由 $f$ 全忠实（见《叠的性质》引理
\ref{stacks-properties-lemma-monomorphism}）以及《范畴》第
\ref{categories-section-inertia} 节中惯性叠的定义立即得到。
具体地，$\mathcal{I}_\mathcal{X}$ 在概形 $T$ 上的对象等同于一对
$(x, \alpha)$，其中 $x$ 是 $\mathcal{X}$ 在 $T$ 上的对象，而
$\alpha : x \to x$ 是 $\mathcal{X}$ 在 $T$ 上的纤维范畴中的态射。
由于 $f$ 全忠实，$\alpha$ 等同于 $\mathcal{Y}$ 在 $T$ 上的
纤维范畴中的态射 $\beta : f(x) \to f(x)$。
因此可把 $\mathcal{I}_\mathcal{X}$ 在 $T$ 上的对象视为三元组
$((y, \beta), x, \gamma)$，其中 $y$ 是 $\mathcal{Y}$ 在 $T$ 上的对象，
$\beta : y \to y$ 属于 $\mathcal{Y}_T$，而
$\gamma : y \to f(x)$ 是 $T$ 上的同构；换言之，这是
$\mathcal{I}_\mathcal{Y} \times_\mathcal{Y} \mathcal{X}$
在 $T$ 上的对象。
\end{proof}

\begin{lemma}
\label{stacks-morphisms-lemma-presentation-inertia}
设 $\mathcal{X}$ 是代数叠。设 $[U/R] \to \mathcal{X}$ 是一个表现。
设 $G/U$ 是与群胚 $(U, R, s, t, c)$ 相关联的稳定子群代数空间。则
$$
\xymatrix{
G \ar[d] \ar[r] & U \ar[d] \\
\mathcal{I}_\mathcal{X} \ar[r] & \mathcal{X}
}
$$
是纤维积图。
\end{lemma}

\begin{proof}
由《空间中的群胚》引理
\ref{spaces-groupoids-lemma-2-cartesian-inertia} 立即得到。
\end{proof}









\section{高阶对角态射}
\label{stacks-morphisms-section-higher-diagonals}

\noindent
设 $f : \mathcal{X} \to \mathcal{Y}$ 是代数叠的态射。在这种情形下，
不仅可以考虑对角态射
$$
\Delta_f : \mathcal{X} \to \mathcal{X} \times_\mathcal{Y} \mathcal{X}
$$
还可以考虑对角态射的对角态射，即态射
$$
\Delta_{\Delta_f} :
\mathcal{X}
\longrightarrow
\mathcal{X} \times_{(\mathcal{X} \times_\mathcal{Y} \mathcal{X})} \mathcal{X}
$$
因此我们有时采用如下术语：称 $\Delta_{f, 0} = f$ 为{\it 第零对角态射}，
称 $\Delta_{f, 1} = \Delta_f$ 为{\it 第一对角态射}，并称
$\Delta_{f, 2} = \Delta_{\Delta_f}$ 为{\it 第二对角态射}。
注意，$\Delta_{f, 1}$ 可由代数空间表示且局部有限型；见引理
\ref{stacks-morphisms-lemma-properties-diagonal}。因此 $\Delta_{f, 2}$ 可表、是单态射、
局部有限型、分离且局部拟有限；见引理
\ref{stacks-morphisms-lemma-properties-diagonal-representable}。

\medskip\noindent
可以用相对惯性叠描述第二对角态射。具体而言，由《范畴》引理
\ref{categories-lemma-inertia-fibred-category}，纤维积
$\mathcal{X}
\times_{(\mathcal{X} \times_\mathcal{Y} \mathcal{X})} \mathcal{X}$
等价于相对惯性叠 $\mathcal{I}_{\mathcal{X}/\mathcal{Y}}$。
此外，在这一等同下，第二对角态射成为相对惯性叠的{\it 中性截面}
$$
\Delta_{f, 2} = e : \mathcal{X} \to \mathcal{I}_{\mathcal{X}/\mathcal{Y}}
$$
。类比代数空间中群胚的情形（《空间中的群胚》引理
\ref{spaces-groupoids-lemma-diagonal}），可得如下等价关系。

\begin{lemma}
\label{stacks-morphisms-lemma-diagonal-diagonal}
设 $f : \mathcal{X} \to \mathcal{Y}$ 是代数叠的态射。
\begin{enumerate}
\item
下列条件等价：
\begin{enumerate}
\item $\mathcal{I}_{\mathcal{X}/\mathcal{Y}} \to \mathcal{X}$
分离；
\item $\Delta_{f, 1} = \Delta_f :
\mathcal{X} \to \mathcal{X} \times_\mathcal{Y} \mathcal{X}$
分离；以及
\item $\Delta_{f, 2} = e :
\mathcal{X} \to \mathcal{I}_{\mathcal{X}/\mathcal{Y}}$
是闭浸入。
\end{enumerate}
\item
下列条件等价：
\begin{enumerate}
\item $\mathcal{I}_{\mathcal{X}/\mathcal{Y}} \to \mathcal{X}$
拟分离；
\item $\Delta_{f, 1} = \Delta_f :
\mathcal{X} \to \mathcal{X} \times_\mathcal{Y} \mathcal{X}$
拟分离；以及
\item $\Delta_{f, 2} = e :
\mathcal{X} \to \mathcal{I}_{\mathcal{X}/\mathcal{Y}}$
拟紧。
\end{enumerate}
\item
下列条件等价：
\begin{enumerate}
\item $\mathcal{I}_{\mathcal{X}/\mathcal{Y}} \to \mathcal{X}$
局部分离；
\item $\Delta_{f, 1} = \Delta_f :
\mathcal{X} \to \mathcal{X} \times_\mathcal{Y} \mathcal{X}$
局部分离；以及
\item $\Delta_{f, 2} = e :
\mathcal{X} \to \mathcal{I}_{\mathcal{X}/\mathcal{Y}}$
是浸入。
\end{enumerate}
\item
下列条件等价：
\begin{enumerate}
\item $\mathcal{I}_{\mathcal{X}/\mathcal{Y}} \to \mathcal{X}$
非分歧；
\item $f$ 是 DM。
\end{enumerate}
\item
下列条件等价：
\begin{enumerate}
\item $\mathcal{I}_{\mathcal{X}/\mathcal{Y}} \to \mathcal{X}$
局部拟有限；
\item $f$ 是拟 DM。
\end{enumerate}
\end{enumerate}
\end{lemma}

\begin{proof}
先证明 (1)、(2) 和 (3)。取代数空间 $U$ 以及满光滑态射
$U \to \mathcal{X}$。则
$G = U \times_\mathcal{X} \mathcal{I}_{\mathcal{X}/\mathcal{Y}}$
是 $U$ 上的代数空间（引理 \ref{stacks-morphisms-lemma-inertia}）。事实上，由注记
\ref{stacks-morphisms-remark-inertia-is-group-in-spaces} 中在相对惯性上构造的群律，
$G$ 是 $U$ 上的群代数空间。此外，作为 $U \to \mathcal{X}$ 的基变换，
$G \to \mathcal{I}_{\mathcal{X}/\mathcal{Y}}$ 满且光滑。最后，态射
$e : \mathcal{X} \to \mathcal{I}_{\mathcal{X}/\mathcal{Y}}$
沿 $G \to \mathcal{I}_{\mathcal{X}/\mathcal{Y}}$ 的基变换是 $G/U$ 的
单位截面 $U \to G$。因此，(a) 与 (c) 的等价性由《空间中的群胚》引理
\ref{spaces-groupoids-lemma-group-scheme-separated} 得到。
由于 $\Delta_{f, 2}$ 是 $\Delta_f$ 的对角态射，根据定义有
(b) $\Leftrightarrow$ (c)。

\medskip\noindent
再证明 (4) 和 (5)。回顾，(4)(b) 意味着 $\Delta_f$ 非分歧，
(5)(b) 意味着 $\Delta_f$ 局部拟有限。取概形 $Z$ 以及态射
$a : Z \to \mathcal{X} \times_\mathcal{Y} \mathcal{X}$。于是
$a = (x_1, x_2, \alpha)$，其中 $x_i : Z \to \mathcal{X}$，而
$\alpha : f \circ x_1  \to f \circ x_2$ 是 $2$-态射。回顾
$$
\vcenter{
\xymatrix{
\mathit{Isom}_{\mathcal{X}/\mathcal{Y}}^\alpha(x_1, x_2)
\ar[d] \ar[r] &
Z \ar[d] \\
\mathcal{X} \ar[r]^{\Delta_f} &
\mathcal{X} \times_\mathcal{Y} \mathcal{X}
}
}
\quad\text{且}\quad
\vcenter{
\xymatrix{
\mathit{Isom}_{\mathcal{X}/\mathcal{Y}}(x_2, x_2)
\ar[d] \ar[r] &
Z \ar[d]^{x_2} \\
\mathcal{I}_{\mathcal{X}/\mathcal{Y}} \ar[r] &
\mathcal{X}
}
}
$$
是笛卡尔方块。由引理 \ref{stacks-morphisms-lemma-isom-pseudo-torsor-aut-over-space}，
代数空间 $\mathit{Isom}_{\mathcal{X}/\mathcal{Y}}^\alpha(x_1, x_2)$
是 $\mathit{Isom}_{\mathcal{X}/\mathcal{Y}}(x_2, x_2)$ 在 $Z$ 上的伪挠子。
因此，(4) 和 (5) 中的等价性由《空间中的群胚》引理
\ref{spaces-groupoids-lemma-pseudo-torsor-implications} 得到。
\end{proof}

\begin{lemma}
\label{stacks-morphisms-lemma-second-diagonal}
设 $f : \mathcal{X} \to \mathcal{Y}$ 是代数叠的态射。下列条件等价：
\begin{enumerate}
\item 态射 $f$ 可由代数空间表示；
\item $f$ 的第二对角态射是同构；
\item 群叠 $ \mathcal{I}_{\mathcal{X}/\mathcal{Y}}$
在 $\mathcal X$ 上平凡；以及
\item 对概形 $T$ 和态射 $x : T \to \mathcal{X}$，
$\mathit{Isom}_\mathcal{X}(x, x) \to
\mathit{Isom}_\mathcal{Y}(f(x), f(x))$ 的核平凡。
\end{enumerate}
\end{lemma}

\begin{proof}
先证明 (1) 与 (2) 的等价性。事实上，$f$ 可由代数空间表示当且仅当
$f$ 忠实；见《代数叠》引理
\ref{algebraic-lemma-characterize-representable-by-algebraic-spaces}。
另一方面，$f$ 忠实当且仅当对概形 $T$ 上 $\mathcal{X}$ 的每个对象 $x$，
函子 $f$ 诱导单射
$\mathit{Isom}_\mathcal{X}(x, x) \to
\mathit{Isom}_\mathcal{Y}(f(x), f(x))$；这又当且仅当其核 $K$ 平凡，
亦即当且仅当对每个 $x : T \to \mathcal{X}$，$e : T \to K$ 是同构。
如上所述，$K = T \times_{x, \mathcal{X}} \mathcal{I}_{\mathcal{X}/\mathcal{Y}}$，
故这证明了 (1) 与 (2) 的等价性。为证明 (2) 与 (3) 的等价性，根据上述讨论，
只须注意群叠平凡当且仅当其单位截面是同构。最后，(3) 与 (4) 的等价性由定义
得到：在引理 \ref{stacks-morphisms-lemma-inertia} 的证明中，我们已经看到 (4) 中的核对应于
$T$ 上的纤维积
$T \times_{x, \mathcal{X}} \mathcal{I}_{\mathcal{X}/\mathcal{Y}}$。
\end{proof}

\noindent
该引理给出代数叠态射的如下层级关系。

\begin{lemma}
\label{stacks-morphisms-lemma-hierarchy}
代数叠的态射 $f : \mathcal{X} \to \mathcal{Y}$ 满足：
\begin{enumerate}
\item 该态射是单态射当且仅当 $\Delta_{f, 1}$ 是同构；并且
\item 该态射可由代数空间表示当且仅当 $\Delta_{f, 1}$ 是单态射。
\end{enumerate}
此外，第二对角态射 $\Delta_{f, 2}$ 总是单态射。
\end{lemma}

\begin{proof}
回顾《叠的性质》引理 \ref{stacks-properties-lemma-monomorphism}：
代数叠的态射是单态射，当且仅当其对角态射是叠的同构。因此，引理
\ref{stacks-morphisms-lemma-second-diagonal} 可以改述为：若一个态射的对角态射是单态射，
则该态射可由代数空间表示。特别地，这说明引理
\ref{stacks-morphisms-lemma-properties-diagonal-representable} 的条件 (3) 实际上是充要条件；
也就是说，代数叠的态射可由代数空间表示，当且仅当其对角态射是单态射。
\end{proof}

\begin{lemma}
\label{stacks-morphisms-lemma-first-diagonal-separated-second-diagonal-closed}
设 $f : \mathcal{X} \to \mathcal{Y}$ 是代数叠的态射。则
\begin{enumerate}
\item $\Delta_{f, 1}$ 分离 $\Leftrightarrow$
$\Delta_{f, 2}$ 是闭浸入 $\Leftrightarrow$
$\Delta_{f, 2}$ 为真 $\Leftrightarrow$
$\Delta_{f, 2}$ 泛闭；
\item $\Delta_{f, 1}$ 拟分离 $\Leftrightarrow$
$\Delta_{f, 2}$ 有限型 $\Leftrightarrow$ $\Delta_{f, 2}$ 拟紧；以及
\item $\Delta_{f, 1}$ 局部分离 $\Leftrightarrow$
$\Delta_{f, 2}$ 是浸入。
\end{enumerate}
\end{lemma}

\begin{proof}
将引理 \ref{stacks-morphisms-lemma-representable-separated-diagonal-closed}、
\ref{stacks-morphisms-lemma-representable-quasi-separated-diagonal-quasi-compact} 和
\ref{stacks-morphisms-lemma-representable-locally-separated-diagonal-immersion}
应用于 $\Delta_{f, 1}$ 即得。
\end{proof}

\noindent
下面这个引理颇为巧妙，并可能启发我们把这些条件推广到高阶代数叠。

\begin{lemma}
\label{stacks-morphisms-lemma-definition-separated}
设 $f : \mathcal{X} \to \mathcal{Y}$ 是代数叠的态射。则
\begin{enumerate}
\item $f$ 分离当且仅当 $\Delta_{f, 1}$ 和 $\Delta_{f, 2}$ 均泛闭；
\item $f$ 拟分离当且仅当 $\Delta_{f, 1}$ 和 $\Delta_{f, 2}$ 均拟紧；
\item $f$ 是拟 DM 当且仅当 $\Delta_{f, 1}$ 和 $\Delta_{f, 2}$
均局部拟有限；
\item $f$ 是 DM 当且仅当 $\Delta_{f, 1}$ 和 $\Delta_{f, 2}$ 均非分歧。
\end{enumerate}
\end{lemma}

\begin{proof}
证明 (1)。假设 $\Delta_{f, 2}$ 和 $\Delta_{f, 1}$ 均泛闭。由引理
\ref{stacks-morphisms-lemma-first-diagonal-separated-second-diagonal-closed}，
$\Delta_{f, 1}$ 分离且泛闭。再由《空间的态射》引理
\ref{spaces-morphisms-lemma-universally-closed-quasi-compact} 和
《代数叠》引理
\ref{algebraic-lemma-representable-transformations-property-implication}，
可知 $\Delta_{f, 1}$ 拟紧。因此它拟紧、分离、泛闭且局部有限型
（由引理 \ref{stacks-morphisms-lemma-properties-diagonal}），故为真。这证明了 (1) 的
``$\Leftarrow$''。反方向蕴涵的证明略去。

\medskip\noindent
证明 (2)。这立即由引理
\ref{stacks-morphisms-lemma-first-diagonal-separated-second-diagonal-closed} 得到。

\medskip\noindent
证明 (3)。把引理 \ref{stacks-morphisms-lemma-properties-diagonal-representable}
应用于 $\Delta_f = \Delta_{f, 1}$ 可知 $\Delta_{f, 2}$ 总是局部拟有限，
结论由此得到。

\medskip\noindent
证明 (4)。把引理 \ref{stacks-morphisms-lemma-properties-diagonal-representable}
应用于 $\Delta_f = \Delta_{f, 1}$，可知 $\Delta_{f, 2}$ 局部有限型且为
单态射，因而 $\Delta_{f, 2}$ 总是非分歧；结论由此得到。见《空间态射进阶》引理
\ref{spaces-more-morphisms-lemma-universally-injective-unramified}。
\end{proof}

\begin{lemma}
\label{stacks-morphisms-lemma-separated-implies-isom}
设 $f : \mathcal{X} \to \mathcal{Y}$ 是代数叠的分离
（相应地，拟分离、拟 DM、DM）态射。则
\begin{enumerate}
\item 给定代数空间 $T_i$（$i = 1, 2$）以及态射
$x_i : T_i \to \mathcal{X}$，令 $y_i = f \circ x_i$，则态射
$$
T_1 \times_{x_1, \mathcal{X}, x_2} T_2 \longrightarrow
T_1 \times_{y_1, \mathcal{Y}, y_2} T_2
$$
为真（相应地，拟紧且拟分离、局部拟有限、非分歧）；
\item 给定代数空间 $T$ 以及态射
$x_i : T \to \mathcal{X}$（$i = 1, 2$），令 $y_i = f \circ x_i$，则态射
$$
\mathit{Isom}_\mathcal{X}(x_1, x_2) \longrightarrow
\mathit{Isom}_\mathcal{Y}(y_1, y_2)
$$
为真（相应地，拟紧且拟分离、局部拟有限、非分歧）。
\end{enumerate}
\end{lemma}

\begin{proof}
证明 (1)。注意，图
$$
\xymatrix{
T_1 \times_{x_1, \mathcal{X}, x_2} T_2 \ar[d] \ar[r] &
T_1 \times_{y_1, \mathcal{Y}, y_2} T_2 \ar[d] \\
\mathcal{X} \ar[r] & \mathcal{X} \times_\mathcal{Y} \mathcal{X}
}
$$
是笛卡尔图。因此，结论由下述事实得到：$f$ 分离
（相应地，拟分离、拟 DM、DM）当且仅当其对角态射为真
（相应地，拟紧且拟分离、局部拟有限、非分歧）。

\medskip\noindent
证明 (2)。这是因为
$$
\mathit{Isom}_\mathcal{X}(x_1, x_2) =
(T \times_{x_1, \mathcal{X}, x_2} T) \times_{T \times T, \Delta_T} T
$$
，故 (2) 中的态射是 (1) 中态射的基变换。
\end{proof}








\section{拟紧态射}
\label{stacks-morphisms-section-quasi-compact}

\noindent
设 $f$ 是可由代数空间表示的代数叠态射。我们已在《叠的性质》第
\ref{stacks-properties-section-properties-morphisms} 节定义了 $f$ 拟紧的含义。
下面给出另一种刻画。

\begin{lemma}
\label{stacks-morphisms-lemma-characterize-representable-quasi-compact}
设 $f : \mathcal{X} \to \mathcal{Y}$ 是可由代数空间表示的代数叠态射。
下列条件等价：
\begin{enumerate}
\item $f$ 拟紧（定义见《叠的性质》第
\ref{stacks-properties-section-properties-morphisms} 节）；以及
\item 对每个拟紧代数叠 $\mathcal{Z}$ 及任意态射
$\mathcal{Z} \to \mathcal{Y}$，代数叠
$\mathcal{Z} \times_\mathcal{Y} \mathcal{X}$ 拟紧。
\end{enumerate}
\end{lemma}

\begin{proof}
假设 (1)，并设 $\mathcal{Z} \to \mathcal{Y}$ 是代数叠的态射，其中
$\mathcal{Z}$ 拟紧。由《叠的性质》引理
\ref{stacks-properties-lemma-quasi-compact-stack}，存在拟紧概形 $U$
及满光滑态射 $U \to \mathcal{Z}$。由于 $f$ 可由代数空间表示且拟紧，
根据定义，$U \times_\mathcal{Y} \mathcal{X}$ 是代数空间，并且
$U \times_\mathcal{Y} \mathcal{X} \to U$ 拟紧。因此
$U \times_\mathcal{Y} \mathcal{X}$ 是拟紧代数空间。态射
$U \times_\mathcal{Y} \mathcal{X} \to
\mathcal{Z} \times_\mathcal{Y} \mathcal{X}$
光滑且满（它是光滑满态射 $U \to \mathcal{Z}$ 的基变换）。
再次应用《叠的性质》引理
\ref{stacks-properties-lemma-quasi-compact-stack}，可知
$\mathcal{Z} \times_\mathcal{Y} \mathcal{X}$ 拟紧

\medskip\noindent
假设 (2)。设 $Z \to \mathcal{Y}$ 是态射，其中 $Z$ 是概形。
我们必须证明代数空间的态射
$p : Z \times_\mathcal{Y} \mathcal{X} \to Z$ 拟紧。
设 $U \subset Z$ 是仿射开子概形。则
$p^{-1}(U) = U \times_\mathcal{Y} Z$
，而代数空间 $U \times_\mathcal{Y} Z$ 由假设 (2) 拟紧。因此 $p$ 拟紧；
见《空间的态射》引理 \ref{spaces-morphisms-lemma-quasi-compact-local}。
\end{proof}

\noindent
这引出如下定义。

\begin{definition}
\label{stacks-morphisms-definition-quasi-compact}
设 $f : \mathcal{X} \to \mathcal{Y}$ 是代数叠的态射。若对每个拟紧代数叠
$\mathcal{Z}$ 及态射 $\mathcal{Z} \to \mathcal{Y}$，纤维积
$\mathcal{Z} \times_\mathcal{Y} \mathcal{X}$ 都拟紧，则称 $f$ {\it 拟紧}。
\end{definition}

\noindent
由上面的引理 \ref{stacks-morphisms-lemma-characterize-representable-quasi-compact}，
对可由代数空间表示的代数叠态射，这与已有概念一致。特别地，
这一概念与已经为代数叠和概形之间的态射定义的概念一致。

\begin{lemma}
\label{stacks-morphisms-lemma-base-change-quasi-compact}
代数叠的拟紧态射沿任意代数叠态射的基变换仍拟紧。
\end{lemma}

\begin{proof}
略。
\end{proof}

\begin{lemma}
\label{stacks-morphisms-lemma-composition-quasi-compact}
两个代数叠拟紧态射的复合拟紧。
\end{lemma}

\begin{proof}
略。
\end{proof}

\begin{lemma}
\label{stacks-morphisms-lemma-closed-immersion-quasi-compact}
代数叠的闭浸入拟紧。
\end{lemma}

\begin{proof}
这由浸入总是可表以及代数空间闭浸入的相应事实得到。
\end{proof}

\begin{lemma}
\label{stacks-morphisms-lemma-surjection-from-quasi-compact}
设
$$
\xymatrix{
\mathcal{X} \ar[rr]_f \ar[rd]_p & &
\mathcal{Y} \ar[dl]^q \\
& \mathcal{Z}
}
$$
是代数叠态射的 $2$-交换图。若 $f$ 满且 $p$ 拟紧，则 $q$ 拟紧。
\end{lemma}

\begin{proof}
设 $\mathcal{T}$ 是拟紧代数叠，并设 $\mathcal{T} \to \mathcal{Z}$ 是态射。
由《叠的性质》引理
\ref{stacks-properties-lemma-base-change-surjective}，态射
$\mathcal{T} \times_\mathcal{Z} \mathcal{X} \to
\mathcal{T} \times_\mathcal{Z} \mathcal{Y}$
满，而根据假设，$\mathcal{T} \times_\mathcal{Z} \mathcal{X}$ 拟紧。
因此，由《叠的性质》引理
\ref{stacks-properties-lemma-quasi-compact-stack}，
$\mathcal{T} \times_\mathcal{Z} \mathcal{Y}$ 拟紧。
\end{proof}

\begin{lemma}
\label{stacks-morphisms-lemma-quasi-compact-permanence}
设 $f : \mathcal{X} \to \mathcal{Y}$ 和
$g : \mathcal{Y} \to \mathcal{Z}$ 是代数叠的态射。若 $g \circ f$ 拟紧且
$g$ 拟分离，则 $f$ 拟紧。
\end{lemma}

\begin{proof}
这是因为 $f$ 等于如下复合：
$(1, f) : \mathcal{X} \to \mathcal{X} \times_\mathcal{Z} \mathcal{Y} \to
\mathcal{Y}$.
第一个态射是拟分离态射
$\mathcal{X} \times_\mathcal{Z} \mathcal{Y} \to \mathcal{X}$
的截面，故由引理 \ref{stacks-morphisms-lemma-section-immersion} 拟紧；该拟分离态射是 $g$
的基变换，见引理 \ref{stacks-morphisms-lemma-base-change-separated}。第二个态射是 $f$ 的
基变换，故拟紧；见引理 \ref{stacks-morphisms-lemma-base-change-quasi-compact}。
拟紧态射的复合拟紧；见引理 \ref{stacks-morphisms-lemma-composition-quasi-compact}。
\end{proof}

\begin{lemma}
\label{stacks-morphisms-lemma-quasi-compact-quasi-separated-permanence}
设 $f : \mathcal{X} \to \mathcal{Y}$ 是代数叠的态射。
\begin{enumerate}
\item 若 $\mathcal{X}$ 拟紧且 $\mathcal{Y}$ 拟分离，则 $f$ 拟紧。
\item 若 $\mathcal{X}$ 拟紧且拟分离，并且 $\mathcal{Y}$ 拟分离，
则 $f$ 拟紧且拟分离。
\item 拟紧且拟分离的代数叠的纤维积拟紧且拟分离。
\end{enumerate}
\end{lemma}

\begin{proof}
(1) 由引理 \ref{stacks-morphisms-lemma-quasi-compact-permanence} 得到。
(2) 由 (1) 和引理 \ref{stacks-morphisms-lemma-compose-after-separated} 得到。
对 (3)，设 $\mathcal{X} \to \mathcal{Y}$ 和 $\mathcal{Z} \to \mathcal{Y}$
是拟紧且拟分离代数叠的态射。利用 (2) 以及引理
\ref{stacks-morphisms-lemma-base-change-quasi-compact} 和
\ref{stacks-morphisms-lemma-base-change-separated}，态射
$\mathcal{X} \times_\mathcal{Y} \mathcal{Z} \to \mathcal{Z}$
作为 $\mathcal{X} \to \mathcal{Y}$ 的基变换拟紧且拟分离。因此，
$\mathcal{X} \times_\mathcal{Y} \mathcal{Z}$ 是 $\mathcal{Z}$ 上拟紧且
拟分离的代数叠，从而拟紧且拟分离；见引理
\ref{stacks-morphisms-lemma-separated-over-separated} 和
\ref{stacks-morphisms-lemma-composition-quasi-compact}。
\end{proof}

\begin{lemma}
\label{stacks-morphisms-lemma-reach-points-scheme-theoretic-image}
设 $f : \mathcal{X} \to \mathcal{Y}$ 是代数叠的拟紧态射。设
$y \in |\mathcal{Y}|$ 是 $|f|$ 的像的闭包中的一点。则存在以 $K$ 为分式域的
赋值环 $A$ 以及交换图
$$
\xymatrix{
\Spec(K) \ar[r] \ar[d] & \mathcal{X} \ar[d] \\
\Spec(A) \ar[r] & \mathcal{Y}
}
$$
使得 $\Spec(A)$ 的闭点映到 $y$。
\end{lemma}

\begin{proof}
取仿射概形 $V$、点 $v \in V$ 以及把 $v$ 映到 $y$ 的光滑态射
$V \to \mathcal{Y}$。考虑基变换图
$$
\xymatrix{
V \times_\mathcal{Y} \mathcal{X} \ar[r] \ar[d]_g & \mathcal{X} \ar[d]^f \\
V \ar[r] & \mathcal{Y}
}
$$
回顾，$|V \times_\mathcal{Y} \mathcal{X}| \to
|V| \times_{|\mathcal{Y}|} |\mathcal{X}|$ 满（《叠的性质》引理
\ref{stacks-properties-lemma-points-cartesian}）。由于
$|V| \to |\mathcal{Y}|$ 是开映射（《叠的性质》引理
\ref{stacks-properties-lemma-topology-points}），可知 $v$ 属于 $|g|$ 的像的闭包。
因此，只须对拟紧态射 $g$ 证明此引理（引理
\ref{stacks-morphisms-lemma-base-change-quasi-compact}）；下一段完成这一点。

\medskip\noindent
假设 $\mathcal{Y} = Y$ 是仿射概形。由于 $f$ 拟紧，$\mathcal{X}$ 拟紧
（定义 \ref{stacks-morphisms-definition-quasi-compact}）。取仿射概形 $W$ 以及满光滑态射
$W \to \mathcal{X}$。于是 $|f|$ 的像就是 $W \to Y$ 的像。
由《态射》引理 \ref{morphisms-lemma-reach-points-scheme-theoretic-image}，
可取图
$$
\xymatrix{
\Spec(K) \ar[r] \ar[d] & W \ar[d] \ar[r] & \mathcal{X} \ar[d] \\
\Spec(A) \ar[r] & Y \ar[r] & Y
}
$$
使得 $\Spec(A)$ 的闭点映到 $y$。再与 $W \to \mathcal{X}$ 复合即得所求。
\end{proof}

\begin{lemma}
\label{stacks-morphisms-lemma-check-quasi-compact-covering}
设 $f : \mathcal{X} \to \mathcal{Y}$ 是代数叠的态射。设
$W \to \mathcal{Y}$ 满、平坦且局部有限表示，其中 $W$ 是代数空间。
若基变换 $W \times_\mathcal{Y} \mathcal{X} \to W$ 拟紧，则 $f$ 拟紧。
\end{lemma}

\begin{proof}
假设 $W \times_\mathcal{Y} \mathcal{X} \to W$ 拟紧。设
$\mathcal{Z} \to \mathcal{Y}$ 是态射，其中 $\mathcal{Z}$ 是拟紧代数叠。
取概形 $U$ 以及满光滑态射
$U \to W \times_\mathcal{Y} \mathcal{Z}$。由于
$U \to \mathcal{Z}$ 平坦、满且局部有限表示，而 $\mathcal{Z}$ 拟紧，
可以找到拟紧开子概形 $U' \subset U$，使 $U' \to \mathcal{Z}$ 满。于是
$U' \times_\mathcal{Y} \mathcal{X} =
U' \times_W (W \times_\mathcal{Y} \mathcal{X})$
由假设拟紧，并满射到 $\mathcal{Z} \times_\mathcal{Y} \mathcal{X}$。
因此 $\mathcal{Z} \times_\mathcal{Y} \mathcal{X}$ 拟紧，正合所需。
\end{proof}











\section{诺特代数叠}
\label{stacks-morphisms-section-noetherian}

\noindent
我们已在《叠的性质》第 \ref{stacks-properties-section-types-properties} 节
定义了局部诺特代数叠。

\begin{definition}
\label{stacks-morphisms-definition-noetherian}
设 $\mathcal{X}$ 是代数叠。若 $\mathcal{X}$ 拟紧、拟分离且局部诺特，
则称 $\mathcal{X}$ {\it 诺特}。
\end{definition}

\noindent
注意，诺特代数叠 $\mathcal{X}$ 不仅拟紧且局部诺特，而且还拟分离。
用第 \ref{stacks-morphisms-section-higher-diagonals} 节的语言，若把“绝对”结构态射
$p : \mathcal{X} \to \Spec(\mathbf{Z})$（即把 $\mathcal{X}$ 看作
$\mathbf{Z}$ 上代数叠时的结构态射），则
$$
\mathcal{X}\text{ Noetherian}
\Leftrightarrow
\mathcal{X}\text{ locally Noetherian and }
\Delta_{p, 0}, \Delta_{p, 1}, \Delta_{p, 2}
\text{ quasi-compact}.
$$
这在后面将意味着：诺特代数叠上的有限型代数叠并不自动是诺特的。

\begin{lemma}
\label{stacks-morphisms-lemma-locally-closed-in-noetherian}
设 $j : \mathcal{X} \to \mathcal{Y}$ 是代数叠的浸入。
\begin{enumerate}
\item 若 $\mathcal{Y}$ 局部诺特，则 $\mathcal{X}$ 局部诺特且 $j$ 拟紧。
\item 若 $\mathcal{Y}$ 诺特，则 $\mathcal{X}$ 诺特。
\end{enumerate}
\end{lemma}

\begin{proof}
取概形 $V$ 以及满光滑态射 $V \to \mathcal{Y}$。则
$U = \mathcal{X} \times_\mathcal{Y} V$ 是概形，而 $V \to U$ 是浸入；
见《叠的性质》定义 \ref{stacks-properties-definition-immersion}。
回顾，$\mathcal{Y}$ 局部诺特当且仅当 $V$ 局部诺特。在这种情形下，
$U$ 也局部诺特（《态射》引理
\ref{morphisms-lemma-immersion-locally-finite-type} 和
\ref{morphisms-lemma-finite-type-noetherian}），且 $U \to V$ 拟紧
（《性质》引理 \ref{properties-lemma-immersion-into-noetherian}）。
这表明 $j$ 拟紧（引理 \ref{stacks-morphisms-lemma-check-quasi-compact-covering}），并且
$\mathcal{X}$ 局部诺特。最后，若 $\mathcal{Y}$ 诺特，则由上述论证可知
$\mathcal{X}$ 拟紧且局部诺特。为完成证明，注意 $j$ 分离；又因为
$\mathcal{Y}$ 拟分离，引理 \ref{stacks-morphisms-lemma-separated-over-separated} 表明
$\mathcal{X}$ 拟分离。
\end{proof}

\begin{lemma}
\label{stacks-morphisms-lemma-Noetherian-topology}
设 $\mathcal{X}$ 是代数叠。
\begin{enumerate}
\item 若 $\mathcal{X}$ 局部诺特，则 $|\mathcal{X}|$ 是局部诺特拓扑空间。
\item 若 $\mathcal{X}$ 拟紧且局部诺特，则 $|\mathcal{X}|$ 是诺特拓扑空间。
\end{enumerate}
\end{lemma}

\begin{proof}
假设 $\mathcal{X}$ 局部诺特。取概形 $U$ 以及满光滑态射
$U \to \mathcal{X}$。由于 $\mathcal{X}$ 局部诺特，$U$ 也局部诺特。
由《性质》引理 \ref{properties-lemma-Noetherian-topology}，这意味着
$|U|$ 是局部诺特拓扑空间。由于 $|U| \to |\mathcal{X}|$ 开且满，
由《拓扑》引理 \ref{topology-lemma-image-Noetherian}，
$|\mathcal{X}|$ 局部诺特。这证明了 (1)。若 $\mathcal{X}$ 拟紧且局部诺特，
则 $|\mathcal{X}|$ 拟紧且局部诺特。因此，由《拓扑》引理
\ref{topology-lemma-quasi-compact-locally-Noetherian-Noetherian}，
$|\mathcal{X}|$ 诺特。
\end{proof}

\begin{lemma}
\label{stacks-morphisms-lemma-locally-Noetherian-quasi-sober}
设 $\mathcal{X}$ 是局部诺特代数叠。则 $|\mathcal{X}|$ 拟清醒
（《拓扑》定义 \ref{topology-definition-generic-point}）。
\end{lemma}

\begin{proof}
我们必须证明每个不可约闭子集 $T \subset |\mathcal{X}|$ 都有一般点。
取仿射概形 $U$ 以及光滑态射 $f : U \to \mathcal{X}$，使
$f^{-1}(T) \subset |U|$ 非空。由于 $U$ 诺特，闭子集 $f^{-1}(T)$
只有有限多个不可约分支（《拓扑》引理 \ref{topology-lemma-Noetherian}）。
设 $f^{-1}(T) = Z_1 \cup \ldots \cup Z_n$ 是其不可约分支分解。
由于 $f$ 是开映射，$f|_{f^{-1}(T)} : f^{-1}(T) \to T$ 的像包含 $T$
的一个非空开子集。由于 $T$ 不可约，这意味着 $f(f^{-1}(T))$ 稠密。
仍由 $T$ 不可约，存在某个 $i$ 使 $f(Z_i)$ 稠密。若
$\xi_i \in Z_i$ 是一般点，则 $f(\xi_i)$ 是 $T$ 的一般点。
\end{proof}





\section{仿射态射}
\label{stacks-morphisms-section-affine}

\noindent
代数叠的仿射态射定义如下。

\begin{definition}
\label{stacks-morphisms-definition-affine}
若代数叠的态射可表，并且在《叠的性质》第
\ref{stacks-properties-section-properties-morphisms} 节的意义下仿射，
则称它{\it 仿射}。
\end{definition}

\noindent
对我们而言，把代数叠的仿射态射看作可由代数空间表示，并且在《叠的性质》第
\ref{stacks-properties-section-properties-morphisms} 节的意义下仿射的代数叠态射，
会稍方便一些。（回顾，在 Stacks 项目中，“可表”默认意为可由概形表示。）
这显然与刚才定义的概念等价，故下文将直接使用这一刻画，不再说明。
下面证明关于此概念的几个简单引理。

\begin{lemma}
\label{stacks-morphisms-lemma-base-change-affine}
设 $\mathcal{X} \to \mathcal{Y}$ 是代数叠的态射，
$\mathcal{Z} \to \mathcal{Y}$ 是代数叠的仿射态射。则
$\mathcal{Z} \times_\mathcal{Y} \mathcal{X} \to \mathcal{X}$
是代数叠的仿射态射。
\end{lemma}

\begin{proof}
这由《叠的性质》第
\ref{stacks-properties-section-properties-morphisms} 节的讨论得到。
\end{proof}

\begin{lemma}
\label{stacks-morphisms-lemma-composition-affine}
代数叠仿射态射的复合仿射。
\end{lemma}

\begin{proof}
这由《叠的性质》第
\ref{stacks-properties-section-properties-morphisms} 节的讨论以及
《空间的态射》引理 \ref{spaces-morphisms-lemma-composition-affine} 得到。
\end{proof}

\begin{lemma}
\label{stacks-morphisms-lemma-affine-permanence}
设
$$
\xymatrix{
\mathcal{X} \ar[rr]_f \ar[rd]_a & & \mathcal{Y} \ar[dl]^b \\
& \mathcal{Z}
}
$$
是代数叠态射的交换图。若 $a$ 仿射且 $\Delta_b$ 仿射，则 $f$ 仿射。
\end{lemma}

\begin{proof}
态射 $a$ 的基变换
$\text{pr}_2 : \mathcal{X} \times_\mathcal{Z} \mathcal{Y} \to \mathcal{Y}$
由引理 \ref{stacks-morphisms-lemma-base-change-affine} 仿射。态射
$(1, f) : \mathcal{X} \to \mathcal{X} \times_\mathcal{Z} \mathcal{Y}$
是
$\Delta_b : \mathcal{Y} \to \mathcal{Y} \times_\mathcal{Z} \mathcal{Y}$
沿态射 $\mathcal{X} \times_\mathcal{Z} \mathcal{Y} \to
\mathcal{Y} \times_\mathcal{Z} \mathcal{Y}$ 的基变换（见《范畴》第
\ref{categories-section-2-fibre-products} 节的材料）。因此，由引理
\ref{stacks-morphisms-lemma-base-change-affine}，它仿射。仿射态射的复合
$f = \text{pr}_2 \circ (1, f)$ 由引理 \ref{stacks-morphisms-lemma-composition-affine} 仿射，
证明完成。
\end{proof}





\section{整态射与有限态射}
\label{stacks-morphisms-section-integral}

\noindent
代数叠的整态射与有限态射定义如下。

\begin{definition}
\label{stacks-morphisms-definition-integral}
设 $f : \mathcal{X} \to \mathcal{Y}$ 是代数叠的态射。
\begin{enumerate}
\item 若 $f$ 可表，并且在《叠的性质》第
\ref{stacks-properties-section-properties-morphisms} 节的意义下为整，
则称 $f$ {\it 整}。
\item 若 $f$ 可表，并且在《叠的性质》第
\ref{stacks-properties-section-properties-morphisms} 节的意义下有限，
则称 $f$ {\it 有限}。
\end{enumerate}
\end{definition}

\noindent
对我们而言，把代数叠的整态射（相应地，有限态射）看作可由代数空间表示，
并且在《叠的性质》第
\ref{stacks-properties-section-properties-morphisms} 节的意义下为整
（相应地，有限）的代数叠态射，会稍方便一些。（回顾，在 Stacks 项目中，
“可表”默认意为可由概形表示。）这显然与刚才定义的概念等价，故下文将
直接使用这一刻画，不再说明。下面证明关于此概念的几个简单引理。

\begin{lemma}
\label{stacks-morphisms-lemma-base-change-integral}
设 $\mathcal{X} \to \mathcal{Y}$ 是代数叠的态射。设
$\mathcal{Z} \to \mathcal{Y}$ 是代数叠的整（或有限）态射。则
$\mathcal{Z} \times_\mathcal{Y} \mathcal{X} \to \mathcal{X}$
是代数叠的整（或有限）态射。
\end{lemma}

\begin{proof}
这由《叠的性质》第
\ref{stacks-properties-section-properties-morphisms} 节的讨论得到。
\end{proof}

\begin{lemma}
\label{stacks-morphisms-lemma-composition-integral}
代数叠整态射（相应地，有限态射）的复合为整（相应地，有限）。
\end{lemma}

\begin{proof}
这由《叠的性质》第
\ref{stacks-properties-section-properties-morphisms} 节的讨论以及
《空间的态射》引理 \ref{spaces-morphisms-lemma-composition-integral} 得到。
\end{proof}




\section{开态射}
\label{stacks-morphisms-section-open}

\noindent
设 $f$ 是可由代数空间表示的代数叠态射。我们已在《叠的性质》第
\ref{stacks-properties-section-properties-morphisms} 节定义了 $f$ 泛开的含义。
下面给出另一种刻画。

\begin{lemma}
\label{stacks-morphisms-lemma-characterize-representable-universally-open}
设 $f : \mathcal{X} \to \mathcal{Y}$ 是可由代数空间表示的代数叠态射。
下列条件等价：
\begin{enumerate}
\item $f$ 泛开（定义见《叠的性质》第
\ref{stacks-properties-section-properties-morphisms} 节）；以及
\item 对代数叠的每个态射 $\mathcal{Z} \to \mathcal{Y}$，拓扑空间的态射
$|\mathcal{Z} \times_\mathcal{Y} \mathcal{X}| \to |\mathcal{Z}|$ 是开映射。
\end{enumerate}
\end{lemma}

\begin{proof}
假设 (1)，并设 $\mathcal{Z} \to \mathcal{Y}$ 如 (2) 所述。取概形 $V$
及满光滑态射 $V \to \mathcal{Z}$。根据假设，代数空间的态射
$V \times_\mathcal{Y} \mathcal{X} \to V$ 泛开；特别地，映射
$|V \times_\mathcal{Y} \mathcal{X}| \to |V|$ 是开映射。由《叠的性质》第
\ref{stacks-properties-section-points} 节，在交换图
$$
\xymatrix{
|V \times_\mathcal{Y} \mathcal{X}| \ar[r] \ar[d] &
|\mathcal{Z} \times_\mathcal{Y} \mathcal{X}| \ar[d] \\
|V| \ar[r] & |\mathcal{Z}|
}
$$
中，水平箭头开且满，并且
$$
|V \times_\mathcal{Y} \mathcal{X}| \longrightarrow
|V| \times_{|\mathcal{Z}|} |\mathcal{Z} \times_\mathcal{Y} \mathcal{X}|
$$
满。因此，由于左侧竖直箭头是开映射，右侧竖直箭头也是开映射。
这证明了 (2)。蕴涵 (2) $\Rightarrow$ (1) 由定义得到。
\end{proof}

\noindent
因此可以采用如下自然定义。

\begin{definition}
\label{stacks-morphisms-definition-open}
设 $f : \mathcal{X} \to \mathcal{Y}$ 是代数叠的态射。
\begin{enumerate}
\item 若拓扑空间的映射 $|\mathcal{X}| \to |\mathcal{Y}|$ 是开映射，
则称 $f$ {\it 开}。
\item 若对代数叠的每个态射 $\mathcal{Z} \to \mathcal{Y}$，
拓扑空间的态射
$$
|\mathcal{Z} \times_\mathcal{Y} \mathcal{X}| \to |\mathcal{Z}|
$$
都是开映射，即基变换
$\mathcal{Z} \times_\mathcal{Y} \mathcal{X} \to \mathcal{Z}$ 开，
则称 $f$ {\it 泛开}。
\end{enumerate}
\end{definition}

\begin{lemma}
\label{stacks-morphisms-lemma-base-change-universally-open}
代数叠的泛开态射沿任意代数叠态射的基变换仍泛开。
\end{lemma}

\begin{proof}
这立即由定义得到。
\end{proof}

\begin{lemma}
\label{stacks-morphisms-lemma-composition-universally-open}
两个代数叠（泛）开态射的复合（泛）开。
\end{lemma}

\begin{proof}
略。
\end{proof}







\section{商映射态射}
\label{stacks-morphisms-section-submersive}

\noindent
设 $f$ 是可由代数空间表示的代数叠态射。我们已在《叠的性质》第
\ref{stacks-properties-section-properties-morphisms} 节定义了 $f$ 泛商映射的含义。
下面给出另一种刻画。

\begin{lemma}
\label{stacks-morphisms-lemma-characterize-representable-universally-submersive}
设 $f : \mathcal{X} \to \mathcal{Y}$ 是可由代数空间表示的代数叠态射。
下列条件等价：
\begin{enumerate}
\item $f$ 泛商映射（定义见《叠的性质》第
\ref{stacks-properties-section-properties-morphisms} 节）；以及
\item 对代数叠的每个态射 $\mathcal{Z} \to \mathcal{Y}$，拓扑空间的态射
$|\mathcal{Z} \times_\mathcal{Y} \mathcal{X}| \to |\mathcal{Z}|$ 是商映射。
\end{enumerate}
\end{lemma}

\begin{proof}
假设 (1)，并设 $\mathcal{Z} \to \mathcal{Y}$ 如 (2) 所述。取概形 $V$
及满光滑态射 $V \to \mathcal{Z}$。根据假设，代数空间的态射
$V \times_\mathcal{Y} \mathcal{X} \to V$ 泛商映射；特别地，映射
$|V \times_\mathcal{Y} \mathcal{X}| \to |V|$ 是商映射。由《叠的性质》第
\ref{stacks-properties-section-points} 节，在交换图
$$
\xymatrix{
|V \times_\mathcal{Y} \mathcal{X}| \ar[r] \ar[d] &
|\mathcal{Z} \times_\mathcal{Y} \mathcal{X}| \ar[d] \\
|V| \ar[r] & |\mathcal{Z}|
}
$$
中，水平箭头开且满，并且
$$
|V \times_\mathcal{Y} \mathcal{X}| \longrightarrow
|V| \times_{|\mathcal{Z}|} |\mathcal{Z} \times_\mathcal{Y} \mathcal{X}|
$$
满。因此，由于左侧竖直箭头是商映射，右侧竖直箭头也是商映射。
这证明了 (2)。蕴涵 (2) $\Rightarrow$ (1) 由定义得到。
\end{proof}

\noindent
因此可以采用如下自然定义。

\begin{definition}
\label{stacks-morphisms-definition-submersive}
设 $f : \mathcal{X} \to \mathcal{Y}$ 是代数叠的态射。
\begin{enumerate}
\item 若连续映射 $|\mathcal{X}| \to |\mathcal{Y}|$ 是商映射，
则称 $f$ {\it 是商映射}\footnote{这与微分流形的浸没概念截然不同。}；
见《拓扑》定义 \ref{topology-definition-submersive}。
\item 若对代数叠的每个态射 $\mathcal{Y}' \to \mathcal{Y}$，
基变换 $\mathcal{Y}' \times_\mathcal{Y} \mathcal{X} \to \mathcal{Y}'$
都是商映射，则称 $f$ {\it 泛商映射}。
\end{enumerate}
\end{definition}

\noindent
注意，商映射态射特别是满的。

\begin{lemma}
\label{stacks-morphisms-lemma-base-change-universally-submersive}
代数叠的泛商映射态射沿任意代数叠态射的基变换仍泛商映射。
\end{lemma}

\begin{proof}
这立即由定义得到。
\end{proof}

\begin{lemma}
\label{stacks-morphisms-lemma-composition-universally-submersive}
两个代数叠（泛）商映射态射的复合仍（泛）商映射。
\end{lemma}

\begin{proof}
略。
\end{proof}












\section{泛闭态射}
\label{stacks-morphisms-section-universally-closed}

\noindent
设 $f$ 是可由代数空间表示的代数叠态射。我们已在《叠的性质》第
\ref{stacks-properties-section-properties-morphisms} 节定义了 $f$ 泛闭的含义。
下面给出另一种刻画。

\begin{lemma}
\label{stacks-morphisms-lemma-characterize-representable-universally-closed}
设 $f : \mathcal{X} \to \mathcal{Y}$ 是可由代数空间表示的代数叠态射。
下列条件等价：
\begin{enumerate}
\item $f$ 泛闭（定义见《叠的性质》第
\ref{stacks-properties-section-properties-morphisms} 节）；以及
\item 对代数叠的每个态射 $\mathcal{Z} \to \mathcal{Y}$，拓扑空间的态射
$|\mathcal{Z} \times_\mathcal{Y} \mathcal{X}| \to |\mathcal{Z}|$ 是闭映射。
\end{enumerate}
\end{lemma}

\begin{proof}
假设 (1)，并设 $\mathcal{Z} \to \mathcal{Y}$ 如 (2) 所述。取概形 $V$
及满光滑态射 $V \to \mathcal{Z}$。根据假设，代数空间的态射
$V \times_\mathcal{Y} \mathcal{X} \to V$ 泛闭；特别地，映射
$|V \times_\mathcal{Y} \mathcal{X}| \to |V|$ 是闭映射。由《叠的性质》第
\ref{stacks-properties-section-points} 节，在交换图
$$
\xymatrix{
|V \times_\mathcal{Y} \mathcal{X}| \ar[r] \ar[d] &
|\mathcal{Z} \times_\mathcal{Y} \mathcal{X}| \ar[d] \\
|V| \ar[r] & |\mathcal{Z}|
}
$$
中，水平箭头开且满，并且
$$
|V \times_\mathcal{Y} \mathcal{X}| \longrightarrow
|V| \times_{|\mathcal{Z}|} |\mathcal{Z} \times_\mathcal{Y} \mathcal{X}|
$$
满。因此，由于左侧竖直箭头是闭映射，右侧竖直箭头也是闭映射。
这证明了 (2)。蕴涵 (2) $\Rightarrow$ (1) 由定义得到。
\end{proof}

\noindent
因此可以采用如下自然定义。

\begin{definition}
\label{stacks-morphisms-definition-closed}
设 $f : \mathcal{X} \to \mathcal{Y}$ 是代数叠的态射。
\begin{enumerate}
\item 若拓扑空间的映射 $|\mathcal{X}| \to |\mathcal{Y}|$ 是闭映射，
则称 $f$ {\it 闭}。
\item 若对代数叠的每个态射 $\mathcal{Z} \to \mathcal{Y}$，
拓扑空间的态射
$$
|\mathcal{Z} \times_\mathcal{Y} \mathcal{X}| \to |\mathcal{Z}|
$$
都是闭映射，即基变换
$\mathcal{Z} \times_\mathcal{Y} \mathcal{X} \to \mathcal{Z}$ 闭，
则称 $f$ {\it 泛闭}。
\end{enumerate}
\end{definition}

\begin{lemma}
\label{stacks-morphisms-lemma-base-change-universally-closed}
代数叠的泛闭态射沿任意代数叠态射的基变换仍泛闭。
\end{lemma}

\begin{proof}
这立即由定义得到。
\end{proof}

\begin{lemma}
\label{stacks-morphisms-lemma-composition-universally-closed}
两个代数叠（泛）闭态射的复合（泛）闭。
\end{lemma}

\begin{proof}
略。
\end{proof}

\begin{lemma}
\label{stacks-morphisms-lemma-universally-closed-local}
设 $f : \mathcal{X} \to \mathcal{Y}$ 是代数叠的态射。下列条件等价：
\begin{enumerate}
\item $f$ 泛闭；
\item 对每个概形 $Z$ 和每个态射 $Z \to \mathcal{Y}$，投影
$|Z \times_\mathcal{Y} \mathcal{X}| \to |Z|$ 是闭映射；
\item 对每个仿射概形 $Z$ 和每个态射 $Z \to \mathcal{Y}$，投影
$|Z \times_\mathcal{Y} \mathcal{X}| \to |Z|$ 是闭映射；以及
\item 存在代数空间 $V$ 及满光滑态射 $V \to \mathcal{Y}$，使得
$V \times_\mathcal{Y} \mathcal{X} \to V$ 是代数叠的泛闭态射。
\end{enumerate}
\end{lemma}

\begin{proof}
略去 (1) 蕴含 (2) 以及 (2) 蕴含 (3) 的证明。

\medskip\noindent
假设 (3)。取满光滑态射 $V \to \mathcal{Y}$。我们将证明
$V \times_\mathcal{Y} \mathcal{X} \to V$ 是代数叠的泛闭态射。
设 $\mathcal{Z} \to V$ 是从代数叠到 $V$ 的态射。设
$W \to \mathcal{Z}$ 是满光滑态射，其中 $W = \coprod W_i$ 是仿射概形的
不交并。于是有交换图
$$
\xymatrix{
\coprod_i |W_i \times_\mathcal{Y} \mathcal{X}| \ar@{=}[r] \ar[d] &
|W \times_\mathcal{Y} \mathcal{X}| \ar[r] \ar[d] &
|\mathcal{Z} \times_\mathcal{Y} \mathcal{X}| \ar[d] \ar@{=}[r] &
|\mathcal{Z} \times_V (V \times_\mathcal{Y} \mathcal{X})| \ar[ld] \\
\coprod |W_i| \ar@{=}[r] &
|W| \ar[r] &
|\mathcal{Z}|
}
$$
我们必须证明东南方向的箭头是闭映射。中间的水平箭头满且开
（《叠的性质》引理 \ref{stacks-properties-lemma-topology-points}）。
由假设 (3) 以及 $W_i$ 仿射这一事实，左侧竖直箭头都是闭映射，
从而右侧竖直箭头也是闭映射。

\medskip\noindent
假设 (4)。我们将证明 $f$ 泛闭。设 $\mathcal{Z} \to \mathcal{Y}$
是代数叠的态射。考虑图
$$
\xymatrix{
|(V \times_\mathcal{Y} \mathcal{Z})
\times_V (V \times_\mathcal{Y} \mathcal{X})| \ar@{=}[r] \ar[rd] &
|V \times_\mathcal{Y} \mathcal{X}| \ar[r] \ar[d] &
|Z \times_\mathcal{Y} \mathcal{X}| \ar[d] \\
 &
|V \times_\mathcal{Y} \mathcal{Z}| \ar[r] &
|\mathcal{Z}|
}
$$
根据假设，西南方向的箭头是闭映射。水平箭头满且开，因为相应的代数叠态射
满且光滑（见上面的文献）。因此，右侧竖直箭头是闭映射。
\end{proof}










\section{泛单射态射}
\label{stacks-morphisms-section-universally-injective}

\noindent
设 $f$ 是可由代数空间表示的代数叠态射。我们已在《叠的性质》第
\ref{stacks-properties-section-properties-morphisms} 节定义了 $f$ 泛单射的含义。
下面给出另一种刻画。

\begin{lemma}
\label{stacks-morphisms-lemma-characterize-representable-universally-injective}
设 $f : \mathcal{X} \to \mathcal{Y}$ 是可由代数空间表示的代数叠态射。
下列条件等价：
\begin{enumerate}
\item $f$ 泛单射（定义见《叠的性质》第
\ref{stacks-properties-section-properties-morphisms} 节）；以及
\item 对代数叠的每个态射 $\mathcal{Z} \to \mathcal{Y}$，映射
$|\mathcal{Z} \times_\mathcal{Y} \mathcal{X}| \to |\mathcal{Z}|$ 单射。
\end{enumerate}
\end{lemma}

\begin{proof}
假设 (1)，并设 $\mathcal{Z} \to \mathcal{Y}$ 如 (2) 所述。取概形 $V$
及满光滑态射 $V \to \mathcal{Z}$。根据假设，代数空间的态射
$V \times_\mathcal{Y} \mathcal{X} \to V$ 泛单射；特别地，映射
$|V \times_\mathcal{Y} \mathcal{X}| \to |V|$ 单射。由《叠的性质》第
\ref{stacks-properties-section-points} 节，在交换图
$$
\xymatrix{
|V \times_\mathcal{Y} \mathcal{X}| \ar[r] \ar[d] &
|\mathcal{Z} \times_\mathcal{Y} \mathcal{X}| \ar[d] \\
|V| \ar[r] & |\mathcal{Z}|
}
$$
中，水平箭头开且满，并且
$$
|V \times_\mathcal{Y} \mathcal{X}| \longrightarrow
|V| \times_{|\mathcal{Z}|} |\mathcal{Z} \times_\mathcal{Y} \mathcal{X}|
$$
满。因此，由于左侧竖直箭头单射，右侧竖直箭头也单射。
这证明了 (2)。蕴涵 (2) $\Rightarrow$ (1) 由定义得到。
\end{proof}

\noindent
因此可以采用如下自然定义。

\begin{definition}
\label{stacks-morphisms-definition-universally-injective}
设 $f : \mathcal{X} \to \mathcal{Y}$ 是代数叠的态射。若对代数叠的每个态射
$\mathcal{Z} \to \mathcal{Y}$，映射
$$
|\mathcal{Z} \times_\mathcal{Y} \mathcal{X}| \to |\mathcal{Z}|
$$
均单射，则称 $f$ {\it 泛单射}。
\end{definition}

\begin{lemma}
\label{stacks-morphisms-lemma-base-change-universally-injective}
代数叠的泛单射态射沿任意代数叠态射的基变换仍泛单射。
\end{lemma}

\begin{proof}
这立即由定义得到。
\end{proof}

\begin{lemma}
\label{stacks-morphisms-lemma-composition-universally-injective}
两个代数叠泛单射态射的复合泛单射。
\end{lemma}

\begin{proof}
略。
\end{proof}

\begin{lemma}
\label{stacks-morphisms-lemma-universally-injective}
设 $f : \mathcal{X} \to \mathcal{Y}$ 是代数叠的态射。下列条件等价：
\begin{enumerate}
\item $f$ 泛单射；
\item $\Delta : \mathcal{X} \to \mathcal{X} \times_\mathcal{Y} \mathcal{X}$
满；以及
\item 对代数闭域、态射 $x_1, x_2 : \Spec(k) \to \mathcal{X}$ 以及
$2$-箭头 $\beta : f \circ x_1 \to f \circ x_2$，存在 $2$-箭头
$\alpha : x_1 \to x_2$，使得 $\beta = \text{id}_f \star \alpha$。
\end{enumerate}
\end{lemma}

\begin{proof}
(1) $\Rightarrow$ (2)。若 $f$ 泛单射，则第一投影
$|\mathcal{X} \times_\mathcal{Y} \mathcal{X}| \to |\mathcal{X}|$
单射，这蕴含 $|\Delta|$ 满。

\medskip\noindent
(2) $\Rightarrow$ (1)。假设 $\Delta$ 满。则 $\Delta$ 的任意基变换都满
（见《叠的性质》第 \ref{stacks-properties-section-surjective} 节）。
由于 $f$ 的基变换的对角态射是 $\Delta$ 的基变换，只须证明
$|\mathcal{X}| \to |\mathcal{Y}|$ 单射。否则，由《叠的性质》引理
\ref{stacks-properties-lemma-points-cartesian}，第一投影
$|\mathcal{X} \times_\mathcal{Y} \mathcal{X}| \to |\mathcal{X}|$
不单射。这当然意味着 $|\Delta|$ 不满。

\medskip\noindent
(3) $\Rightarrow$ (2)。设
$t \in |\mathcal{X} \times_\mathcal{Y} \mathcal{X}|$。可以用态射
$t : \Spec(k) \to \mathcal{X} \times_\mathcal{Y} \mathcal{X}$
表示 $t$，其中 $k$ 是代数闭域。根据我们对 $2$-纤维积的构造，
可以用 $(x_1, x_2, \beta)$ 表示 $t$，其中
$x_1, x_2 : \Spec(k) \to \mathcal{X}$，而
$\beta : f \circ x_1 \to f \circ x_2$ 是 $2$-态射。于是 (3) 蕴含存在
映到 $\beta$ 的 $2$-态射 $\alpha : x_1 \to x_2$。这恰好意味着
$\Delta(x_1) = (x_1, x_1, \text{id})$ 与 $t$ 同构。因此 (2) 成立。

\medskip\noindent
(2) $\Rightarrow$ (3)。设 $x_1, x_2 : \Spec(k) \to \mathcal{X}$ 是态射，
其中 $k$ 是代数闭域。设
$\beta : f \circ x_1 \to f \circ x_2$ 是 $2$-态射。如上一段所述，得到态射
$t = (x_1, x_2, \beta) :
\Spec(k) \to \mathcal{X} \times_\mathcal{Y} \mathcal{X}$.
由引理 \ref{stacks-morphisms-lemma-properties-diagonal}，
$$
T = \mathcal{X}
\times_{\Delta, \mathcal{X} \times_\mathcal{Y} \mathcal{X}, t} \Spec(k)
$$
是 $\Spec(k)$ 上局部有限型的代数空间。条件 (2) 蕴含 $T$ 非空。
由于 $k$ 代数闭，$T$ 中有一个 $k$-点。展开定义，这意味着
$\Mor(\Spec(k), \mathcal{X})$ 中存在态射
$\alpha : x_1 \to x_2$，使得 $\beta = \text{id}_f \star \alpha$。
\end{proof}

\begin{lemma}
\label{stacks-morphisms-lemma-universally-injective-point}
设 $f : \mathcal{X} \to \mathcal{Y}$ 是代数叠的泛单射态射。设
$y : \Spec(k) \to \mathcal{Y}$ 是态射，其中 $k$ 是代数闭域。
若 $y$ 属于 $|\mathcal{X}| \to |\mathcal{Y}|$ 的像，则存在态射
$x : \Spec(k) \to \mathcal{X}$，使得 $y = f \circ x$。
\end{lemma}

\begin{proof}
先指出此引理并非显然，因为 $y$ 属于 $|f|$ 的像这一假设只意味着：
在可能用扩域替换 $k$ 之后，可以把 $y$ 提升为到 $\mathcal{X}$ 的态射。
为证明该引理，可以沿 $y$ 基变换 $f$；因而可以假设有非空代数叠
$\mathcal{X}$ 和泛单射态射 $\mathcal{X} \to \Spec(k)$，并希望找到
$\mathcal{X}$ 的一个 $k$-值点。可以用约化替换 $\mathcal{X}$。
可以选取域 $k'$ 以及满、平坦、局部有限型态射
$\Spec(k') \to \mathcal{X}$；见《叠的性质》引理
\ref{stacks-properties-lemma-unique-point}。由于
$\mathcal{X} \to \Spec(k)$ 泛单射，
$$
\Spec(k') \times_\mathcal{X} \Spec(k') \to \Spec(k' \otimes_k k')
$$
作为满态射
$\Delta : \mathcal{X} \to \mathcal{X} \times_{\Spec(k)} \mathcal{X}$
的基变换，它是满的（引理 \ref{stacks-morphisms-lemma-universally-injective}）。
由于 $k$ 代数闭，$k' \otimes_k k'$ 是整环（《代数》引理
\ref{algebra-lemma-geometrically-integral-any-integral-base-change}）。
设 $\xi \in \Spec(k') \times_\mathcal{X} \Spec(k')$ 是映到
$\Spec(k' \otimes_k k')$ 一般点的点。令 $U$ 为
$\Spec(k') \times_\mathcal{X} \Spec(k')$ 中含有 $\xi$ 的连通分支上
所赋的约化诱导闭子概形结构。于是两个投影 $U \to \Spec(k')$ 都局部有限型，
因为投影 $\Spec(k') \times_\mathcal{X} \Spec(k') \to \Spec(k')$
作为态射 $\Spec(k') \to \mathcal{X}$ 的基变换也具有这一性质。
应用《簇》命题 \ref{varieties-proposition-unique-base-field}，可知 $k'$ 在
$\Gamma(U, \mathcal{O}_U)$ 中的两个像的整闭包相等。在 $\kappa(\xi)$
中考察，这意味着每个形如 $\lambda \otimes 1$ 的元素都在子域
$$
1 \otimes k' \subset 
(\text{fraction field of }k' \otimes_k k') \subset
\kappa(\xi).
$$
上代数相关。由于 $k$ 代数闭，这只有在 $k' = k$ 时才可能，证明完成。
\end{proof}

\begin{lemma}
\label{stacks-morphisms-lemma-universally-injective-local}
设 $f : \mathcal{X} \to \mathcal{Y}$ 是代数叠的态射。下列条件等价：
\begin{enumerate}
\item $f$ 泛单射；
\item 对每个仿射概形 $Z$ 及任意态射 $Z \to \mathcal{Y}$，态射
$Z \times_\mathcal{Y} \mathcal{X} \to Z$ 泛单射；以及
\item 此处添加更多内容。
\end{enumerate}
\end{lemma}

\begin{proof}
蕴涵 (1) $\Rightarrow$ (2) 是立即的。假设 (2) 成立。我们将证明
$\Delta_f : \mathcal{X} \to \mathcal{X} \times_\mathcal{Y} \mathcal{X}$
满；由引理 \ref{stacks-morphisms-lemma-universally-injective}，这蕴含 (1)。考虑仿射概形 $V$
以及光滑态射 $V \to \mathcal{Y}$。由于
$g : V \times_\mathcal{Y} \mathcal{X} \to V$
由 (2) 泛单射，可知 $\Delta_g$ 满。然而，$\Delta_g$ 是 $\Delta_f$
沿光滑态射 $V \to \mathcal{Y}$ 的基变换。由于这些态射
$V \to \mathcal{Y}$ 合起来满，故 $\Delta_f$ 满。
\end{proof}

\begin{lemma}
\label{stacks-morphisms-lemma-check-universally-injective-covering}
设 $f : \mathcal{X} \to \mathcal{Y}$ 是代数叠的态射。设
$W \to \mathcal{Y}$ 满、平坦且局部有限表示，其中 $W$ 是代数空间。
若基变换 $W \times_\mathcal{Y} \mathcal{X} \to W$ 泛单射，
则 $f$ 泛单射。
\end{lemma}

\begin{proof}
注意，态射
$g : W \times_\mathcal{Y} \mathcal{X} \to W$
的对角态射 $\Delta_g$ 是 $\Delta_f$ 沿 $W \to \mathcal{Y}$ 的基变换。
因此，若 $\Delta_g$ 满，则由《叠的性质》引理
\ref{stacks-properties-lemma-check-property-covering}，$\Delta_f$ 也满。
于是该引理由引理 \ref{stacks-morphisms-lemma-universally-injective} 中的刻画 (2) 得到。
\end{proof}








\section{泛同胚}
\label{stacks-morphisms-section-universal-homeomorphisms}

\noindent
设 $f$ 是可由代数空间表示的代数叠态射。我们已在《叠的性质》第
\ref{stacks-properties-section-properties-morphisms} 节定义了 $f$ 为泛同胚的含义。
下面给出另一种刻画。

\begin{lemma}
\label{stacks-morphisms-lemma-characterize-representable-universal-homeomorphism}
设 $f : \mathcal{X} \to \mathcal{Y}$ 是可由代数空间表示的代数叠态射。
下列条件等价：
\begin{enumerate}
\item $f$ 是泛同胚（《叠的性质》第
\ref{stacks-properties-section-properties-morphisms} 节）；以及
\item 对代数叠的每个态射 $\mathcal{Z} \to \mathcal{Y}$，拓扑空间的映射
$|\mathcal{Z} \times_\mathcal{Y} \mathcal{X}| \to |\mathcal{Z}|$
是同胚。
\end{enumerate}
\end{lemma}

\begin{proof}
假设 (1)，并设 $\mathcal{Z} \to \mathcal{Y}$ 如 (2) 所述。取概形 $V$
及满光滑态射 $V \to \mathcal{Z}$。根据假设，代数空间的态射
$V \times_\mathcal{Y} \mathcal{X} \to V$ 是泛同胚；特别地，映射
$|V \times_\mathcal{Y} \mathcal{X}| \to |V|$ 是同胚。由《叠的性质》第
\ref{stacks-properties-section-points} 节，在交换图
$$
\xymatrix{
|V \times_\mathcal{Y} \mathcal{X}| \ar[r] \ar[d] &
|\mathcal{Z} \times_\mathcal{Y} \mathcal{X}| \ar[d] \\
|V| \ar[r] & |\mathcal{Z}|
}
$$
中，水平箭头开且满，并且
$$
|V \times_\mathcal{Y} \mathcal{X}| \longrightarrow
|V| \times_{|\mathcal{Z}|} |\mathcal{Z} \times_\mathcal{Y} \mathcal{X}|
$$
满。因此，由于左侧竖直箭头是同胚，右侧竖直箭头也是同胚。
这证明了 (2)。蕴涵 (2) $\Rightarrow$ (1) 由定义得到。
\end{proof}

\noindent
因此可以采用如下自然定义。

\begin{definition}
\label{stacks-morphisms-definition-universal-homeomorphism}
设 $f : \mathcal{X} \to \mathcal{Y}$ 是代数叠的态射。若对代数叠的每个态射
$\mathcal{Z} \to \mathcal{Y}$，拓扑空间的映射
$$
|\mathcal{Z} \times_\mathcal{Y} \mathcal{X}| \to |\mathcal{Z}|
$$
都是同胚，则称 $f$ 是{\it 泛同胚}。
\end{definition}

\begin{lemma}
\label{stacks-morphisms-lemma-base-change-universal-homeomorphism}
代数叠的泛同胚沿任意代数叠态射的基变换仍是泛同胚。
\end{lemma}

\begin{proof}
这立即由定义得到。
\end{proof}

\begin{lemma}
\label{stacks-morphisms-lemma-composition-universal-homeomorphism}
两个代数叠泛同胚的复合是泛同胚。
\end{lemma}

\begin{proof}
略。
\end{proof}

\begin{lemma}
\label{stacks-morphisms-lemma-check-universal-homeomorphism-covering}
设 $f : \mathcal{X} \to \mathcal{Y}$ 是代数叠的态射。设
$W \to \mathcal{Y}$ 满、平坦且局部有限表示，其中 $W$ 是代数空间。
若基变换 $W \times_\mathcal{Y} \mathcal{X} \to W$ 是泛同胚，
则 $f$ 是泛同胚。
\end{lemma}

\begin{proof}
假设 $g : W \times_\mathcal{Y} \mathcal{X} \to W$ 是泛同胚。
则 $g$ 泛单射，因而由引理
\ref{stacks-morphisms-lemma-check-universally-injective-covering}，$f$ 泛单射。
另一方面，设 $\mathcal{Z} \to \mathcal{Y}$ 是态射，其中 $\mathcal{Z}$
是代数叠。取概形 $U$ 及满光滑态射
$U \to W \times_\mathcal{Y} \mathcal{Z}$。考虑图
$$
\xymatrix{
W \times_\mathcal{Y} \mathcal{X} \ar[d]^g &
U \times_\mathcal{Y} \mathcal{X} \ar[d] \ar[l] \ar[r] &
\mathcal{Z} \times_\mathcal{Y} \mathcal{X} \ar[d] \\
W &
U \ar[l] \ar[r] &
\mathcal{Z}
}
$$
根据对 $g$ 的假设，中间的竖直箭头在拓扑空间上诱导同胚。态射
$U \to \mathcal{Z}$ 和
$U \times_\mathcal{Y} \mathcal{X} \to
\mathcal{Z} \times_\mathcal{Y} \mathcal{X}$
满、平坦且局部有限表示，故在拓扑空间上诱导开映射。由此可知
$|\mathcal{Z} \times_\mathcal{Y} \mathcal{X}| \to |\mathcal{Z}|$
是开映射。满性容易证明；我们略去其证明。
\end{proof}










\section{在源与靶上光滑局部的态射类型}
\label{stacks-morphisms-section-local-source-target}

\noindent
给定代数空间态射的一个{\it 在源与靶上光滑局部}的性质（见
《空间上的下降》定义 \ref{spaces-descent-definition-local-source-target}），
可以用它定义代数叠态射的相应性质，具体做法是要求下面引理中的任一等价条件。

\begin{lemma}
\label{stacks-morphisms-lemma-local-source-target}
设 $\mathcal{P}$ 是代数空间态射的一个在源与靶上光滑局部的性质。
设 $f : \mathcal{X} \to \mathcal{Y}$ 是代数叠的态射。考虑交换图
$$
\xymatrix{
U \ar[d]_a \ar[r]_h & V \ar[d]^b \\
\mathcal{X} \ar[r]^f & \mathcal{Y}
}
$$
，其中 $U$ 和 $V$ 是代数空间，竖直箭头光滑。下列条件等价：
\begin{enumerate}
\item 对任何如上并且 $U \to \mathcal{X} \times_\mathcal{Y} V$ 也光滑的图，
态射 $h$ 具有性质 $\mathcal{P}$；以及
\item 对某个如上的图，若 $a : U \to \mathcal{X}$ 满，则态射 $h$
具有性质 $\mathcal{P}$。
\end{enumerate}
若 $\mathcal{X}$ 和 $\mathcal{Y}$ 可由代数空间表示，则这些条件还等价于
$f$（作为代数空间的态射）具有性质 $\mathcal{P}$。若 $\mathcal{P}$
还在任意基变换下保持，并且在基上为 fppf 局部的，则对于可由代数空间表示的
态射 $f$，这些条件还等价于 $f$ 在《叠的性质》第
\ref{stacks-properties-section-properties-morphisms} 节的意义下具有性质
$\mathcal{P}$。
\end{lemma}

\begin{proof}
先证明蕴涵 (1) $\Rightarrow$ (2)。取代数空间 $V$ 以及满光滑态射
$V \to \mathcal{Y}$。取代数空间 $U$ 以及满光滑态射
$U \to \mathcal{X} \times_\mathcal{Y} V$。注意，$U \to \mathcal{X}$
作为基变换 $\mathcal{X} \times_\mathcal{Y} V \to \mathcal{X}$ 与所选态射
$U \to \mathcal{X} \times_\mathcal{Y} V$ 的复合，也满且光滑。因此得到
如 (1) 中的图。于是，若 (1) 成立，则 $h : U \to V$ 具有性质
$\mathcal{P}$；由于 $U \to \mathcal{X}$ 满，这意味着 (2) 成立。

\medskip\noindent
反之，假设 (2) 成立，并设 $U, V, a, b, h$ 如 (2) 所述。再设
$U', V', a', b', h'$ 是 (1) 中的任意图。图示为
$$
\xymatrix{
U \ar[d] \ar[r]_h & V \ar[d] \\
\mathcal{X} \ar[r]^f & \mathcal{Y}
}
\quad\quad
\xymatrix{
U' \ar[d] \ar[r]_{h'} & V' \ar[d] \\
\mathcal{X} \ar[r]^f & \mathcal{Y}
}
$$
为证明 (2) 蕴含 (1)，必须证明 $h'$ 具有 $\mathcal{P}$。为此考虑交换图
$$
\xymatrix{
U \ar[dd]^h &
U \times_\mathcal{X} U' \ar[d] \ar[l] \ar@/^6ex/[dd]^{(h, h')} \ar[r] &
U' \ar[dd]^{h'} \\
& U \times_\mathcal{Y} V' \ar[lu] \ar[d] & \\
V &
V \times_\mathcal{Y} V' \ar[l] \ar[r] &
V'
}
$$
，其中各对象都是代数空间。注意，水平箭头分别作为光滑态射
$V \to \mathcal{Y}$、$V' \to \mathcal{Y}$、$U \to \mathcal{X}$ 和
$U' \to \mathcal{X}$ 的基变换而光滑。还要注意，图
$$
\xymatrix{
U \times_\mathcal{X} U' \ar[d] \ar[r] & U' \ar[d] \\
U \times_\mathcal{Y} V' \ar[r] & \mathcal{X} \times_\mathcal{Y} V'
}
$$
是笛卡尔图；由于 $U', V', a', b', h'$ 如 (1) 所述，故左侧竖直箭头光滑。
由《空间上的下降》引理
\ref{spaces-descent-lemma-local-source-target-implies} 的 (2)，性质
$\mathcal{P}$ 在靶上光滑局部，因此基变换
$U \times_\mathcal{Y} V' \to V \times_\mathcal{Y} V'$ 具有
$\mathcal{P}$。由《空间上的下降》引理
\ref{spaces-descent-lemma-local-source-target-implies} 的 (1)，
$\mathcal{P}$ 在源上光滑局部，
所以可前复合光滑态射
$U \times_\mathcal{X} U' \to U \times_\mathcal{Y} V'$，并得出
$(h, h')$ 具有 $\mathcal{P}$。由于 $V \times_\mathcal{Y} V' \to V'$
光滑，由《空间上的下降》引理
\ref{spaces-descent-lemma-local-source-target-implies} 的 (3)，可知
$U \times_\mathcal{X} U' \to V'$ 具有 $\mathcal{P}$。最后，由于
$U \times_X U' \to U'$ 满且光滑，并且 $\mathcal{P}$ 在源上光滑局部
（同一引理），可知 $h'$ 具有 $\mathcal{P}$。这就完成了 (1) 与 (2)
等价性的证明。

\medskip\noindent
若 $\mathcal{X}$ 和 $\mathcal{Y}$ 可表，则应用《空间上的下降》引理
\ref{spaces-descent-lemma-local-source-target-characterize} 可知，
(1) 和 (2) 等价于 $f$ 具有 $\mathcal{P}$。

\medskip\noindent
最后，假设 $f$ 可表，$U, V, a, b, h$ 如引理的 (2) 所述，并且
$\mathcal{P}$ 在任意基变换下保持。必须证明：对任意概形 $Z$ 和态射
$Z \to \mathcal{Y}$，基变换
$Z \times_\mathcal{Y} \mathcal{X} \to Z$
具有性质 $\mathcal{P}$。考虑图
$$
\xymatrix{
Z \times_\mathcal{Y} U \ar[d] \ar[r] &
Z \times_\mathcal{Y} V \ar[d] \\
Z \times_\mathcal{Y} \mathcal{X} \ar[r] &
Z
}
$$
注意，上方水平箭头是 $h$ 的基变换，因而具有性质 $\mathcal{P}$。
左侧竖直箭头光滑且满，右侧竖直箭头光滑。因此，应用《空间上的下降》引理
\ref{spaces-descent-lemma-local-source-target-characterize} 可知
$Z \times_\mathcal{Y} \mathcal{X} \to Z$ 具有性质 $\mathcal{P}$。
\end{proof}

\begin{definition}
\label{stacks-morphisms-definition-P}
设 $\mathcal{P}$ 是代数空间态射的一个在源与靶上光滑局部的性质。
若引理 \ref{stacks-morphisms-lemma-local-source-target} 中的等价条件成立，则称代数叠态射
$f : \mathcal{X} \to \mathcal{Y}$ {\it 具有性质 $\mathcal{P}$}。
\end{definition}

\begin{remark}
\label{stacks-morphisms-remark-composition}
设 $\mathcal{P}$ 是代数空间态射的一个在源与靶上光滑局部且在复合下稳定的
性质。则定义 \ref{stacks-morphisms-definition-P} 所定义的代数叠态射性质在复合下稳定。
事实上，设 $f : \mathcal{X} \to \mathcal{Y}$ 和
$g : \mathcal{Y} \to \mathcal{Z}$ 是具有性质 $\mathcal{P}$ 的代数叠态射。
取代数空间 $W$ 以及满光滑态射 $W \to \mathcal{Z}$。取代数空间 $V$
以及满光滑态射 $V \to \mathcal{Y} \times_\mathcal{Z} W$。最后，
取代数空间 $U$ 以及满光滑态射
$U \to \mathcal{X} \times_\mathcal{Y} V$。根据定义，态射
$V \to W$ 和 $U \to V$ 具有性质 $\mathcal{P}$。由于假设
$\mathcal{P}$ 在复合下稳定，故 $U \to W$ 具有性质 $\mathcal{P}$。
再次由定义可知，$g \circ f : \mathcal{X} \to \mathcal{Z}$ 具有性质
$\mathcal{P}$。
\end{remark}

\begin{remark}
\label{stacks-morphisms-remark-base-change}
设 $\mathcal{P}$ 是代数空间态射的一个在源与靶上光滑局部且在基变换下稳定的
性质。则定义 \ref{stacks-morphisms-definition-P} 所定义的代数叠态射性质在基变换下稳定。
事实上，设 $f : \mathcal{X} \to \mathcal{Y}$ 和
$g : \mathcal{Y}' \to \mathcal{Y}$ 是代数叠的态射，并假设 $f$ 具有性质
$\mathcal{P}$。取代数空间 $V$ 以及满光滑态射 $V \to \mathcal{Y}$。
取代数空间 $U$ 以及满光滑态射
$U \to \mathcal{X} \times_\mathcal{Y} V$。最后，取代数空间 $V'$
以及满光滑态射 $V' \to \mathcal{Y}' \times_\mathcal{Y} V$。
根据定义，态射 $U \to V$ 具有性质 $\mathcal{P}$。由于假设
$\mathcal{P}$ 在基变换下稳定，故 $V' \times_V U \to V'$ 具有性质
$\mathcal{P}$。考虑图
$$
\xymatrix{
V' \times_V U \ar[r] \ar[d] &
\mathcal{Y}' \times_\mathcal{Y} \mathcal{X} \ar[r] \ar[d] &
\mathcal{X} \ar[d] \\
V' \ar[r] & \mathcal{Y}' \ar[r] & \mathcal{Y}
}
$$
可见左上方的水平箭头光滑且满，因此根据定义，投影
$\mathcal{Y}' \times_\mathcal{Y} \mathcal{X} \to \mathcal{Y}'$ 具有性质
$\mathcal{P}$。
\end{remark}

\begin{remark}
\label{stacks-morphisms-remark-implication}
设 $\mathcal{P}, \mathcal{P}'$ 是代数空间态射的两个在源与靶上光滑局部的
性质。假设对代数空间态射有
$\mathcal{P} \Rightarrow \mathcal{P}'$。则对于定义
\ref{stacks-morphisms-definition-P} 中用 $\mathcal{P}$ 和 $\mathcal{P}'$ 定义的代数叠态射
性质，同样有 $\mathcal{P} \Rightarrow \mathcal{P}'$。这由定义显然。
\end{remark}









\section{有限型态射}
\label{stacks-morphisms-section-finite-type}

\noindent
代数空间态射的“局部有限型”性质在源与靶上光滑局部；见《空间上的下降》
注记 \ref{spaces-descent-remark-list-local-source-target}。它还在基变换下稳定，
并且在靶上为 fpqc 局部的；见《空间的态射》引理
\ref{spaces-morphisms-lemma-base-change-finite-type} 和《空间上的下降》引理
\ref{spaces-descent-lemma-descending-property-locally-finite-type}。
因此，由上面的引理 \ref{stacks-morphisms-lemma-local-source-target}，可以如下定义代数空间的态射
局部有限型的含义；当态射可由代数空间表示时，这与《叠的性质》第
\ref{stacks-properties-section-properties-morphisms} 节定义的已有概念一致。

\begin{definition}
\label{stacks-morphisms-definition-locally-finite-type}
设 $f : \mathcal{X} \to \mathcal{Y}$ 是代数叠的态射。
\begin{enumerate}
\item 若引理 \ref{stacks-morphisms-lemma-local-source-target} 中的等价条件对
$\mathcal{P} = \text{locally of finite type}$ 成立，则称 $f$
{\it 局部有限型}。
\item 若 $f$ 局部有限型且拟紧，则称它{\it 有限型}。
\end{enumerate}
\end{definition}

\begin{lemma}
\label{stacks-morphisms-lemma-composition-finite-type}
有限型态射的复合为有限型。局部有限型也同样成立。
\end{lemma}

\begin{proof}
结合注记 \ref{stacks-morphisms-remark-composition} 与《空间的态射》引理
\ref{spaces-morphisms-lemma-composition-finite-type}。
\end{proof}

\begin{lemma}
\label{stacks-morphisms-lemma-base-change-finite-type}
有限型态射的基变换有限型。局部有限型也同样成立。
\end{lemma}

\begin{proof}
结合注记 \ref{stacks-morphisms-remark-base-change} 与《空间的态射》引理
\ref{spaces-morphisms-lemma-base-change-finite-type}。
\end{proof}

\begin{lemma}
\label{stacks-morphisms-lemma-immersion-locally-finite-type}
浸入局部有限型。
\end{lemma}

\begin{proof}
结合注记 \ref{stacks-morphisms-remark-implication} 与《空间的态射》引理
\ref{spaces-morphisms-lemma-immersion-locally-finite-type}。
\end{proof}

\begin{lemma}
\label{stacks-morphisms-lemma-locally-finite-type-locally-noetherian}
设 $f : \mathcal{X} \to \mathcal{Y}$ 是代数叠的态射。若 $f$ 局部有限型且
$\mathcal{Y}$ 局部诺特，则 $\mathcal{X}$ 局部诺特。
\end{lemma}

\begin{proof}
设
$$
\xymatrix{
U \ar[d] \ar[r] & V \ar[d] \\
\mathcal{X} \ar[r] & \mathcal{Y}
}
$$
是交换图，其中 $U$、$V$ 是概形，$V \to \mathcal{Y}$ 满且光滑，
$U \to V \times_\mathcal{Y} \mathcal{X}$ 满且光滑。则 $U \to V$
局部有限型。若 $\mathcal{Y}$ 局部诺特，则 $V$ 局部诺特。由《态射》引理
\ref{morphisms-lemma-finite-type-noetherian}，$U$ 局部诺特，
这意味着 $\mathcal{X}$ 局部诺特。
\end{proof}

\noindent
下面两个引理将在后文得到改进（在讨论局部有限表示的代数叠态射之后）。

\begin{lemma}
\label{stacks-morphisms-lemma-check-finite-type-covering}
设 $f : \mathcal{X} \to \mathcal{Y}$ 是代数叠的态射。设
$W \to \mathcal{Y}$ 满、平坦且局部有限表示，其中 $W$ 是代数空间。
若基变换 $W \times_\mathcal{Y} \mathcal{X} \to W$ 局部有限型，
则 $f$ 局部有限型。
\end{lemma}

\begin{proof}
取代数空间 $V$ 以及满光滑态射 $V \to \mathcal{Y}$。取代数空间 $U$
以及满光滑态射 $U \to V \times_\mathcal{Y} \mathcal{X}$。
必须证明 $U \to V$ 局部有限表示。现在把所有对象沿
$W \to \mathcal{Y}$ 基变换：令
$U' = W \times_\mathcal{Y} U$,
$V' = W \times_\mathcal{Y} V$,
$\mathcal{X}' = W \times_\mathcal{Y} \mathcal{X}$,
以及 $\mathcal{Y}' = W \times_\mathcal{Y} \mathcal{Y} = W$。
由基变换，$U' \to V' \times_{\mathcal{Y}'} \mathcal{X}'$ 仍光滑。
因此，根据代数叠局部有限型态射的定义以及
$\mathcal{X}' \to \mathcal{Y}'$ 局部有限型这一假设，
$U' \to V'$ 局部有限型。又因为 $V' \to V$ 作为
$W \to \mathcal{Y}$ 的基变换满、平坦且局部有限表示，由《空间上的下降》引理
\ref{spaces-descent-lemma-descending-property-locally-finite-type}，
$U \to V$ 局部有限型，证明完成。
\end{proof}

\begin{lemma}
\label{stacks-morphisms-lemma-check-finite-type-precompose}
设 $\mathcal{X} \to \mathcal{Y} \to \mathcal{Z}$ 是代数叠的态射。
假设 $\mathcal{X} \to \mathcal{Z}$ 局部有限型，并且
$\mathcal{X} \to \mathcal{Y}$ 可由代数空间表示、满、平坦且局部有限表示。
则 $\mathcal{Y} \to \mathcal{Z}$ 局部有限型。
\end{lemma}

\begin{proof}
取代数空间 $W$ 以及满光滑态射 $W \to \mathcal{Z}$。取代数空间 $V$
以及满光滑态射 $V \to W \times_\mathcal{Z} \mathcal{Y}$。令
$U = V \times_\mathcal{Y} \mathcal{X}$，它是代数空间。已知
$U \to V$ 满、平坦且局部有限表示，而 $U \to W$ 局部有限型。
因此，该引理归结为代数空间态射的情形，而该情形就是《空间上的下降》引理
\ref{spaces-descent-lemma-locally-finite-type-fppf-local-source}。
\end{proof}

\begin{lemma}
\label{stacks-morphisms-lemma-finite-type-permanence}
设 $f : \mathcal{X} \to \mathcal{Y}$、
$g : \mathcal{Y} \to \mathcal{Z}$ 是代数叠的态射。若
$g \circ f : \mathcal{X} \to \mathcal{Z}$ 局部有限型，
则 $f : \mathcal{X} \to \mathcal{Y}$ 局部有限型。
\end{lemma}

\begin{proof}
可以找到图
$$
\xymatrix{
U \ar[r] \ar[d] & V \ar[r] \ar[d] & W \ar[d] \\
\mathcal{X} \ar[r] & \mathcal{Y} \ar[r] & \mathcal{Z}
}
$$
，其中 $U$、$V$、$W$ 是概形，竖直箭头 $W \to \mathcal{Z}$ 满且光滑，
箭头 $V \to \mathcal{Y} \times_\mathcal{Z} W$ 满且光滑，并且箭头
$U \to \mathcal{X} \times_\mathcal{Y} V$ 满且光滑。于是
$U \to \mathcal{X} \times_\mathcal{Z} V$ 也满且光滑（它是满光滑态射与
此类态射的一个基变换的复合）。根据定义，$U \to W$ 局部有限型。
因此，由《态射》引理 \ref{morphisms-lemma-permanence-finite-type}，
$U \to V$ 局部有限型；根据定义，这又意味着
$\mathcal{X} \to \mathcal{Y}$ 局部有限型。
\end{proof}










\section{有限型点}
\label{stacks-morphisms-section-points-finite-type}

\noindent
设 $\mathcal{X}$ 是代数叠。有限型点 $x \in |\mathcal{X}|$ 是可以由局部有限型
态射 $\Spec(k) \to \mathcal{X}$ 表示的点。对于代数空间和代数叠，
有限型点适合作为闭点的替代。例如，它们总是“足够多”。

\begin{lemma}
\label{stacks-morphisms-lemma-point-finite-type}
设 $\mathcal{X}$ 是代数叠，$x \in |\mathcal{X}|$。下列条件等价：
\begin{enumerate}
\item 存在表示 $x$ 的局部有限型态射 $\Spec(k) \to \mathcal{X}$。
\item 存在概形 $U$、闭点 $u \in U$ 以及光滑态射
$\varphi : U \to \mathcal{X}$，使得 $\varphi(u) = x$。
\end{enumerate}
\end{lemma}

\begin{proof}
设 $u \in U$ 和 $U \to \mathcal{X}$ 如 (2) 所述。则
$\Spec(\kappa(u)) \to U$ 有限型，而 $U \to \mathcal{X}$ 可表且局部有限型
（由《空间的态射》引理
\ref{spaces-morphisms-lemma-etale-locally-finite-presentation} 和
\ref{spaces-morphisms-lemma-finite-presentation-finite-type}）。因此，
由引理 \ref{stacks-morphisms-lemma-composition-finite-type}，(1) 成立。

\medskip\noindent
反之，假设 $\Spec(k) \to \mathcal{X}$ 局部有限型并表示 $x$。设
$U \to \mathcal{X}$ 是满光滑态射，其中 $U$ 是概形。根据假设，
$U \times_\mathcal{X} \Spec(k) \to U$ 是局部有限型的代数空间态射。
取 $U \times_\mathcal{X} \Spec(k)$ 的一个有限型点 $v$（至少存在一个；
见《空间的态射》引理
\ref{spaces-morphisms-lemma-identify-finite-type-points}）。由《空间的态射》引理
\ref{spaces-morphisms-lemma-finite-type-points-morphism}，$v$ 的像
$u \in U$ 是 $U$ 的有限型点。因此，由《态射》引理
\ref{morphisms-lemma-identify-finite-type-points}，缩小 $U$ 后可以假设 $u$
是 $U$ 的闭点，即 (2) 成立。
\end{proof}

\begin{definition}
\label{stacks-morphisms-definition-finite-type-point}
设 $\mathcal{X}$ 是代数叠。若引理 \ref{stacks-morphisms-lemma-point-finite-type} 中的等价条件
成立，则称点 $x \in |\mathcal{X}|$ 是{\it 有限型点}\footnote{这略微滥用了
术语，因为称为“局部有限型点”或许更准确。}。以
$\mathcal{X}_{\text{ft-pts}}$ 表示 $\mathcal{X}$ 的有限型点集。
\end{definition}

\noindent
有限型点集可以描述如下。

\begin{lemma}
\label{stacks-morphisms-lemma-identify-finite-type-points}
设 $\mathcal{X}$ 是代数叠。则
$$
\mathcal{X}_{\text{ft-pts}} =
\bigcup\nolimits_{\varphi : U \to \mathcal{X}\text{ smooth}} |\varphi|(U_0)
$$
，其中 $U_0$ 是 $U$ 的闭点集。这里可以让 $U$ 遍历 $\mathcal{X}$ 上所有
光滑概形，或只遍历 $\mathcal{X}$ 上所有光滑仿射概形。
\end{lemma}

\begin{proof}
立即由引理 \ref{stacks-morphisms-lemma-point-finite-type} 得到。
\end{proof}

\begin{lemma}
\label{stacks-morphisms-lemma-finite-type-points-morphism}
设 $f : \mathcal{X} \to \mathcal{Y}$ 是代数叠的态射。若 $f$ 局部有限型，则
$f(\mathcal{X}_{\text{ft-pts}}) \subset \mathcal{Y}_{\text{ft-pts}}$.
\end{lemma}

\begin{proof}
取 $x \in \mathcal{X}_{\text{ft-pts}}$。用局部有限型态射
$x : \Spec(k) \to \mathcal{X}$ 表示 $x$。则由引理
\ref{stacks-morphisms-lemma-composition-finite-type}，$f \circ x$ 局部有限型。因此
$f(x) \in \mathcal{Y}_{\text{ft-pts}}$。
\end{proof}

\begin{lemma}
\label{stacks-morphisms-lemma-finite-type-points-surjective-morphism}
设 $f : \mathcal{X} \to \mathcal{Y}$ 是代数叠的态射。若 $f$ 局部有限型且满，则
$f(\mathcal{X}_{\text{ft-pts}}) = \mathcal{Y}_{\text{ft-pts}}$.
\end{lemma}

\begin{proof}
由引理 \ref{stacks-morphisms-lemma-finite-type-points-morphism}，有
$f(\mathcal{X}_{\text{ft-pts}}) \subset \mathcal{Y}_{\text{ft-pts}}$。
设 $y \in |\mathcal{Y}|$ 是有限型点。用局部有限型态射
$\Spec(k) \to \mathcal{Y}$ 表示 $y$。由于 $f$ 满，代数叠
$\mathcal{X}_k = \Spec(k) \times_\mathcal{Y} \mathcal{X}$ 非空，
故由引理 \ref{stacks-morphisms-lemma-identify-finite-type-points}，它有有限型点
$x \in |\mathcal{X}_k|$。态射 $\mathcal{X}_k \to \mathcal{X}$
是 $\Spec(k) \to \mathcal{Y}$ 的基变换，因而局部有限型
（引理 \ref{stacks-morphisms-lemma-base-change-finite-type}）。所以由引理
\ref{stacks-morphisms-lemma-finite-type-points-morphism}，$x$ 在 $\mathcal{X}$ 中的像是有限型点，
并且根据构造映到 $y$。
\end{proof}

\begin{lemma}
\label{stacks-morphisms-lemma-enough-finite-type-points}
设 $\mathcal{X}$ 是代数叠。对任意局部闭子集 $T \subset |\mathcal{X}|$，有
$$
T \not = \emptyset
\Rightarrow
T \cap \mathcal{X}_{\text{ft-pts}} \not = \emptyset.
$$
特别地，对任意闭子集 $T \subset |\mathcal{X}|$，
$T \cap \mathcal{X}_{\text{ft-pts}}$ 在 $T$ 中稠密。
\end{lemma}

\begin{proof}
令 $i : \mathcal{Z} \to \mathcal{X}$ 为 $T$ 上的约化诱导子叠结构；
见《叠的性质》注记
\ref{stacks-properties-remark-stack-structure-locally-closed-subset}。
浸入局部有限型；见引理 \ref{stacks-morphisms-lemma-immersion-locally-finite-type}。
因此，由引理 \ref{stacks-morphisms-lemma-finite-type-points-morphism}，有
$\mathcal{Z}_{\text{ft-pts}} \subset \mathcal{X}_{\text{ft-pts}} \cap T$.
最后，任何带有到 $\mathcal{Z}$ 的光滑态射的非空仿射概形 $U$ 至少有一个闭点，
故由引理 \ref{stacks-morphisms-lemma-identify-finite-type-points}，$\mathcal{Z}$ 至少有一个
有限型点。引理得证。
\end{proof}

\noindent
下面给出代数叠上有限型点的另一种、更具技术性的刻画。特别地，它告诉我们：
只要 $x$ 是有限型点，$\mathcal{X}$ 在 $x$ 处的剩余胚就存在！

\begin{lemma}
\label{stacks-morphisms-lemma-point-finite-type-monomorphism}
设 $\mathcal{X}$ 是代数叠，$x \in |\mathcal{X}|$。下列条件等价：
\begin{enumerate}
\item $x$ 是有限型点；
\item 存在代数叠 $\mathcal{Z}$，其底拓扑空间 $|\mathcal{Z}|$ 是单点集，
并且存在局部有限型态射 $f : \mathcal{Z} \to \mathcal{X}$，使得
$\{x\} = |f|(|\mathcal{Z}|)$；以及
\item $\mathcal{X}$ 在 $x$ 处的剩余胚 $\mathcal{Z}_x$ 存在，并且包含态射
$\mathcal{Z}_x \to \mathcal{X}$ 局部有限型。
\end{enumerate}
\end{lemma}

\begin{proof}
（本段出现的所有态射都可由代数空间表示，因此《叠的性质》第
\ref{stacks-properties-section-properties-morphisms} 节的约定和结果均适用。）
假设 $x$ 是有限型点。取仿射概形 $U$、闭点 $u \in U$ 以及光滑态射
$\varphi : U \to \mathcal{X}$，使得 $\varphi(u) = x$；见引理
\ref{stacks-morphisms-lemma-identify-finite-type-points}。照例令 $u = \Spec(\kappa(u))$。
令 $R = u \times_\mathcal{X} u$，于是得到代数空间中的群胚
$(u, R, s, t, c)$；见《代数叠》引理
\ref{algebraic-lemma-map-space-into-stack}。投影态射 $R \to u$ 是如下复合：
$$
R = u \times_\mathcal{X} u \to
u \times_\mathcal{X} U \to
u \times_\mathcal{X} X = u
$$
其中第一个箭头有限型（它是概形闭浸入 $u \to U$ 的基变换），
第二个箭头光滑（它是光滑态射 $U \to \mathcal{X}$ 的基变换）。因此，
$s, t : R \to u$ 作为复合局部有限型；见《空间的态射》引理
\ref{spaces-morphisms-lemma-composition-finite-type}。由于 $u$ 是域的谱，
$s, t$ 平坦且局部有限表示（由《空间的态射》引理
\ref{spaces-morphisms-lemma-noetherian-finite-type-finite-presentation}）。
由《可表性判据》定理
\ref{criteria-theorem-flat-groupoid-gives-algebraic-stack}，
$\mathcal{Z} = [u/R]$ 是代数叠。
由《代数叠》引理 \ref{algebraic-lemma-map-space-into-stack}，得到典范态射
$$
f : \mathcal{Z} \longrightarrow \mathcal{X}
$$
，它全忠实。因此，该态射可由代数空间表示；见《代数叠》引理
\ref{algebraic-lemma-characterize-representable-by-algebraic-spaces}；
并且是单态射；见《叠的性质》引理
\ref{stacks-properties-lemma-monomorphism}。由此，$\mathcal{X}$ 在 $x$ 处的
剩余胚 $\mathcal{Z}_x \subset \mathcal{X}$ 存在，并且 $f$ 通过一个等价
$\mathcal{Z} \to \mathcal{Z}_x$ 分解；见《叠的性质》引理
\ref{stacks-properties-lemma-residual-gerbe-unique}。根据构造，图
$$
\xymatrix{
u \ar[d] \ar[r] & U \ar[d] \\
\mathcal{Z} \ar[r]^f & \mathcal{X}
}
$$
交换。由《可表性判据》引理
\ref{criteria-lemma-flat-quotient-flat-presentation}，左侧竖直箭头满、
平坦且局部有限表示。考虑图
$$
\xymatrix{
u \times_\mathcal{X} U \ar[d] \ar[r] &
\mathcal{Z} \times_\mathcal{X} U \ar[r] \ar[d] & U \ar[d] \\
u \ar[r] & \mathcal{Z} \ar[r]^f & \mathcal{X}
}
$$
由于 $u \to \mathcal{X}$ 局部有限型，其基变换
$u \times_\mathcal{X} U \to U$ 局部有限型。此外，
$u \times_\mathcal{X} U \to \mathcal{Z} \times_\mathcal{X} U$
作为 $u \to \mathcal{Z}$ 的基变换满、平坦且局部有限表示。因此
$\{u \times_\mathcal{X} U \to \mathcal{Z} \times_\mathcal{X} U\}$
是代数空间的 fppf 覆盖，故由《空间上的下降》引理
\ref{spaces-descent-lemma-locally-finite-presentation-fppf-local-source}，
$\mathcal{Z} \times_\mathcal{X} U \to U$ 局部有限型。根据定义，
这意味着 $f$ 局部有限型（因为竖直箭头
$\mathcal{Z} \times_\mathcal{X} U \to \mathcal{Z}$ 作为
$U \to \mathcal{X}$ 的基变换光滑，并且由于 $\mathcal{Z}$ 只有一个点而满）。
由于 $\mathcal{Z} = \mathcal{Z}_x$，(3) 成立。

\medskip\noindent
(3) 蕴含 (2) 是显然的。若 (2) 成立，则由引理
\ref{stacks-morphisms-lemma-finite-type-points-morphism}，并用引理
\ref{stacks-morphisms-lemma-enough-finite-type-points} 看出
$\mathcal{Z}_{\text{ft-pts}}$ 非空，即 $\mathcal{Z}$ 的唯一点是
$\mathcal{Z}$ 的有限型点，可知 $x$ 是 $\mathcal{X}$ 的有限型点。
\end{proof}







\section{自同构群}
\label{stacks-morphisms-section-automorphism-groups}

\noindent
设 $\mathcal{X}$ 是代数叠。设 $x \in |\mathcal{X}|$ 对应于
$x : \Spec(k) \to \mathcal{X}$。在这种情形下，我们常说“设 $G_x/k$
是 $x$ 的自同构群代数空间”。这只是指
$$
G_x = \mathit{Isom}_\mathcal{X}(x, x) =
\Spec(k) \times_\mathcal{X} \mathcal{I}_\mathcal{X}
$$
是 $x$ 的自同构所成的群代数空间。它是 $\Spec(k)$ 上的群代数空间。
若 $k'/k$ 是域扩张，则所诱导态射 $x' : \Spec(k') \to \mathcal{X}$
的自同构群代数空间是 $G_x$ 到 $\Spec(k')$ 的基变换。

\begin{lemma}
\label{stacks-morphisms-lemma-automorphism-group-scheme}
在上述情形下，若下列条件之一成立，则 $G_x$ 是概形：
\begin{enumerate}
\item $\Delta : \mathcal{X} \to \mathcal{X} \times \mathcal{X}$
拟分离；
\item $\Delta : \mathcal{X} \to \mathcal{X} \times \mathcal{X}$
局部分离；
\item $\mathcal{X}$ 是拟 DM；
\item $\mathcal{I}_\mathcal{X} \to \mathcal{X}$
拟分离；
\item $\mathcal{I}_\mathcal{X} \to \mathcal{X}$
局部分离；或
\item $\mathcal{I}_\mathcal{X} \to \mathcal{X}$
局部拟有限。
\end{enumerate}
\end{lemma}

\begin{proof}
由引理 \ref{stacks-morphisms-lemma-diagonal-diagonal}，注意到 (1) $\Rightarrow$ (4)、
(2) $\Rightarrow$ (5) 且 (3) $\Rightarrow$ (6)。在情形 (4) 中，$G_x$
是拟分离代数空间；在情形 (5) 中，$G_x$ 是局部分离代数空间。
在两种情形下，$G_x$ 都是适宜代数空间（《适宜代数空间》第
\ref{decent-spaces-section-reasonable-decent} 节以及引理
\ref{decent-spaces-lemma-locally-separated-decent}）。于是，由《空间中群胚进阶》
引理 \ref{spaces-more-groupoids-lemma-group-scheme-over-field-separated}，
$G_x$ 分离；继而由《空间中群胚进阶》命题
\ref{spaces-more-groupoids-proposition-group-space-scheme-over-field}，
$G_x$ 是概形。在情形 (6) 中，$G_x \to \Spec(k)$ 局部拟有限，
故由《域上的空间》引理
\ref{spaces-over-fields-lemma-locally-quasi-finite-over-field}，$G_x$ 是概形。
\end{proof}

\begin{lemma}
\label{stacks-morphisms-lemma-property-automorphism-groups}
设 $\mathcal{X}$ 是代数叠，$x \in |\mathcal{X}|$ 是一点。设 $P$
是域上代数空间的一个在基域扩张下不变的性质；例如
$P(X/k) = X \to \Spec(k)\text{ is finite}$.
下列条件等价：
\begin{enumerate}
\item 对 $x$ 的类中的某个态射 $x : \Spec(k) \to \mathcal{X}$，
自同构群代数空间 $G_x/k$ 具有 $P$；以及
\item 对 $x$ 的类中的任意态射 $x : \Spec(k) \to \mathcal{X}$，
自同构群代数空间 $G_x/k$ 具有 $P$。
\end{enumerate}
\end{lemma}

\begin{proof}
略。
\end{proof}

\begin{remark}
\label{stacks-morphisms-remark-property-automorphism-groups}
设 $P$ 是域上代数空间的一个在基域扩张下不变的性质。给定代数叠
$\mathcal{X}$ 和 $x \in |\mathcal{X}|$，若引理
\ref{stacks-morphisms-lemma-property-automorphism-groups} 中的等价条件成立，则称
$\mathcal{X}$ 在 $x$ 处的自同构群具有 $P$。例如，若每当
$x : \Spec(k) \to \mathcal{X}$ 是 $x$ 的代表时，
$G_x \to \Spec(k)$ 都有限，则称{\it $\mathcal{X}$ 在 $x$ 处的自同构群有限}。
光滑、真等性质也同样处理。（这里显然有术语滥用，但我们相信它不会造成
混淆或不精确。）
\end{remark}

\begin{lemma}
\label{stacks-morphisms-lemma-iso-automorphism-groups}
设 $f : \mathcal{X} \to \mathcal{Y}$ 是代数叠的态射，
$x \in |\mathcal{X}|$ 是一点。下列条件等价：
\begin{enumerate}
\item 对 $x$ 的类中的某个态射 $x : \Spec(k) \to \mathcal{X}$，
令 $y = f \circ x$，则自同构群代数空间的映射
$G_x \to G_y$ 是同构；以及
\item 对 $x$ 的类中的任意态射 $x : \Spec(k) \to \mathcal{X}$，
令 $y = f \circ x$，则自同构群代数空间的映射
$G_x \to G_y$ 是同构。
\end{enumerate}
\end{lemma}

\begin{proof}
这归结为同构性在靶上为 fpqc 局部这一事实；见《空间上的下降》引理
\ref{spaces-descent-lemma-descending-property-isomorphism}。具体地，
假设 $k'/k$ 是域扩张，把相应的复合记作
$x' : \Spec(k') \to \mathcal{X}$，并令 $y' = f \circ x'$。则态射
$G_{x'} \to G_{y'}$ 是 $G_x \to G_y$ 沿
$\Spec(k') \to \Spec(k)$ 的基变换。因此，$G_x \to G_y$ 是同构当且仅当
$G_{x'} \to G_{y'}$ 是同构。由此可见，只要此性质对一个代表成立，
它就会传遍整个等价类。
\end{proof}

\begin{remark}
\label{stacks-morphisms-remark-identify-automorphism-groups}
设 $f : \mathcal{X} \to \mathcal{Y}$ 是代数叠的态射，
$x \in |\mathcal{X}|$ 是一点。文献中用“{\it $f$ 在 $x$ 处保持稳定子}”
或“{\it $f$ 在 $x$ 处反映不动点}”来表示引理
\ref{stacks-morphisms-lemma-iso-automorphism-groups} 的等价条件对 $f$ 和 $x$ 成立。
我们更愿意说“{\it $f$ 在 $x$ 与 $f(x)$ 处的自同构群之间诱导同构}”。
\end{remark}





\section{表现与代数叠的性质}
\label{stacks-morphisms-section-presentations-properties}

\noindent
设 $(U, R, s, t, c)$ 是代数空间中的群胚。若
$s, t : R \to U$ 平坦且局部有限表现，则商叠 $[U/R]$
是代数叠；参见《可表示性判据》定理
\ref{criteria-theorem-flat-groupoid-gives-algebraic-stack}.
本节研究 $(U, R, s, t, c)$ 的哪些性质会蕴含代数叠
$[U/R]$ 的相应性质。

\begin{lemma}
\label{stacks-morphisms-lemma-properties-diagonal-from-presentation}
设 $(U, R, s, t, c)$ 是代数空间中的群胚，且
$s, t : R \to U$ 平坦并局部有限表现。考虑代数叠
$\mathcal{X} = [U/R]$（见上文）。
\begin{enumerate}
\item 若 $R \to U \times U$ 分离，则 $\Delta_\mathcal{X}$ 分离。
\item 若 $U$、$R$ 分离，则 $\Delta_\mathcal{X}$ 分离。
\item 若 $R \to U \times U$ 局部拟有限，则 $\mathcal{X}$ 为拟 DM。
\item 若 $s, t : R \to U$ 局部拟有限，则 $\mathcal{X}$ 为拟 DM。
\item 若 $R \to U \times U$ 为真态射，则 $\mathcal{X}$ 分离。
\item 若 $s, t : R \to U$ 为真态射且 $U$ 分离，则
$\mathcal{X}$ 分离。
\item 此处添加更多内容。
\end{enumerate}
\end{lemma}

\begin{proof}
由《可表示性判据》引理
\ref{criteria-lemma-flat-quotient-flat-presentation}，
态射 $U \to \mathcal{X}$ 满、平坦且局部有限表现。因此
$U \times U \to \mathcal{X} \times \mathcal{X}$ 也有这些性质。
有笛卡儿图
$$
\xymatrix{
R  = U \times_\mathcal{X} U \ar[r] \ar[d] & U \times U \ar[d] \\
\mathcal{X} \ar[r] & \mathcal{X} \times \mathcal{X}
}
$$
（参见《空间中的群胚》引理
\ref{spaces-groupoids-lemma-quotient-stack-2-cartesian}).
于是，$\Delta_\mathcal{X}$ 具有《代数叠的性质》一节
\ref{stacks-properties-section-properties-morphisms}
所列性质之一，当且仅当态射 $R \to U \times U$ 具有该性质；参见
《代数叠的性质》引理
\ref{stacks-properties-lemma-check-property-covering}。
这说明了 (1)、(3) 和 (5)。条件 (2) 蕴含
$R \to U \times U$ 分离，故 (2) 由 (1) 得出。条件 (4)
蕴含条件 (3)，故 (4) 由 (3) 得出。由《空间的态射》引理
\ref{spaces-morphisms-lemma-universally-closed-permanence}
可知，条件 (6) 蕴含条件 (5)，故 (6) 由 (5) 得出。
\end{proof}

\begin{lemma}
\label{stacks-morphisms-lemma-points-presentation}
设 $(U, R, s, t, c)$ 是代数空间中的群胚，且
$s, t : R \to U$ 平坦并局部有限表现。考虑代数叠
$\mathcal{X} = [U/R]$（见上文）。则 $|R| \to |U| \times |U|$
的像是一个等价关系，而 $|\mathcal{X}|$ 是 $|U|$ 对此等价关系的商。
\end{lemma}

\begin{proof}
诱导态射 $p : U \to \mathcal{X}$ 满、平坦且局部有限表现；参见
《可表示性判据》引理
\ref{criteria-lemma-flat-quotient-flat-presentation}.
因此，由《代数叠的性质》引理
\ref{stacks-properties-lemma-characterize-surjective}.
可知 $|U| \to |\mathcal{X}|$ 满。注意到
$R = U \times_\mathcal{X} U$；参见《空间中的群胚》引理
\ref{spaces-groupoids-lemma-quotient-stack-2-cartesian}。于是，
《代数叠的性质》引理 \ref{stacks-properties-lemma-points-cartesian}
蕴含映射
$$
|R| \longrightarrow |U| \times_{|\mathcal{X}|} |U|
$$
满。因此，$|R| \to |U| \times |U|$ 的像恰为所有满足如下条件的
$(u_1, u_2) \in |U| \times |U|$ 所成的集合：$u_1$ 与 $u_2$
在 $|\mathcal{X}|$ 中有相同的像。合并这两个陈述即得引理。
\end{proof}






\section{代数叠的特殊表现}
\label{stacks-morphisms-section-presentations}

\noindent
本节证明两个重要定理。第一个把拟 DM 叠 $\mathcal{X}$ 刻画为如下形式的叠：
$\mathcal{X} = [U/R]$，其中 $s, t : R \to U$ 局部拟有限（并且平坦且
局部有限表现）。第二个定理断言，DM 代数叠是 Deligne--Mumford 叠。

\medskip\noindent
下述引理给出一个判据，用以判断表现的一个“截片”何时在该代数叠上仍然平坦。

\begin{lemma}
\label{stacks-morphisms-lemma-slice}
设 $\mathcal{X}$ 是代数叠。考虑笛卡儿图
$$
\xymatrix{
U \ar[d] & F \ar[l]^p \ar[d] \\
\mathcal{X} & \Spec(k) \ar[l]
}
$$
其中 $U$ 是代数空间，$k$ 是域，且 $U \to \mathcal{X}$ 平坦并局部有限表现。
设
$f_1, \ldots, f_r \in \Gamma(U, \mathcal{O}_U)$
且 $z \in |F|$，使 $f_1, \ldots, f_r$ 在局部环
$\mathcal{O}_{F, \overline{z}}$ 中的像构成正则序列。则以一个包含
$p(z)$ 的开子空间替换 $U$ 后，态射
$$
V(f_1, \ldots, f_r) \longrightarrow \mathcal{X}
$$
平坦且局部有限表现。
\end{lemma}

\begin{proof}
选取概形 $W$ 以及满光滑态射 $W \to \mathcal{X}$。选取域扩张
$k'/k$ 和态射 $w : \Spec(k') \to W$，使
$\Spec(k') \to W \to \mathcal{X}$ 与
$\Spec(k') \to \Spec(k) \to \mathcal{X}$ 是 $2$-同构的。
这是可能的，因为 $W \to \mathcal{X}$ 满。考虑交换图
$$
\xymatrix{
U \ar[d] &
U \times_\mathcal{X} W \ar[l]^-{\text{pr}_0} \ar[d] &
F' \ar[l]^-{p'} \ar[d] \\
\mathcal{X} &
W \ar[l] &
\Spec(k') \ar[l]
}
$$
其中两个方块都是笛卡儿的。由 $w$ 的选取可知
$F' = F \times_{\Spec(k)} \Spec(k')$。因此 $F' \to F$ 满，故可选取
映到 $z$ 的点 $z' \in |F'|$。由于 $F' \to F$ 平坦，
$\mathcal{O}_{F, \overline{z}} \to \mathcal{O}_{F', \overline{z}'}$
平坦；参见《空间的态射》引理
\ref{spaces-morphisms-lemma-flat-at-point-etale-local-rings}。
因此，$f_1, \ldots, f_r$ 在 $\mathcal{O}_{F', \overline{z}'}$ 中的像
构成正则序列；参见《代数》引理
\ref{algebra-lemma-flat-increases-depth}。注意，
$U \times_\mathcal{X} W \to W$ 是平坦且局部有限表现的代数空间态射。
因此，由《空间的态射进阶》引理
\ref{spaces-more-morphisms-lemma-slice} 可知，存在
$U \times_\mathcal{X} W$ 的开子空间 $U'$，它包含 $p'(z')$，并且交
$U' \cap (V(f_1, \ldots, f_r) \times_\mathcal{X} W)$ 在 $W$ 上平坦且
局部有限表现。注意，$\text{pr}_0(U')$ 是 $U$ 中包含 $p(z)$ 的开子空间，
因为 $\text{pr}_0$ 光滑，因而是开的。现在可知，态射
$U' \cap (V(f_1, \ldots, f_r) \times_\mathcal{X} W) \to \mathcal{X}$
作为复合
$$
U' \cap (V(f_1, \ldots, f_r) \times_\mathcal{X} W) \to W \to \mathcal{X}.
$$
平坦且局部有限表现。因此，《代数叠的性质》引理
\ref{stacks-properties-lemma-check-property-after-precomposing}
蕴含 $\text{pr}_0(U') \cap V(f_1, \ldots, f_r) \to \mathcal{X}$
平坦且局部有限表现，正合所需。
\end{proof}

\begin{lemma}
\label{stacks-morphisms-lemma-quasi-finite-at-point}
设 $\mathcal{X}$ 是代数叠。考虑笛卡儿图
$$
\xymatrix{
U \ar[d] & F \ar[l]^p \ar[d] \\
\mathcal{X} & \Spec(k) \ar[l]
}
$$
其中 $U$ 是代数空间，$k$ 是域，且 $U \to \mathcal{X}$ 局部有限型。
设 $z \in |F|$ 且 $\dim_z(F) = 0$。则以一个包含 $p(z)$ 的开子空间
替换 $U$ 后，态射
$$
U \longrightarrow \mathcal{X}
$$
局部拟有限。
\end{lemma}

\begin{proof}
由于 $f : U \to \mathcal{X}$ 局部有限型，存在极大开集
$W(f) \subset U$，使限制 $f|_{W(f)} : W(f) \to \mathcal{X}$
局部拟有限；参见《代数叠的性质》评注
\ref{stacks-properties-remark-local-source-apply}
(\ref{stacks-properties-item-loc-quasi-finite}).
因此，只需证明 $p(z)$ 是 $W(f)$ 的点。此外，上述评注还说明，
$W(f)$ 的构造与沿可由代数空间表示的态射所作的任意基变换交换。
故只需证明态射 $F \to \Spec(k)$ 在 $z$ 处局部拟有限。这立即由
《空间的态射》引理
\ref{spaces-morphisms-lemma-locally-quasi-finite-rel-dimension-0} 得出。
\end{proof}

\noindent
拟 DM 叠具有由概形给出的局部拟有限“覆盖”。

\begin{theorem}
\label{stacks-morphisms-theorem-quasi-DM}
设 $\mathcal{X}$ 是代数叠。下列条件等价：
\begin{enumerate}
\item $\mathcal{X}$ 是拟 DM 叠；
\item 存在概形 $W$ 以及满、平坦、局部有限表现且局部拟有限的态射
$W \to \mathcal{X}$。
\end{enumerate}
\end{theorem}

\begin{proof}
(2) $\Rightarrow$ (1) 即引理
\ref{stacks-morphisms-lemma-properties-covering-imply-diagonal}。假设 (1)。
设 $x \in |\mathcal{X}|$ 是有限型点。我们将构造一个在 $\mathcal{X}$
上、并在 $x$ 的某个邻域内“起作用”的概形。证明末尾取所有这些概形的
不交并即可得出结论。

\medskip\noindent
设 $U$ 是仿射概形，$U \to \mathcal{X}$ 是光滑态射，而
$u \in U$ 是映到 $x$ 的闭点；参见引理
\ref{stacks-morphisms-lemma-point-finite-type}。照例记 $u = \Spec(\kappa(u))$。
考虑交换图
$$
\xymatrix{
u \ar[d] & R \ar[l] \ar[d] \\
U \ar[d] & F \ar[d] \ar[l]^p \\
\mathcal{X} & u \ar[l]
}
$$
其中两个方块都是纤维积方块；特别地，$R = u \times_\mathcal{X} u$。
在引理 \ref{stacks-morphisms-lemma-point-finite-type-monomorphism} 的证明中，已经看到
$(u, R, s, t, c)$ 是代数空间中的群胚，且 $s, t$ 局部有限型。
令 $G \to u$ 为稳定子群代数空间（参见《空间中的群胚》定义
\ref{spaces-groupoids-definition-stabilizer-groupoid}).
注意
$$
G = R \times_{(u \times u)} u =
(u \times_\mathcal{X} u) \times_{(u \times u)} u =
\mathcal{X} \times_{\mathcal{X} \times \mathcal{X}} u.
$$
由于 $\mathcal{X}$ 为拟 DM，$G$ 在 $u$ 上局部拟有限。由
《空间中的群胚进阶》引理
\ref{spaces-more-groupoids-lemma-groupoid-on-field-dimension-equal-stabilizer}
可得 $\dim(R) = 0$。

\medskip\noindent
令 $e : u \to R$ 为该群胚的单位元。复合 $u \to R \to u$ 等于
$u$ 的恒等态射。由于 $u \subset U$ 是闭子概形，$R \subset F$
是闭子空间。因此，也可把 $e$ 看作 $F$ 的点，并将其记为 $z$。
考虑 \'etale 局部环之间的映射
$$
\mathcal{O}_{U, u}
\xrightarrow{p^\sharp}
\mathcal{O}_{F, \overline{z}}
\longrightarrow
\mathcal{O}_{R, \overline{e}}
$$
由第一段的结果，$\mathcal{O}_{R, \overline{e}}$ 的维数为 $0$。
另一方面，由于 $R$ 在 $F$ 中由 $\mathfrak m_u$ 截出，第二个箭头的核为
$p^\sharp(\mathfrak m_u)\mathcal{O}_{F, \overline{z}}$。因此
$$
\mathfrak m_{\overline{z}} =
\sqrt{p^\sharp(\mathfrak m_u)\mathcal{O}_{F, \overline{z}}}
$$
另一方面，由于态射 $U \to \mathcal{X}$ 光滑，$F \to u$ 是代数空间的
光滑态射。这意味着 $F$ 是正则代数空间（《域上的空间》引理
\ref{spaces-over-fields-lemma-smooth-regular}）。因此
$\mathcal{O}_{F, \overline{z}}$ 是正则局部环（《空间的性质》引理
\ref{spaces-properties-lemma-regular}）。注意，正则局部环是
Cohen--Macaulay 环（《代数》引理
\ref{algebra-lemma-regular-ring-CM}）。令
$d = \dim(\mathcal{O}_{F, \overline{z}})$。由《代数》引理
\ref{algebra-lemma-find-sequence-image-regular}，可找到
$f_1, \ldots, f_d \in \mathcal{O}_{U, u}$，使其像
$\varphi(f_1), \ldots, \varphi(f_d)$ 在
$\mathcal{O}_{F, \overline{z}}$ 中构成正则序列。由引理
\ref{stacks-morphisms-lemma-slice}，缩小 $U$ 后可假设
$Z = V(f_1, \ldots, f_d) \to \mathcal{X}$ 平坦且局部有限表现。
注意，按构造，$F_Z = Z \times_\mathcal{X} u$ 是
$F = U \times_\mathcal{X} u$ 的闭子空间，$z$ 是该闭子空间的点，并且
$$
\dim(\mathcal{O}_{F_Z, \overline{z}}) = 0.
$$
由于 $z$ 相对于 $u$ 的超越次数为零，由《空间的态射》引理
\ref{spaces-morphisms-lemma-dimension-fibre-at-a-point}
可得 $\dim_z(F_Z) = 0$。因而由引理
\ref{stacks-morphisms-lemma-quasi-finite-at-point} 可知，必要时再次缩小 $U$ 后，态射
$Z \to \mathcal{X}$ 局部拟有限。

\medskip\noindent
综上，对 $\mathcal{X}$ 的每个有限型点 $x$，存在局部拟有限、平坦且
局部有限表现的态射 $f_x : Z_x \to \mathcal{X}$，使 $x$ 属于
$|f_x|$ 的像。令 $W = \coprod_x Z_x$ 及 $f = \coprod f_x$。
则 $f$ 平坦、局部有限表现且局部拟有限。特别地，$|f|$ 的像是开集；
参见《代数叠的性质》引理
\ref{stacks-properties-lemma-topology-points}。按构造，该像包含
$\mathcal{X}$ 的所有有限型点，故由引理
\ref{stacks-morphisms-lemma-enough-finite-type-points}（以及《代数叠的性质》引理
\ref{stacks-properties-lemma-characterize-surjective}）可知 $f$ 满。
\end{proof}

\begin{lemma}
\label{stacks-morphisms-lemma-DM-residual-gerbe}
设 $\mathcal{Z}$ 是 DM、局部诺特且约化的代数叠，并且
$|\mathcal{Z}|$ 是单点集。则存在域 $k$ 以及满的 \'etale 态射
$\Spec(k) \to \mathcal{Z}$。
\end{lemma}

\begin{proof}
由《代数叠的性质》引理
\ref{stacks-properties-lemma-unique-point-better}，存在域 $k$ 以及满、平坦且
局部有限表现的态射 $\Spec(k) \to \mathcal{Z}$。令
$U = \Spec(k)$ 及 $R = U \times_\mathcal{Z} U$，于是得到代数空间中的群胚
$(U, R, s, t, c)$；参见《代数叠》引理
\ref{algebraic-lemma-criterion-map-representable-spaces-fibred-in-groupoids}.
注意，由《代数叠》评注
\ref{algebraic-remark-flat-fp-presentation}，有等价
$$
f_{can} : [U/R] \longrightarrow \mathcal{Z}
$$
投影 $s, t : R \to U$ 局部有限表现。由于 $\mathcal{Z}$ 为 DM，
稳定子群代数空间
$$
G = U \times_{U \times U} R = U \times_{U \times U} (U \times_\mathcal{Z} U) =
U \times_{\mathcal{Z} \times \mathcal{Z}, \Delta_\mathcal{Z}} \mathcal{Z}
$$
在 $U$ 上非分歧。特别地，$\dim(G) = 0$；由《空间中的群胚进阶》引理
\ref{spaces-more-groupoids-lemma-groupoid-on-field-dimension-equal-stabilizer}
可得 $\dim(R) = 0$。这蕴含 $R$ 是概形；参见《域上的空间》引理
\ref{spaces-over-fields-lemma-locally-finite-type-dim-zero}。由《簇》引理
\ref{varieties-lemma-algebraic-scheme-dim-0} 可知，$R$（以及 $G$）是若干
Artin 局部环之谱的不交并，而这些环经 $s$ 或 $t$ 在 $k$ 上有限。令
$P = \Spec(A) \subset R$ 为既开又闭的子概形，其底层点是群胚概形
$(U, R, s, t, c)$ 的单位元 $e$。由于
$s \circ e = t \circ e = \text{id}_{\Spec(k)}$，$A$ 是 Artin 局部环，
其剩余域经 $s^\sharp : k \to A$ 或 $t^\sharp : k \to A$ 与 $k$ 等同。
注意，$s, t : \Spec(A) \to \Spec(k)$ 都是有限态射（由上述引理）。
由于 $G \to \Spec(k)$ 非分歧，
$$
G \cap P = P \times_{U \times U} U = \Spec(A \otimes_{k \otimes k} k)
$$
在 $k$ 上非分歧。另一方面，$A \otimes_{k \otimes k} k$ 作为 $A$ 的商环
是局部环，并满射到 $k$。于是
$A \otimes_{k \otimes k} k = k$。由此，$P \to U \times U$ 泛单射
（因为 $P$ 只有一个点，且其剩余域为 $k$）、非分歧（由上面对唯一像点
之上的纤维的计算）且有限型（因为 $s, t$ 有限型），因而是单态射（参见
《\'Etale 态射》引理
\ref{etale-lemma-universally-injective-unramified}).
因此，$s|_P, t|_P : P \to U$ 定义有限平坦等价关系。应用《群胚》命题
\ref{groupoids-proposition-finite-flat-equivalence} 可知 $U/P$ 存在，且为概形
$\overline{U}$。此外，$U \to \overline{U}$ 有限局部自由，并且
$P = U \times_{\overline{U}} U$。事实上，
$\overline{U} = \Spec(k_0)$，其中 $k_0 \subset k$ 是 $R$-不变函数所成的环。
由于 $k$ 是域，由定义《群胚》等式
(\ref{groupoids-equation-invariants}) 可知 $k_0$ 是域。

\medskip\noindent
我们断言
\begin{equation}
\label{stacks-morphisms-equation-etale-covering}
\Spec(k_0) = \overline{U} = U/P \to [U/R] = \mathcal{Z}
\end{equation}
是所需的满 \'etale 态射。由《代数叠的性质》引理
\ref{stacks-properties-lemma-flat-cover-by-field}，该态射满。因此只需证明
(\ref{stacks-morphisms-equation-etale-covering}) 为 \'etale\footnote{我们诚请读者自行给出
这一事实的证明。事实上，该论证与《自举》定理
\ref{bootstrap-theorem-final-bootstrap} 的证明末尾的论证有很多共同之处，
因而或许应当在某处将其单独提炼为一个引理。}。我们不直接证明 \'etale 性，
而是先应用《自举》引理 \ref{bootstrap-lemma-divide-subgroupoid}，可知存在群胚概形
$(\overline{U}, \overline{R}, \overline{s}, \overline{t}, \overline{c})$
使 $(U, R, s, t, c)$ 是
$(\overline{U}, \overline{R}, \overline{s}, \overline{t}, \overline{c})$
经商态射 $U \to \overline{U}$ 所作的限制。（上面已验证该引理的全部假设，
只有 $j : R \to U \times U$ 分离且局部拟有限这一断言尚未验证；它由
$R$ 是在 $k$ 上局部拟有限的分离概形这一事实得出。）由于
$U \to \overline{U}$ 有限局部自由，$[U/R] \to [\overline{U}/\overline{R}]$
是等价；参见《空间中的群胚》引理
\ref{spaces-groupoids-lemma-quotient-stack-restrict-equivalence}。

\medskip\noindent
注意，$s, t$ 是态射 $\overline{s}, \overline{t}$ 沿
$U \to \overline{U}$ 的基变换。由于 $\{U \to \overline{U}\}$ 是 fppf 覆盖，
可知 $\overline{s}, \overline{t}$ 平坦、局部有限表现且局部拟有限；参见
《下降》引理 \ref{descent-lemma-descending-property-flat}、
\ref{descent-lemma-descending-property-locally-finite-presentation} 及
\ref{descent-lemma-descending-property-quasi-finite}。
考虑交换图
$$
\xymatrix{
U \times_{\overline{U}} U \ar@{=}[r] \ar[rd] & P \ar[r] \ar[d] & R \ar[d] \\
& \overline{U} \ar[r]^{\overline{e}} & \overline{R}
}
$$
关于限制的一般事实说明，外侧四个角构成笛卡儿图。由该等式可知内侧方块
也是笛卡儿的。由于 $P$ 在 $R$ 中为开集，由《下降》引理
\ref{descent-lemma-descending-property-open-immersion} 可知
$\overline{e}$ 是开浸入。

\medskip\noindent
当然，若 $\overline{e}$ 是开浸入，且 $\overline{s}, \overline{t}$ 平坦并
局部有限表现，则态射 $\overline{t}, \overline{s}$ 为 \'etale。例如，应用
《群胚进阶》引理 \ref{more-groupoids-lemma-sheaf-differentials} 可看出这一点：
该引理说明 $\Omega_{\overline{R}/\overline{U}} = 0$ 蕴含
$\overline{s}, \overline{t} : \overline{R} \to \overline{U}$ 非分歧（参见
《态射》引理 \ref{morphisms-lemma-unramified-omega-zero}），进而蕴含
$\overline{s}, \overline{t}$ 为 \'etale（参见《态射》引理
\ref{morphisms-lemma-flat-unramified-etale}）。因此，
$\mathcal{Z} = [\overline{U}/\overline{R}]$ 是代数叠 $\mathcal{Z}$ 的一个
\'etale 表现，并由《代数叠的性质》引理
\ref{stacks-properties-lemma-check-property-covering}
可知 $\overline{U} \to \mathcal{Z}$ 为 \'etale。
\end{proof}

\begin{lemma}
\label{stacks-morphisms-lemma-etale-at-point}
设 $\mathcal{X}$ 是代数叠。考虑笛卡儿图
$$
\xymatrix{
U \ar[d] & F \ar[l]^p \ar[d] \\
\mathcal{X} & \Spec(k) \ar[l]
}
$$
其中 $U$ 是代数空间，$k$ 是域，且 $U \to \mathcal{X}$ 平坦并局部有限表现。
设 $z \in |F|$，使 $F \to \Spec(k)$ 在 $z$ 处非分歧。则以一个包含
$p(z)$ 的开子空间替换 $U$ 后，态射
$$
U \longrightarrow \mathcal{X}
$$
为 \'etale。
\end{lemma}

\begin{proof}
由于 $f : U \to \mathcal{X}$ 平坦且局部有限表现，存在极大开集
$W(f) \subset U$，使限制 $f|_{W(f)} : W(f) \to \mathcal{X}$ 为
\'etale；参见《代数叠的性质》评注
\ref{stacks-properties-remark-local-source-apply}
(\ref{stacks-properties-item-etale}).
因此，只需证明 $p(z)$ 是 $W(f)$ 的点。此外，上述评注还说明，
$W(f)$ 的构造与沿可由代数空间表示的态射所作的任意基变换交换。
故只需证明态射 $F \to \Spec(k)$ 在 $z$ 处为 \'etale。它作为
$U \to \mathcal{X}$ 的基变换是平坦且局部有限表现的，并且按假设
$F \to \Spec(k)$ 在 $z$ 处非分歧，故结论由《空间的态射》引理
\ref{spaces-morphisms-lemma-unramified-flat-lfp-etale} 得出。
\end{proof}

\noindent
DM 叠是 Deligne--Mumford 叠。

\begin{theorem}
\label{stacks-morphisms-theorem-DM}
设 $\mathcal{X}$ 是代数叠。下列条件等价：
\begin{enumerate}
\item $\mathcal{X}$ 为 DM；
\item $\mathcal{X}$ 为 Deligne--Mumford；
\item 存在概形 $W$ 以及满 \'etale 态射 $W \to \mathcal{X}$。
\end{enumerate}
\end{theorem}

\begin{proof}
回忆，(3) 是 (2) 的定义；参见《代数叠》定义
\ref{algebraic-definition-deligne-mumford}。(3) $\Rightarrow$ (1) 即引理
\ref{stacks-morphisms-lemma-properties-covering-imply-diagonal}。假设 (1)。设
$x \in |\mathcal{X}|$ 是有限型点。我们将构造一个在 $\mathcal{X}$ 上、
并在 $x$ 的某个邻域内“起作用”的概形。证明末尾取所有这些概形的不交并
即可得出结论。

\medskip\noindent
由引理 \ref{stacks-morphisms-lemma-point-finite-type-monomorphism}，$\mathcal{X}$ 在 $x$
处的剩余胚 $\mathcal{Z}_x$ 存在，且 $\mathcal{Z}_x \to \mathcal{X}$
局部有限型。由引理 \ref{stacks-morphisms-lemma-separation-properties-residual-gerbe}，
代数叠 $\mathcal{Z}_x$ 为 DM。由引理
\ref{stacks-morphisms-lemma-DM-residual-gerbe}，存在域 $k$ 以及满 \'etale 态射
$z : \Spec(k) \to \mathcal{Z}_x$。特别地，复合
$x : \Spec(k) \to \mathcal{X}$ 局部有限型（由《空间的态射》引理
\ref{spaces-morphisms-lemma-composition-finite-type} 和
\ref{spaces-morphisms-lemma-etale-locally-finite-type}）。

\medskip\noindent
选取概形 $U$ 以及光滑态射 $U \to \mathcal{X}$，使 $x$ 属于
$|U| \to |\mathcal{X}|$ 的像。考虑纤维积方块
$$
\xymatrix{
U \ar[d] & F  \ar[l] \ar[d] \\
\mathcal{X} & \Spec(k) \ar[l]_-x
}
$$
换言之，$F = U \times_{\mathcal{X}, x} \Spec(k)$。由《代数叠的性质》引理
\ref{stacks-properties-lemma-points-cartesian} 可知 $F$ 非空。
由于 $\mathcal{Z}_x \to \mathcal{X}$ 是单态射，有
$$
\Spec(k) \times_{z, \mathcal{Z}_x, z} \Spec(k)
=
\Spec(k) \times_{x, \mathcal{X}, x} \Spec(k)
$$
按 $z$ 的构造，等式两端到 $\Spec(k)$ 的投影为 \'etale。由于
$$
F \times_U F =
(\Spec(k) \times_\mathcal{X} \Spec(k))
\times_{\Spec(k)} F
$$
可知投影映射 $F \times_U F \to F$ 也为 \'etale。因此
$\Delta_{F/U} : F \to F \times_U F$ 为 \'etale（参见《空间的态射》引理
\ref{spaces-morphisms-lemma-etale-permanence}）。由《空间的态射》引理
\ref{spaces-morphisms-lemma-etale-universally-injective-open}
可知 $\Delta_{F/U}$ 是开浸入，进而由《空间的态射》引理
\ref{spaces-morphisms-lemma-diagonal-unramified-morphism}
可知 $F \to U$ 非分歧。

\medskip\noindent
选取非空仿射概形 $V$ 以及 \'etale 态射 $V \to F$。（也可直接使用 $F$
而避免这一步，但这样做似乎更容易说明所发生的情形。）图示如下：
$$
\xymatrix{
U \ar[d] & F  \ar[l] \ar[d] & V \ar[l] \ar[ld] \\
\mathcal{X} & \Spec(k) \ar[l]_-x
}
$$
于是 $V \to \Spec(k)$ 是概形的光滑态射，而 $V \to U$ 是概形的非分歧态射
（参见《空间的态射》引理
\ref{spaces-morphisms-lemma-composition-smooth} 和
\ref{spaces-morphisms-lemma-composition-unramified}）。
选取闭点 $v \in V$，使 $k \subset \kappa(v)$ 有限可分；参见《簇》引理
\ref{varieties-lemma-smooth-separable-closed-points-dense}。令 $u \in U$
为其像点。局部环 $\mathcal{O}_{V, v}$ 正则（参见《簇》引理
\ref{varieties-lemma-smooth-regular}），且来自态射 $V \to U$ 的局部环同态
$$
\varphi : \mathcal{O}_{U, u} \longrightarrow \mathcal{O}_{V, v}
$$
满足 $\varphi(\mathfrak m_u)\mathcal{O}_{V, v} = \mathfrak m_v$；参见
《态射》引理 \ref{morphisms-lemma-unramified-at-point}。因此可找到
$f_1, \ldots, f_d \in \mathcal{O}_{U, u}$，使各像
$\varphi(f_1), \ldots, \varphi(f_d)$ 构成
$\mathfrak m_v/\mathfrak m_v^2$ 在 $\kappa(v)$ 上的一组基。由于
$\mathcal{O}_{V, v}$ 是正则局部环，这蕴含
$\varphi(f_1), \ldots, \varphi(f_d)$ 在 $\mathcal{O}_{V, v}$ 中构成
正则序列（参见《代数》引理 \ref{algebra-lemma-regular-ring-CM}）。
以 $u$ 的一个开邻域替换 $U$ 后，可假设
$f_1, \ldots, f_d \in \Gamma(U, \mathcal{O}_U)$。必要时再以 $u$ 的一个
更小开邻域替换 $U$，由引理 \ref{stacks-morphisms-lemma-slice} 可假设
$V(f_1, \ldots, f_d) \to \mathcal{X}$ 平坦且局部有限表现。按构造，
$$
V(f_1, \ldots, f_d) \times_\mathcal{X} \Spec(k)
\longleftarrow
V(f_1, \ldots, f_d) \times_U V
$$
为 \'etale，且 $V(f_1, \ldots, f_d) \times_U V$ 是由
$f_1|_V, \ldots, f_d|_V$ 截出的闭子概形 $T \subset V$。因此按构造
$v \in T$，并且
$$
\mathcal{O}_{T, v} =
\mathcal{O}_{V, v}/(\varphi(f_1), \ldots, \varphi(f_d)) = \kappa(v)
$$
是 $k$ 的有限可分扩张。由此 $T \to \Spec(k)$ 在 $v$ 处非分歧；参见
《态射》引理 \ref{morphisms-lemma-unramified-at-point}。根据代数空间的
非分歧态射的定义，这意味着
$V(f_1, \ldots, f_d) \times_\mathcal{X} \Spec(k) \to \Spec(k)$
在 $v$ 于 $V(f_1, \ldots, f_d) \times_\mathcal{X} \Spec(k)$ 中的像处
非分歧。应用引理 \ref{stacks-morphisms-lemma-etale-at-point} 可知，再将 $U$ 缩小为
$u$ 的另一个开邻域后，态射 $V(f_1, \ldots, f_d) \to \mathcal{X}$
为 \'etale。

\medskip\noindent
综上，对 $\mathcal{X}$ 的每个有限型点 $x$，存在 \'etale 态射
$f_x : W_x \to \mathcal{X}$，使 $x$ 属于 $|f_x|$ 的像。令
$W = \coprod_x W_x$ 及 $f = \coprod f_x$。则 $f$ 为 \'etale。特别地，
$|f|$ 的像是开集；参见《代数叠的性质》引理
\ref{stacks-properties-lemma-topology-points}。按构造，该像包含
$\mathcal{X}$ 的所有有限型点，故由引理
\ref{stacks-morphisms-lemma-enough-finite-type-points}（以及《代数叠的性质》引理
\ref{stacks-properties-lemma-characterize-surjective}）可知 $f$ 满。
\end{proof}

\noindent
下面这个有用推论说明，代数叠的 DM 态射的“纤维”是 Deligne--Mumford 的。

\begin{lemma}
\label{stacks-morphisms-lemma-DM}
设 $f : \mathcal{X} \to \mathcal{Y}$ 是代数叠的 DM 态射。则
\begin{enumerate}
\item 对每个 DM 代数叠 $\mathcal{Z}$ 和态射
$\mathcal{Z} \to \mathcal{Y}$，存在一个概形以及满 \'etale 态射
$U \to \mathcal{X} \times_\mathcal{Y} \mathcal{Z}$.
\item 对每个代数空间 $Z$ 和态射 $Z \to \mathcal{Y}$，存在一个概形以及
满 \'etale 态射
$U \to \mathcal{X} \times_\mathcal{Y} Z$.
\end{enumerate}
\end{lemma}

\begin{proof}
(1) 的证明。由于 $f$ 为 DM，由引理
\ref{stacks-morphisms-lemma-base-change-separated}，基变换
$\mathcal{X} \times_\mathcal{Y} \mathcal{Z} \to \mathcal{Z}$ 为 DM。
又因 $\mathcal{Z}$ 为 DM，引理 \ref{stacks-morphisms-lemma-separated-over-separated}
蕴含 $\mathcal{X} \times_\mathcal{Y} \mathcal{Z}$ 为 DM。因此存在概形
$U$ 以及满 \'etale 态射
$U \to \mathcal{X} \times_\mathcal{Y} \mathcal{Z}$；参见定理
\ref{stacks-morphisms-theorem-DM}。第 (2) 项是 (1) 的特例，因为代数空间在视为代数叠时
由引理 \ref{stacks-morphisms-lemma-trivial-implications} 可知是 DM。
\end{proof}






\section{Deligne--Mumford 轨迹}
\label{stacks-morphisms-section-DM}

\noindent
每个代数叠都有一个最大的开子叠，并且该子叠是 Deligne--Mumford 叠；
这基本显然，但我们仍在下面写出证明。当然，该子叠可能为空，例如当
$X = [\Spec(\mathbf{Z})/\mathbf{G}_{m, \mathbf{Z}}]$ 时如此。
下面将刻画 DM 轨迹中的点。

\begin{lemma}
\label{stacks-morphisms-lemma-open-DM-locus}
设 $\mathcal{X}$ 是代数叠。存在开子叠
$$
\mathcal{X}'' \subset \mathcal{X}' \subset \mathcal{X}
$$
使 $\mathcal{X}''$ 为 DM，$\mathcal{X}'$ 为拟 DM，并且它们分别是具有
这些性质的最大开子叠。
\end{lemma}

\begin{proof}
这里真正要说的只是：若 $\mathcal{U} \subset \mathcal{X}$ 和
$\mathcal{V} \subset \mathcal{X}$ 是 DM 开子叠，则满足
$|\mathcal{W}| = |\mathcal{U}| \cup |\mathcal{V}|$ 的开子叠
$\mathcal{W} \subset \mathcal{X}$ 也为 DM。（拟 DM 的情形类似。）
尽管这有些取巧，我们用定理 \ref{stacks-morphisms-theorem-DM} 来证明。由该定理，
可选取概形 $U$、$V$ 以及满 \'etale 态射
$U \to \mathcal{U}$、$V \to \mathcal{V}$。于是当然
$U \amalg V \to \mathcal{W}$ 满且为 \'etale。拟 DM 的情形使用定理
\ref{stacks-morphisms-theorem-quasi-DM} 以完全相同的方法证明。
\end{proof}

\begin{lemma}
\label{stacks-morphisms-lemma-points-DM-locus}
设 $\mathcal{X}$ 是代数叠。设 $x \in |\mathcal{X}|$ 对应于
$x : \Spec(k) \to \mathcal{X}$。令 $G_x/k$ 为 $x$ 的自同构群代数空间。
则
\begin{enumerate}
\item $x$ 属于 $\mathcal{X}$ 的 DM 轨迹，当且仅当
$G_x \to \Spec(k)$ 非分歧；
\item $x$ 属于 $\mathcal{X}$ 的拟 DM 轨迹，当且仅当
$G_x \to \Spec(k)$ 局部拟有限。
\end{enumerate}
\end{lemma}

\begin{proof}
(2) 的证明。选取概形 $U$ 以及满光滑态射 $U \to \mathcal{X}$。
考虑纤维积
$$
\xymatrix{
G \ar[r] \ar[d] & \mathcal{I}_\mathcal{X} \ar[d] \\
U \ar[r] & \mathcal{X}
}
$$
回忆，$G$ 是 $U \to \mathcal{X}$ 的自同构群代数空间。由
《空间中的群胚》引理
\ref{spaces-groupoids-lemma-open-over-which-unramified-or-lqf}
可知，存在最大开子概形 $U' \subset U$，使 $G_{U'} \to U'$ 局部拟有限。
此外，$U'$ 的构造与任意基变换交换。特别地，$U'$ 在
$R = U \times_\mathcal{X} U$ 中的两个逆像是 $R$ 的同一个开子空间
（毕竟，两个映射 $R \to \mathcal{X}$ 同构，因而有同构的自同构群空间）。
因此，由《代数叠的性质》引理
\ref{stacks-properties-lemma-substacks-presentation}
可知，$U'$ 是某个开子叠 $\mathcal{X}' \subset \mathcal{X}$ 的逆像，
并且有笛卡儿图
$$
\xymatrix{
G_{U'} \ar[r] \ar[d] & \mathcal{I}_{\mathcal{X}'} \ar[d] \\
U' \ar[r] & \mathcal{X}'
}
$$
因此，态射 $\mathcal{I}_{\mathcal{X}'} \to \mathcal{X}'$ 局部拟有限，
由引理 \ref{stacks-morphisms-lemma-diagonal-diagonal} 第 (5) 项可知 $\mathcal{X}'$ 为拟 DM。
另一方面，若 $\mathcal{W} \subset \mathcal{X}$ 是拟 DM 开子叠，则由
$U'$ 的构造，$\mathcal{W}$ 的逆像 $W \subset U$ 必须包含于
$U'$，因为
$\mathcal{I}_\mathcal{W} =
\mathcal{W} \times_\mathcal{X} \mathcal{I}_\mathcal{X}$
在 $\mathcal{W}$ 上局部拟有限。因此，$\mathcal{X}'$ 是拟 DM 轨迹。
最后，选取域扩张 $K/k$ 以及 $2$-交换图
$$
\xymatrix{
\Spec(K) \ar[r] \ar[d] & \Spec(k) \ar[d]^x \\
U \ar[r] & \mathcal{X}
}
$$
于是得到 $K$ 上群代数空间的同构
$G_x \times_{\Spec(k)} \Spec(K) \cong G \times_U \Spec(K)$
。因此，$G_x$ 在 $k$ 上局部拟有限，当且仅当
$\Spec(K) \to U$ 的像落在 $U'$ 中（为此使用 $U'$ 的构造与基变换交换，
以及对 $\Spec(K) \to \Spec(k)$ 和 $G_x$ 应用《空间中的群胚》引理
\ref{spaces-groupoids-lemma-open-over-which-unramified-or-lqf}
）。这就完成了 (2) 的证明。(1) 的证明完全相同。
\end{proof}









\section{局部拟有限态射}
\label{stacks-morphisms-section-locally-quasi-finite}

\noindent
代数空间态射的“局部拟有限”性质不是在源与靶上光滑局部的，故不能利用
第 \ref{stacks-morphisms-section-local-source-target} 节的材料来定义代数叠的局部拟有限态射。
对于可由代数空间表示的代数叠态射，我们已经知道局部拟有限的含义；参见
《代数叠的性质》一节
\ref{stacks-properties-section-properties-morphisms}.
为了找到适用于一般态射的条件，我们作如下观察。

\begin{lemma}
\label{stacks-morphisms-lemma-representable-by-spaces-quasi-finite}
设 $f : \mathcal{X} \to \mathcal{Y}$ 是代数叠的态射。假设 $f$ 可由
代数空间表示。下列条件等价：
\begin{enumerate}
\item $f$ 局部拟有限（其含义见《代数叠的性质》一节
\ref{stacks-properties-section-properties-morphisms}）；
\item $f$ 局部有限型，且对每个态射 $\Spec(k) \to \mathcal{Y}$，
其中 $k$ 为域，空间 $|\Spec(k) \times_\mathcal{Y} \mathcal{X}|$ 离散。
\end{enumerate}
\end{lemma}

\begin{proof}
假设 (1)。在此情形，代数空间的态射
$\mathcal{X}_k \to \Spec(k)$ 作为 $f$ 的基变换局部拟有限。因此由
《空间的态射》引理 \ref{spaces-morphisms-lemma-locally-quasi-finite}，
$|\mathcal{X}_k|$ 离散。反之，假设 (2)。选取满光滑态射
$V \to \mathcal{Y}$，其中 $V$ 是概形。只需证明代数空间态射
$V \times_\mathcal{Y} \mathcal{X} \to V$ 局部拟有限；参见
《代数叠的性质》引理
\ref{stacks-properties-lemma-check-property-covering}.
按假设，态射 $V \times_\mathcal{Y} \mathcal{X} \to V$ 局部有限型。
对任意态射 $\Spec(k) \to V$，其中 $k$ 为域，有
$$
\Spec(k) \times_V (V \times_\mathcal{Y} \mathcal{X}) =
\Spec(k) \times_\mathcal{Y} \mathcal{X}
$$
按假设，其点空间离散。因此，由《空间的态射》引理
\ref{spaces-morphisms-lemma-locally-quasi-finite} 可知
$V \times_\mathcal{Y} \mathcal{X} \to V$ 局部拟有限。
\end{proof}

\noindent
可由代数空间表示的代数叠态射为拟 DM；参见引理
\ref{stacks-morphisms-lemma-trivial-implications}。结合上述引理可知，在可由代数空间表示的
态射这一情形，下述定义与已有概念并不冲突。

\begin{definition}
\label{stacks-morphisms-definition-locally-quasi-finite}
设 $f : \mathcal{X} \to \mathcal{Y}$ 是代数叠的态射。若 $f$ 为拟 DM、
局部有限型，并且对每个态射 $\Spec(k) \to \mathcal{Y}$，其中 $k$ 为域，
空间 $|\mathcal{X}_k|$ 离散，则称 $f$ {\it 局部拟有限}。
\end{definition}

\noindent
要求 $f$ 为拟 DM 是自然的。例如，设 $k$ 为域，并考虑态射
$\pi : [\Spec(k)/\mathbf{G}_m] \to \Spec(k)$
。它具有单点纤维，并且局部有限型。稍后将看到，该态射光滑且相对维数为
$-1$，而我们希望局部拟有限态射的相对维数为 $0$。还要注意，截面
$\Spec(k) \to [\Spec(k)/\mathbf{G}_m]$ 不具有离散纤维，故不局部拟有限；
而我们希望局部拟有限态射具有如下恒久性质：若
$f : \mathcal{X} \to \mathcal{X}'$ 是代数叠的态射，并且在代数叠
$\mathcal{Y}$ 上局部拟有限，则 $f$ 局部拟有限（事实上有稍强的结论，
参见引理 \ref{stacks-morphisms-lemma-locally-quasi-finite-permanence}）。

\medskip\noindent
上述定义的另一个理由是下面的引理
\ref{stacks-morphisms-lemma-characterize-locally-quasi-finite}；它用 $\mathcal{X}$ 和
$\mathcal{Y}$ 的适当“表现”或“覆盖”的存在性刻画局部拟有限性。

\begin{lemma}
\label{stacks-morphisms-lemma-base-change-locally-quasi-finite}
局部拟有限态射的基变换局部拟有限。
\end{lemma}

\begin{proof}
对拟 DM 态射，这已见于引理 \ref{stacks-morphisms-lemma-base-change-separated}；
对局部有限型态射，这已见于引理
\ref{stacks-morphisms-lemma-base-change-finite-type}。关于纤维的条件显然由基变换继承。
\end{proof}

\begin{lemma}
\label{stacks-morphisms-lemma-locally-quasi-finite-over-field}
设 $\mathcal{X} \to \Spec(k)$ 是局部拟有限态射，其中 $\mathcal{X}$
为代数叠，$k$ 为域。设 $f : V \to \mathcal{X}$ 是局部拟有限态射，
其中 $V$ 为概形。则 $V \to \Spec(k)$ 局部拟有限。
\end{lemma}

\begin{proof}
由引理 \ref{stacks-morphisms-lemma-composition-finite-type} 可知
$V \to \Spec(k)$ 局部有限型。反设 $V \to \Spec(k)$ 不局部拟有限，
以求矛盾。则 $V$ 的点之间存在非平凡特化 $v \leadsto v'$；参见《态射》引理
\ref{morphisms-lemma-quasi-finite-at-point-characterize}。特别地，
$\text{trdeg}_k(\kappa(v)) > \text{trdeg}_k(\kappa(v'))$；参见《态射》引理
\ref{morphisms-lemma-dimension-fibre-specialization}。由于
$|\mathcal{X}|$ 离散，有 $|f|(v) = |f|(v')$。考虑
$R = V \times_\mathcal{X} V$。则 $R$ 是代数空间，而投影
$s, t : R \to V$ 作为 $V \to \mathcal{X}$ 的基变换局部拟有限
（后者可由代数空间表示，故这由《代数叠的性质》一节中的讨论得出
\ref{stacks-properties-section-properties-morphisms}）。由《代数叠的性质》引理
\ref{stacks-properties-lemma-points-cartesian} 可知，存在 $r \in |R|$，
使 $s(r) = v$ 且 $t(r) = v'$。由《空间的态射》引理
\ref{spaces-morphisms-lemma-compare-tr-deg} 可知，$v/k$、$r/k$ 与
$v'/k$ 的超越次数相等。此矛盾证明了引理。
\end{proof}

\begin{lemma}
\label{stacks-morphisms-lemma-composition-locally-quasi-finite}
局部拟有限态射的复合局部拟有限。
\end{lemma}

\begin{proof}
对拟 DM 态射，这已见于引理 \ref{stacks-morphisms-lemma-composition-separated}；
对局部有限型态射，这已见于引理
\ref{stacks-morphisms-lemma-composition-finite-type}。设 $\mathcal{X} \to \mathcal{Y}$ 和
$\mathcal{Y} \to \mathcal{Z}$ 局部拟有限。设 $k$ 为域，并设
$\Spec(k) \to \mathcal{Z}$ 为态射。只需证明
$|\mathcal{X}_k|$ 离散。由引理
\ref{stacks-morphisms-lemma-base-change-locally-quasi-finite}，态射
$\mathcal{X}_k \to \mathcal{Y}_k$ 和
$\mathcal{Y}_k \to \Spec(k)$ 局部拟有限。特别地，$\mathcal{Y}_k$
是拟 DM 代数叠；参见引理
\ref{stacks-morphisms-lemma-separated-implies-morphism-separated}。由定理
\ref{stacks-morphisms-theorem-quasi-DM}，可找到概形 $V$ 以及满、平坦、局部有限表现且
局部拟有限的态射 $V \to \mathcal{Y}_k$。由引理
\ref{stacks-morphisms-lemma-locally-quasi-finite-over-field} 可知，$V$ 在 $k$ 上局部拟有限；
特别地，$|V|$ 离散。态射
$V \times_{\mathcal{Y}_k} \mathcal{X}_k \to \mathcal{X}_k$ 满、平坦且
局部有限表现，故
$|V \times_{\mathcal{Y}_k} \mathcal{X}_k| \to |\mathcal{X}_k|$
满且开。因此只需证明 $|V \times_{\mathcal{Y}_k} \mathcal{X}_k|$ 离散。
注意，$V$ 是若干 Artin 局部 $k$-代数 $A_i$ 之谱的不交并，其剩余域为
$k_i$；参见《簇》引理
\ref{varieties-lemma-algebraic-scheme-dim-0}。因此，只需证明每个
$$
|\Spec(A_i) \times_{\mathcal{Y}_k} \mathcal{X}_k| =
|\Spec(k_i) \times_{\mathcal{Y}_k} \mathcal{X}_k| =
|\Spec(k_i) \times_\mathcal{Y} \mathcal{X}|
$$
离散；这由 $\mathcal{X} \to \mathcal{Y}$ 局部拟有限这一假设得出。
\end{proof}

\noindent
在用覆盖刻画局部拟有限态射之前，我们先对拟 DM 态射这样做。

\begin{lemma}
\label{stacks-morphisms-lemma-characterize-quasi-DM}
设 $f : \mathcal{X} \to \mathcal{Y}$ 是代数叠的态射。下列条件等价：
\begin{enumerate}
\item $f$ 为拟 DM；
\item 对任意态射 $V \to \mathcal{Y}$，其中 $V$ 为代数空间，存在满、平坦、
局部有限表现且局部拟有限的态射
$U \to \mathcal{X} \times_\mathcal{Y} V$，其中 $U$ 为代数空间；
\item 存在代数空间 $U$、$V$，满、平坦且局部有限表现的态射
$V \to \mathcal{Y}$，以及满、平坦、局部有限表现且局部拟有限的态射
$U \to \mathcal{X} \times_\mathcal{Y} V$。
\end{enumerate}
\end{lemma}

\begin{proof}
(2) $\Rightarrow$ (3) 是显然的。

\medskip\noindent
假设 (1)，并设 $V \to \mathcal{Y}$ 如 (2) 所述。则
$\mathcal{X} \times_\mathcal{Y} V \to V$ 为拟 DM；参见引理
\ref{stacks-morphisms-lemma-base-change-separated}。由引理
\ref{stacks-morphisms-lemma-trivial-implications}，代数空间 $V$ 为 DM，因而为拟 DM。
于是，由引理 \ref{stacks-morphisms-lemma-separated-over-separated}，
$\mathcal{X} \times_\mathcal{Y} V$ 为拟 DM。因此可以应用定理
\ref{stacks-morphisms-theorem-quasi-DM} 得到 (2) 所述的态射
$U \to \mathcal{X} \times_\mathcal{Y} V$。

\medskip\noindent
假设 (3)。设 $V \to \mathcal{Y}$ 和
$U \to \mathcal{X} \times_\mathcal{Y} V$ 如 (3) 所述。为证明 $f$ 为拟 DM，
只需证明 $\mathcal{X} \times_\mathcal{Y} V \to V$ 为拟 DM；参见引理
\ref{stacks-morphisms-lemma-check-separated-covering}。由引理
\ref{stacks-morphisms-lemma-properties-covering-imply-diagonal} 可知
$\mathcal{X} \times_\mathcal{Y} V$ 为拟 DM。于是由引理
\ref{stacks-morphisms-lemma-separated-implies-morphism-separated}，
$\mathcal{X} \times_\mathcal{Y} V \to V$ 为拟 DM，故 (1) 成立。
这就完成了引理的证明。
\end{proof}

\begin{lemma}
\label{stacks-morphisms-lemma-characterize-locally-quasi-finite}
设 $f : \mathcal{X} \to \mathcal{Y}$ 是代数叠的态射。下列条件等价：
\begin{enumerate}
\item $f$ 局部拟有限；
\item $f$ 为拟 DM，并且对任意代数空间 $V$ 给出的态射
$V \to \mathcal{Y}$，以及任意代数空间 $U$ 给出的局部拟有限态射
$U \to \mathcal{X} \times_\mathcal{Y} V$，态射 $U \to V$ 局部拟有限；
\item 对代数空间 $V$ 给出的任意态射 $V \to \mathcal{Y}$，存在满、平坦、
局部有限表现且局部拟有限的态射
$U \to \mathcal{X} \times_\mathcal{Y} V$，其中 $U$ 是代数空间，
并且 $U \to V$ 局部拟有限；
\item 存在代数空间 $U$、$V$，满、平坦且局部有限表现的态射
$V \to \mathcal{Y}$，以及满、平坦、局部有限表现且局部拟有限的态射
$U \to \mathcal{X} \times_\mathcal{Y} V$，并且 $U \to V$ 局部拟有限。
\end{enumerate}
\end{lemma}

\begin{proof}
假设 (1)。按假设，$f$ 为拟 DM。设 $V \to \mathcal{Y}$ 和
$U \to \mathcal{X} \times_\mathcal{Y} V$ 如 (2) 所述。由引理
\ref{stacks-morphisms-lemma-composition-locally-quasi-finite}，复合
$U \to \mathcal{X} \times_\mathcal{Y} V \to V$ 局部拟有限。
因此 (1) 蕴含 (2)。

\medskip\noindent
假设 (2)。设 $V \to \mathcal{Y}$ 如 (3) 所述。由引理
\ref{stacks-morphisms-lemma-characterize-quasi-DM}，可找到代数空间 $U$ 以及满、平坦、
局部有限表现且局部拟有限的态射
$U \to \mathcal{X} \times_\mathcal{Y} V$。由 (2)，复合
$U \to V$ 局部拟有限。因此 (2) 蕴含 (3)。

\medskip\noindent
(3) 蕴含 (4) 是显然的。

\medskip\noindent
假设 (4)。我们将证明 (1) 成立，从而完成证明。由引理
\ref{stacks-morphisms-lemma-characterize-quasi-DM} 可知 $f$ 为拟 DM。为证明 $f$ 局部有限型，
只需证明 $g : \mathcal{X} \times_\mathcal{Y} V \to V$ 局部有限型；参见引理
\ref{stacks-morphisms-lemma-check-finite-type-covering}。进而只需检验 $g$ 与
$h : U \to \mathcal{X} \times_\mathcal{Y} V$ 预复合后局部有限型；参见引理
\ref{stacks-morphisms-lemma-check-finite-type-precompose}。由于按假设
$g \circ h : U \to V$ 局部拟有限，这一点成立，故 $f$ 局部有限型。
最后，设 $k$ 为域，并设 $\Spec(k) \to \mathcal{Y}$ 为态射。则
$V \times_\mathcal{Y} \Spec(k)$ 是非空代数空间，且在 $k$ 上局部有限表现。
因此可找到有限扩张 $k'/k$ 以及态射 $\Spec(k') \to V$，使
$$
\xymatrix{
\Spec(k') \ar[r] \ar[d] & V \ar[d] \\
\Spec(k) \ar[r] & \mathcal{Y}
}
$$
交换（细节略）。于是 $\mathcal{X}_{k'} \to \mathcal{X}_k$ 可由概形表示、
满且有限局部自由。特别地，
$|\mathcal{X}_{k'}| \to |\mathcal{X}_k|$ 满且开。因此只需证明
$|\mathcal{X}_{k'}|$ 离散。由于
$$
U \times_V \Spec(k') =
U \times_{\mathcal{X} \times_\mathcal{Y} V} \mathcal{X}_{k'}
$$
可知 $U \times_V \Spec(k') \to \mathcal{X}_{k'}$ 满、平坦且局部有限表现
（作为 $U \to \mathcal{X} \times_\mathcal{Y} V$ 的基变换）。因此
$|U \times_V \Spec(k')| \to |\mathcal{X}_{k'}|$ 满且开。故只需证明
$|U \times_V \Spec(k')|$ 离散。这由 $U \to V$ 局部拟有限这一事实得出
（或者使用上面的定义，或者由代数空间态射的原定义并借助《空间的态射》引理
\ref{spaces-morphisms-lemma-locally-quasi-finite} 得出）。
\end{proof}

\begin{lemma}
\label{stacks-morphisms-lemma-locally-quasi-finite-permanence}
设 $\mathcal{X} \to \mathcal{Y} \to \mathcal{Z}$ 是代数叠的态射。
假设 $\mathcal{X} \to \mathcal{Z}$ 局部拟有限，且
$\mathcal{Y} \to \mathcal{Z}$ 为拟 DM。则
$\mathcal{X} \to \mathcal{Y}$ 局部拟有限。
\end{lemma}

\begin{proof}
把 $\mathcal{X} \to \mathcal{Y}$ 写成复合
$$
\mathcal{X} \longrightarrow
\mathcal{X} \times_\mathcal{Z} \mathcal{Y} \longrightarrow
\mathcal{Y}
$$
第二个箭头作为 $\mathcal{X} \to \mathcal{Z}$ 的基变换局部拟有限；参见引理
\ref{stacks-morphisms-lemma-base-change-locally-quasi-finite}。由于
$\mathcal{Y} \to \mathcal{Z}$ 为拟 DM，由引理
\ref{stacks-morphisms-lemma-semi-diagonal}，第一个箭头局部拟有限。因此由引理
\ref{stacks-morphisms-lemma-composition-locally-quasi-finite}，
$\mathcal{X} \to \mathcal{Y}$ 局部拟有限。
\end{proof}

















\section{拟有限态射}
\label{stacks-morphisms-section-quasi-finite}

\noindent
我们已在第 \ref{stacks-morphisms-section-locally-quasi-finite} 节定义了代数叠的“局部拟有限”
态射，并在第 \ref{stacks-morphisms-section-quasi-compact} 节定义了代数叠的“拟紧”态射。
根据定义，代数空间的态射拟有限，当且仅当它既局部拟有限又拟紧
（《空间的态射》定义
\ref{spaces-morphisms-definition-locally-quasi-finite}),
）。因此，可如下定义代数叠态射拟有限的含义；当该态射可由代数空间表示时，
它与《代数叠的性质》一节
\ref{stacks-properties-section-properties-morphisms} 中已经定义的概念一致。

\begin{definition}
\label{stacks-morphisms-definition-quasi-finite}
\begin{reference}
\cite{rydh_approx}
\end{reference}
设 $f : \mathcal{X} \to \mathcal{Y}$ 是代数叠的态射。若 $f$ 局部拟有限
（定义 \ref{stacks-morphisms-definition-locally-quasi-finite}）且拟紧
（定义 \ref{stacks-morphisms-definition-quasi-compact}），则称 $f$ {\it 拟有限}。
\end{definition}

\begin{lemma}
\label{stacks-morphisms-lemma-composition-quasi-finite}
拟有限态射的复合拟有限。
\end{lemma}

\begin{proof}
合并引理 \ref{stacks-morphisms-lemma-composition-locally-quasi-finite} 和
\ref{stacks-morphisms-lemma-composition-quasi-compact}。
\end{proof}

\begin{lemma}
\label{stacks-morphisms-lemma-base-change-quasi-finite}
拟有限态射的基变换拟有限。
\end{lemma}

\begin{proof}
合并引理 \ref{stacks-morphisms-lemma-base-change-locally-quasi-finite} 和
\ref{stacks-morphisms-lemma-base-change-quasi-compact}。
\end{proof}

\begin{lemma}
\label{stacks-morphisms-lemma-quasi-finite-permanence}
设 $f : \mathcal{X} \to \mathcal{Y}$ 和
$g : \mathcal{Y} \to \mathcal{Z}$ 是代数叠的态射。若 $g \circ f$ 拟有限，
且 $g$ 拟分离并为拟 DM，则 $f$ 拟有限。
\end{lemma}

\begin{proof}
合并引理 \ref{stacks-morphisms-lemma-locally-quasi-finite-permanence} 和
\ref{stacks-morphisms-lemma-quasi-compact-permanence}。
\end{proof}












\section{平坦态射}
\label{stacks-morphisms-section-flat}

\noindent
代数空间态射的“平坦”性质在源与靶上光滑局部；参见《空间上的下降》评注
\ref{spaces-descent-remark-list-local-source-target}。它还在基变换下稳定，
并且在靶上是 fpqc 局部的；参见《空间的态射》引理
\ref{spaces-morphisms-lemma-base-change-flat}
以及《空间上的下降》引理
\ref{spaces-descent-lemma-descending-property-flat}.
因此，由上述引理 \ref{stacks-morphisms-lemma-local-source-target}，可如下定义代数空间态射
何谓平坦；当该态射可由代数空间表示时，这与《代数叠的性质》一节
\ref{stacks-properties-section-properties-morphisms} 中已有的定义一致。

\begin{definition}
\label{stacks-morphisms-definition-flat}
设 $f : \mathcal{X} \to \mathcal{Y}$ 是代数叠的态射。若引理
\ref{stacks-morphisms-lemma-local-source-target} 的等价条件对
$\mathcal{P} = \text{flat}$ 成立，则称 $f$ {\it 平坦}。
\end{definition}

\begin{lemma}
\label{stacks-morphisms-lemma-composition-flat}
平坦态射的复合平坦。
\end{lemma}

\begin{proof}
合并评注 \ref{stacks-morphisms-remark-composition} 与《空间的态射》引理
\ref{spaces-morphisms-lemma-composition-flat}。
\end{proof}

\begin{lemma}
\label{stacks-morphisms-lemma-base-change-flat}
平坦态射的基变换平坦。
\end{lemma}

\begin{proof}
合并评注 \ref{stacks-morphisms-remark-base-change} 与《空间的态射》引理
\ref{spaces-morphisms-lemma-base-change-flat}。
\end{proof}

\begin{lemma}
\label{stacks-morphisms-lemma-descent-flat}
设 $f : \mathcal{X} \to \mathcal{Y}$ 是代数叠的态射。设
$\mathcal{Z} \to \mathcal{Y}$ 是代数叠的满平坦态射。若基变换
$\mathcal{Z} \times_\mathcal{Y} \mathcal{X} \to \mathcal{Z}$ 平坦，
则 $f$ 平坦。
\end{lemma}

\begin{proof}
选取代数空间 $W$ 以及满光滑态射 $W \to \mathcal{Z}$。则
$W \to \mathcal{Z}$ 满且平坦（《空间的态射》引理
\ref{spaces-morphisms-lemma-smooth-flat}），故 $W \to \mathcal{Y}$ 满且平坦
（由《代数叠的性质》引理
\ref{stacks-properties-lemma-composition-surjective}
和引理 \ref{stacks-morphisms-lemma-composition-flat}）。由于
$\mathcal{Z} \times_\mathcal{Y} \mathcal{X} \to \mathcal{Z}$
沿 $W \to \mathcal{Z}$ 的基变换是平坦态射（引理
\ref{stacks-morphisms-lemma-base-change-flat}），可用 $W$ 替换 $\mathcal{Z}$。

\medskip\noindent
选取代数空间 $V$ 以及满光滑态射 $V \to \mathcal{Y}$。选取代数空间 $U$
以及满光滑态射 $U \to V \times_\mathcal{Y} \mathcal{X}$。需要证明
$U \to V$ 平坦。现在沿 $W \to \mathcal{Y}$ 对所有对象作基变换：令
$U' = W \times_\mathcal{Y} U$，
$V' = W \times_\mathcal{Y} V$,
$\mathcal{X}' = W \times_\mathcal{Y} \mathcal{X}$,
以及 $\mathcal{Y}' = W \times_\mathcal{Y} \mathcal{Y} = W$。由基变换，
$U' \to V' \times_{\mathcal{Y}'} \mathcal{X}'$ 仍然光滑。因此，由我们对
代数叠的平坦态射的定义，以及 $\mathcal{X}' \to \mathcal{Y}'$ 平坦这一假设，
可知 $U' \to V'$ 平坦。又因 $V' \to V$ 作为
$W \to \mathcal{Y}$ 的基变换满，由《空间的态射》引理
\ref{spaces-morphisms-lemma-base-change-module-flat} (2)
可知 $U \to V$ 平坦，得证。
\end{proof}

\begin{lemma}
\label{stacks-morphisms-lemma-flat-permanence}
设 $\mathcal{X} \to \mathcal{Y} \to \mathcal{Z}$ 是代数叠的态射。
若 $\mathcal{X} \to \mathcal{Z}$ 平坦，且
$\mathcal{X} \to \mathcal{Y}$ 满且平坦，则
$\mathcal{Y} \to \mathcal{Z}$ 平坦。
\end{lemma}

\begin{proof}
选取代数空间 $W$ 以及满光滑态射 $W \to \mathcal{Z}$。选取代数空间 $V$
以及满光滑态射 $V \to W \times_\mathcal{Z} \mathcal{Y}$。选取代数空间
$U$ 以及满光滑态射 $U \to V \times_\mathcal{Y} \mathcal{X}$。
我们知道 $U \to V$ 平坦，且 $U \to W$ 平坦。此外，由于
$\mathcal{X} \to \mathcal{Y}$ 满，$U \to V$ 也满（作为满态射的复合）。
因此，引理归结为代数空间态射的情形。该情形即《空间的态射》引理
\ref{spaces-morphisms-lemma-flat-permanence}.
\end{proof}

\begin{lemma}
\label{stacks-morphisms-lemma-lift-valuation-ring-through-flat-morphism}
设 $f : \mathcal{X} \to \mathcal{Y}$ 是代数叠的平坦态射。设
$\Spec(A) \to \mathcal{Y}$ 为态射，其中 $A$ 是赋值环。若
$\Spec(A)$ 的闭点映到 $|\mathcal{Y}|$ 中属于
$|\mathcal{X|} \to |\mathcal{Y}|$ 之像的点，则存在交换图
$$
\xymatrix{
\Spec(A') \ar[r] \ar[d] & \mathcal{X} \ar[d] \\
\Spec(A) \ar[r] & \mathcal{Y}
}
$$
其中 $A \to A'$ 是赋值环扩张（《代数进阶》定义
\ref{more-algebra-definition-extension-valuation-rings}).
\end{lemma}

\begin{proof}
基变换 $\mathcal{X}_A \to \Spec(A)$ 平坦（引理
\ref{stacks-morphisms-lemma-base-change-flat}），且 $\Spec(A)$ 的闭点属于
$|\mathcal{X}_A| \to |\Spec(A)|$ 的像（《代数叠的性质》引理
\ref{stacks-properties-lemma-points-cartesian}）。因此可假设
$\mathcal{Y} = \Spec(A)$。设 $U \to \mathcal{X}$ 为满光滑态射，
其中 $U$ 为概形。于是对态射 $U \to \Spec(A)$ 应用《空间的态射》引理
\ref{spaces-morphisms-lemma-lift-valuation-ring-through-flat-morphism}
即得结论。
\end{proof}







\section{在一点处平坦}
\label{stacks-morphisms-section-flat-at-point}

\noindent
我们仍需发展一般机制，才能说明代数叠的态射在一点处具有某个给定性质的
含义。目前，下述引理已经足够。

\begin{lemma}
\label{stacks-morphisms-lemma-flat-at-point}
设 $f : \mathcal{X} \to \mathcal{Y}$ 是代数叠的态射。设
$x \in |\mathcal{X}|$。考虑交换图
$$
\vcenter{
\xymatrix{
U \ar[d]_a \ar[r]_h & V \ar[d]^b \\
\mathcal{X} \ar[r]^f & \mathcal{Y}
}
}
\quad\text{带标记点}
\vcenter{
\xymatrix{
u \in |U| \ar[d] \\
x \in |\mathcal{X}|
}
}
$$
其中 $U$、$V$ 是代数空间，$b$ 平坦，并且
$(a, h) : U \to \mathcal{X} \times_\mathcal{Y} V$
平坦。下列条件等价：
\begin{enumerate}
\item 对某一个上述图，$h$ 在 $u$ 处平坦；
\item 对每一个上述图，$h$ 在 $u$ 处平坦。
\end{enumerate}
\end{lemma}

\begin{proof}
假设给定引理中所述的第二个图 $U', V', u', a', b', h'$。于是可以考虑
$$
\xymatrix{
U \ar[d] & U \times_\mathcal{X} U' \ar[l] \ar[d] \ar[r] & U' \ar[d] \\
V & V \times_\mathcal{Y} V' \ar[l] \ar[r] & V'
}
$$
由《代数叠的性质》引理
\ref{stacks-properties-lemma-points-cartesian}，存在点
$u'' \in |U \times_\mathcal{X} U'|$ 映到 $u$ 和 $u'$。若 $h$ 在 $u$ 处平坦，
则基变换
$U \times_V (V \times_\mathcal{Y} V') \to V \times_\mathcal{Y} V'$
在 $u$ 上方的任一点处平坦；参见《空间的态射》引理
\ref{spaces-morphisms-lemma-base-change-module-flat}。另一方面，态射
$$
U \times_\mathcal{X} U' \to
U \times_\mathcal{X} (\mathcal{X} \times_\mathcal{Y} V') =
U \times_\mathcal{Y} V' =
U \times_V (V \times_\mathcal{Y} V')
$$
作为 $(a', h')$ 的基变换平坦；参见引理
\ref{stacks-morphisms-lemma-base-change-flat}。作复合并应用《空间的态射》引理
\ref{spaces-morphisms-lemma-composition-module-flat}，可知
$U \times_\mathcal{X} U' \to V \times_\mathcal{Y} V'$ 在 $u''$ 处平坦。
再与平坦映射 $V \times_\mathcal{Y} V' \to V'$ 复合，可知
$U \times_\mathcal{X} U' \to V'$ 在 $u''$ 处平坦。最后，由于
$U \times_\mathcal{X} U' \to U'$ 在 $u''$ 处平坦，且 $u''$ 映到 $u'$，
由《空间的态射》引理 \ref{spaces-morphisms-lemma-flat-permanence} 可知
$U' \to V'$ 在 $u'$ 处平坦。
\end{proof}

\begin{definition}
\label{stacks-morphisms-definition-flat-at-point}
设 $f : \mathcal{X} \to \mathcal{Y}$ 是代数叠的态射。设
$x \in |\mathcal{X}|$。若引理 \ref{stacks-morphisms-lemma-flat-at-point} 的等价条件成立，
则称 $f$ {\it 在 $x$ 处平坦}。
\end{definition}





\section{有限表现态射}
\label{stacks-morphisms-section-finite-presentation}

\noindent
代数空间态射的“局部有限表现”性质在源与靶上光滑局部；参见
《空间上的下降》评注
\ref{spaces-descent-remark-list-local-source-target}。它还在基变换下稳定，
并且在靶上是 fpqc 局部的；参见《空间的态射》引理
\ref{spaces-morphisms-lemma-base-change-finite-presentation} 以及
《空间上的下降》引理
\ref{spaces-descent-lemma-descending-property-locally-finite-presentation}.
因此，由上述引理 \ref{stacks-morphisms-lemma-local-source-target}，可如下定义代数叠态射
何谓局部有限表现；当该态射可由代数空间表示时，这与《代数叠的性质》一节
\ref{stacks-properties-section-properties-morphisms} 中已有的定义一致。

\begin{definition}
\label{stacks-morphisms-definition-locally-finite-presentation}
设 $f : \mathcal{X} \to \mathcal{Y}$ 是代数叠的态射。
\begin{enumerate}
\item 若引理 \ref{stacks-morphisms-lemma-local-source-target} 的等价条件对
$\mathcal{P} = \text{locally of finite presentation}$ 成立，则称 $f$
{\it 局部有限表现}。
\item 若该态射局部有限表现、拟紧且拟分离，则称 $f$ {\it 有限表现}。
\end{enumerate}
\end{definition}

\noindent
注意，有限表现态射{\bf 不只是}一个局部有限表现的拟紧态射。

\begin{lemma}
\label{stacks-morphisms-lemma-composition-finite-presentation}
有限表现态射的复合有限表现。对局部有限表现态射也有同样结论。
\end{lemma}

\begin{proof}
合并评注 \ref{stacks-morphisms-remark-composition} 与《空间的态射》引理
\ref{spaces-morphisms-lemma-composition-finite-presentation}。
\end{proof}

\begin{lemma}
\label{stacks-morphisms-lemma-base-change-finite-presentation}
有限表现态射的基变换有限表现。对局部有限表现态射也有同样结论。
\end{lemma}

\begin{proof}
合并评注 \ref{stacks-morphisms-remark-base-change} 与《空间的态射》引理
\ref{spaces-morphisms-lemma-base-change-finite-presentation}。
\end{proof}

\begin{lemma}
\label{stacks-morphisms-lemma-finite-presentation-finite-type}
局部有限表现态射局部有限型。有限表现态射有限型。
\end{lemma}

\begin{proof}
合并评注 \ref{stacks-morphisms-remark-implication} 与《空间的态射》引理
\ref{spaces-morphisms-lemma-finite-presentation-finite-type}。
\end{proof}

\begin{lemma}
\label{stacks-morphisms-lemma-noetherian-finite-type-finite-presentation}
设 $f : \mathcal{X} \to \mathcal{Y}$ 是代数叠的态射。
\begin{enumerate}
\item 若 $\mathcal{Y}$ 局部诺特且 $f$ 局部有限型，则 $f$ 局部有限表现。
\item 若 $\mathcal{Y}$ 局部诺特，且 $f$ 有限型并拟分离，则 $f$ 有限表现。
\end{enumerate}
\end{lemma}

\begin{proof}
假设 $f : \mathcal{X} \to \mathcal{Y}$ 局部有限型，且 $\mathcal{Y}$
局部诺特。这意味着存在引理 \ref{stacks-morphisms-lemma-local-source-target} 中所述的图，
其中 $h$ 局部有限型，竖直箭头 $a$ 满。由《空间的态射》引理
\ref{spaces-morphisms-lemma-noetherian-finite-type-finite-presentation}
可知 $h$ 局部有限表现。因此按定义，
$\mathcal{X} \to \mathcal{Y}$ 局部有限表现。这证明了 (1)。若 $f$
有限型且拟分离，则它也拟紧且拟分离，故 (2) 立即得出。
\end{proof}

\begin{lemma}
\label{stacks-morphisms-lemma-finite-presentation-permanence}
设 $f : \mathcal{X} \to \mathcal{Y}$ 和
$g : \mathcal{Y} \to \mathcal{Z}$ 是代数叠的态射。若 $g \circ f$
局部有限表现，且 $g$ 局部有限型，则 $f$ 局部有限表现。
\end{lemma}

\begin{proof}
选取代数空间 $W$ 以及满光滑态射 $W \to \mathcal{Z}$。选取代数空间 $V$
以及满光滑态射 $V \to \mathcal{Y} \times_\mathcal{Z} W$。选取代数空间
$U$ 以及满光滑态射 $U \to \mathcal{X} \times_\mathcal{Y} V$。
对态射 $U \to V \to W$ 应用《空间的态射》引理
\ref{spaces-morphisms-lemma-finite-presentation-permanence}
即得引理。
\end{proof}

\begin{lemma}
\label{stacks-morphisms-lemma-diagonal-morphism-finite-type}
设 $f : \mathcal{X} \to \mathcal{Y}$ 是代数叠的态射，其对角态射为
$\Delta : \mathcal{X} \to \mathcal{X} \times_\mathcal{Y} \mathcal{X}$。
若 $f$ 局部有限型，则 $\Delta$ 局部有限表现。若 $f$ 拟分离且局部有限型，
则 $\Delta$ 有限表现。
\end{lemma}

\begin{proof}
注意，$\Delta$ 是 $\mathcal{X}$ 上的态射（经第二投影）。若 $f$ 局部有限型，
则 $\mathcal{X}$ 在 $\mathcal{X}$ 上有限表现，且由引理
\ref{stacks-morphisms-lemma-base-change-finite-type}，
$\text{pr}_2 : \mathcal{X} \times_\mathcal{Y} \mathcal{X} \to \mathcal{X}$
局部有限型。因此，第一个陈述由引理
\ref{stacks-morphisms-lemma-finite-presentation-permanence} 得出。第二个陈述由第一个陈述
和定义得出（因为按定义，$f$ 拟分离意味着 $\Delta_f$ 拟紧且拟分离）。
\end{proof}

\begin{lemma}
\label{stacks-morphisms-lemma-open-immersion-locally-finite-presentation}
开浸入局部有限表现。
\end{lemma}

\begin{proof}
由《代数叠的性质》定义
\ref{stacks-properties-definition-immersion}
以及《空间的态射》引理
\ref{spaces-morphisms-lemma-open-immersion-locally-finite-presentation}
即得。
\end{proof}

\begin{lemma}
\label{stacks-morphisms-lemma-check-property-after-fppf-base-change}
设 $P$ 是代数空间态射的一个性质，它在靶上为 fppf 局部的，并在任意基变换
下保持。设 $f : \mathcal{X} \to \mathcal{Y}$ 是可由代数空间表示的代数叠态射。
设 $\mathcal{Z} \to \mathcal{Y}$ 是满、平坦且局部有限表现的代数叠态射。
令 $\mathcal{W} = \mathcal{Z} \times_\mathcal{Y} \mathcal{X}$。则
$$
(f\text{ 具有 }P) \Leftrightarrow
(\text{the projection }\mathcal{W} \to \mathcal{Z}\text{ 具有 }P).
$$
该陈述的含义见《代数叠的性质》一节
\ref{stacks-properties-section-properties-morphisms}.
\end{lemma}

\begin{proof}
选取代数空间 $W$ 以及满、平坦且局部有限表现的态射
$W \to \mathcal{Z}$。由《代数叠的性质》引理
\ref{stacks-properties-lemma-composition-surjective}
以及引理 \ref{stacks-morphisms-lemma-composition-flat} 和
\ref{stacks-morphisms-lemma-composition-finite-presentation}
可知，复合 $W \to \mathcal{Y}$ 也满、平坦且局部有限表现。记
$V = W \times_\mathcal{Z} \mathcal{W} = V \times_\mathcal{Y} \mathcal{X}$.
由《代数叠的性质》引理
\ref{stacks-properties-lemma-check-property-covering}
可知，$f$ 具有 $\mathcal{P}$ 当且仅当 $V \to W$ 具有该性质，
而 $\mathcal{W} \to \mathcal{Z}$ 具有 $\mathcal{P}$ 当且仅当
$V \to W$ 具有该性质。引理得证。
\end{proof}

\begin{lemma}
\label{stacks-morphisms-lemma-descent-property}
设 $\mathcal{P}$ 是代数空间态射的一个性质，它在源与靶上光滑局部，
并且在靶上为 fppf 局部。设 $f : \mathcal{X} \to \mathcal{Y}$ 是代数叠的态射。
设 $\mathcal{Z} \to \mathcal{Y}$ 是代数叠的满、平坦且局部有限表现的态射。
若基变换 $\mathcal{Z} \times_\mathcal{Y} \mathcal{X} \to \mathcal{Z}$
具有 $\mathcal{P}$，则 $f$ 具有 $\mathcal{P}$。
\end{lemma}

\begin{proof}
假设 $\mathcal{Z} \times_\mathcal{Y} \mathcal{X} \to \mathcal{Z}$
具有 $\mathcal{P}$。选取代数空间 $W$ 以及满光滑态射
$W \to \mathcal{Z}$。注意
$W \times_\mathcal{Z} \mathcal{Z} \times_\mathcal{Y} \mathcal{X} =
W \times_\mathcal{Y} \mathcal{X}$。根据
$\mathcal{Z} \times_\mathcal{Y} \mathcal{X} \to \mathcal{Z}$ 具有
$\mathcal{P}$ 的定义本身（参见定义 \ref{stacks-morphisms-definition-P} 和引理
\ref{stacks-morphisms-lemma-local-source-target}），可知
$W \times_\mathcal{Y} \mathcal{X} \to W$ 具有 $\mathcal{P}$。另一方面，
$W \to \mathcal{Z}$ 满、平坦且局部有限表现（《空间的态射》引理
\ref{spaces-morphisms-lemma-smooth-flat} 和
\ref{spaces-morphisms-lemma-smooth-locally-finite-presentation}）
），故 $W \to \mathcal{Y}$ 满、平坦且局部有限表现（由《代数叠的性质》引理
\ref{stacks-properties-lemma-composition-surjective}
以及引理 \ref{stacks-morphisms-lemma-composition-flat} 和
\ref{stacks-morphisms-lemma-composition-finite-presentation}）。
因此可用 $W$ 替换 $\mathcal{Z}$。

\medskip\noindent
选取代数空间 $V$ 以及满光滑态射 $V \to \mathcal{Y}$。选取代数空间 $U$
以及满光滑态射 $U \to V \times_\mathcal{Y} \mathcal{X}$。需要证明
$U \to V$ 具有 $\mathcal{P}$。现在沿 $W \to \mathcal{Y}$ 对所有对象
作基变换：令
$U' = W \times_\mathcal{Y} U$,
$V' = W \times_\mathcal{Y} V$,
$\mathcal{X}' = W \times_\mathcal{Y} \mathcal{X}$,
以及 $\mathcal{Y}' = W \times_\mathcal{Y} \mathcal{Y} = W$。由基变换，
$U' \to V' \times_{\mathcal{Y}'} \mathcal{X}'$ 仍然光滑。因此，在定义
$\mathcal{X}' \to \mathcal{Y}' = W$ 具有 $\mathcal{P}$ 时应用引理
\ref{stacks-morphisms-lemma-local-source-target}，可知 $U' \to V'$ 具有 $\mathcal{P}$。
又因 $V' \to V$ 作为 $W \to \mathcal{Y}$ 的基变换满、平坦且局部有限表现，
而 $\mathcal{P}$ 在靶上对 fppf 拓扑是局部的，故
$U \to V$ 具有 $\mathcal{P}$。
\end{proof}

\begin{lemma}
\label{stacks-morphisms-lemma-descent-finite-presentation}
设 $f : \mathcal{X} \to \mathcal{Y}$ 是代数叠的态射。设
$\mathcal{Z} \to \mathcal{Y}$ 是代数叠的满、平坦且局部有限表现的态射。
若基变换 $\mathcal{Z} \times_\mathcal{Y} \mathcal{X} \to \mathcal{Z}$
局部有限表现，则 $f$ 局部有限表现。
\end{lemma}

\begin{proof}
“局部有限表现”性质满足引理 \ref{stacks-morphisms-lemma-descent-property} 的条件。
在源与靶上光滑局部这一点已见于本节引言，而在靶上为 fppf 局部这一点即
《空间上的下降》引理
\ref{spaces-descent-lemma-descending-property-locally-finite-presentation}.
\end{proof}

\begin{lemma}
\label{stacks-morphisms-lemma-flat-finite-presentation-permanence}
设 $\mathcal{X} \to \mathcal{Y} \to \mathcal{Z}$ 是代数叠的态射。
若 $\mathcal{X} \to \mathcal{Z}$ 局部有限表现，且
$\mathcal{X} \to \mathcal{Y}$ 满、平坦且局部有限表现，则
$\mathcal{Y} \to \mathcal{Z}$ 局部有限表现。
\end{lemma}

\begin{proof}
选取代数空间 $W$ 以及满光滑态射 $W \to \mathcal{Z}$。选取代数空间 $V$
以及满光滑态射 $V \to W \times_\mathcal{Z} \mathcal{Y}$。选取代数空间
$U$ 以及满光滑态射 $U \to V \times_\mathcal{Y} \mathcal{X}$。我们知道
$U \to V$ 平坦且局部有限表现，而 $U \to W$ 局部有限表现。此外，
由于 $\mathcal{X} \to \mathcal{Y}$ 满，$U \to V$ 也满（作为满态射的复合）。
因此，引理归结为代数空间态射的情形。该情形即《空间上的下降》引理
\ref{spaces-descent-lemma-locally-finite-presentation-fppf-local-source}.
\end{proof}

\begin{lemma}
\label{stacks-morphisms-lemma-surjective-flat-locally-finite-presentation}
设 $f : \mathcal{X} \to \mathcal{Y}$ 是满、平坦且局部有限表现的代数叠态射。
则对每个概形 $U$ 以及 $\mathcal{Y}$ 在 $U$ 上的对象 $y$，存在 fppf 覆盖
$\{U_i \to U\}$ 以及 $\mathcal{X}$ 在 $U_i$ 上的对象 $x_i$，使在
$\mathcal{Y}_{U_i}$ 中有 $f(x_i) \cong y|_{U_i}$。
\end{lemma}

\begin{proof}
可把 $y$ 看作态射 $U \to \mathcal{Y}$。由《代数叠的性质》引理
\ref{stacks-properties-lemma-base-change-surjective}
以及引理 \ref{stacks-morphisms-lemma-base-change-finite-presentation} 和
\ref{stacks-morphisms-lemma-base-change-flat}
可知 $\mathcal{X} \times_\mathcal{Y} U \to U$ 满、平坦且局部有限表现。
设 $V$ 为概形，并设 $V \to \mathcal{X} \times_\mathcal{Y} U$ 光滑且满。
则 $V \to \mathcal{X} \times_\mathcal{Y} U$ 也满、平坦且局部有限表现
（参见《空间的态射》引理
\ref{spaces-morphisms-lemma-smooth-flat} 和
\ref{spaces-morphisms-lemma-smooth-locally-finite-presentation}）。
因此 $V \to U$ 也满、平坦且局部有限表现；参见《代数叠的性质》引理
\ref{stacks-properties-lemma-composition-surjective}
以及引理 \ref{stacks-morphisms-lemma-composition-finite-presentation} 和
\ref{stacks-morphisms-lemma-composition-flat}.
因此，$\{V \to U\}$ 是所需的 fppf 覆盖，而
$x : V \to \mathcal{X}$ 是所需对象。
\end{proof}

\begin{lemma}
\label{stacks-morphisms-lemma-surjective-family-flat-locally-finite-presentation}
设 $f_j : \mathcal{X}_j \to \mathcal{X}$（$j \in J$）是一族代数叠态射，
每个都平坦且局部有限表现，并且它们联合为满，即
$|\mathcal{X}| = \bigcup |f_j|(|\mathcal{X}_j|)$.
则对每个概形 $U$ 以及 $\mathcal{X}$ 在 $U$ 上的对象 $x$，存在 fppf 覆盖
$\{U_i \to U\}_{i \in I}$、映射 $a : I \to J$，以及
$\mathcal{X}_{a(i)}$ 在 $U_i$ 上的对象 $x_i$，使在
$\mathcal{X}_{U_i}$ 中有 $f_{a(i)}(x_i) \cong y|_{U_i}$。
\end{lemma}

\begin{proof}
对态射 $\coprod_{j \in J} \mathcal{X}_j \to \mathcal{X}$ 应用引理
\ref{stacks-morphisms-lemma-surjective-flat-locally-finite-presentation}。（由于我们的设置，
这里有一个轻微的集合论问题；我们将其忽略。）最后注意，态射
$x_i : U_i \to \coprod_{j \in J} \mathcal{X}_j$ 由不交并分解
$U_i = \coprod U_{i, j}$ 和态射 $U_{i, j} \to \mathcal{X}_j$ 给出。
于是 fppf 覆盖 $\{U_{i, j} \to U\}$ 以及态射
$U_{i, j} \to \mathcal{X}_j$ 满足要求。
\end{proof}

\begin{lemma}
\label{stacks-morphisms-lemma-fppf-open}
设 $f : \mathcal{X} \to \mathcal{Y}$ 平坦且局部有限表现。则
$|f| : |\mathcal{X}| \to |\mathcal{Y}|$ 是开映射。
\end{lemma}

\begin{proof}
选取概形 $V$ 以及满光滑态射 $V \to \mathcal{Y}$。选取概形 $U$ 以及
满光滑态射 $U \to V \times_\mathcal{Y} \mathcal{X}$。按假设，概形态射
$U \to V$ 平坦且局部有限表现。因此，由《态射》引理
\ref{morphisms-lemma-fppf-open}，$U \to V$ 是开映射。根据
$|\mathcal{Y}|$ 上拓扑的构造，映射 $|V| \to |\mathcal{Y}|$ 是开的。
映射 $|U| \to |\mathcal{X}|$ 满。由初等拓扑，这些事实蕴含结论。
\end{proof}

\begin{lemma}
\label{stacks-morphisms-lemma-descent-quasi-compact}
设 $f : \mathcal{X} \to \mathcal{Y}$ 是代数叠的态射。设
$\mathcal{Z} \to \mathcal{Y}$ 是代数叠的满、平坦且局部有限表现的态射。
若基变换 $\mathcal{Z} \times_\mathcal{Y} \mathcal{X} \to \mathcal{Z}$
拟紧，则 $f$ 拟紧。
\end{lemma}

\begin{proof}
需要证明：对任意态射 $\mathcal{Y}' \to \mathcal{Y}$，若
$\mathcal{Y}'$ 拟紧，则 $\mathcal{Y}' \times_\mathcal{Y} \mathcal{X}$ 拟紧。
记 $\mathcal{Z}' = \mathcal{Z} \times_\mathcal{Y} \mathcal{Y}'$。
由引理 \ref{stacks-morphisms-lemma-fppf-open}，
$|\mathcal{Z}'| \to |\mathcal{Y}'|$ 是开映射。因此可找到映满
$\mathcal{Y}'$ 的拟紧开子叠 $\mathcal{W} \subset \mathcal{Z}'$。由于
$\mathcal{Z} \times_\mathcal{Y} \mathcal{X} \to \mathcal{Z}$
拟紧，可知
$$
\mathcal{W} \times_\mathcal{Z} \mathcal{Z} \times_\mathcal{Y} \mathcal{X} =
\mathcal{W} \times_\mathcal{Y} \mathcal{X}
$$
拟紧。而映射
$\mathcal{W} \times_\mathcal{Y} \mathcal{X} \to
\mathcal{Y}' \times_\mathcal{Y} \mathcal{X}$
满，故得证。部分细节略。
\end{proof}

\begin{lemma}
\label{stacks-morphisms-lemma-check-separated-on-ui-cover}
设 $f : \mathcal{X} \to \mathcal{Y}$、$g : \mathcal{Y} \to \mathcal{Z}$
是可复合的代数叠态射，其复合为
$h = g \circ f : \mathcal{X} \to \mathcal{Z}$。若 $f$ 满、平坦、
局部有限表现且泛单射，并且 $h$ 分离，则 $g$ 分离。
\end{lemma}

\begin{proof}
考虑图
$$
\xymatrix{
\mathcal{X} \ar[r]_\Delta \ar[rd] &
\mathcal{X} \times_\mathcal{Y} \mathcal{X} \ar[r] \ar[d] &
\mathcal{X} \times_\mathcal{Z} \mathcal{X} \ar[d] \\
& \mathcal{Y} \ar[r] & \mathcal{Y} \times_\mathcal{Z} \mathcal{Y}
}
$$
该方块是笛卡儿的。需要证明底部水平箭头为真态射。我们已经知道它可由
代数空间表示且局部有限型（引理 \ref{stacks-morphisms-lemma-properties-diagonal}）。
由于右侧竖直箭头满、平坦且局部有限表现，只需证明右上方水平箭头为真态射
（引理 \ref{stacks-morphisms-lemma-check-property-after-fppf-base-change}）。由于 $h$ 分离，
顶部水平箭头的复合为真态射。

\medskip\noindent
由于 $f$ 泛单射，$\Delta$ 满（引理
\ref{stacks-morphisms-lemma-universally-injective}）。由于 $\Delta$ 与投影
$\mathcal{X} \times_\mathcal{Y} \mathcal{X} \to \mathcal{X}$ 的复合为恒等，
可知 $\Delta$ 泛闭。由《空间的态射》引理
\ref{spaces-morphisms-lemma-image-universally-closed-separated}
可知 $\mathcal{X} \times_\mathcal{Y} \mathcal{X} \to
\mathcal{X} \times_\mathcal{Z} \mathcal{X}$ 分离，因为
$\mathcal{X} \to \mathcal{X} \times_\mathcal{Z} \mathcal{X}$ 分离。
这里使用了如下事实：代数空间态射的性质之间的蕴含可转移为可由代数空间表示的
代数叠态射的相应性质之间的蕴含；这在《代数叠的性质》一节中讨论
\ref{stacks-properties-section-properties-morphisms}.
最后，使用同一原则，并由《空间的态射》引理
\ref{spaces-morphisms-lemma-image-proper-is-proper}
可知 $\mathcal{X} \times_\mathcal{Y} \mathcal{X} \to
\mathcal{X} \times_\mathcal{Z} \mathcal{X}$ 为真态射。
\end{proof}














\section{胚}
\label{stacks-morphisms-section-gerbes}

\noindent
一类重要的代数叠是形如 $[B/G]$ 的叠，其中 $B$ 是代数空间，而 $G$ 是
$B$ 上平坦且局部有限表现的群代数空间（在 $B$ 上平凡作用）；参见
《可表示性判据》引理 \ref{criteria-lemma-BG-algebraic}。事实证明，若一个
代数叠在 fppf 拓扑下局部地具有这种形式，则它是一个胚；参见引理
\ref{stacks-morphisms-lemma-local-structure-gerbe}。本节简要讨论这个概念及其相应的相对概念。

\begin{definition}
\label{stacks-morphisms-definition-gerbe}
设 $f : \mathcal{X} \to \mathcal{Y}$ 是代数叠的态射。若把源与靶
视为 $(\Sch/S)_{fppf}$ 上的群胚叠时，
$\mathcal{X}$ 是 $\mathcal{Y}$ 上的胚，则称 $\mathcal{X}$ 是
$\mathcal{Y}$ {\it 上的胚}；参见《叠》定义
\ref{stacks-definition-gerbe-over-stack-in-groupoids}。若存在态射
$\mathcal{X} \to X$，其中 $X$ 是代数空间，并且该态射使
$\mathcal{X}$ 成为 $X$ 上的胚，则称代数叠 $\mathcal{X}$ 是一个{\it 胚}。
\end{definition}

\noindent
$\mathcal{X}$ 是 $\mathcal{Y}$ 上的胚这一条件，纯粹用给定代数叠背后的
拓扑和范畴论来定义；但稍后将看到，该条件具有几何后果。例如，它蕴含
$\mathcal{X} \to \mathcal{Y}$ 满、平坦且局部有限表现；参见引理
\ref{stacks-morphisms-lemma-gerbe-fppf}。绝对概念更难解析，因为起初可能不清楚 $X$
是否被良好确定。事实上，它确实被良好确定。

\begin{lemma}
\label{stacks-morphisms-lemma-gerbe-over-iso-classes}
设 $\mathcal{X}$ 是代数叠。若 $\mathcal{X}$ 是胚，则预层
$$
(\Sch/S)_{fppf}^{opp} \to \textit{Sets}, \quad
U \mapsto \Ob(\mathcal{X}_U)/\!\!\cong
$$
的层化是代数空间，并且 $\mathcal{X}$ 是该空间上的胚。
\end{lemma}

\begin{proof}
（在这个证明中，第 \ref{stacks-morphisms-section-conventions} 节引入的术语滥用确实很有用。）
选取态射 $\pi : \mathcal{X} \to X$，其中 $X$ 是代数空间，并且该态射使
$\mathcal{X}$ 成为 $X$ 上的胚。只需证明 $X$ 是引理中所示预层
$\mathcal{F}$ 的层化。显然有映射 $c : \mathcal{F} \to X$。我们将使用
《叠》引理 \ref{stacks-lemma-when-gerbe} 的性质 (2)(a) 和 (2)(b)，
说明映射 $c^\# : \mathcal{F}^\# \to X$ 既满又单，因而是同构；参见
《站点》引理 \ref{sites-lemma-mono-epi-sheaves}。

满射性：设 $T$ 是概形，并设 $f : T \to X$。由性质 (2)(a)，存在 fppf 覆盖
$\{h_i : T_i \to T\}$ 和态射 $x_i : T_i \to \mathcal{X}$，使
$f \circ h_i$ 对应于 $\pi \circ x_i$。因此，$f|_{T_i}$ 属于 $c$ 的像。

单射性：设 $T$ 是概形，并设 $x, y : T \to \mathcal{X}$ 为态射，且
$c \circ x = c \circ y$。由 (2)(b)，可找到覆盖 $\{T_i \to T\}$ 以及
纤维范畴 $\mathcal{X}_{T_i}$ 中的态射
$x|_{T_i} \to y|_{T_i}$。因此，限制 $x|_{T_i}, y|_{T_i}$ 在
$\mathcal{F}(T_i)$ 中相等。这证明了 $x, y$ 按所需给出
$\mathcal{F}^\#$ 在 $T$ 上的同一截面。
\end{proof}

\begin{lemma}
\label{stacks-morphisms-lemma-base-change-gerbe}
设
$$
\xymatrix{
\mathcal{X}' \ar[r] \ar[d] & \mathcal{X} \ar[d] \\
\mathcal{Y}' \ar[r] & \mathcal{Y}
}
$$
为代数叠的纤维积。若 $\mathcal{X}$ 是 $\mathcal{Y}$ 上的胚，则
$\mathcal{X}'$ 是 $\mathcal{Y}'$ 上的胚。
\end{lemma}

\begin{proof}
立即由定义和《叠》引理 \ref{stacks-lemma-base-change-gerbe} 得出。
\end{proof}

\begin{lemma}
\label{stacks-morphisms-lemma-composition-gerbe}
设 $\mathcal{X} \to \mathcal{Y}$ 和 $\mathcal{Y} \to \mathcal{Z}$
是代数叠的态射。若 $\mathcal{X}$ 是 $\mathcal{Y}$ 上的胚，且
$\mathcal{Y}$ 是 $\mathcal{Z}$ 上的胚，则
$\mathcal{X}$ 是 $\mathcal{Z}$ 上的胚。
\end{lemma}

\begin{proof}
立即由《叠》引理 \ref{stacks-lemma-composition-gerbe} 得出。
\end{proof}

\begin{lemma}
\label{stacks-morphisms-lemma-gerbe-descent}
设
$$
\xymatrix{
\mathcal{X}' \ar[r] \ar[d] & \mathcal{X} \ar[d] \\
\mathcal{Y}' \ar[r] & \mathcal{Y}
}
$$
为代数叠的纤维积。若 $\mathcal{Y}' \to \mathcal{Y}$ 满、平坦且局部有限表现，
并且 $\mathcal{X}'$ 是 $\mathcal{Y}'$ 上的胚，则
$\mathcal{X}$ 是 $\mathcal{Y}$ 上的胚。
\end{lemma}

\begin{proof}
立即由引理 \ref{stacks-morphisms-lemma-surjective-flat-locally-finite-presentation} 和
《叠》引理 \ref{stacks-lemma-gerbe-descent} 得出。
\end{proof}

\begin{lemma}
\label{stacks-morphisms-lemma-gerbe-with-section}
设 $\pi : \mathcal{X} \to U$ 是从代数叠到代数空间的态射，并设
$x : U \to \mathcal{X}$ 是 $\pi$ 的截面。令
$G = \mathit{Isom}_\mathcal{X}(x, x)$；参见定义
\ref{stacks-morphisms-definition-isom}。若 $\mathcal{X}$ 是 $U$ 上的胚，则
\begin{enumerate}
\item 有群胚叠的典范等价
$$
x_{can} : [U/G] \longrightarrow \mathcal{X}.
$$
其中 $[U/G]$ 是 $G$ 在 $U$ 上平凡作用所对应的商叠；
\item $G \to U$ 平坦且局部有限表现；
\item $U \to \mathcal{X}$ 满、平坦且局部有限表现。
\end{enumerate}
\end{lemma}

\begin{proof}
令 $R = U \times_{x, \mathcal{X}, x} U$。态射 $R \to U \times U$
经过对角态射 $\Delta_U : U \to U \times U$ 分解，因为它经过
$U \times_U U = U$ 分解。因此 $R = G$，因为
\begin{align*}
G & = \mathit{Isom}_\mathcal{X}(x, x) \\
& = U \times_{x, \mathcal{X}} \mathcal{I}_\mathcal{X} \\
& = U \times_{x, \mathcal{X}}
(\mathcal{X}
\times_{\Delta, \mathcal{X} \times_S \mathcal{X}, \Delta}
\mathcal{X}) \\
& = (U \times_{x, \mathcal{X}, x} U) \times_{U \times U, \Delta_U} U \\
& = R \times_{U \times U, \Delta_U} U \\
& = R
\end{align*}
第四个等号使用《范畴》引理 \ref{categories-lemma-diagonal-2}。
令 $t, s : R \to U$ 为投影。《代数叠》引理
\ref{algebraic-lemma-map-space-into-stack} 在 $R$ 上构造的复合法则
$c : R \times_{s, U, t} R \to R$ 与 $G$ 上的群法则一致（证明略）。
因此，《代数叠》引理 \ref{algebraic-lemma-map-space-into-stack}
说明我们得到典范的全忠实 $1$-态射
$$
x_{can} : [U/G] \longrightarrow \mathcal{X}
$$
，它是 $(\Sch/S)_{fppf}$ 上的群胚叠之间的态射。为说明它是等价，只需证明
它本质满。为此，只需证明 $\mathcal{X}$ 在任一概形 $T$ 上的对象，在
fppf 局部下经某个态射 $T \to U$ 来自 $x$；参见《叠》引理
\ref{stacks-lemma-characterize-essentially-surjective-when-ff}。
而这由 $\pi$ 使 $\mathcal{X}$ 成为 $U$ 上的胚这一条件得出；参见
《叠》引理 \ref{stacks-lemma-when-gerbe} 的性质 (2)(a)。

\medskip\noindent
由《可表示性判据》引理 \ref{criteria-lemma-BG-algebraic} 可知
$G \to U$ 平坦且局部有限表现。最后，由《可表示性判据》引理
\ref{criteria-lemma-flat-quotient-flat-presentation}
可知 $U \to \mathcal{X}$ 满、平坦且局部有限表现。
\end{proof}

\begin{lemma}
\label{stacks-morphisms-lemma-local-structure-gerbe}
设 $\pi : \mathcal{X} \to \mathcal{Y}$ 是代数叠的态射。下列条件等价：
\begin{enumerate}
\item $\mathcal{X}$ 是 $\mathcal{Y}$ 上的胚；
\item 存在代数空间 $U$、在 $U$ 上平坦且局部有限表现的群代数空间 $G$，
以及满、平坦且局部有限表现的态射 $U \to \mathcal{Y}$，使在 $U$ 上有
$\mathcal{X} \times_\mathcal{Y} U \cong [U/G]$。
\end{enumerate}
\end{lemma}

\begin{proof}
假设 (2)。由引理 \ref{stacks-morphisms-lemma-gerbe-descent}，为证明 (1)，只需说明
$[U/G]$ 是 $U$ 上的胚。这立即由《空间中的群胚》引理
\ref{spaces-groupoids-lemma-group-quotient-gerbe} 得出。

\medskip\noindent
假设 (1)。$\pi$ 的任意基变换都是胚；参见引理
\ref{stacks-morphisms-lemma-base-change-gerbe}。第一步，选取概形 $V$ 以及满光滑态射
$V \to \mathcal{Y}$。于是可假设 $\pi : \mathcal{X} \to V$ 是概形上的胚。
这意味着存在 fppf 覆盖 $\{V_i \to V\}$，使纤维范畴
$\mathcal{X}_{V_i}$ 非空；参见《叠》引理
\ref{stacks-lemma-when-gerbe} (2)(a)。注意，
$U = \coprod V_i \to V$ 满、平坦且局部有限表现。因此，可用 $U$ 替换 $V$，
并假设 $\pi : \mathcal{X} \to U$ 是概形 $U$ 上的胚，且存在
$\mathcal{X}$ 在 $U$ 上的对象 $x$。由引理
\ref{stacks-morphisms-lemma-gerbe-with-section} 可知，在 $U$ 上有
$\mathcal{X} = [U/G]$，其中 $G$ 是 $U$ 上平坦且局部有限表现的群代数空间。
\end{proof}

\begin{lemma}
\label{stacks-morphisms-lemma-gerbe-fppf}
设 $\pi : \mathcal{X} \to \mathcal{Y}$ 是代数叠的态射。若
$\mathcal{X}$ 是 $\mathcal{Y}$ 上的胚，则 $\pi$ 满、平坦且局部有限表现。
\end{lemma}

\begin{proof}
由《代数叠的性质》引理
\ref{stacks-properties-lemma-descent-surjective}
以及引理 \ref{stacks-morphisms-lemma-descent-flat} 和
\ref{stacks-morphisms-lemma-descent-finite-presentation}
可知，只需在用满、平坦且局部有限表现的态射
$\mathcal{Y}' \to \mathcal{Y}$ 对 $\pi$ 作基变换后证明引理。由引理
\ref{stacks-morphisms-lemma-local-structure-gerbe}，可假设 $\mathcal{Y} = U$ 是代数空间，
并且在 $U$ 上有 $\mathcal{X} = [U/G]$。此时
$U \to [U/G]$ 满、平坦且局部有限表现；参见引理
\ref{stacks-morphisms-lemma-gerbe-with-section}。于是由《代数叠的性质》引理
\ref{stacks-properties-lemma-surjective-permanence} 以及引理
\ref{stacks-morphisms-lemma-flat-permanence} 和
\ref{stacks-morphisms-lemma-flat-finite-presentation-permanence}
可知 $\pi$ 满、平坦且局部有限表现。
\end{proof}

\begin{proposition}
\label{stacks-morphisms-proposition-when-gerbe}
设 $\mathcal{X}$ 是代数叠。下列条件等价：
\begin{enumerate}
\item $\mathcal{X}$ 是胚；
\item $\mathcal{I}_\mathcal{X} \to \mathcal{X}$ 平坦且局部有限表现。
\end{enumerate}
\end{proposition}

\begin{proof}
假设 (1)。选取到代数空间 $X$ 的态射 $\mathcal{X} \to X$，使
$\mathcal{X}$ 成为 $X$ 上的胚。设 $X' \to X$ 是满、平坦且局部有限表现的态射，
并令 $\mathcal{X}' = X' \times_X \mathcal{X}$。注意，由引理
\ref{stacks-morphisms-lemma-base-change-gerbe}，$\mathcal{X}'$ 是 $X'$ 上的胚。下图中的两个
方块都是纤维积方块：
$$
\xymatrix{
\mathcal{I}_{\mathcal{X}'} \ar[r] \ar[d] &
\mathcal{X}' \ar[r] \ar[d] & X' \ar[d] \\
\mathcal{I}_\mathcal{X} \ar[r] &
\mathcal{X} \ar[r] & X
}
$$
；参见引理 \ref{stacks-morphisms-lemma-cartesian-square-inertia}。因此，由引理
\ref{stacks-morphisms-lemma-descent-flat} 和
\ref{stacks-morphisms-lemma-descent-finite-presentation}，为证明
$\mathcal{I}_\mathcal{X} \to \mathcal{X}$ 平坦且局部有限表现，
只需在这样的基变换后证明。于是可应用引理
\ref{stacks-morphisms-lemma-local-structure-gerbe}，假设 $\mathcal{X} = [U/G]$。由引理
\ref{stacks-morphisms-lemma-gerbe-with-section} 可知，$G$ 在 $U$ 上平坦且局部有限表现，
而 $x : U \to [U/G]$ 满、平坦且局部有限表现。此外，
$\mathcal{I}_\mathcal{X}$ 沿 $x$ 的拉回为 $G$，再次由下降，即由引理
\ref{stacks-morphisms-lemma-descent-flat} 和
\ref{stacks-morphisms-lemma-descent-finite-presentation}，可知 (2) 成立。

\medskip\noindent
反之，假设 (2)。选取光滑表现 $\mathcal{X} = [U/R]$；参见《代数叠》一节
\ref{algebraic-section-stack-to-presentation}。记 $G \to U$ 为群胚
$(U, R, s, t, c, e, i)$ 的稳定子群代数空间；参见《空间中的群胚》定义
\ref{spaces-groupoids-definition-stabilizer-groupoid}.
由引理 \ref{stacks-morphisms-lemma-presentation-inertia} 可知，$G \to U$ 作为
$\mathcal{I}_\mathcal{X} \to \mathcal{X}$ 的基变换平坦且局部有限表现；
参见引理 \ref{stacks-morphisms-lemma-base-change-flat} 和
\ref{stacks-morphisms-lemma-base-change-finite-presentation}.
考虑如下 $G$ 在 $R$ 上的作用：
$$
a : G \times_{U, t} R \to R, \quad (g, r) \mapsto c(g, r)
$$
。由于 $R$ 是群胚，该作用在任意概形 $T$ 的 $T$-值点上自由。因此，
由《自举》引理 \ref{bootstrap-lemma-quotient-free-action}，$R' = R/G$
是代数空间，而商态射 $\pi : R \to R'$ 满、平坦且局部有限表现。
投影 $s, t : R \to U$ 是 $G$-不变的，故得到态射
$s' , t' : R' \to U$，使 $s = s' \circ \pi$ 且
$t = t' \circ \pi$。由于 $s, t : R \to U$ 平坦且局部有限表现，
由《空间的态射》引理
\ref{spaces-morphisms-lemma-flat-permanence} 和《空间上的下降》引理
\ref{spaces-descent-lemma-locally-finite-presentation-fppf-local-source}
可知 $s', t'$ 平坦且局部有限表现。考虑态射
$$
j' = (t', s') : R' \longrightarrow U \times U.
$$
我们断言它是单态射。事实上，设 $T$ 是概形，并设
$a, b : T \to R'$ 是在 $U \times U$ 中有相同像的态射。根据商
$R' = R/G$ 的定义，存在 fppf 覆盖 $\{h_j : T_j \to T\}$，使对某些态射
$a_j, b_j : T_j \to R$ 有 $a \circ h_j = \pi \circ a_j$ 且
$b \circ h_j = \pi \circ b_j$。由于 $a_j, b_j$ 在 $U \times U$ 中有相同像，
$g_j = c(a_j, i(b_j))$ 是 $G$ 的一个 $T_j$-值点，并且
$c(g_j, b_j) = a_j$。换言之，$a_j$ 和 $b_j$ 在 $R'$ 中有相同的像，
故断言得证。由于 $j : R \to U \times U$ 是预等价关系（参见
《空间中的群胚》引理
\ref{spaces-groupoids-lemma-groupoid-pre-equivalence}），且
$R \to R'$ 作为层的映射为满，可知
$j' : R' \to U \times U$ 是等价关系。因此，《自举》定理
\ref{bootstrap-theorem-final-bootstrap} 说明 $X = U/R'$ 是代数空间。
最后，我们断言态射
$$
\mathcal{X} = [U/R] \longrightarrow X = U/R'
$$
使 $\mathcal{X}$ 成为 $X$ 上的胚。这由《空间中的群胚》引理
\ref{spaces-groupoids-lemma-when-gerbe} 得出，因为 $R \to R'$ 满、平坦且
局部有限表现（必要时使用《自举》引理
\ref{bootstrap-lemma-surjective-flat-locally-finite-presentation}
以看出这蕴含所需假设）。
\end{proof}

\begin{lemma}
\label{stacks-morphisms-lemma-gerbe-isom-fppf}
设 $f : \mathcal{X} \to \mathcal{Y}$ 是使 $\mathcal{X}$ 成为
$\mathcal{Y}$ 上之胚的代数叠态射。则
\begin{enumerate}
\item $\mathcal{I}_{\mathcal{X}/\mathcal{Y}} \to \mathcal{X}$
平坦且局部有限表现；
\item $\mathcal{X} \to \mathcal{X} \times_\mathcal{Y} \mathcal{X}$
满、平坦且局部有限表现；
\item 给定代数空间 $T_i$（$i = 1, 2$）和态射
$x_i : T_i \to \mathcal{X}$，令 $y_i = f \circ x_i$，则态射
$$
T_1 \times_{x_1, \mathcal{X}, x_2} T_2 \longrightarrow
T_1 \times_{y_1, \mathcal{Y}, y_2} T_2
$$
满、平坦且局部有限表现；
\item 给定代数空间 $T$ 和态射 $x_i : T \to \mathcal{X}$（$i = 1, 2$），
令 $y_i = f \circ x_i$，则态射
$$
\mathit{Isom}_\mathcal{X}(x_1, x_2) \longrightarrow
\mathit{Isom}_\mathcal{Y}(y_1, y_2)
$$
满、平坦且局部有限表现。
\end{enumerate}
\end{lemma}

\begin{proof}
(1) 的证明。选取概形 $Y$ 以及满光滑态射 $Y \to \mathcal{Y}$。令
$\mathcal{X}' = \mathcal{X} \times_\mathcal{Y} Y$。由引理
\ref{stacks-morphisms-lemma-cartesian-square-inertia}，得到笛卡儿方块
$$
\xymatrix{
\mathcal{I}_{\mathcal{X}'} \ar[r] \ar[d] &
\mathcal{X}' \ar[r] \ar[d] & Y \ar[d] \\
\mathcal{I}_{\mathcal{X}/\mathcal{Y}} \ar[r] &
\mathcal{X} \ar[r] & \mathcal{Y}
}
$$
由引理 \ref{stacks-morphisms-lemma-descent-flat} 和
\ref{stacks-morphisms-lemma-descent-finite-presentation}
可知，只需证明 $\mathcal{I}_{\mathcal{X}'} \to \mathcal{X}'$
平坦且局部有限表现。这由命题 \ref{stacks-morphisms-proposition-when-gerbe} 得出
（因为由引理 \ref{stacks-morphisms-lemma-base-change-gerbe}，$\mathcal{X}'$ 是 $Y$ 上的胚）。

\medskip\noindent
(2) 的证明。沿用上述记号，注意可假设
$\mathcal{X}' = [Y/G]$，其中 $G$ 是 $Y$ 上平坦且局部有限表现的
群代数空间；参见引理 \ref{stacks-morphisms-lemma-local-structure-gerbe}。态射
$\Delta : \mathcal{X} \to \mathcal{X} \times_\mathcal{Y} \mathcal{X}$
在 $\mathcal{Y}$ 上沿态射 $Y \to \mathcal{Y}$ 的基变换是态射
$\Delta' : \mathcal{X}' \to \mathcal{X}' \times_Y \mathcal{X}'$.
因此，只需证明 $\Delta'$ 满、平坦且局部有限表现（参见引理
\ref{stacks-morphisms-lemma-descent-flat} 和
\ref{stacks-morphisms-lemma-descent-finite-presentation}).
换言之，需要证明
$$
[Y/G] \longrightarrow [Y/G \times_Y G]
$$
满、平坦且局部有限表现。这是因为沿满、平坦且局部有限表现的态射
$Y \to [Y/G \times_Y G]$ 作基变换后得到态射 $G \to Y$。

\medskip\noindent
(3) 的证明。注意图
$$
\xymatrix{
T_1 \times_{x_1, \mathcal{X}, x_2} T_2 \ar[d] \ar[r] &
T_1 \times_{y_1, \mathcal{Y}, y_2} T_2 \ar[d] \\
\mathcal{X} \ar[r] & \mathcal{X} \times_\mathcal{Y} \mathcal{X}
}
$$
是笛卡儿的。因此 (3) 由 (2) 得出。

\medskip\noindent
(4) 的证明。由于
$$
\mathit{Isom}_\mathcal{X}(x_1, x_2) =
(T \times_{x_1, \mathcal{X}, x_2} T) \times_{T \times T, \Delta_T} T
$$
，(4) 中的态射是 (3) 中态射的基变换，故结论成立。
\end{proof}

\begin{proposition}
\label{stacks-morphisms-proposition-when-gerbe-over}
设 $f : \mathcal{X} \to \mathcal{Y}$ 是代数叠的态射。下列条件等价：
\begin{enumerate}
\item $\mathcal{X}$ 是 $\mathcal{Y}$ 上的胚；
\item $f : \mathcal{X} \to \mathcal{Y}$ 与
$\Delta : \mathcal{X} \to \mathcal{X} \times_\mathcal{Y} \mathcal{X}$
均满、平坦且局部有限表现。
\end{enumerate}
\end{proposition}

\begin{proof}
(1) $\Rightarrow$ (2) 由引理 \ref{stacks-morphisms-lemma-gerbe-fppf} 和
\ref{stacks-morphisms-lemma-gerbe-isom-fppf} 得出。

\medskip\noindent
假设 (2)。只需对 $f$ 沿满、平坦且局部有限表现的态射
$\mathcal{Y}' \to \mathcal{Y}$ 所作的基变换证明 (1)；参见引理
\ref{stacks-morphisms-lemma-gerbe-descent}（注意，$f$ 的对角态射的基变换就是基变换的
对角态射）。因此可假设 $\mathcal{Y}$ 是概形 $Y$。在此情形，
$\mathcal{I}_\mathcal{X} \to \mathcal{X}$ 是 $\Delta$ 的基变换，故由命题
\ref{stacks-morphisms-proposition-when-gerbe} 可知 $\mathcal{X}$ 是胚。还需证明
$\mathcal{X}$ 是 $Y$ 上的胚。令 $\mathcal{X} \to X$ 为引理
\ref{stacks-morphisms-lemma-gerbe-over-iso-classes} 的态射，它使 $\mathcal{X}$ 成为代数空间
$X$ 上的胚，而后者分类 $\mathcal{X}$ 的对象的同构类。显然，
$f : \mathcal{X} \to Y$ 分解为 $\mathcal{X} \to X \to Y$。由于 $f$
满、平坦且局部有限表现，可知 $X \to Y$ 作为 fppf 层的映射为满
（例如使用引理 \ref{stacks-morphisms-lemma-surjective-flat-locally-finite-presentation}）。
另一方面，$X \to Y$ 也单：对任意概形 $T$ 以及 $X$ 的任意两个
$T$-值点 $x_1, x_2$，若它们映到 $Y$ 的同一点，则首先可在 $T$ 上
fppf 局部地把 $x_1, x_2$ 提升为 $\mathcal{X}$ 在 $T$ 上的对象
$\xi_1, \xi_2$；其次，由
$\Delta : \mathcal{X} \to \mathcal{X} \times_Y \mathcal{X}$
满、平坦且局部有限表现这一假设，可推出 $\xi_1$ 和 $\xi_2$ 在 fppf
局部同构。于是按 $X$ 的构造，$x_1 = x_2$。因此 $X = Y$，证明完成。
\end{proof}

\noindent
至此，我们已经发展出足够的机制，可以证明剩余胚在存在时确实是胚。

\begin{lemma}
\label{stacks-morphisms-lemma-noetherian-singleton-stack-gerbe}
设 $\mathcal{Z}$ 是约化、局部诺特的代数叠，并且 $|\mathcal{Z}|$ 是单点集。
则 $\mathcal{Z}$ 是某个约化、局部诺特的代数空间 $Z$ 上的胚，且
$|Z|$ 是单点集。
\end{lemma}

\begin{proof}
由《代数叠的性质》引理
\ref{stacks-properties-lemma-unique-point-better}，存在满、平坦且局部有限表现的
态射 $\Spec(k) \to \mathcal{Z}$，其中 $k$ 为域。于是
$\mathcal{I}_Z \times_\mathcal{Z} \Spec(k) \to \Spec(k)$ 可由代数空间表示，
并且局部有限型（作为 $\mathcal{I}_\mathcal{Z} \to \mathcal{Z}$ 的基变换；
参见引理 \ref{stacks-morphisms-lemma-inertia} 和 \ref{stacks-morphisms-lemma-base-change-finite-type}）。
因此，它局部有限表现；参见《空间的态射》引理
\ref{spaces-morphisms-lemma-noetherian-finite-type-finite-presentation}.
当然，因为 $k$ 是域，它也平坦。因此可以应用引理
\ref{stacks-morphisms-lemma-descent-flat} 和
\ref{stacks-morphisms-lemma-descent-finite-presentation}
，可知 $\mathcal{I}_\mathcal{Z} \to \mathcal{Z}$ 平坦且局部有限表现。
由命题 \ref{stacks-morphisms-proposition-when-gerbe}，$\mathcal{Z}$ 是胚。设
$\pi : \mathcal{Z} \to Z$ 为到代数空间的态射，使 $\mathcal{Z}$ 成为
$Z$ 上的胚。由引理 \ref{stacks-morphisms-lemma-gerbe-fppf}，$\pi$ 满、平坦且局部有限表现。
因此，作为复合，$\Spec(k) \to Z$ 满、平坦且局部有限表现；参见
《代数叠的性质》引理
\ref{stacks-properties-lemma-composition-surjective} 以及引理
\ref{stacks-morphisms-lemma-composition-flat} 和
\ref{stacks-morphisms-lemma-composition-finite-presentation}.
于是由《代数叠的性质》引理
\ref{stacks-properties-lemma-unique-point-better} 可知，$|Z|$ 是单点集，
且 $Z$ 局部诺特并约化。
\end{proof}

\begin{lemma}
\label{stacks-morphisms-lemma-gerbe-bijection-points}
设 $f : \mathcal{X} \to \mathcal{Y}$ 是代数叠的态射。若
$\mathcal{X}$ 是 $\mathcal{Y}$ 上的胚，则 $f$ 是泛同胚。
\end{lemma}

\begin{proof}
由引理 \ref{stacks-morphisms-lemma-base-change-gerbe}，关于 $f$ 的假设在基变换下保持。
因此，只需说明映射 $|\mathcal{X}| \to |\mathcal{Y}|$ 是拓扑空间的同胚。
设 $k$ 为域，并设 $y$ 是 $\mathcal{Y}$ 在 $\Spec(k)$ 上的对象。由
《叠》引理 \ref{stacks-lemma-when-gerbe} 的性质 (2)(a)，存在 fppf 覆盖
$\{T_i \to \Spec(k)\}$ 以及 $\mathcal{X}$ 在 $T_i$ 上的对象 $x_i$，
使 $f(x_i) \cong y|_{T_i}$。选取满足 $T_i \not = \emptyset$ 的 $i$。
对某个域 $K$，选取态射 $\Spec(K) \to T_i$。则 $k \subset K$，而
$x_i|_K$ 是位于 $y|_K$ 上方的 $\mathcal{X}$ 的对象。因此可知
$|\mathcal{Y}| \to |\mathcal{X}|$ 满。映射
$|\mathcal{Y}| \to |\mathcal{X}|$ 也单。事实上，若 $x, x'$ 是
$\mathcal{X}$ 在 $\Spec(k)$ 上的对象，并且它们的像 $f(x), f(x')$
在 $\mathcal{Y}$ 中经扩张后同构，则《叠》引理
\ref{stacks-lemma-when-gerbe} 的性质 (2)(b) 保证存在 $k$ 的一个扩张，
使 $x$ 和 $x'$ 同构（细节略）。因此
$|\mathcal{X}| \to |\mathcal{Y}|$ 连续且双射，只需再证明它是开的。
这由引理 \ref{stacks-morphisms-lemma-gerbe-fppf} 和 \ref{stacks-morphisms-lemma-fppf-open} 得出。
\end{proof}

\begin{lemma}
\label{stacks-morphisms-lemma-gerbe-diagonal-quasi-compact}
设 $f : \mathcal{X} \to \mathcal{Y}$ 是代数叠的态射，并且
$\mathcal{X}$ 是 $\mathcal{Y}$ 上的胚。若 $\Delta_\mathcal{X}$ 拟紧，
则 $\Delta_\mathcal{Y}$ 也拟紧。
\end{lemma}

\begin{proof}
考虑图
$$
\xymatrix{
\mathcal{X} \ar[r] &
\mathcal{X} \times_\mathcal{Y} \mathcal{X} \ar[r] \ar[d] &
\mathcal{X} \times \mathcal{X} \ar[d] \\
&
\mathcal{Y} \ar[r] &
\mathcal{Y} \times \mathcal{Y}
}
$$
由命题 \ref{stacks-morphisms-proposition-when-gerbe-over} 可知左上方箭头满。由于顶部水平箭头的
复合拟紧，由引理 \ref{stacks-morphisms-lemma-surjection-from-quasi-compact} 可知右上方箭头拟紧。
该方块是笛卡儿的，而右侧竖直箭头满、平坦且局部有限表现。因此由引理
\ref{stacks-morphisms-lemma-descent-quasi-compact} 得出结论。
\end{proof}

\noindent
下述引理说明，对任何本身为胚的代数叠，其所有点的剩余胚都存在。

\begin{lemma}
\label{stacks-morphisms-lemma-gerbe-residual-gerbe-exists}
设 $\mathcal{X}$ 是代数叠。若 $\mathcal{X}$ 是胚，则对每个
$x \in |\mathcal{X}|$，$\mathcal{X}$ 在 $x$ 处的剩余胚都存在。
\end{lemma}

\begin{proof}
设 $\pi : \mathcal{X} \to X$ 是从 $\mathcal{X}$ 到代数空间 $X$
的态射，并且它使 $\mathcal{X}$ 成为 $X$ 上的胚。设
$Z_x \to X$ 是 $X$ 在 $x$ 处的剩余空间，见《适宜代数空间》定义
\ref{decent-spaces-definition-residual-space}。令
$\mathcal{Z} = \mathcal{X} \times_X Z_x$。由引理
\ref{stacks-morphisms-lemma-base-change-gerbe}，代数叠 $\mathcal{Z}$ 是 $Z_x$ 上的胚。
因此 $|\mathcal{Z}| = |Z_x|$（引理
\ref{stacks-morphisms-lemma-gerbe-bijection-points}）是单点集。由于
$\mathcal{Z} \to Z_x$ 是 $\pi$ 的基变换，故其局部有限表现（见引理
\ref{stacks-morphisms-lemma-gerbe-fppf} 和 \ref{stacks-morphisms-lemma-base-change-finite-presentation}）。
于是 $\mathcal{Z}$ 局部诺特，见引理
\ref{stacks-morphisms-lemma-locally-finite-type-locally-noetherian}。
所以 $\mathcal{X}$ 在 $x$ 处的剩余胚 $\mathcal{Z}_x$ 存在，并且等于
$\mathcal{Z}_x = \mathcal{Z}_{red}$，即代数叠 $\mathcal{Z}$ 的约化。
事实上，上文已经说明 $|\mathcal{Z}_{red}|$ 是映到
$x \in |\mathcal{X}|$ 的单点集；它按构造是约化的，并且局部诺特
（因为局部诺特代数叠的约化仍局部诺特，细节略）。
\end{proof}









\section{按胚分层}
\label{stacks-morphisms-section-stratify}

\noindent
本节的目标是证明，许多代数叠 $\mathcal{X}$ 都有由局部闭子叠
$\mathcal{X}_i \subset \mathcal{X}$ 构成的“分层”，其中每个
$\mathcal{X}_i$ 都是胚。这说明在某种意义下，胚是构成任意代数叠的基本构件。
注意，在本节中，我们所谓的分层仅指
$$
|\mathcal{X}| = \bigcup\nolimits_i |\mathcal{X}_i|
$$
是与 $\mathcal{X}$ 关联的拓扑空间的分层，并无其他含义。因此，为研究这种分层，
可以无妨地用 $\mathcal{X}$ 的约化替换它（见《叠的性质》一节
\ref{stacks-properties-section-reduced}）。

\medskip\noindent
下述命题说明，$\mathcal{X}$ 的约化中（几乎总是）存在一个稠密开子叠。

\begin{proposition}
\label{stacks-morphisms-proposition-open-stratum}
设 $\mathcal{X}$ 是约化代数叠，且
$\mathcal{I}_\mathcal{X} \to \mathcal{X}$ 拟紧。则存在一个稠密开子叠
$\mathcal{U} \subset \mathcal{X}$，它是胚。
\end{proposition}

\begin{proof}
由命题 \ref{stacks-morphisms-proposition-when-gerbe}，只需找到一个稠密开子叠
$\mathcal{U}$，使得 $\mathcal{I}_\mathcal{U} \to \mathcal{U}$
平坦且局部有限表现。注意
$\mathcal{I}_\mathcal{U} =
\mathcal{I}_\mathcal{X} \times_\mathcal{X} \mathcal{U}$，见引理
\ref{stacks-morphisms-lemma-cartesian-square-inertia}。

\medskip\noindent
选取一个表现 $\mathcal{X} = [U/R]$。设 $G \to U$ 是群胚 $R$ 的稳定子群
代数空间。由引理 \ref{stacks-morphisms-lemma-presentation-inertia}，$G \to U$ 是
$\mathcal{I}_\mathcal{X} \to \mathcal{X}$ 的基变换，因而按假设拟紧，并由引理
\ref{stacks-morphisms-lemma-inertia} 局部有限型。设 $W \subset U$ 是最大的开子概形
（可能为空），使得限制 $G_W \to W$ 平坦且局部有限表现
（我们略去 $W$ 存在性的证明；提示：使用这些性质是局部的这一事实）。
由《代数空间的态射》命题
\ref{spaces-morphisms-proposition-generic-flatness-reduced}，
$W \subset U$ 稠密。注意，由《空间中群胚的进一步内容》引理
\ref{spaces-more-groupoids-lemma-property-G-invariant}，
$W \subset U$ 是 $R$-不变的。因此，由《叠的性质》引理
\ref{stacks-properties-lemma-substacks-presentation}，$W$ 对应于一个开子叠
$\mathcal{U} \subset \mathcal{X}$。由于 $|U| \to |\mathcal{X}|$ 是开映射，
且 $|W| \subset |U|$ 稠密，可知 $\mathcal{U}$ 在 $\mathcal{X}$ 中稠密。
最后，态射 $\mathcal{I}_\mathcal{U} \to \mathcal{U}$ 平坦且局部有限表现，
因为沿满光滑态射 $W \to \mathcal{U}$ 的基变换是态射 $G_W \to W$，
而后者按构造平坦且局部有限表现。见引理 \ref{stacks-morphisms-lemma-descent-flat} 和
\ref{stacks-morphisms-lemma-descent-finite-presentation}。
\end{proof}

\noindent
上述命题立即蕴含：在惯性拟紧的代数叠上，任意点都有剩余胚；我们将在引理
\ref{stacks-morphisms-lemma-every-point-residual-gerbe} 中证明这一点。另一方面，由胚构成的
有限分层并非总是存在。下面给出一个例子。

\begin{example}
\label{stacks-morphisms-example-infinite-stratification}
设 $k$ 是域。取 $U = \Spec(k[x_0, x_1, x_2, \ldots])$，并令
$\mathbf{G}_m$ 按下式作用：
$t(x_0, x_1, x_2, \ldots) = (tx_0, t^p x_1, t^{p^2} x_2, \ldots)$
，其中 $p$ 是素数。令 $\mathcal{X} = [U/\mathbf{G}_m]$。这是代数叠。
$\mathcal{X}$ 有如下分层：
\begin{enumerate}
\item $\mathcal{X}_0$ 是 $x_0$ 非零的部分；
\item $\mathcal{X}_1$ 是 $x_0$ 为零而 $x_1$ 非零的部分；
\item $\mathcal{X}_2$ 是 $x_0, x_1$ 为零而 $x_2$ 非零的部分；
\item 依此类推；并且
\item $\mathcal{X}_{\infty}$ 是所有 $x_i$ 都为零的部分。
\end{enumerate}
每一层都是某个概形上的胚：对 $\mathcal{X}_i$，其群为 $\mu_{p^i}$；
对 $\mathcal{X}_{\infty}$，其群为 $\mathbf{G}_m$。这些层是约化局部闭子叠。
不存在具有相同性质而更粗的分层。
\end{example}

\noindent
尽管如此，利用超限归纳与命题 \ref{stacks-morphisms-proposition-open-stratum}，我们可以找到
可能无限的胚分层……！

\begin{lemma}
\label{stacks-morphisms-lemma-every-point-in-a-stratum}
设 $\mathcal{X}$ 是代数叠，并且
$\mathcal{I}_\mathcal{X} \to \mathcal{X}$ 拟紧。则存在一个良序指标集 $I$，
并且对每个 $i \in I$，存在一个约化局部闭子叠
$\mathcal{U}_i \subset \mathcal{X}$，使得
\begin{enumerate}
\item 每个 $\mathcal{U}_i$ 都是胚；
\item 有 $|\mathcal{X}| = \bigcup_{i \in I} |\mathcal{U}_i|$；
\item $T_i = |\mathcal{X}| \setminus \bigcup_{i' < i} |\mathcal{U}_{i'}|$
对所有 $i \in I$ 都是 $|\mathcal{X}|$ 中的闭集；并且
\item $|\mathcal{U}_i|$ 是 $T_i$ 中的开集。
\end{enumerate}
此外，还可以作如下安排：或者 (a) $|\mathcal{U}_i| \subset T_i$ 稠密，
或者 (b) $\mathcal{U}_i$ 拟紧。在情形 (a) 中，若把 $\mathcal{U}_i$
取得尽可能大（细节见证明），则该分层是典范的。
\end{lemma}

\begin{proof}
设 $T \subset |\mathcal{X}|$ 是非空闭子集。对每个这样的 $T$，我们要找出
（分别地，选取）一个约化局部闭子叠
$\mathcal{U}(T) \subset \mathcal{X}$，使得
$|\mathcal{U}(T)| \subset T$ 开且稠密（分别地，非空且拟紧）。事实上，
由《叠的性质》引理 \ref{stacks-properties-lemma-reduced-closed-substack}，
存在唯一的约化闭子叠 $\mathcal{X}' \subset \mathcal{X}$，使得
$T = |\mathcal{X}'|$。注意，由引理
\ref{stacks-morphisms-lemma-monomorphism-cartesian-square-inertia}，
$\mathcal{I}_{\mathcal{X}'} =
\mathcal{I}_\mathcal{X} \times_\mathcal{X} \mathcal{X}'$。
因此，作为一个基变换，$\mathcal{I}_{\mathcal{X}'} \to \mathcal{X}'$ 拟紧；
见引理 \ref{stacks-morphisms-lemma-base-change-quasi-compact}。所以命题
\ref{stacks-morphisms-proposition-open-stratum} 蕴含：存在一个稠密且极大的开子叠
$\mathcal{U} \subset \mathcal{X}'$（见该命题的证明），它是胚。在情形 (a)
中令 $\mathcal{U}(T) = \mathcal{U}$（这是典范的）；在情形 (b) 中，
只需选取一个非空拟紧开子叠
$\mathcal{U}(T) \subset \mathcal{U}$，见《叠的性质》引理
\ref{stacks-properties-lemma-space-locally-quasi-compact}
（借助选择公理，可以对所有 $T$ 同时这样做）。

\medskip\noindent
利用超限递归，我们对每个序数 $\alpha$ 构造一个闭子集
$T_\alpha \subset |\mathcal{X}|$。对 $\alpha = 0$，令
$T_0 = |\mathcal{X}|$。给定 $T_\alpha$ 后，令
$$
T_{\alpha + 1} = T_\alpha \setminus |\mathcal{U}(T_\alpha)|.
$$
若 $\beta$ 是极限序数，则令
$$
T_\beta = \bigcap\nolimits_{\alpha < \beta} T_\alpha.
$$
我们断言，对所有充分大的 $\alpha$，都有 $T_\alpha = \emptyset$。事实上，
假设对所有 $\alpha$ 都有 $T_\alpha \not = \emptyset$。于是可以得到从
全体序数这个类到 $|\mathcal{X}|$ 的子集所成集合的单射，矛盾。

\medskip\noindent
该断言蕴含本引理。事实上，令
$$
I = \{\alpha \mid \mathcal{U}_\alpha \not = \emptyset \}.
$$
由该断言，这是良序集。对 $i = \alpha \in I$，令
$\mathcal{U}_i = \mathcal{U}_\alpha$。于是 $\mathcal{U}_i$
是约化局部闭子叠且是胚，即 (1) 成立。按构造，若
$i = \alpha \in I$，则 $T_i = T\alpha$，所以 (3) 成立。由我们对
$\mathcal{U}(T)$ 的选取，(4) 以及 (a) 或 (b) 也成立。最后，为证明 (2)，
取 $x \in |\mathcal{X}|$。存在使 $x \not \in T_\beta$ 的最小序数
$\beta$（因为序数是良序的）。由极限序数情形下 $T_\beta$ 的定义，
此时 $\beta$ 必为后继序数。因此 $\beta = \alpha + 1$，且
$x \in |\mathcal{U}(T_\alpha)|$，得证。
\end{proof}

\begin{remark}
\label{stacks-morphisms-remark-order-type}
我们可以问，引理 \ref{stacks-morphisms-lemma-every-point-in-a-stratum} 所构造的 (a) 型分层
给出的典范分层具有怎样的序型。一个自然猜想是，良序集 $I$ 的
{\it 基数}至多为 $\aleph_0$。我们不知道这是真是假。若你知道，请发送邮件至
\href{mailto:stacks.project@gmail.com}{stacks.project@gmail.com}.
\end{remark}







\section{代数叠的拓扑空间}
\label{stacks-morphisms-section-topology}

\noindent
本节把前面的结果应用于代数叠所关联的拓扑空间 $|\mathcal{X}|$。

\begin{lemma}
\label{stacks-morphisms-lemma-spectral-qc-diagonal-qc}
设 $\mathcal{X}$ 是拟紧代数叠，且其对角态射 $\Delta$ 拟紧。
则 $|\mathcal{X}|$ 是谱拓扑空间。
\end{lemma}

\begin{proof}
选取仿射概形 $U$ 以及满光滑态射 $U \to \mathcal{X}$，见《叠的性质》引理
\ref{stacks-properties-lemma-quasi-compact-stack}。于是
$|U| \to |\mathcal{X}|$ 连续、开且满，见《叠的性质》引理
\ref{stacks-properties-lemma-topology-points}。因此，$|\mathcal{X}|$
的拟紧开子集构成其拓扑的一组基。对拟紧开子集
$W_1, W_2 \subset |\mathcal{X}|$，可以选取 $U$ 的拟紧开子集
$V_1, V_2$，它们分别映到 $W_1$ 和 $W_2$。由于 $\Delta$ 拟紧，可知
$$
V_1 \times_\mathcal{X} V_2 =
(V_1 \times V_2) \times_{\mathcal{X} \times \mathcal{X}, \Delta} \mathcal{X}
$$
拟紧。由《叠的性质》引理
\ref{stacks-properties-lemma-points-cartesian}，
$|V_1 \times_\mathcal{X} V_2|$ 在 $|\mathcal{X}|$ 中的像是
$W_1 \cap W_2$。所以 $W_1 \cap W_2$ 拟紧。为完成证明，只需证明
$|\mathcal{X}|$ 是清醒的，见《拓扑学》定义
\ref{topology-definition-spectral-space}。

\medskip\noindent
设 $T \subset |\mathcal{X}|$ 是不可约闭子集。我们必须证明 $T$
有唯一的一般点。设 $\mathcal{Z} \subset \mathcal{X}$ 是与 $T$
对应的约化诱导闭子叠，见《叠的性质》定义
\ref{stacks-properties-definition-reduced-induced-stack}。
由于 $\mathcal{Z} \to \mathcal{X}$ 是闭浸入，可知
$\Delta_\mathcal{Z}$ 拟紧：首先，把
$\mathcal{Z} \to \mathcal{X} \times \mathcal{X}$ 写成
$\mathcal{Z} \to \mathcal{X}$ 与 $\Delta$ 的复合，从而证明它拟紧；
再把 $\mathcal{Z} \to \mathcal{X} \times \mathcal{X}$ 写成
$\Delta_\mathcal{Z}$ 与
$\mathcal{Z} \times \mathcal{Z} \to \mathcal{X} \times \mathcal{X}$
的复合，并使用引理 \ref{stacks-morphisms-lemma-quasi-compact-permanence} 以及
$\mathcal{Z} \times \mathcal{Z} \to \mathcal{X} \times \mathcal{X}$
分离这一事实。因此，问题归结为下一段讨论的情形。

\medskip\noindent
假设 $\mathcal{X}$ 拟紧，$\Delta$ 拟紧，$\mathcal{X}$ 约化，且
$|\mathcal{X}|$ 不可约。我们必须证明 $|\mathcal{X}|$ 有唯一的一般点。
由于 $\mathcal{I}_\mathcal{X} \to \mathcal{X}$ 是 $\Delta$ 的基变换，
可知 $\mathcal{I}_\mathcal{X} \to \mathcal{X}$ 拟紧
（引理 \ref{stacks-morphisms-lemma-base-change-quasi-compact}）。因此，由命题
\ref{stacks-morphisms-proposition-open-stratum}，存在一个稠密开子叠
$\mathcal{U} \subset \mathcal{X}$，它是胚。换言之，
$|\mathcal{U}| \subset |\mathcal{X}|$ 开且稠密。于是可以假设
$\mathcal{X}$ 是胚。设 $\mathcal{X} \to X$ 使 $\mathcal{X}$
成为代数空间 $X$ 上的胚。由引理 \ref{stacks-morphisms-lemma-gerbe-bijection-points}，
$|\mathcal{X}| \cong |X|$。特别地，$X$ 拟紧。由引理
\ref{stacks-morphisms-lemma-gerbe-diagonal-quasi-compact}，$X$ 的对角态射拟紧；也就是说，
$X$ 是拟分离代数空间。于是，由《代数空间的性质》引理
\ref{spaces-properties-lemma-quasi-compact-quasi-separated-spectral}，
$|X|$ 是谱空间，这就给出了所需结论。
\end{proof}

\begin{lemma}
\label{stacks-morphisms-lemma-spectral-qcqs}
设 $\mathcal{X}$ 是拟紧拟分离代数叠。则 $|\mathcal{X}|$ 是谱拓扑空间。
\end{lemma}

\begin{proof}
这是引理 \ref{stacks-morphisms-lemma-spectral-qc-diagonal-qc} 的特殊情形。
\end{proof}

\begin{lemma}
\label{stacks-morphisms-lemma-sober-qs}
设 $\mathcal{X}$ 是对角态射拟紧的代数叠（例如，若 $\mathcal{X}$ 拟分离）。
则存在开覆盖 $|\mathcal{X}| = \bigcup U_i$，其中每个 $U_i$ 都是谱空间。
特别地，$|\mathcal{X}|$ 是清醒拓扑空间。
\end{lemma}

\begin{proof}
这是引理 \ref{stacks-morphisms-lemma-spectral-qc-diagonal-qc} 的直接推论。
\end{proof}






\section{剩余胚的存在性}
\label{stacks-morphisms-section-existence-residual-gerbes}

\noindent
代数叠上一点的剩余胚定义见《叠的性质》定义
\ref{stacks-properties-definition-residual-gerbe}。我们已经证明，有限型点
（引理 \ref{stacks-morphisms-lemma-point-finite-type-monomorphism}）以及胚的任意一点
（引理 \ref{stacks-morphisms-lemma-gerbe-residual-gerbe-exists}）都有剩余胚。本节证明，
许多代数叠上都存在剩余胚。首先给出前面所承诺的命题
\ref{stacks-morphisms-proposition-open-stratum} 的应用。

\begin{lemma}
\label{stacks-morphisms-lemma-every-point-residual-gerbe}
设 $\mathcal{X}$ 是代数叠，并且
$\mathcal{I}_\mathcal{X} \to \mathcal{X}$ 拟紧。则对每个
$x \in |\mathcal{X}|$，$\mathcal{X}$ 在 $x$ 处的剩余胚都存在。
\end{lemma}

\begin{proof}
设 $T = \overline{\{x\}} \subset |\mathcal{X}|$ 是 $x$ 的闭包。
由《叠的性质》引理 \ref{stacks-properties-lemma-reduced-closed-substack}，
存在约化闭子叠 $\mathcal{X}' \subset \mathcal{X}$，使得
$T = |\mathcal{X}'|$。注意，由引理
\ref{stacks-morphisms-lemma-monomorphism-cartesian-square-inertia}，
$\mathcal{I}_{\mathcal{X}'} =
\mathcal{I}_\mathcal{X} \times_\mathcal{X} \mathcal{X}'$。因此，作为一个
基变换，$\mathcal{I}_{\mathcal{X}'} \to \mathcal{X}'$ 拟紧，见引理
\ref{stacks-morphisms-lemma-base-change-quasi-compact}。所以由命题
\ref{stacks-morphisms-proposition-open-stratum}，存在稠密开子叠
$\mathcal{U} \subset \mathcal{X}'$，它是胚。注意，
$x \in |\mathcal{U}|$，因为 $\{x\} \subset T$ 也是稠密子集。于是由引理
\ref{stacks-morphisms-lemma-gerbe-residual-gerbe-exists}，$\mathcal{U}$ 在 $x$ 处的剩余胚
$\mathcal{Z}_x \subset \mathcal{U}$ 存在。由定义立即可知，
$\mathcal{Z}_x \to \mathcal{X}$ 是 $\mathcal{X}$ 在 $x$ 处的剩余胚。
\end{proof}

\noindent
若叠是拟 DM 的，则剩余胚也存在。特别地，Deligne--Mumford 叠上总存在剩余胚。

\begin{lemma}
\label{stacks-morphisms-lemma-every-point-residual-gerbe-quasi-DM}
设 $\mathcal{X}$ 是拟 DM 代数叠。则对每个 $x \in |\mathcal{X}|$，
$\mathcal{X}$ 在 $x$ 处的剩余胚都存在。
\end{lemma}

\begin{proof}
选取概形 $U$ 以及满、平坦、局部有限表现且局部拟有限的态射
$U \to \mathcal{X}$，见定理 \ref{stacks-morphisms-theorem-quasi-DM}。令
$R = U \times_\mathcal{X} U$。作为态射 $U \to \mathcal{X}$ 的基变换，
投影 $s, t : R \to U$ 满、平坦、局部有限表现且局部拟有限。存在典范态射
$[U/R] \to \mathcal{X}$（见《代数叠》引理
\ref{algebraic-lemma-map-space-into-stack}）；由于
$U \to \mathcal{X}$ 满、平坦且局部有限表现，它是等价，见《代数叠》评注
\ref{algebraic-remark-flat-fp-presentation}。因此，可以假设
$\mathcal{X} = [U/R]$，其中 $(U, R, s, t, c)$ 是代数空间中的群胚，
且 $s, t : R \to U$ 满、平坦、局部有限表现且局部拟有限。令
$$
U' = \coprod\nolimits_{u \in U\text{ 位于下列对象之上： }x} \Spec(\kappa(u)).
$$
典范态射 $U' \to U$ 是单态射。令
$$
R' = U' \times_\mathcal{X} U' =
R \times_{(U \times U)} (U' \times U')
$$
由于 $U' \to U$ 是单态射，可知两个投影 $s', t' : R' \to U'$
都分解为一个单态射与一个局部拟有限态射的复合。因此，由于 $U'$
是域的谱的不交并，应用《域上的代数空间》引理
\ref{spaces-over-fields-lemma-mono-towards-locally-quasi-finite-over-field}，
可知态射 $s', t' : R' \to U'$ 局部拟有限。再次利用 $U'$ 是域的谱的不交并，
态射 $s', t'$ 也平坦。最后，$s', t'$ 局部拟有限蕴含 $s', t'$
局部有限型，进而 $s', t'$ 局部有限表现（因为 $U'$ 是域的谱的不交并，特别地局部诺特，
所以可以应用《代数空间的态射》引理
\ref{spaces-morphisms-lemma-noetherian-finite-type-finite-presentation}）。
于是，由《可表示性判据》定理
\ref{criteria-theorem-flat-groupoid-gives-algebraic-stack}，
$\mathcal{Z} = [U'/R']$ 是代数叠。由于 $R'$ 是 $R$ 沿 $U' \to U$
的限制，由《空间中的群胚》引理
\ref{spaces-groupoids-lemma-quotient-stack-restrict} 和《叠的性质》引理
\ref{stacks-properties-lemma-monomorphism}，可知
$\mathcal{Z} \to \mathcal{X}$ 是单态射。由于
$\mathcal{Z} \to \mathcal{X}$ 是单态射，因此
$|\mathcal{Z}| \to |\mathcal{X}|$ 单，见《叠的性质》引理
\ref{stacks-properties-lemma-monomorphism-injective-points}。由《叠的性质》引理
\ref{stacks-properties-lemma-points-cartesian}，可知
$$
|U'| = |\mathcal{Z} \times_\mathcal{X} U'|
\longrightarrow
|\mathcal{Z}| \times_{|\mathcal{X}|} |U'|
$$
满；结合我们对 $U'$ 的选取，这蕴含
$|\mathcal{Z}| \to |\mathcal{X}|$ 的像为 $\{x\}$。所以
$|\mathcal{Z}|$ 是单点集。最后，按构造，$U'$ 局部诺特且约化；也就是说，
$\mathcal{Z}$ 约化且局部诺特。这说明 $\mathcal{Z} \to \mathcal{X}$
的本质像是 $\mathcal{X}$ 在 $x$ 处的剩余胚，见《叠的性质》引理
\ref{stacks-properties-lemma-residual-gerbe-unique}。
\end{proof}

\begin{lemma}
\label{stacks-morphisms-lemma-every-point-residual-gerbe-locally-Noetherian}
设 $\mathcal{X}$ 是局部诺特代数叠。则对每个
$x \in |\mathcal{X}|$，$\mathcal{X}$ 在 $x$ 处的剩余胚都存在。
\end{lemma}

\begin{proof}
选取仿射概形 $U$ 和光滑态射 $U \to \mathcal{X}$，使 $x$ 位于开连续映射
$|U| \to |\mathcal{X}|$ 的像中。我们可以并且确实用与
$|U| \to |\mathcal{X}|$ 的像对应的开子叠替换 $\mathcal{X}$，见
《叠的性质》引理 \ref{stacks-properties-lemma-open-substacks}。因此，可以假设
$\mathcal{X} = [U/R]$，其中 $(U, R, s, t, c)$ 是代数空间中的光滑群胚，
且 $U$ 是诺特仿射概形；见《代数叠》一节
\ref{algebraic-section-stack-to-presentation}。

\medskip\noindent
设 $E \subset |U|$ 是 $\{x\} \subset |\mathcal{X}|$ 的逆像。当然，
$E \not = \emptyset$。由于 $|U|$ 是诺特拓扑空间，可以选取
$u \in E$，使得 $\overline{\{u\}} \cap E = \{u\}$。像往常一样，把
$u = \Spec(\kappa(u))$ 看作其剩余域的谱。记
$$
F = u \times_{U, t} R = u \times_\mathcal{X} U
\quad\text{且}\quad
R' = (u \times u) \times_{(U \times U), (t, s)} R = u \times_\mathcal{X} u
$$
。此外，记 $Z = \overline{\{u\}} \subset U$，赋予它约化诱导概形结构。
记 $p : F \to U$ 为第二投影诱导的态射（在 $F$ 的第一个纤维积描述中，
使用 $s : R \to U$）。则 $E$ 是 $p$ 的集合论像。态射 $R' \to F$
是单态射，并且经由 $Z$ 的逆像 $p^{-1}(Z)$ 分解。该逆像
$p^{-1}(Z) \subset F$ 是闭子概形；由于上面对 $u \in E$ 的谨慎选取，
限制 $p|_{p^{-1}(Z)} : p^{-1}(Z) \to Z$ 的集合论像包含于
$\{u\} \subset Z$。由于 $u = \lim W$，其中极限遍历不可约约化概形
$Z$ 的非空仿射开子概形，可知态射
$p|_{p^{-1}(Z)} : p^{-1}(Z) \to Z$ 经由态射 $u \to Z$ 分解。这显然蕴含
$R' = p^{-1}(Z)$。特别地，态射 $t' : R' \to u$ 局部有限表现，因为它是
局部诺特代数空间的闭浸入 $p^{-1}(Z) \to F$ 与光滑态射
$\text{pr}_1 : F \to u$ 的复合；使用《代数空间的态射》诸引理
\ref{spaces-morphisms-lemma-locally-finite-type-locally-noetherian}、
\ref{spaces-morphisms-lemma-closed-immersion-finite-presentation} 和
\ref{spaces-morphisms-lemma-composition-finite-presentation}。因此，
$(U, R, s, t, c)$ 沿 $u \to U$ 的限制 $(u, R', s', t', c')$
是代数空间中的群胚，其中 $s'$ 和 $t'$ 平坦且局部有限表现。所以由
《可表示性判据》定理 \ref{criteria-theorem-flat-groupoid-gives-algebraic-stack}，
$\mathcal{Z} = [u/R']$ 是代数叠。由于 $R'$ 是 $R$ 沿 $u \to U$
的限制，由《空间中的群胚》引理
\ref{spaces-groupoids-lemma-quotient-stack-restrict} 和《叠的性质》引理
\ref{stacks-properties-lemma-monomorphism}，可知
$\mathcal{Z} \to \mathcal{X}$ 是单态射。再由《叠的性质》一节
\ref{stacks-properties-section-residual-gerbe} 的内容，
$\mathcal{Z}$（同构于）该剩余胚。
\end{proof}













\section{\'Etale 局部结构}
\label{stacks-morphisms-section-etale-local}

\noindent
本节开始讨论关于代数叠的 \'etale 局部结构可以说些什么。

\begin{lemma}
\label{stacks-morphisms-lemma-quotient-etale}
设 $Y$ 是代数空间。设 $(U, R, s, t, c)$ 是 $Y$ 上代数空间中的群胚。
假设 $U \to Y$ 平坦且局部有限表现，并且
$R \to U \times_Y U$ 是开浸入。则 $X = [U/R] = U/R$ 是代数空间，
且 $X \to Y$ 是 \'etale 的。
\end{lemma}

\begin{proof}
由《可表示性判据》定理
\ref{criteria-theorem-flat-groupoid-gives-algebraic-stack}，
商叠 $[U/R]$ 是代数叠。另一方面，由于
$R \to U \times U$ 是单态射，它是代数空间（这里沿用我们的语言滥用，并见
《代数叠》命题
\ref{algebraic-proposition-algebraic-stack-no-automorphisms}），当然也等于
代数空间 $U/R$（《引导》定理 \ref{bootstrap-theorem-final-bootstrap}）。
由于 $U \to X$ 满、平坦且局部有限表现
（《引导》引理 \ref{bootstrap-lemma-covering-quotient}），由《代数空间的态射》
引理 \ref{spaces-morphisms-lemma-flat-permanence} 和《代数空间上的下降》引理
\ref{spaces-descent-lemma-flat-finitely-presented-permanence}，可知
$X \to Y$ 平坦且局部有限表现。最后，考虑笛卡儿图
$$
\xymatrix{
R \ar[d] \ar[r] & U \times_Y U \ar[d] \\
X \ar[r] & X \times_Y X
}
$$
由于右侧竖直箭头满、平坦且局部有限表现（略去一个小细节），而顶部水平箭头
按假设是开浸入，可知 $X \to X \times_Y X$ 是开浸入（使用《叠的性质》引理
\ref{stacks-properties-lemma-check-property-covering}）。所以由《代数空间的态射》
引理 \ref{spaces-morphisms-lemma-diagonal-unramified-morphism}，
$X \to Y$ 非分歧。最后应用《代数空间的态射》引理
\ref{spaces-morphisms-lemma-unramified-flat-lfp-etale} 即得结论。
\end{proof}

\begin{lemma}
\label{stacks-morphisms-lemma-quasi-splitting-etale}
设 $S$ 是概形。设 $(U, R, s, t, c)$ 是 $S$ 上代数空间中的群胚。
假设 $s, t$ 平坦且局部有限表现。设 $P \subset R$ 是开子空间，使得
$(U, P, s|_P, t|_P, c|_{P \times_{s, U, t} P})$
构成 $S$ 上代数空间中的群胚。则
$$
[U/P] \longrightarrow [U/R]
$$
是代数叠的态射，可由代数空间表示、满且 \'etale。
\end{lemma}

\begin{proof}
由于 $P \subset R$ 是开的，$s|_P$ 和 $t|_P$ 平坦且局部有限表现。因此，
由《可表示性判据》定理
\ref{criteria-theorem-flat-groupoid-gives-algebraic-stack}，
$[U/R]$ 和 $[U/P]$ 都是代数叠。为证明该态射可由代数空间表示，只需证明
$[U/P] \to [U/R]$ 在纤维范畴上忠实，见《代数叠》引理
\ref{algebraic-lemma-characterize-representable-by-algebraic-spaces}。
这立即由 $P \to R$ 是单态射以及《空间中的群胚》引理
\ref{spaces-groupoids-lemma-quotient-stack-objects} 给出的商叠对象的显式描述得出。
至此，由《代数叠》定义
\ref{algebraic-definition-relative-representable-property}，我们知道
$[U/P] \to [U/R]$ 满和 \'etale 分别意味着什么。例如，满性由
《可表示性判据》引理 \ref{criteria-lemma-representable-by-spaces-surjective}
以及域的谱上的商叠对象描述（见上引引理）得出。还需证明该态射是 \'etale 的。

\medskip\noindent
为此，只需证明 $U \times_{[U/R]} [U/P] \to U$ 是 \'etale 的，见
《叠的性质》引理 \ref{stacks-properties-lemma-check-property-covering}。
由《空间中的群胚》引理
\ref{spaces-groupoids-lemma-cartesian-square-of-morphism}，该纤维积等于
$[R/P \times_{s, U, t} R]$，而到 $U$ 的态射由
$s : R \to U$ 诱导。映射
$s', t' : P \times_{s, U, t} R \to R$ 分别由
$s' : (p, r) \mapsto r$ 和 $t' : (p, r) \mapsto c(p, r)$ 给出。由于
$P \subset R$ 是开的，可知
$(t', s') : P \times_{s, U, t} R \to R \times_{s, U, s} R$
是开浸入。因此应用引理 \ref{stacks-morphisms-lemma-quotient-etale} 即得结论。
\end{proof}

\begin{lemma}
\label{stacks-morphisms-lemma-etale-local-quasi-DM}
设 $\mathcal{X}$ 是代数叠。假设 $\mathcal{X}$ 拟 DM 且对角态射分离
（等价地，$\mathcal{I}_\mathcal{X} \to \mathcal{X}$
局部拟有限且分离）。设 $x \in |\mathcal{X}|$。则存在代数叠的态射
$$
\mathcal{U} \longrightarrow \mathcal{X}
$$
，具有下列性质：
\begin{enumerate}
\item 存在映到 $x$ 的点 $u \in |\mathcal{U}|$；
\item $\mathcal{U} \to \mathcal{X}$ 可由代数空间表示且 \'etale；
\item $\mathcal{U} = [U/R]$，其中 $(U, R, s, t, c)$ 是概形群胚，
$U$、$R$ 仿射，而 $s, t$ 有限、平坦且局部有限表现。
\end{enumerate}
\end{lemma}

\begin{proof}
（括号中的陈述由引理 \ref{stacks-morphisms-lemma-diagonal-diagonal} 中的等价条件得出。）
选取仿射概形 $U$ 以及平坦、局部有限表现、局部拟有限的态射
$U \to \mathcal{X}$，使得 $x$ 是某个点 $u \in U$ 的像。由定理
\ref{stacks-morphisms-theorem-quasi-DM} 以及 $\mathcal{X}$ 拟 DM 的假设，这可以做到。
令 $(U, R, s, t, c)$ 为取 $R = U \times_\mathcal{X} U$
所得的代数空间中的群胚，见《代数叠》引理
\ref{algebraic-lemma-map-space-into-stack}。设
$\mathcal{X}' \subset \mathcal{X}$ 是与 $|U| \to |\mathcal{X}|$
的开像对应的开子叠（《叠的性质》诸引理
\ref{stacks-properties-lemma-topology-points} 和
\ref{stacks-properties-lemma-open-substacks}）。显然，可以用开子叠
$\mathcal{X}'$ 替换 $\mathcal{X}$。于是可以假设
$U \to \mathcal{X}$ 满，此时《代数叠》评注
\ref{algebraic-remark-flat-fp-presentation} 给出
$\mathcal{X} = [U/R]$。注意，$s, t : R \to U$ 平坦、局部有限表现且
局部拟有限。由于
$R = U \times U \times_{(\mathcal{X} \times \mathcal{X})} \mathcal{X}$，
且 $\mathcal{X}$ 的对角态射分离，可知 $R$ 分离。因此
$s, t : R \to U$ 分离。将《代数空间的态射》命题
\ref{spaces-morphisms-proposition-locally-quasi-finite-separated-over-scheme}
应用于 $s : R \to U$，可知 $R$ 是概形。

\medskip\noindent
上文已经验证，$(U, R, s, t, c)$ 与 $u$ 满足
《空间中群胚的进一步内容》引理
\ref{spaces-more-groupoids-lemma-quasi-splitting-affine-scheme}
的所有假设。因此，可以找到基本 \'etale 邻域
$(U', u') \to (U, u)$，使得 $R$ 在 $U'$ 上的限制 $R'$ 在 $u$ 上拟分裂。
注意，$R' = U' \times_\mathcal{X} U'$（略去一个小细节；提示：使用纤维积的
传递性）。用 $(U', R', s', t', c')$ 替换 $(U, R, s, t, c)$，并像上文那样
缩小 $\mathcal{X}$，可以假设 $(U, R, s, t, c)$ 在 $u$ 上有一个拟分裂
（从《空间中群胚的进一步内容》定义
\ref{spaces-more-groupoids-definition-split-at-point} 的脚注可见，
从此以后点 $u$ 已无关紧要）。设 $P \subset R$ 是 $R$ 在 $u$ 上的一个拟分裂。
应用引理 \ref{stacks-morphisms-lemma-quasi-splitting-etale} 可知
$$
\mathcal{U} = [U/P] \longrightarrow [U/R] = \mathcal{X}
$$
具有所有所需性质。
\end{proof}

\begin{lemma}
\label{stacks-morphisms-lemma-etale-local-quasi-DM-at-x}
设 $\mathcal{X}$ 是代数叠。假设 $\mathcal{X}$ 拟 DM 且对角态射分离
（等价地，$\mathcal{I}_\mathcal{X} \to \mathcal{X}$ 局部拟有限且分离）。
设 $x \in |\mathcal{X}|$。假设 $\mathcal{X}$ 在 $x$ 处的自同构群有限
（评注 \ref{stacks-morphisms-remark-property-automorphism-groups}）。则存在代数叠的态射
$$
g : \mathcal{U} \longrightarrow \mathcal{X}
$$
，具有下列性质：
\begin{enumerate}
\item 存在映到 $x$ 的点 $u \in |\mathcal{U}|$，且 $g$ 诱导 $u$ 与 $x$
处自同构群之间的同构（评注 \ref{stacks-morphisms-remark-identify-automorphism-groups}）；
\item $\mathcal{U} \to \mathcal{X}$ 可由代数空间表示且 \'etale；
\item $\mathcal{U} = [U/R]$，其中 $(U, R, s, t, c)$ 是概形群胚，
$U$、$R$ 仿射，而 $s, t$ 有限、平坦且局部有限表现。
\end{enumerate}
\end{lemma}

\begin{proof}
由引理 \ref{stacks-morphisms-lemma-automorphism-group-scheme}，$G_x$ 是群概形。证明的第一部分
与引理 \ref{stacks-morphisms-lemma-etale-local-quasi-DM} 的证明第一部分{\bf 完全}相同。因此，
可以假设 $\mathcal{X} = [U/R]$，其中 $(U, R, s, t, c)$ 以及映到 $x$ 的
$u \in U$ 满足《空间中群胚的进一步内容》引理
\ref{spaces-more-groupoids-lemma-quasi-splitting-affine-scheme}
的全部假设。关于 $G_x$ 的假设蕴含 $G_u$ 在 $u$ 上有限。因此，
《空间中群胚的进一步内容》引理
\ref{spaces-more-groupoids-lemma-splitting-affine-scheme} 的全部假设都成立。
所以可以找到基本 \'etale 邻域 $(U', u') \to (U, u)$，使得 $R$ 在 $U'$
上的限制 $R'$ 在 $u$ 上分裂。注意，
$R' = U' \times_\mathcal{X} U'$（略去一个小细节；提示：使用纤维积的传递性）。
用 $(U', R', s', t', c')$ 替换 $(U, R, s, t, c)$，并像上文那样缩小
$\mathcal{X}$，可以假设 $(U, R, s, t, c)$ 在 $u$ 上有一个分裂。
设 $P \subset R$ 是 $R$ 在 $u$ 上的一个分裂。应用引理
\ref{stacks-morphisms-lemma-quasi-splitting-etale} 可知
$$
\mathcal{U} = [U/P] \longrightarrow [U/R] = \mathcal{X}
$$
可由代数空间表示且 \'etale。按构造，$G_u$ 包含于 $P$，所以该态射在
$u$ 处的自同构群上给出所需的同构。
\end{proof}

\begin{lemma}
\label{stacks-morphisms-lemma-etale-local-quasi-DM-at-x-inertia}
设 $\mathcal{X}$ 是代数叠。假设 $\mathcal{X}$ 拟 DM 且对角态射分离
（等价地，$\mathcal{I}_\mathcal{X} \to \mathcal{X}$ 局部拟有限且分离）。
设 $x \in |\mathcal{X}|$。假设 $x$ 可以由拟紧态射
$\Spec(k) \to \mathcal{X}$ 表示。则存在代数叠的态射
$$
g : \mathcal{U} \longrightarrow \mathcal{X}
$$
，具有下列性质：
\begin{enumerate}
\item 存在映到 $x$ 的点 $u \in |\mathcal{U}|$，且 $g$ 诱导 $u$ 与 $x$
处剩余胚之间的同构；
\item $\mathcal{U} \to \mathcal{X}$ 可由代数空间表示且 \'etale；
\item $\mathcal{U} = [U/R]$，其中 $(U, R, s, t, c)$ 是概形群胚，
$U$、$R$ 仿射，而 $s, t$ 有限、平坦且局部有限表现。
\end{enumerate}
\end{lemma}

\begin{proof}
证明的第一部分与引理 \ref{stacks-morphisms-lemma-etale-local-quasi-DM} 的证明第一部分
{\bf 完全}相同。因此，可以假设 $\mathcal{X} = [U/R]$，其中
$(U, R, s, t, c)$ 以及映到 $x$ 的 $u \in U$ 满足
《空间中群胚的进一步内容》引理
\ref{spaces-more-groupoids-lemma-quasi-splitting-affine-scheme}
的全部假设。注意，$u = \Spec(\kappa(u)) \to \mathcal{X}$ 拟紧，见
《叠的性质》引理
\ref{stacks-properties-lemma-UR-quasi-compact-above-x}。考虑笛卡儿图
$$
\xymatrix{
F \ar[d] \ar[r] & U \ar[d] \\
u \ar[r]^u & \mathcal{X}
}
$$
由于 $U$ 是仿射概形且 $F \to U$ 拟紧，可知 $F$ 拟紧。由于
$U \to \mathcal{X}$ 局部拟有限，可知 $F \to u$ 局部拟有限。因此
$F \to u$ 拟有限，并且 $F$ 是仿射概形，其底拓扑空间有限离散
（《域上的代数空间》引理
\ref{spaces-over-fields-lemma-locally-quasi-finite-over-field}）。
注意，我们有单态射 $u \times_\mathcal{X} u \to F$。特别地，集合
$\{r \in R : s(r) = u, t(r) = u\}$ 是
$|u \times_\mathcal{X} u| \to |R|$ 的像，故为有限集。于是，
《空间中群胚的进一步内容》引理
\ref{spaces-more-groupoids-lemma-strong-splitting-affine-scheme}
的全部假设成立。

\medskip\noindent
所以可以找到基本 \'etale 邻域 $(U', u') \to (U, u)$，使得 $R$ 在 $U'$
上的限制 $R'$ 在 $u'$ 上强分裂。注意，
$R' = U' \times_\mathcal{X} U'$（略去一个小细节；提示：使用纤维积的传递性）。
用 $(U', R', s', t', c')$ 替换 $(U, R, s, t, c)$，并像上文那样缩小
$\mathcal{X}$，可以假设 $(U, R, s, t, c)$ 在 $u$ 上有一个强分裂。
设 $P \subset R$ 是 $R$ 在 $u$ 上的一个强分裂。应用引理
\ref{stacks-morphisms-lemma-quasi-splitting-etale} 可知
$$
\mathcal{U} = [U/P] \longrightarrow [U/R] = \mathcal{X}
$$
可由代数空间表示且 \'etale。由于 $P \subset R$ 是开的，并且按构造包含
$\{r \in R : s(r) = u, t(r) = u\}$，可知
$u \times_\mathcal{U} u \to u \times_\mathcal{X} u$ 是同构。关于剩余胚的
陈述遂由《叠的性质》引理
\ref{stacks-properties-lemma-residual-gerbe-isomorphic} 得出
（注意，由引理 \ref{stacks-morphisms-lemma-every-point-residual-gerbe-quasi-DM}，
所涉及的剩余胚确实存在）。
\end{proof}







\section{光滑态射}
\label{stacks-morphisms-section-smooth}

\noindent
代数空间态射“为光滑”的性质在源与靶上光滑局部，见《代数空间上的下降》评注
\ref{spaces-descent-remark-list-local-source-target}。该性质在基变换下也稳定，
并且在靶上 fpqc 局部，见《代数空间的态射》引理
\ref{spaces-morphisms-lemma-base-change-smooth} 和《代数空间上的下降》引理
\ref{spaces-descent-lemma-descending-property-smooth}。因此，由上面的引理
\ref{stacks-morphisms-lemma-local-source-target}，可以如下定义代数空间的态射何时光滑；
当该态射可由代数空间表示时，这一定义与《叠的性质》一节
\ref{stacks-properties-section-properties-morphisms} 中已有的概念一致。

\begin{definition}
\label{stacks-morphisms-definition-smooth}
设 $f : \mathcal{X} \to \mathcal{Y}$ 是代数叠的态射。若引理
\ref{stacks-morphisms-lemma-local-source-target} 的等价条件对
$\mathcal{P} = \text{smooth}$ 成立，则称 $f$ {\it 光滑}。
\end{definition}

\begin{lemma}
\label{stacks-morphisms-lemma-composition-smooth}
光滑态射的复合仍光滑。
\end{lemma}

\begin{proof}
结合评注 \ref{stacks-morphisms-remark-composition} 与《代数空间的态射》引理
\ref{spaces-morphisms-lemma-composition-smooth}。
\end{proof}

\begin{lemma}
\label{stacks-morphisms-lemma-base-change-smooth}
光滑态射的基变换仍光滑。
\end{lemma}

\begin{proof}
结合评注 \ref{stacks-morphisms-remark-base-change} 与《代数空间的态射》引理
\ref{spaces-morphisms-lemma-base-change-smooth}。
\end{proof}

\begin{lemma}
\label{stacks-morphisms-lemma-descent-smooth}
设 $f : \mathcal{X} \to \mathcal{Y}$ 是代数叠的态射。设
$\mathcal{Z} \to \mathcal{Y}$ 是代数叠的满、平坦、局部有限表现态射。
若基变换
$\mathcal{Z} \times_\mathcal{Y} \mathcal{X} \to \mathcal{Z}$
光滑，则 $f$ 光滑。
\end{lemma}

\begin{proof}
“光滑”这一性质满足引理 \ref{stacks-morphisms-lemma-descent-property} 的条件。在源与靶上
光滑局部这一点已在本节导言中看到，而在靶上 fppf 局部是
《代数空间上的下降》引理
\ref{spaces-descent-lemma-descending-property-smooth}。
\end{proof}

\begin{lemma}
\label{stacks-morphisms-lemma-smooth-locally-finite-presentation}
代数叠的光滑态射局部有限表现。
\end{lemma}

\begin{proof}
略。
\end{proof}

\begin{lemma}
\label{stacks-morphisms-lemma-where-smooth}
设 $f : \mathcal{X} \to \mathcal{Y}$ 是代数叠的态射。存在最大的开子叠
$\mathcal{U} \subset \mathcal{X}$，使得
$f|_\mathcal{U} : \mathcal{U} \to \mathcal{Y}$ 光滑。此外，
该开子叠的构造与下列操作交换：
\begin{enumerate}
\item 前复合光滑态射；
\item 沿平坦且局部有限表现的态射作基变换；
\item 在 $f$ 局部有限表现时，沿平坦态射作基变换。
\end{enumerate}
\end{lemma}

\begin{proof}
选取交换图
$$
\xymatrix{
U \ar[d]_a \ar[r]_h & V \ar[d]^b \\
\mathcal{X} \ar[r]^f & \mathcal{Y}
}
$$
，其中 $U$ 和 $V$ 是代数空间，竖直箭头光滑，而
$a : U \to \mathcal{X}$ 满。存在最大的开子空间 $U' \subset U$，
使得 $h_{U'} : U' \to V$ 光滑，见《代数空间的态射》引理
\ref{spaces-morphisms-lemma-where-smooth}。设
$\mathcal{U} \subset \mathcal{X}$ 是与 $|U'| \to |\mathcal{X}|$
的像对应的开子叠（《叠的性质》诸引理
\ref{stacks-properties-lemma-topology-points} 和
\ref{stacks-properties-lemma-open-substacks}）。由引理
\ref{stacks-morphisms-lemma-local-source-target} 中的等价性，可知
$f|_\mathcal{U} : \mathcal{U} \to \mathcal{Y}$ 光滑，而且
$\mathcal{U}$ 是具有该性质的最大开子叠。

\medskip\noindent
断言 (1) 由光滑性在源上光滑局部这一事实得出（代数叠的光滑态射正是利用该性质
定义的）。断言 (2) 和 (3) 由代数空间的情形得出，见《代数空间的态射》引理
\ref{spaces-morphisms-lemma-where-smooth}。
\end{proof}

\begin{lemma}
\label{stacks-morphisms-lemma-smooth-quotient-stack}
设 $X \to Y$ 是代数空间的光滑态射。设 $G$ 是 $Y$ 上的群代数空间，
它在 $Y$ 上平坦且局部有限表现。设 $G$ 在 $Y$ 上作用于 $X$。
则商叠 $[X/G]$ 在 $Y$ 上光滑。
\end{lemma}

\noindent
即使 $G$ 在 $Y$ 上不光滑，这也成立！

\begin{proof}
由《可表示性判据》定理
\ref{criteria-theorem-flat-groupoid-gives-algebraic-stack}，
商 $[X/G]$ 是代数叠。$[X/G]$ 在 $Y$ 上的光滑性由光滑性可沿 fppf 覆盖
下降这一事实得出：选取满光滑态射 $U \to [X/G]$，其中 $U$ 是概形。
$[X/G]$ 在 $Y$ 上光滑，等价于 $U$ 在 $Y$ 上光滑。注意，
$U \times_{[X/G]} X$ 在 $X$ 上光滑，因而也在 $Y$ 上光滑
（因为光滑态射的复合仍光滑）。另一方面，
$U \times_{[X/G]} X \to U$ 局部有限表现、平坦且满
（因为它是具有这些性质的 $X \to [X/G]$ 的基变换；例如见
《可表示性判据》引理
\ref{criteria-lemma-flat-quotient-flat-presentation}）。因此，可以应用
《代数空间上的下降》引理 \ref{spaces-descent-lemma-smooth-permanence}。
\end{proof}

\begin{lemma}
\label{stacks-morphisms-lemma-gerbe-smooth}
设 $\pi : \mathcal{X} \to \mathcal{Y}$ 是代数叠的态射。若
$\mathcal{X}$ 是 $\mathcal{Y}$ 上的胚，则 $\pi$ 满且光滑。
\end{lemma}

\begin{proof}
满性已经在引理 \ref{stacks-morphisms-lemma-gerbe-fppf} 中得到。由引理
\ref{stacks-morphisms-lemma-descent-smooth}，只需沿满、平坦、局部有限表现的态射
$\mathcal{Y}' \to \mathcal{Y}$ 作基变换，用所得基变换替换 $\pi$ 后证明本引理。
由引理 \ref{stacks-morphisms-lemma-local-structure-gerbe}，可以假设
$\mathcal{Y} = U$ 是代数空间，而 $\mathcal{X} = [U/G]$ 是 $U$ 上的叠，
其中 $G \to U$ 平坦且局部有限表现。此时由引理
\ref{stacks-morphisms-lemma-smooth-quotient-stack} 得证。
\end{proof}






\section{在源与靶上 \'etale--光滑局部的态射类型}
\label{stacks-morphisms-section-etale-smooth-local-source-target}

\noindent
给定代数空间态射的一个{\it 在源与靶上 \'etale--光滑局部}的性质
（见《代数空间上的下降》定义
\ref{spaces-descent-definition-etale-smooth-local-source-target}），
可以用它定义代数叠的 DM 态射的相应性质，即要求下述引理的任一等价条件成立。

\begin{lemma}
\label{stacks-morphisms-lemma-etale-smooth-local-source-target}
设 $\mathcal{P}$ 是代数空间态射的一个在源与靶上 \'etale--光滑局部的性质。
设 $f : \mathcal{X} \to \mathcal{Y}$ 是代数叠的 DM 态射。考虑交换图
$$
\xymatrix{
U \ar[d]_a \ar[r]_h & V \ar[d]^b \\
\mathcal{X} \ar[r]^f & \mathcal{Y}
}
$$
，其中 $U$ 和 $V$ 是代数空间，$V \to \mathcal{Y}$ 光滑，而
$U \to \mathcal{X} \times_\mathcal{Y} V$ 是 \'etale 的。下列条件等价：
\begin{enumerate}
\item 对任意上述交换图，态射 $h$ 具有性质 $\mathcal{P}$；
\item 对某个满足 $a : U \to \mathcal{X}$ 为满的上述交换图，态射 $h$
具有性质 $\mathcal{P}$。
\end{enumerate}
若 $\mathcal{X}$ 和 $\mathcal{Y}$ 可由代数空间表示，则这些条件还等价于
$f$（作为代数空间的态射）具有性质 $\mathcal{P}$。若 $\mathcal{P}$
还在任意基变换下保持，并且在基上 fppf 局部，则对可由代数空间表示的态射
$f$，这些条件还等价于 $f$ 在《叠的性质》一节
\ref{stacks-properties-section-properties-morphisms} 的意义下具有性质
$\mathcal{P}$。
\end{lemma}

\begin{proof}
先证明 (1) $\Rightarrow$ (2)。选取代数空间 $V$ 以及满光滑态射
$V \to \mathcal{Y}$。由于 $f$ 是 DM 的，存在概形 $U$ 以及满 \'etale 态射
$U \to V \times_\mathcal{Y} \mathcal{X}$，见引理 \ref{stacks-morphisms-lemma-DM}。
于是可见 (2) 成立。注意，$U \to \mathcal{X}$ 也满且光滑，因为它是基变换
$\mathcal{X} \times_\mathcal{Y} V \to \mathcal{X}$ 与所选映射
$U \to \mathcal{X} \times_\mathcal{Y} V$ 的复合。因此得到了 (1)
中所述的交换图。所以若 (1) 成立，则 $h : U \to V$ 具有性质
$\mathcal{P}$；又因 $U \to \mathcal{X}$ 满，这正说明 (2) 成立。

\medskip\noindent
反之，假设 (2) 成立，并设 $U, V, a, b, h$ 如 (2) 所述。再任取 (1)
中所述的交换图 $U', V', a', b', h'$。图示为
$$
\xymatrix{
U \ar[d] \ar[r]_h & V \ar[d] \\
\mathcal{X} \ar[r]^f & \mathcal{Y}
}
\quad\quad
\xymatrix{
U' \ar[d] \ar[r]_{h'} & V' \ar[d] \\
\mathcal{X} \ar[r]^f & \mathcal{Y}
}
$$
为证明 (2) 蕴含 (1)，必须证明 $h'$ 具有 $\mathcal{P}$。为此，考虑代数空间的
交换图
$$
\xymatrix{
U \ar[d]^h &
U \times_\mathcal{X} U' \ar[l] \ar[d]^{(h, h')} \ar[r] &
U' \ar[d]^{h'} \\
V &
V \times_\mathcal{Y} V' \ar[l] \ar[r] &
V'
}
$$
。注意，水平箭头都是光滑态射
$V \to \mathcal{Y}$、$V' \to \mathcal{Y}$、$U \to \mathcal{X}$
和 $U' \to \mathcal{X}$ 的基变换，因而光滑。还要注意，下列方块
$$
\xymatrix{
U \ar[d] & U \times_\mathcal{X} U' \ar[l] \ar[d] &
U \times_\mathcal{X} U' \ar[d] \ar[r] & U' \ar[d] \\
V \times_\mathcal{Y} \mathcal{X} &
V \times_\mathcal{Y} U' \ar[l] &
U \times_\mathcal{Y} V' \ar[r] &
\mathcal{X} \times_\mathcal{Y} V'
}
$$
是笛卡儿的，故由对 $U', V', a', b', h'$ 和 $U, V, a, b, h$ 的假设，
竖直箭头都是 \'etale 的。由《代数空间上的下降》引理
\ref{spaces-descent-lemma-etale-smooth-local-source-target-implies} 的 (2)，
$\mathcal{P}$ 在靶上光滑局部，因此 $h$ 的基变换
$t : U \times_\mathcal{Y} V' \to V \times_\mathcal{Y} V'$
具有 $\mathcal{P}$。由《代数空间上的下降》引理
\ref{spaces-descent-lemma-etale-smooth-local-source-target-implies} 的 (1)，
$\mathcal{P}$ 在源上 \'etale 局部；
又因 $s : U \times_\mathcal{X} U' \to U \times_\mathcal{Y} V'$ 是
\'etale 的，可知态射 $(h, h') = t \circ s$ 具有 $\mathcal{P}$。考虑交换图
$$
\xymatrix{
U \times_\mathcal{X} U' \ar[r]_{(h, h')} \ar[d] &
V \times_\mathcal{Y} V' \ar[d] \\
U' \ar[r]^{h'} & V'
}
$$
左侧竖直箭头满，右侧竖直箭头光滑，而诱导态射
$$
U \times_\mathcal{X} U'
\longrightarrow
U' \times_{V'} (V \times_\mathcal{Y} V') = V \times_\mathcal{Y} U'
$$
如上所见是 \'etale 的。因此，由《代数空间上的下降》定义
\ref{spaces-descent-definition-etale-smooth-local-source-target} 的 (3)，
可知 $h'$ 具有 $\mathcal{P}$。这就证明了 (1) 与 (2) 的等价性。

\medskip\noindent
若 $\mathcal{X}$ 和 $\mathcal{Y}$ 可表，则《代数空间上的下降》引理
\ref{spaces-descent-lemma-etale-smooth-local-source-target-characterize}
适用，并说明 (1) 和 (2) 等价于 $f$ 具有 $\mathcal{P}$。

\medskip\noindent
最后，假设 $f$ 可表，$U, V, a, b, h$ 如本引理 (2) 所述，并且
$\mathcal{P}$ 在任意基变换下保持。必须证明，对任意概形 $Z$ 和态射
$Z \to \mathcal{X}$，基变换
$Z \times_\mathcal{Y} \mathcal{X} \to Z$
具有性质 $\mathcal{P}$。考虑交换图
$$
\xymatrix{
Z \times_\mathcal{Y} U \ar[d] \ar[r] &
Z \times_\mathcal{Y} V \ar[d] \\
Z \times_\mathcal{Y} \mathcal{X} \ar[r] &
Z
}
$$
注意，顶部水平箭头是 $h$ 的基变换，因而具有性质 $\mathcal{P}$。
左侧竖直箭头满，诱导态射
$$
Z \times_\mathcal{Y} U \longrightarrow
(Z \times_\mathcal{Y} \mathcal{X}) \times_Z (Z \times_\mathcal{Y} V)
$$
是 \'etale 的，而右侧竖直箭头光滑。因此，《代数空间上的下降》引理
\ref{spaces-descent-lemma-etale-smooth-local-source-target-characterize}
适用，并说明 $Z \times_\mathcal{Y} \mathcal{X} \to Z$ 具有性质
$\mathcal{P}$。
\end{proof}

\begin{definition}
\label{stacks-morphisms-definition-etale-smooth-P}
设 $\mathcal{P}$ 是代数空间态射的一个在源与靶上 \'etale--光滑局部的性质。
若引理 \ref{stacks-morphisms-lemma-local-source-target} 的等价条件成立，则称代数叠的 DM 态射
$f : \mathcal{X} \to \mathcal{Y}$ {\it 具有性质 $\mathcal{P}$}。
\end{definition}

\begin{remark}
\label{stacks-morphisms-remark-etale-smooth-composition}
设 $\mathcal{P}$ 是代数空间态射的一个在源与靶上 \'etale--光滑局部并且
在复合下稳定的性质。则定义 \ref{stacks-morphisms-definition-etale-smooth-P} 所定义的
代数叠 DM 态射的相应性质在复合下稳定。事实上，设
$f : \mathcal{X} \to \mathcal{Y}$ 和 $g : \mathcal{Y} \to \mathcal{Z}$
是具有性质 $\mathcal{P}$ 的代数叠 DM 态射。由引理
\ref{stacks-morphisms-lemma-composition-separated}，复合 $g \circ f$ 是 DM 的。选取代数空间
$W$ 以及满光滑态射 $W \to \mathcal{Z}$。选取代数空间 $V$ 以及满
\'etale 态射 $V \to \mathcal{Y} \times_\mathcal{Z} W$
（引理 \ref{stacks-morphisms-lemma-DM}）。再选取代数空间 $U$ 以及满 \'etale 态射
$U \to \mathcal{X} \times_\mathcal{Y} V$。按定义，态射
$V \to W$ 和 $U \to V$ 具有性质 $\mathcal{P}$。由于假设
$\mathcal{P}$ 在复合下稳定，故 $U \to W$ 具有性质 $\mathcal{P}$。
再次由定义可知，$g \circ f : \mathcal{X} \to \mathcal{Z}$ 具有性质
$\mathcal{P}$。
\end{remark}

\begin{remark}
\label{stacks-morphisms-remark-etale-smooth-base-change}
设 $\mathcal{P}$ 是代数空间态射的一个在源与靶上 \'etale--光滑局部并且
在基变换下稳定的性质。则定义 \ref{stacks-morphisms-definition-etale-smooth-P} 所定义的
代数叠 DM 态射的相应性质在任意基变换下稳定。事实上，设
$f : \mathcal{X} \to \mathcal{Y}$ 是代数叠的 DM 态射，而
$g : \mathcal{Y}' \to \mathcal{Y}$ 是代数叠的态射，并假设 $f$
具有性质 $\mathcal{P}$。则基变换
$\mathcal{Y}' \times_\mathcal{Y} \mathcal{X} \to \mathcal{Y}'$
由引理 \ref{stacks-morphisms-lemma-base-change-separated} 是 DM 态射。选取代数空间 $V$
以及满光滑态射 $V \to \mathcal{Y}$。选取代数空间 $U$ 以及满 \'etale 态射
$U \to \mathcal{X} \times_\mathcal{Y} V$（引理 \ref{stacks-morphisms-lemma-DM}）。
最后，选取代数空间 $V'$ 以及满光滑态射
$V' \to \mathcal{Y}' \times_\mathcal{Y} V$。按定义，态射
$U \to V$ 具有性质 $\mathcal{P}$。由于假设 $\mathcal{P}$ 在基变换下稳定，
$V' \times_V U \to V'$ 也具有性质 $\mathcal{P}$。考虑交换图
$$
\xymatrix{
V' \times_V U \ar[r] \ar[d] &
\mathcal{Y}' \times_\mathcal{Y} \mathcal{X} \ar[r] \ar[d] &
\mathcal{X} \ar[d] \\
V' \ar[r] & \mathcal{Y}' \ar[r] & \mathcal{Y}
}
$$
可知左上方水平箭头满，并且
$$
V' \times_V U \to V' \times_\mathcal{Y}
(\mathcal{Y}' \times_{\mathcal{Y}'} \mathcal{X}) =
V' \times_V (\mathcal{X} \times_\mathcal{Y} V)
$$
作为 $U \to \mathcal{X} \times_\mathcal{Y} V$ 的基变换是 \'etale 的。
所以由定义可知，投影
$\mathcal{Y}' \times_\mathcal{Y} \mathcal{X} \to \mathcal{Y}'$
具有性质 $\mathcal{P}$。
\end{remark}

\begin{remark}
\label{stacks-morphisms-remark-etale-smooth-implication}
设 $\mathcal{P}, \mathcal{P}'$ 是代数空间态射的两个在源与靶上
\'etale--光滑局部的性质。假设对代数空间的态射有
$\mathcal{P} \Rightarrow \mathcal{P}'$。则对利用 $\mathcal{P}$ 和
$\mathcal{P}'$ 在定义 \ref{stacks-morphisms-definition-etale-smooth-P} 中定义的代数叠态射的
性质，也有 $\mathcal{P} \Rightarrow \mathcal{P}'$。这由定义显然成立。
\end{remark}








\section{\'Etale 态射}
\label{stacks-morphisms-section-etale}

\noindent
不应把代数叠的 \'etale 态射定义为相对维数 $0$ 的光滑态射。事实上，态射
$$
[\mathbf{A}^1_k/\mathbf{G}_{m, k}] \longrightarrow \Spec(k)
$$
对群概形 $\mathbf{G}_{m, k}$ 在 $\mathbf{A}^1_k$ 上的任意作用都是相对维数
$0$ 的光滑态射。这与通常认为 \'etale 态射应当认同切空间的观念不符。
上例不是拟有限的，但态射
$$
\mathcal{X} = [\Spec(k)/\mu_{p, k}] \longrightarrow \Spec(k)
$$
光滑且拟有限（一节 \ref{stacks-morphisms-section-locally-quasi-finite}）。然而，若 $k$
的特征为 $p > 0$，则可表态射 $\Spec(k) \to \mathcal{X}$ 不是 \'etale 的，
因为其基变换
$\mu_{p, k} = \Spec(k) \times_\mathcal{X} \Spec(k) \to \Spec(k)$
是从非约化概形到域的谱的态射。因此，若把 \'etale 态射定义为光滑且局部
拟有限的态射，则《代数空间的态射》引理
\ref{spaces-morphisms-lemma-etale-permanence} 的类似结论将不成立。

\medskip\noindent
我们的做法则从下列要求出发：“\'etale 性”应当是在基变换下保持的性质；
并且若 $\mathcal{X} \to X$ 是从代数叠到概形的 \'etale 态射，则
$\mathcal{X}$ 应为 Deligne--Mumford 叠。换言之，我们要求 \'etale 态射为
DM 态射，并用一节 \ref{stacks-morphisms-section-etale-smooth-local-source-target} 的内容
定义代数叠的 \'etale 态射。

\medskip\noindent
在引理 \ref{stacks-morphisms-lemma-characterize-etale} 中，我们将把代数叠的 \'etale 态射刻画为
满足下列条件的态射：(a) 局部有限表现，(b) 平坦，且 (c) 对角态射为 \'etale。

\medskip\noindent
代数空间态射“为 \'etale”的性质在源与靶上 \'etale--光滑局部，见
《代数空间上的下降》评注
\ref{spaces-descent-remark-list-etale-smooth-local-source-target}。该性质在
基变换下也稳定，并且在靶上 fpqc 局部，见《代数空间的态射》引理
\ref{spaces-morphisms-lemma-base-change-etale} 和《代数空间上的下降》引理
\ref{spaces-descent-lemma-descending-property-etale}。因此，由上面的引理
\ref{stacks-morphisms-lemma-etale-smooth-local-source-target}，可以如下定义代数空间的态射何时为
\'etale；当该态射可由代数空间表示时，这一定义与《叠的性质》一节
\ref{stacks-properties-section-properties-morphisms} 中已有的概念一致，因为由引理
\ref{stacks-morphisms-lemma-trivial-implications}，这种态射自动是 DM 的。

\begin{definition}
\label{stacks-morphisms-definition-etale}
设 $f : \mathcal{X} \to \mathcal{Y}$ 是代数叠的态射。若 $f$ 是 DM 的，
并且引理 \ref{stacks-morphisms-lemma-etale-smooth-local-source-target} 的等价条件对
$\mathcal{P} = \etale$ 成立，则称 $f$ {\it 为 \'etale}。
\end{definition}

\noindent
下文将不再特别说明：当 $f$ 可由代数空间表示，或 $\mathcal{X}$ 和
$\mathcal{Y}$ 可由代数空间表示时，这一定义与已有的 \'etale 态射概念一致。

\begin{lemma}
\label{stacks-morphisms-lemma-composition-etale}
\'Etale 态射的复合仍为 \'etale。
\end{lemma}

\begin{proof}
结合评注 \ref{stacks-morphisms-remark-etale-smooth-composition} 与《代数空间的态射》引理
\ref{spaces-morphisms-lemma-composition-etale}。
\end{proof}

\begin{lemma}
\label{stacks-morphisms-lemma-base-change-etale}
\'Etale 态射的基变换仍为 \'etale。
\end{lemma}

\begin{proof}
结合评注 \ref{stacks-morphisms-remark-etale-smooth-base-change} 与《代数空间的态射》引理
\ref{spaces-morphisms-lemma-base-change-etale}。
\end{proof}

\begin{lemma}
\label{stacks-morphisms-lemma-open-immersion-etale}
开浸入是 \'etale 的。
\end{lemma}

\begin{proof}
设 $j : \mathcal{U} \to \mathcal{X}$ 是代数叠的开浸入。由于 $j$ 可表，
由引理 \ref{stacks-morphisms-lemma-trivial-implications}，它是 DM 的。另一方面，若
$X \to \mathcal{X}$ 是满光滑态射，其中 $X$ 是概形，则
$U = \mathcal{U} \times_\mathcal{X} X$ 是 $X$ 的开子概形。因此
$U \to X$ 是 \'etale 的（《态射》引理
\ref{morphisms-lemma-open-immersion-etale}），由定义可知 $j$ 是 \'etale 的。
\end{proof}

\begin{lemma}
\label{stacks-morphisms-lemma-etale}
设 $f : \mathcal{X} \to \mathcal{Y}$ 是代数叠的态射。下列条件等价：
\begin{enumerate}
\item $f$ 是 \'etale 的；
\item $f$ 是 DM 的，并且对任意态射 $V \to \mathcal{Y}$（其中 $V$
是代数空间）以及任意 \'etale 态射
$U \to V \times_\mathcal{Y} \mathcal{X}$（其中 $U$ 是代数空间），
态射 $U \to V$ 都是 \'etale 的；
\item 存在满、局部有限表现、平坦的态射 $W \to \mathcal{Y}$（其中 $W$
是代数空间）以及满 \'etale 态射
$T \to W \times_\mathcal{Y} \mathcal{X}$（其中 $T$ 是代数空间），
使得态射 $T \to W$ 是 \'etale 的。
\end{enumerate}
\end{lemma}

\begin{proof}
假设 (1)。则 $f$ 是 DM 的；由于 \'etale 性在基变换下保持，(2) 成立。

\medskip\noindent
假设 (2)。选取概形 $V$ 以及满 \'etale 态射 $V \to \mathcal{Y}$。
选取概形 $U$ 以及满 \'etale 态射
$U \to V \times_\mathcal{Y} \mathcal{X}$（引理 \ref{stacks-morphisms-lemma-DM}）。
于是 (3) 成立。

\medskip\noindent
假设 $W \to \mathcal{Y}$ 和
$T \to W \times_\mathcal{Y} \mathcal{X}$ 如 (3) 所述。先验证 $f$ 是 DM 的。
只需验证 $W \times_\mathcal{Y} \mathcal{X} \to W$ 是 DM 的，见引理
\ref{stacks-morphisms-lemma-check-separated-covering}。由引理
\ref{stacks-morphisms-lemma-compose-after-separated}，只需验证
$W \times_\mathcal{Y} \mathcal{X}$ 是 DM 叠。这由态射
$T \to W \times_\mathcal{Y} \mathcal{X}$ 的存在性以及定理
\ref{stacks-morphisms-theorem-DM} 的容易方向得出。

\medskip\noindent
假设 $f$ 是 DM 的，且 $W \to \mathcal{Y}$ 和
$T \to W \times_\mathcal{Y} \mathcal{X}$ 如 (3) 所述。设 $V$
是代数空间，$V \to \mathcal{Y}$ 满且光滑；设 $U$ 是代数空间，并令
$U \to V \times_\mathcal{Y} \mathcal{X}$ 满且 \'etale
（引理 \ref{stacks-morphisms-lemma-DM}）。必须验证 $U \to V$ 是 \'etale 的。由
《代数空间上的下降》引理
\ref{spaces-descent-lemma-descending-property-etale}，只需证明
$U \times_\mathcal{Y} W \to V \times_\mathcal{Y} W$ 是 \'etale 的。
可以依次用
$\mathcal{X} \times_\mathcal{Y} W, W, W, T, U \times_\mathcal{Y} W,
V \times_\mathcal{Y} W$ 替换
$\mathcal{X}, \mathcal{Y}, W, T, U, V$（略去一个小细节）。因此可以假设
$Y = \mathcal{Y}$ 是代数空间，存在代数空间 $T$ 以及满 \'etale 态射
$T \to \mathcal{X}$，使得 $T \to Y$ 是 \'etale 的，而 $U$ 和 $V$
仍如上所述。在此情形，有
$$
U \to V\text{ is \'etale}
\Leftrightarrow
\mathcal{X} \to Y\text{ is \'etale}
\Leftrightarrow
T \to Y\text{ is \'etale}
$$
，这是引理 \ref{stacks-morphisms-lemma-etale-smooth-local-source-target} 的性质 (1) 与 (2)
的等价性以及定义 \ref{stacks-morphisms-definition-etale} 所给出的。证明完成。
\end{proof}

\begin{lemma}
\label{stacks-morphisms-lemma-etale-permanence}
设 $\mathcal{X}, \mathcal{Y}$ 是在代数叠 $\mathcal{Z}$ 上 \'etale 的
代数叠。任意 $\mathcal{Z}$ 上的态射
$\mathcal{X} \to \mathcal{Y}$ 都是 \'etale 的。
\end{lemma}

\begin{proof}
由引理 \ref{stacks-morphisms-lemma-compose-after-separated}，态射
$\mathcal{X} \to \mathcal{Y}$ 是 DM 的。设 $W \to \mathcal{Z}$
是源为代数空间的满光滑态射。设
$V \to \mathcal{Y} \times_\mathcal{Z} W$ 是源为代数空间的满
\'etale 态射（引理 \ref{stacks-morphisms-lemma-DM}）。设
$U \to \mathcal{X} \times_\mathcal{Y} V$ 是源为代数空间的满
\'etale 态射（引理 \ref{stacks-morphisms-lemma-DM}）。则
$$
U \longrightarrow \mathcal{X} \times_\mathcal{Z} W
$$
是满 \'etale 态射，因为它是
$U \to \mathcal{X} \times_\mathcal{Y} V$ 与
$V \to \mathcal{Y} \times_\mathcal{Z} W$ 沿
$\mathcal{X} \times_\mathcal{Z} W \to \mathcal{Y} \times_\mathcal{Z} W$
所得基变换的复合。因此，只需证明 $U \to W$ 是 \'etale 的。由于
$\mathcal{X} \to \mathcal{Z}$ 和 $\mathcal{Y} \to \mathcal{Z}$
是 \'etale 的，$U \to W$ 和 $V \to W$ 也都是 \'etale 的；所需结论由
代数空间版本的引理，即《代数空间的态射》引理
\ref{spaces-morphisms-lemma-etale-permanence} 得出。
\end{proof}








\section{非分歧态射}
\label{stacks-morphisms-section-unramified}

\noindent
关于我们为何选择这个非分歧态射的定义，请读者参阅 \'etale 态射一节
\ref{stacks-morphisms-section-etale} 中的讨论。

\medskip\noindent
在引理 \ref{stacks-morphisms-lemma-characterize-unramified} 中，我们将把代数叠的非分歧态射
刻画为局部有限型且对角态射为 \'etale 的态射。

\medskip\noindent
代数空间态射“为非分歧”的性质在源与靶上 \'etale--光滑局部，见
《代数空间上的下降》评注
\ref{spaces-descent-remark-list-etale-smooth-local-source-target}。该性质在基变换下
也稳定，并且在靶上 fpqc 局部，见《代数空间的态射》引理
\ref{spaces-morphisms-lemma-base-change-unramified} 和《代数空间上的下降》引理
\ref{spaces-descent-lemma-descending-property-unramified}。因此，由上面的引理
\ref{stacks-morphisms-lemma-etale-smooth-local-source-target}，可以如下定义代数空间的态射何时
非分歧；当该态射可由代数空间表示时，这一定义与《叠的性质》一节
\ref{stacks-properties-section-properties-morphisms} 中已有的概念一致，因为由引理
\ref{stacks-morphisms-lemma-trivial-implications}，这种态射自动是 DM 的。

\begin{definition}
\label{stacks-morphisms-definition-unramified}
设 $f : \mathcal{X} \to \mathcal{Y}$ 是代数叠的态射。若 $f$ 是 DM 的，
并且引理 \ref{stacks-morphisms-lemma-etale-smooth-local-source-target} 的等价条件对
$\mathcal{P} =$“非分歧”成立，则称 $f$ {\it 非分歧}。
\end{definition}

\noindent
下文将不再特别说明：当 $f$ 可由代数空间表示，或 $\mathcal{X}$ 和
$\mathcal{Y}$ 可由代数空间表示时，这一定义与已有的非分歧态射概念一致。

\begin{lemma}
\label{stacks-morphisms-lemma-composition-unramified}
非分歧态射的复合仍非分歧。
\end{lemma}

\begin{proof}
结合评注 \ref{stacks-morphisms-remark-etale-smooth-composition} 与《代数空间的态射》引理
\ref{spaces-morphisms-lemma-composition-unramified}。
\end{proof}

\begin{lemma}
\label{stacks-morphisms-lemma-base-change-unramified}
非分歧态射的基变换仍非分歧。
\end{lemma}

\begin{proof}
结合评注 \ref{stacks-morphisms-remark-etale-smooth-base-change} 与《代数空间的态射》引理
\ref{spaces-morphisms-lemma-base-change-unramified}。
\end{proof}

\begin{lemma}
\label{stacks-morphisms-lemma-etale-unramified}
\'Etale 态射是非分歧的。
\end{lemma}

\begin{proof}
由评注 \ref{stacks-morphisms-remark-etale-smooth-implication} 和《代数空间的态射》引理
\ref{spaces-morphisms-lemma-etale-unramified} 得出。
\end{proof}

\begin{lemma}
\label{stacks-morphisms-lemma-immersion-unramified}
浸入是非分歧的。
\end{lemma}

\begin{proof}
设 $j : \mathcal{Z} \to \mathcal{X}$ 是代数叠的浸入。由于 $j$ 可表，
由引理 \ref{stacks-morphisms-lemma-trivial-implications}，它是 DM 的。另一方面，若
$X \to \mathcal{X}$ 是满光滑态射，其中 $X$ 是概形，则
$Z = \mathcal{Z} \times_\mathcal{X} X$ 是 $X$ 的局部闭子概形。因此
$Z \to X$ 非分歧（《态射》诸引理
\ref{morphisms-lemma-open-immersion-unramified} 和
\ref{morphisms-lemma-closed-immersion-unramified}），由定义可知 $j$ 非分歧。
\end{proof}

\begin{lemma}
\label{stacks-morphisms-lemma-unramified}
设 $f : \mathcal{X} \to \mathcal{Y}$ 是代数叠的态射。下列条件等价：
\begin{enumerate}
\item $f$ 非分歧；
\item $f$ 是 DM 的，并且对任意态射 $V \to \mathcal{Y}$（其中 $V$
是代数空间）以及任意 \'etale 态射
$U \to V \times_\mathcal{Y} \mathcal{X}$（其中 $U$ 是代数空间），
态射 $U \to V$ 都非分歧；
\item 存在满、局部有限表现、平坦的态射 $W \to \mathcal{Y}$（其中 $W$
是代数空间）以及满 \'etale 态射
$T \to W \times_\mathcal{Y} \mathcal{X}$（其中 $T$ 是代数空间），
使得态射 $T \to W$ 非分歧。
\end{enumerate}
\end{lemma}

\begin{proof}
假设 (1)。则 $f$ 是 DM 的；由于非分歧性在基变换下保持，(2) 成立。

\medskip\noindent
假设 (2)。选取概形 $V$ 以及满 \'etale 态射 $V \to \mathcal{Y}$。
选取概形 $U$ 以及满 \'etale 态射
$U \to V \times_\mathcal{Y} \mathcal{X}$（引理 \ref{stacks-morphisms-lemma-DM}）。
于是 (3) 成立。

\medskip\noindent
假设 $W \to \mathcal{Y}$ 和
$T \to W \times_\mathcal{Y} \mathcal{X}$ 如 (3) 所述。先验证 $f$ 是 DM 的。
只需验证 $W \times_\mathcal{Y} \mathcal{X} \to W$ 是 DM 的，见引理
\ref{stacks-morphisms-lemma-check-separated-covering}。由引理
\ref{stacks-morphisms-lemma-compose-after-separated}，只需验证
$W \times_\mathcal{Y} \mathcal{X}$ 是 DM 叠。这由态射
$T \to W \times_\mathcal{Y} \mathcal{X}$ 的存在性以及定理
\ref{stacks-morphisms-theorem-DM} 的容易方向得出。

\medskip\noindent
假设 $f$ 是 DM 的，且 $W \to \mathcal{Y}$ 和
$T \to W \times_\mathcal{Y} \mathcal{X}$ 如 (3) 所述。设 $V$
是代数空间，$V \to \mathcal{Y}$ 满且光滑；设 $U$ 是代数空间，并令
$U \to V \times_\mathcal{Y} \mathcal{X}$ 满且 \'etale
（引理 \ref{stacks-morphisms-lemma-DM}）。必须验证 $U \to V$ 非分歧。由
《代数空间上的下降》引理
\ref{spaces-descent-lemma-descending-property-unramified}，只需证明
$U \times_\mathcal{Y} W \to V \times_\mathcal{Y} W$ 非分歧。可以依次用
$\mathcal{X} \times_\mathcal{Y} W, W, W, T, U \times_\mathcal{Y} W,
V \times_\mathcal{Y} W$ 替换
$\mathcal{X}, \mathcal{Y}, W, T, U, V$（略去一个小细节）。因此可以假设
$Y = \mathcal{Y}$ 是代数空间，存在代数空间 $T$ 以及满 \'etale 态射
$T \to \mathcal{X}$，使得 $T \to Y$ 非分歧，而 $U$ 和 $V$
仍如上所述。在此情形，有
$$
U \to V\text{ is unramified}
\Leftrightarrow
\mathcal{X} \to Y\text{ is unramified}
\Leftrightarrow
T \to Y\text{ is unramified}
$$
，这是引理 \ref{stacks-morphisms-lemma-etale-smooth-local-source-target} 的性质 (1) 与 (2)
的等价性以及定义 \ref{stacks-morphisms-definition-unramified} 所给出的。证明完成。
\end{proof}

\begin{lemma}
\label{stacks-morphisms-lemma-unramified-quasi-finite}
代数叠的非分歧态射局部拟有限。
\end{lemma}

\begin{proof}
这由引理 \ref{stacks-morphisms-lemma-unramified}（非分歧态射的刻画）、引理
\ref{stacks-morphisms-lemma-characterize-locally-quasi-finite}（局部拟有限态射的刻画），
以及《代数空间的态射》引理
\ref{spaces-morphisms-lemma-unramified-quasi-finite}（代数空间的相应结果）得出。
\end{proof}

\begin{lemma}
\label{stacks-morphisms-lemma-permanence-unramified}
设 $\mathcal{X} \to \mathcal{Y} \to \mathcal{Z}$ 是代数叠的态射。
若 $\mathcal{X} \to \mathcal{Z}$ 非分歧，且
$\mathcal{Y} \to \mathcal{Z}$ 是 DM 的，则
$\mathcal{X} \to \mathcal{Y}$ 非分歧。
\end{lemma}

\begin{proof}
假设 $\mathcal{X} \to \mathcal{Z}$ 非分歧。由引理
\ref{stacks-morphisms-lemma-compose-after-separated}，态射
$\mathcal{X} \to \mathcal{Y}$ 是 DM 的。选取交换图
$$
\xymatrix{
U \ar[d] \ar[r] & V \ar[d] \ar[r] & W \ar[d] \\
\mathcal{X} \ar[r] & \mathcal{Y} \ar[r] & \mathcal{Z}
}
$$
，其中 $U, V, W$ 是代数空间，$W \to \mathcal{Z}$ 满且光滑，
$V \to \mathcal{Y} \times_\mathcal{Z} W$ 满且 \'etale，而
$U \to \mathcal{X} \times_\mathcal{Y} V$ 满且 \'etale
（见引理 \ref{stacks-morphisms-lemma-DM}）。于是
$U \to \mathcal{X} \times_\mathcal{Z} W$ 也满且 \'etale。因此已知
$U \to W$ 非分歧，而必须证明 $U \to V$ 非分歧。这由《代数空间的态射》引理
\ref{spaces-morphisms-lemma-permanence-unramified} 得出。
\end{proof}

\begin{lemma}
\label{stacks-morphisms-lemma-characterize-unramified}
设 $f : \mathcal{X} \to \mathcal{Y}$ 是代数叠的态射。下列条件等价：
\begin{enumerate}
\item $f$ 非分歧；
\item $f$ 局部有限型，且其对角态射为 \'etale。
\end{enumerate}
\end{lemma}

\begin{proof}
假设 $f$ 非分歧。则 $f$ 是 DM 的，因而可以选取代数空间 $U$、$V$，
满光滑态射 $V \to \mathcal{Y}$，以及满 \'etale 态射
$U \to \mathcal{X} \times_\mathcal{Y} V$（引理 \ref{stacks-morphisms-lemma-DM}）。
由于 $f$ 非分歧，诱导态射 $U \to V$ 非分歧。因此 $U \to V$ 局部有限型
（《代数空间的态射》引理
\ref{spaces-morphisms-lemma-unramified-locally-finite-type}），从而 $f$
局部有限型。对角态射
$\Delta : \mathcal{X} \to \mathcal{X} \times_\mathcal{Y} \mathcal{X}$
是 $\mathcal{Y}$ 上代数叠的态射。$\Delta$ 沿满光滑态射
$V \to \mathcal{Y}$ 的基变换，是 $f$ 的基变换，即
$\mathcal{X}_V = \mathcal{X} \times_\mathcal{Y} V \to V$ 的对角态射。
换言之，图
$$
\xymatrix{
\mathcal{X}_V \ar[r] \ar[d] &  \mathcal{X}_V \times_V \mathcal{X}_V \ar[d] \\
\mathcal{X} \ar[r] & \mathcal{X} \times_\mathcal{Y} \mathcal{X}
}
$$
是笛卡儿的。由于右侧竖直箭头满且光滑，由《叠的性质》引理
\ref{stacks-properties-lemma-check-property-weak-covering}，只需证明顶部
水平箭头为 \'etale。考虑交换图
$$
\xymatrix{
U \ar[d] \ar[r] & U \times_V U \ar[d] \\
\mathcal{X}_V \ar[r] & \mathcal{X}_V \times_V \mathcal{X}_V
}
$$
其中所有箭头都可由代数空间表示，竖直箭头为 \'etale，左侧竖直箭头满，
而由《代数空间的态射》引理
\ref{spaces-morphisms-lemma-diagonal-unramified-morphism}，顶部水平箭头是开浸入。
这蕴含所需结论：首先，作为 \'etale 态射的复合，
$U \to \mathcal{X}_V \times_V \mathcal{X}_V$ 是 \'etale 的；然后应用
《叠的性质》引理
\ref{stacks-properties-lemma-check-property-after-precomposing} 可知
$\mathcal{X}_V \to \mathcal{X}_V \times_V \mathcal{X}_V$ 是 \'etale 的，
因为对代数空间的态射而言，\'etale 性在源上对 \'etale 拓扑局部
（《代数空间上的下降》引理
\ref{spaces-descent-lemma-etale-etale-local-source}）。

\medskip\noindent
假设 $f$ 局部有限型且其对角态射为 \'etale。由于代数空间的 \'etale 态射
非分歧，按定义 $f$ 是 DM 的。因此，像上文一样，可以选取代数空间
$U$、$V$，满光滑态射 $V \to \mathcal{Y}$，以及满 \'etale 态射
$U \to \mathcal{X} \times_\mathcal{Y} V$（引理 \ref{stacks-morphisms-lemma-DM}）。
为完成证明，必须证明 $U \to V$ 非分歧。已经知道 $U \to V$ 局部有限型。
像上文一样，得到交换图
$$
\xymatrix{
U \ar[d] \ar[r] & U \times_V U \ar[d] \\
\mathcal{X}_V \ar[r] & \mathcal{X}_V \times_V \mathcal{X}_V
}
$$
，其中所有箭头都可由代数空间表示，竖直箭头为 \'etale，而底部水平箭头作为
$\Delta$ 的基变换也是 \'etale 的。例如由引理
\ref{stacks-morphisms-lemma-etale-permanence}\footnote{也很容易直接从《代数空间的态射》引理
\ref{spaces-morphisms-lemma-etale-permanence} 推出这一点。}，可知
$U \to U \times_V U$ 是 \'etale 的。因此，
$U \to U \times_V U$ 是 \'etale 单态射，故为开浸入
（《代数空间的态射》引理
\ref{spaces-morphisms-lemma-etale-universally-injective-open}）。于是由
《代数空间的态射》引理
\ref{spaces-morphisms-lemma-diagonal-unramified-morphism}，
$U \to V$ 非分歧。
\end{proof}

\begin{lemma}
\label{stacks-morphisms-lemma-characterize-etale}
设 $f : \mathcal{X} \to \mathcal{Y}$ 是代数叠的态射。下列条件等价：
\begin{enumerate}
\item $f$ 是 \'etale 的；
\item $f$ 局部有限表现、平坦且非分歧；
\item $f$ 局部有限表现、平坦，且其对角态射为 \'etale。
\end{enumerate}
\end{lemma}

\begin{proof}
(2) 与 (3) 的等价性立即由引理 \ref{stacks-morphisms-lemma-characterize-unramified} 得出。
所以在每种情形下，态射 $f$ 都是 DM 的。于是可以选取代数空间 $U$、$V$，
满光滑态射 $V \to \mathcal{Y}$，以及满 \'etale 态射
$U \to \mathcal{X} \times_\mathcal{Y} V$（引理 \ref{stacks-morphisms-lemma-DM}）。
为完成证明，必须证明 $U \to V$ 是 \'etale 的，当且仅当它局部有限表现、
平坦且非分歧。这由《代数空间的态射》引理
\ref{spaces-morphisms-lemma-unramified-flat-lfp-etale} 得出
（以及更为平凡的《代数空间的态射》诸引理
\ref{spaces-morphisms-lemma-etale-unramified}、
\ref{spaces-morphisms-lemma-etale-locally-finite-presentation} 和
\ref{spaces-morphisms-lemma-etale-flat}）。
\end{proof}








\section{固有态射}
\label{stacks-morphisms-section-proper}

\noindent
固有态射的概念在代数几何中起着重要作用。下面定义代数叠的固有态射。

\begin{definition}
\label{stacks-morphisms-definition-proper}
设 $f : \mathcal{X} \to \mathcal{Y}$ 是代数叠的态射。若 $f$ 分离、
有限型且泛闭，则称 $f$ {\it 固有}。
\end{definition}

\noindent
这与已有的代数空间固有态射概念并不冲突：代数空间的态射固有，当且仅当它
分离、有限型且泛闭（《代数空间的态射》定义
\ref{spaces-morphisms-definition-proper}）；而我们已经在引理
\ref{stacks-morphisms-lemma-representable-separated-diagonal-closed}、一节
\ref{stacks-morphisms-section-finite-type} 和引理
\ref{stacks-morphisms-lemma-characterize-representable-universally-closed} 中验证了这些概念的
相容性。类似地，若 $f : \mathcal{X} \to \mathcal{Y}$ 是可由代数空间表示的
代数叠态射，则《叠的性质》一节
\ref{stacks-properties-section-properties-morphisms} 已定义 $f$ 何时固有。
然而，该节的讨论说明，这等价于要求 $f$ 分离、有限型且泛闭，而上引同一批结果
给出了相容性。

\begin{lemma}
\label{stacks-morphisms-lemma-base-change-proper}
固有态射的基变换仍固有。
\end{lemma}

\begin{proof}
见引理 \ref{stacks-morphisms-lemma-base-change-separated}、
\ref{stacks-morphisms-lemma-base-change-finite-type} 和
\ref{stacks-morphisms-lemma-base-change-universally-closed}。
\end{proof}

\begin{lemma}
\label{stacks-morphisms-lemma-composition-proper}
固有态射的复合仍固有。
\end{lemma}

\begin{proof}
见引理 \ref{stacks-morphisms-lemma-composition-separated}、
\ref{stacks-morphisms-lemma-composition-finite-type} 和
\ref{stacks-morphisms-lemma-composition-universally-closed}。
\end{proof}

\begin{lemma}
\label{stacks-morphisms-lemma-closed-immersion-proper}
代数叠的闭浸入是代数叠的固有态射。
\end{lemma}

\begin{proof}
按定义，闭浸入可表（《叠的性质》定义
\ref{stacks-properties-definition-immersion}）。因此，结论由《叠的性质》一节
\ref{stacks-properties-section-properties-morphisms} 的讨论以及代数空间态射的
相应结果得出，见《代数空间的态射》引理
\ref{spaces-morphisms-lemma-closed-immersion-proper}。
\end{proof}

\begin{lemma}
\label{stacks-morphisms-lemma-universally-closed-permanence}
考虑代数叠的交换图
$$
\xymatrix{
\mathcal{X} \ar[rr] \ar[rd] & &
\mathcal{Y} \ar[ld] \\
& \mathcal{Z} &
}
$$
\begin{enumerate}
\item 若 $\mathcal{X} \to \mathcal{Z}$ 泛闭，而
$\mathcal{Y} \to \mathcal{Z}$ 分离，则态射
$\mathcal{X} \to \mathcal{Y}$ 泛闭。特别地，$|\mathcal{X}|$
在 $|\mathcal{Y}|$ 中的像是闭的。
\item 若 $\mathcal{X} \to \mathcal{Z}$ 固有，而
$\mathcal{Y} \to \mathcal{Z}$ 分离，则态射
$\mathcal{X} \to \mathcal{Y}$ 固有。
\end{enumerate}
\end{lemma}

\begin{proof}
假设 $\mathcal{X} \to \mathcal{Z}$ 泛闭，且
$\mathcal{Y} \to \mathcal{Z}$ 分离。把该态射分解为
$\mathcal{X} \to \mathcal{X} \times_\mathcal{Z} \mathcal{Y} \to \mathcal{Y}$.
第一个态射固有（引理 \ref{stacks-morphisms-lemma-semi-diagonal}），因而泛闭。投影
$\mathcal{X} \times_\mathcal{Z} \mathcal{Y} \to \mathcal{Y}$
是泛闭态射的基变换，故泛闭，见引理
\ref{stacks-morphisms-lemma-base-change-universally-closed}。所以，作为泛闭态射的复合，
$\mathcal{X} \to \mathcal{Y}$ 泛闭，见引理
\ref{stacks-morphisms-lemma-composition-universally-closed}。这证明了 (1)。为推出 (2)，
把 (1) 与引理 \ref{stacks-morphisms-lemma-compose-after-separated}、
\ref{stacks-morphisms-lemma-quasi-compact-permanence} 和
\ref{stacks-morphisms-lemma-finite-type-permanence} 结合即可。
\end{proof}

\begin{lemma}
\label{stacks-morphisms-lemma-image-proper-is-proper}
设 $\mathcal{Z}$ 是代数叠。设 $f : \mathcal{X} \to \mathcal{Y}$
是 $\mathcal{Z}$ 上代数叠的态射。若 $\mathcal{X}$ 在 $\mathcal{Z}$
上泛闭，且 $f$ 满，则 $\mathcal{Y}$ 在 $\mathcal{Z}$ 上泛闭。
特别地，若 $\mathcal{Y}$ 还在 $\mathcal{Z}$ 上分离且有限型，则
$\mathcal{Y}$ 在 $\mathcal{Z}$ 上固有。
\end{lemma}

\begin{proof}
假设 $\mathcal{X}$ 泛闭且 $f$ 满。记结构态射为
$p : \mathcal{X} \to \mathcal{Z}$ 和
$q : \mathcal{Y} \to \mathcal{Z}$。设
$\mathcal{Z}' \to \mathcal{Z}$ 是代数叠的态射。$f$ 沿
$\mathcal{Z}' \to \mathcal{Z}$ 的基变换
$f' : \mathcal{X}' \to \mathcal{Y}'$ 满
（《叠的性质》引理
\ref{stacks-properties-lemma-base-change-surjective}），而 $p$ 的基变换
$p' : \mathcal{X}' \to \mathcal{Z}'$ 是闭映射。若
$T \subset |\mathcal{Y}'|$ 闭，则
$(f')^{-1}(T) \subset |\mathcal{X}'|$ 闭，因而
$p'((f')^{-1}(T)) = q'(T)$ 闭。所以 $q'$ 是闭映射。
\end{proof}








\section{概形论像}
\label{stacks-morphisms-section-scheme-theoretic-image}

\noindent
定义如下。

\begin{definition}
\label{stacks-morphisms-definition-scheme-theoretic-image}
设 $f : \mathcal{X} \to \mathcal{Y}$ 是代数叠的态射。$f$ 的
{\it 概形论像}是使 $f$ 经由其分解的最小闭子叠
$\mathcal{Z} \subset \mathcal{Y}$\footnote{我们将在引理
\ref{stacks-morphisms-lemma-scheme-theoretic-image-existence} 中看到，概形论像总是存在。}。
\end{definition}

\noindent
我们常也用 $f : \mathcal{X} \to \mathcal{Z}$ 表示 $f$ 的这一分解。
若态射 $f$ 不拟紧，则概形论像的构造通常不与到 $\mathcal{Y}$ 的开子叠的
限制交换。然而，若 $f$ 拟紧，则概形论像与平坦基变换交换
（引理 \ref{stacks-morphisms-lemma-existence-plus-flat-base-change}）。

\begin{lemma}
\label{stacks-morphisms-lemma-cover-upstairs}
设 $f : \mathcal{X} \to \mathcal{Y}$ 是代数叠的态射。设
$g : \mathcal{W} \to \mathcal{X}$ 是代数叠的满、平坦、局部有限表现态射。
则 $f$ 的概形论像存在，当且仅当 $f \circ g$ 的概形论像存在；若存在，
二者相同。
\end{lemma}

\begin{proof}
假设 $\mathcal{Z} \subset \mathcal{Y}$ 是闭子叠，且 $f \circ g$
经由 $\mathcal{Z}$ 分解。为证明本引理，只需证明 $f$ 经由
$\mathcal{Z}$ 分解。考虑概形 $T$ 以及态射 $T \to \mathcal{X}$，
它由 $\mathcal{X}$ 在 $T$ 上的纤维范畴中的对象 $x$ 给出。我们要证明，
$f(x)$ 事实上属于 $\mathcal{Z}$ 在 $T$ 上的纤维范畴。投影
$T \times_\mathcal{X} \mathcal{W} \to T$ 满、平坦且局部有限表现。
因此存在 fppf 覆盖 $\{T_i \to T\}$，使得对所有 $i$，$T_i \to T$
都经由 $T \times_\mathcal{X} \mathcal{W} \to T$ 分解。于是
$T_i \to \mathcal{X}$ 经由 $\mathcal{W}$ 分解，进而
$T_i \to \mathcal{Y}$ 经由 $\mathcal{Z}$ 分解。所以
$f(x)|_{T_i}$ 是 $\mathcal{Z}$ 的对象。由于 $\mathcal{Z}$ 是严格满子叠，
可知 $f(x)$ 是 $\mathcal{Z}$ 的对象，正合所需。
\end{proof}

\begin{lemma}
\label{stacks-morphisms-lemma-scheme-theoretic-image-existence}
设 $f : \mathcal{Y} \to \mathcal{X}$ 是代数叠的态射。则 $f$
的概形论像存在。
\end{lemma}

\begin{proof}
选取概形 $V$ 以及满光滑态射 $V \to \mathcal{Y}$。由引理
\ref{stacks-morphisms-lemma-cover-upstairs}，可以用 $V$ 替换 $\mathcal{Y}$。因此，只需证明：
若 $X \to \mathcal{X}$ 是从概形到代数叠的态射，则其概形论像存在。
选取概形 $U$ 以及满光滑态射 $U \to \mathcal{X}$。令
$R = U \times_\mathcal{X} U$。由《代数叠》引理
\ref{algebraic-lemma-stack-presentation}，有 $\mathcal{X} = [U/R]$。
由《叠的性质》引理
\ref{stacks-properties-lemma-substacks-presentation}，$\mathcal{X}$
的闭子叠 $\mathcal{Z}$ 与 $R$-不变闭子概形 $Z \subset U$ 形成
$1$ 对 $1$ 的对应。
设 $Z_1 \subset U$ 是 $X \times_\mathcal{X} U \to U$ 的概形论像。
注意，$X \to \mathcal{X}$ 经由 $\mathcal{Z}$ 分解，当且仅当
$X \times_\mathcal{X} U \to U$ 经由相应的 $R$-不变闭子概形 $Z$ 分解
（细节略；提示：这是因为 $X \times_\mathcal{X} U \to X$ 满且光滑）。
因此，必须证明存在包含 $Z_1$ 的最小 $R$-不变闭子概形 $Z \subset U$。

\medskip\noindent
设 $\mathcal{I}_1 \subset \mathcal{O}_U$ 是与闭子概形
$Z_1 \subset U$ 对应的拟凝聚理想层。设 $Z_\alpha$（$\alpha \in A$）
为 $U$ 的所有包含 $Z_1$ 的 $R$-不变闭子概形所成的集合。对 $\alpha \in A$，
设 $\mathcal{I}_\alpha \subset \mathcal{O}_U$ 是与闭子概形
$Z_\alpha \subset U$ 对应的拟凝聚理想层。包含关系
$Z_1 \subset Z_\alpha$ 意味着
$\mathcal{I}_\alpha \subset \mathcal{I}_1$。$Z_\alpha$ 的 $R$-不变性意味着
$$
s^{-1}\mathcal{I}_\alpha \cdot \mathcal{O}_R =
t^{-1}\mathcal{I}_\alpha \cdot \mathcal{O}_R
$$
作为（代数空间）$R$ 上的（拟凝聚）理想层相等。考虑像
$$
\mathcal{I} =
\Im\left(
\bigoplus\nolimits_{\alpha \in A} \mathcal{I}_\alpha \to \mathcal{I}_1
\right) =
\Im\left(
\bigoplus\nolimits_{\alpha \in A} \mathcal{I}_\alpha \to \mathcal{O}_X
\right)
$$
由于拟凝聚层的直和仍拟凝聚，拟凝聚层之间映射的像也拟凝聚，可知
$\mathcal{I}$ 拟凝聚。由于拉回是正合的且与直和交换，可知
$$
s^{-1}\mathcal{I} \cdot \mathcal{O}_R =
t^{-1}\mathcal{I} \cdot \mathcal{O}_R
$$
因此，$\mathcal{I}$ 定义一个 $R$-不变闭子概形 $Z \subset U$，
它包含于每个 $Z_\alpha$ 且包含 $Z_1$，正合所需。
\end{proof}

\begin{lemma}
\label{stacks-morphisms-lemma-factor-factor}
设
$$
\xymatrix{
\mathcal{X}_1 \ar[d] \ar[r]_{f_1} & \mathcal{Y}_1 \ar[d] \\
\mathcal{X}_2 \ar[r]^{f_2} & \mathcal{Y}_2
}
$$
是代数叠的交换图。设 $\mathcal{Z}_i \subset \mathcal{Y}_i$（$i = 1, 2$）
是 $f_i$ 的概形论像。则态射
$\mathcal{Y}_1 \to \mathcal{Y}_2$ 诱导态射
$\mathcal{Z}_1 \to \mathcal{Z}_2$ 以及交换图
$$
\xymatrix{
\mathcal{X}_1 \ar[r] \ar[d] &
\mathcal{Z}_1 \ar[d] \ar[r] &
\mathcal{Y}_1 \ar[d] \\
\mathcal{X}_2 \ar[r] &
\mathcal{Z}_2 \ar[r] &
\mathcal{Y}_2
}
$$
\end{lemma}

\begin{proof}
$\mathcal{Z}_2$ 在 $\mathcal{Y}_1$ 中的概形论逆像是
$\mathcal{Y}_1$ 的闭子叠，$f_1$ 经由它分解。因此
$\mathcal{Z}_1$ 包含于其中。本引理得证。
\end{proof}

\begin{lemma}
\label{stacks-morphisms-lemma-existence-plus-flat-base-change}
设 $f : \mathcal{X} \to \mathcal{Y}$ 是代数叠的拟紧态射。
则概形论像的构造与平坦基变换交换。
\end{lemma}

\begin{proof}
设 $\mathcal{Y}' \to \mathcal{Y}$ 是代数叠的平坦态射。选取概形 $V$
以及满光滑态射 $V \to \mathcal{Y}$。选取概形 $V'$ 以及满光滑态射
$V' \to \mathcal{Y}' \times_\mathcal{Y} V$。可以并且确实假设
$V = \coprod_{i \in I} V_i$ 是仿射概形的不交并，并且
$V' = \coprod_{i \in I} \coprod_{j \in J_i} V_{i, j}$
是仿射概形的不交并，其中每个 $V_{i, j}$ 都映入 $V_i$。设
\begin{enumerate}
\item $\mathcal{Z} \subset \mathcal{Y}$ 是 $f$ 的概形论像；
\item $\mathcal{Z}' \subset \mathcal{Y}'$ 是 $f$ 沿
$\mathcal{Y}' \to \mathcal{Y}$ 所得基变换的概形论像；
\item $Z \subset V$ 是 $f$ 沿 $V \to \mathcal{Y}$ 所得基变换的概形论像；
\item $Z' \subset V'$ 是 $f$ 沿 $V' \to \mathcal{Y}$ 所得基变换的概形论像。
\end{enumerate}
若能证明
(a) $Z = V \times_\mathcal{Y} \mathcal{Z}$,
(b) $Z' = V' \times_{\mathcal{Y}'} \mathcal{Z}'$，以及
(c) $Z' = V' \times_V Z$
，则本引理成立：包含态射
$\mathcal{Z}' \to \mathcal{Z} \times_\mathcal{Y} \mathcal{Y}'$
（引理 \ref{stacks-morphisms-lemma-factor-factor}）必为同构，因为沿满光滑态射
$V' \to \mathcal{Y}'$ 作基变换后它是同构。

\medskip\noindent
证明 (a)。令 $R = V \times_\mathcal{Y} V$。由《叠的性质》引理
\ref{stacks-properties-lemma-substacks-presentation}，规则
$\mathcal{Z} \mapsto \mathcal{Z} \times_\mathcal{Y} V$
给出 $\mathcal{Y}$ 的闭子叠与 $V$ 的 $R$-不变闭子空间之间的
$1$ 对 $1$ 对应。
此外，$f : \mathcal{X} \to \mathcal{Y}$ 经由 $\mathcal{Z}$ 分解，
当且仅当其基变换 $g : \mathcal{X} \times_\mathcal{Y} V \to V$
经由 $\mathcal{Z} \times_\mathcal{Y} V$ 分解。我们断言：$g$
的概形论像 $Z \subset V$ 是 $R$-不变的。由刚才所述，该断言蕴含 (a)。

\medskip\noindent
对每个 $i$，态射 $\mathcal{X} \times_\mathcal{Y} V_i \to V_i$ 拟紧，
因而 $\mathcal{X} \times_\mathcal{Y} V_i$ 拟紧。所以可以选取仿射概形
$W_i$ 以及满光滑态射
$W_i \to \mathcal{X} \times_\mathcal{Y} V_i$。注意，
$W = \coprod W_i$ 是概形，带有满光滑态射
$W \to \mathcal{X} \times_\mathcal{Y} V$，并且与 $g$ 复合所得
$W \to V$ 拟紧。设 $Z \to V$ 是 $W \to V$ 的概形论像，见《态射》一节
\ref{morphisms-section-scheme-theoretic-image} 和《代数空间的态射》一节
\ref{spaces-morphisms-section-scheme-theoretic-image}。由引理
\ref{stacks-morphisms-lemma-cover-upstairs}，$Z \subset V$ 是 $g$ 的概形论像。为证明
$Z$ 是 $R$-不变的，我们断言二者
$$
\text{pr}_0^{-1}(Z), \text{pr}_1^{-1}(Z) \subset R = V \times_\mathcal{Y} V
$$
都是 $\mathcal{X} \times_\mathcal{Y} R \to R$ 的概形论像。事实上，
先用《代数空间的态射》引理
\ref{spaces-morphisms-lemma-flat-base-change-scheme-theoretic-image}
可知，$\text{pr}_0^{-1}(Z)$ 是复合
$$
W \times_{V, \text{pr}_0} R = W \times_\mathcal{Y} V \to
\mathcal{X} \times_\mathcal{Y} R \to R
$$
的概形论像。由于这里第一个箭头满且光滑，可知
$\text{pr}_0^{-1}(Z)$ 是
$\mathcal{X} \times_\mathcal{Y} R \to R$ 的概形论像。同一论证适用于
$\text{pr}_1^{-1}(Z)$。所以 $Z$ 是 $R$-不变的。

\medskip\noindent
陈述 (b) 的证明与 (a) 完全相同。

\medskip\noindent
证明 (c)。设 $Z_i \subset V_i$ 是
$\mathcal{X} \times_\mathcal{Y} V_i \to V_i$ 的概形论像，并设
$Z_{i, j} \subset V_{i, j}$ 是
$\mathcal{X} \times_\mathcal{Y} V_{i, j} \to V_{i, j}$ 的概形论像。
显然，只需证明 $Z_i$ 在 $V_{i, j}$ 中的逆像是 $Z_{i, j}$。上文已经看到，
$Z_i$ 是 $W_i \to V_i$ 的概形论像；同理，$Z_{i, j}$ 是
$W_i \times_{V_i} V_{i, j} \to V_{i, j}$ 的概形论像。因此，该等式由
概形的情形（《态射》引理
\ref{morphisms-lemma-flat-base-change-scheme-theoretic-image}）以及
$V_{i, j} \to V_i$ 平坦这一事实得出。
\end{proof}

\begin{lemma}
\label{stacks-morphisms-lemma-topology-scheme-theoretic-image}
设 $f : \mathcal{X} \to \mathcal{Y}$ 是代数叠的拟紧态射。设
$\mathcal{Z} \subset \mathcal{Y}$ 是 $f$ 的概形论像。则
$|\mathcal{Z}|$ 是 $|f|$ 的像的闭包。
\end{lemma}

\begin{proof}
设 $z \in |\mathcal{Z}|$ 是一点。选取仿射概形 $V$、点 $v \in V$
以及把 $v$ 映到 $z$ 的光滑态射 $V \to \mathcal{Y}$。则
$\mathcal{X} \times_\mathcal{Y} V$ 是拟紧代数叠。因此，可以找到仿射概形
$W$ 以及满光滑态射 $W \to \mathcal{X} \times_\mathcal{Y} V$。
由引理 \ref{stacks-morphisms-lemma-existence-plus-flat-base-change}，
$\mathcal{X} \times_\mathcal{Y} V \to V$ 的概形论像为
$Z = \mathcal{Z} \times_\mathcal{Y} V$。所以由《叠的性质》引理
\ref{stacks-properties-lemma-points-cartesian}，$|\mathcal{Z}|$
在 $|V|$ 中的逆像是 $|Z|$。由引理 \ref{stacks-morphisms-lemma-cover-upstairs}，
$Z$ 是 $W \to V$ 的概形论像。由《代数空间的态射》引理
\ref{spaces-morphisms-lemma-quasi-compact-scheme-theoretic-image}，
$|W| \to |Z|$ 的像稠密。因此，
$|\mathcal{X} \times_\mathcal{Y} V| \to |Z|$ 的像稠密。注意
$v \in Z$。由于 $|V| \to |\mathcal{Y}|$ 是开映射，一个拓扑论证说明
$z$ 位于 $|f|$ 的像的闭包中，正合所需。
\end{proof}

\begin{lemma}
\label{stacks-morphisms-lemma-scheme-theoretic-image-of-partial-section}
设 $f : \mathcal{X} \to \mathcal{Y}$ 是可由代数空间表示且分离的代数叠态射。
设 $\mathcal{V} \subset \mathcal{Y}$ 是开子叠，且
$\mathcal{V} \to \mathcal{Y}$ 拟紧。设
$s : \mathcal{V} \to \mathcal{X}$ 是态射，满足
$f \circ s = \text{id}_\mathcal{V}$。设 $\mathcal{Y}'$ 是 $s$
的概形论像。则 $\mathcal{Y}' \to \mathcal{Y}$ 在
$\mathcal{V}$ 上是同构。
\end{lemma}

\begin{proof}
由引理 \ref{stacks-morphisms-lemma-quasi-compact-permanence}，态射
$s : \mathcal{V} \to \mathcal{Y}$ 拟紧。因此，由引理
\ref{stacks-morphisms-lemma-existence-plus-flat-base-change}，$s$ 的概形论像
$\mathcal{Y}'$ 的构造与平坦基变换交换。为证明本引理，可以假设
$\mathcal{Y}$ 可由代数空间表示，从而归结到代数空间的情形，即
《代数空间的态射》引理
\ref{spaces-morphisms-lemma-scheme-theoretic-image-of-partial-section}。
\end{proof}









\section{赋值判据}
\label{stacks-morphisms-section-valuative}

\noindent
陈述赋值判据时必须谨慎。确切地说，陈述中需要谈论交换图，但我们是在
$2$-范畴中工作，还必须确保 $2$-态射也按正确方式复合！

\begin{definition}
\label{stacks-morphisms-definition-fill-in-diagram}
设 $f : \mathcal{X} \to \mathcal{Y}$ 是代数叠的态射。考虑实线部分
$2$-交换的图
\begin{equation}
\label{stacks-morphisms-equation-diagram}
\vcenter{
\xymatrix{
\Spec(K) \ar[r]_-x \ar[d]_j & \mathcal{X} \ar[d]^f \\
\Spec(A) \ar[r]^-y \ar@{..>}[ru] & \mathcal{Y}
}
}
\end{equation}
，其中 $A$ 是分式域为 $K$ 的赋值环。设
$$
\gamma : y \circ j \longrightarrow f \circ x
$$
是见证该图 $2$-交换的 $2$-态射
（记号同《范畴》诸节 \ref{categories-section-formal-cat-cat} 和
\ref{categories-section-2-categories}）。给定
(\ref{stacks-morphisms-equation-diagram}) 与 $\gamma$，一个{\it 虚线箭头}是三元组
$(a, \alpha, \beta)$，由态射 $a : \Spec(A) \to \mathcal{X}$
及 $2$-箭头 $\alpha : a \circ j \to x$、$\beta : y \to f \circ a$
组成，并且满足
$\gamma = (\text{id}_f \star \alpha) \circ (\beta \star \text{id}_j)$；
换言之，下图交换：
$$
\xymatrix{
& f \circ a \circ j \ar[rd]^{\text{id}_f \star \alpha} \\
y \circ j \ar[ru]^{\beta \star \text{id}_j} \ar[rr]^\gamma & &
f \circ x
}
$$
。一个{\it 虚线箭头的态射}
$(a, \alpha, \beta) \to (a', \alpha', \beta')$ 是满足
$\alpha = \alpha' \circ (\theta \star \text{id}_j)$ 以及
$\beta' = (\text{id}_f \star \theta) \circ \beta$ 的
$2$-箭头 $\theta : a \to a'$。
\end{definition}

\noindent
上述定义是《范畴》定义 \ref{categories-definition-dotted-arrows} 的特殊情形。
一般而言，虚线箭头的范畴依赖于 $\gamma$。若 $\mathcal{Y}$ 可由代数空间表示
（或者域上对象的自同构群平凡），则当然至多有一个 $\gamma$，无需检查三角形
的交换性。更一般地，有引理
\ref{stacks-morphisms-lemma-cat-dotted-arrows-independent}。三角形的交换性在证明与基变换的
相容性时很重要，见引理 \ref{stacks-morphisms-lemma-cat-dotted-arrows-base-change} 的证明。

\begin{lemma}
\label{stacks-morphisms-lemma-cat-dotted-arrows}
在定义 \ref{stacks-morphisms-definition-fill-in-diagram} 的情形下，虚线箭头的范畴是群胚。
若 $\Delta_f$ 分离，则它是集合胚。
\end{lemma}

\begin{proof}
由于 $2$-箭头可逆，虚线箭头的范畴显然是群胚。给定虚线箭头
$(a, \alpha, \beta)$，虚线箭头 $(a, \alpha, \beta)$ 的一个自同构
是满足两个条件的 $2$-态射
$\theta : a \to a$。第一个条件
$\beta = (\text{id}_f \star \theta) \circ \beta$ 表明 $\theta$ 定义态射
$(a, \theta) : \Spec(A) \to \mathcal{I}_{\mathcal{X}/\mathcal{Y}}$.
第二个条件
$\alpha = \alpha \circ (\theta \star \text{id}_j)$
蕴含 $(a, \theta)$ 到 $\Spec(K)$ 的限制为恒等态射。图示为
$$
\xymatrix{
\mathcal{I}_{\mathcal{X}/\mathcal{Y}} \ar[d] & &
\Spec(K) \ar[d]^j \ar[ll]_{(a \circ j, \text{id})} \\
\mathcal{X} & & \Spec(A) \ar[ll]_a \ar[llu]_{(a, \theta)}
}
$$
换言之，若 $G \to \Spec(A)$ 是沿 $a$ 拉回相对惯性所得的群代数空间，
则 $\theta$ 定义点 $\theta \in G(A)$，其在 $G(K)$ 中的像平凡。
当然，若单位元 $e : \Spec(A) \to G$ 是闭浸入，则这只有在
$\theta$ 为恒等态射时才可能发生。查看引理
\ref{stacks-morphisms-lemma-diagonal-diagonal} 即得所需结论。
\end{proof}

\begin{lemma}
\label{stacks-morphisms-lemma-cat-dotted-arrows-independent}
在定义 \ref{stacks-morphisms-definition-fill-in-diagram} 中，假设
$\mathcal{I}_\mathcal{Y} \to \mathcal{Y}$ 固有
（例如，若 $\mathcal{Y}$ 分离，或 $\mathcal{Y}$ 在某个代数空间上分离）。
则虚线箭头的范畴不依赖于 $\gamma$ 的选取（精确到非典范等价）；而虚线箭头
对某个 $\gamma$ 存在（从而等价地对所有选择存在），当且仅当存在图
$$
\xymatrix{
\Spec(K) \ar[r]_-x \ar[d]_j & \mathcal{X} \ar[d]^f \\
\Spec(A) \ar[r]^-y \ar[ru]_a & \mathcal{Y}
}
$$
，其两个三角形都 $2$-交换（无需检查 $2$-态射是否正确复合）。
\end{lemma}

\begin{proof}
设 $\gamma, \gamma' : y \circ j \longrightarrow f \circ x$ 是两个
$2$-态射。则 $\gamma^{-1} \circ \gamma'$ 是 $y$ 在 $\Spec(K)$ 上的
自同构。因此，若
$\mathit{Isom}_\mathcal{Y}(y, y) \to \Spec(A)$ 固有，则由固有性的赋值判据
（《代数空间的态射》引理
\ref{spaces-morphisms-lemma-characterize-proper}），可以找到
$\delta : y \to y$，其到 $\Spec(K)$ 的限制为
$\gamma^{-1} \circ \gamma'$。于是可以用 $\delta$ 定义 $\gamma$
的虚线箭头范畴与 $\gamma'$ 的虚线箭头范畴之间的等价：把
$(a, \alpha, \beta)$ 送到 $(a, \alpha, \beta \circ \delta)$。
最后一个陈述显然。
\end{proof}

\begin{lemma}
\label{stacks-morphisms-lemma-cat-dotted-arrows-base-change}
假设给定 $2$-交换图
$$
\xymatrix{
\Spec(K) \ar[r]_-{x'} \ar[d]_j &
\mathcal{X}' \ar[d]^p \ar[r]_q &
\mathcal{X} \ar[d]^f \\
\Spec(A) \ar[r]^-{y'} &
\mathcal{Y}' \ar[r]^g &
\mathcal{Y}
}
$$
，其中右侧方块为 $2$-笛卡儿的。选取 $2$-箭头
$\gamma' : y' \circ j \to p \circ x'$。令
$x = q \circ x'$、$y = g \circ y'$，并令
$\gamma : y \circ j \to f \circ x$ 为 $\gamma'$ 与右侧方块
$2$-交换性中隐含的 $2$-箭头的复合。则左侧方块及 $\gamma'$
的虚线箭头范畴，等价于外部矩形及 $\gamma$ 的虚线箭头范畴。
\end{lemma}

\begin{proof}
（若不用虚线箭头的范畴，而考察如引理
\ref{stacks-morphisms-lemma-cat-dotted-arrows-independent} 中产生两个 $2$-交换三角形的态射的
同构类集合，我们不知道如何证明本引理的类似结论；事实上，该类似结论很可能
是错的。）第一种证明：本引理是《范畴》引理
\ref{categories-lemma-cat-dotted-arrows-base-change} 的特殊情形。
第二种证明：按《范畴》引理
\ref{categories-lemma-2-product-categories-over-C} 所述，可以用 $2$-纤维积
$\mathcal{Y}' \times_\mathcal{Y} \mathcal{X}$ 替换
$\mathcal{X}'$。此时对象 $x'$ 变成三元组
$(y' \circ j, x, \gamma)$。然后，可以把外部矩形的虚线箭头
$(a, \alpha, \beta)$ 送到左侧方块的虚线箭头
$(a', \alpha', \beta')$，其中取 $a' = (y', a, \beta)$、
$\alpha' = (\text{id}_{y' \circ j}, \alpha)$ 以及
$\beta' = \text{id}_{y'}$。细节略。
\end{proof}

\begin{lemma}
\label{stacks-morphisms-lemma-cat-dotted-arrows-composition}
假设给定 $2$-交换图
$$
\xymatrix{
\Spec(K) \ar[r]_-x \ar[dd]_j & \mathcal{X} \ar[d]^f \\
& \mathcal{Y} \ar[d]^g \\
\Spec(A) \ar[r]^-z & \mathcal{Z}
}
$$
选取 $2$-箭头 $\gamma : z \circ j \to g \circ f \circ x$。设
$\mathcal{C}$ 是外部矩形及 $\gamma$ 的虚线箭头范畴。设
$\mathcal{C}'$ 是下列方块
$$
\xymatrix{
\Spec(K) \ar[r]_-{f \circ x} \ar[d]_j & \mathcal{Y} \ar[d]^g \\
\Spec(A) \ar[r]^-z & \mathcal{Z}
}
$$
及 $\gamma$ 的虚线箭头范畴。则 $\mathcal{C}$ 等价于一个范畴
$\mathcal{C}''$，它具有下述性质：存在函子
$\mathcal{C}'' \to \mathcal{C}'$，使 $\mathcal{C}''$ 成为
$\mathcal{C}'$ 上纤维化于群胚的范畴，并且其纤维范畴是某些如下形式方块
$$
\xymatrix{
\Spec(K) \ar[r]_-x \ar[d]_j & \mathcal{X} \ar[d]^f \\
\Spec(A) \ar[r]^-y & \mathcal{Y}
}
$$
及某些 $y \circ j \to f \circ x$ 的选取所对应的虚线箭头范畴。
\end{lemma}

\begin{proof}
本引理是《范畴》引理
\ref{categories-lemma-cat-dotted-arrows-composition} 的特殊情形。
\end{proof}

\begin{definition}
\label{stacks-morphisms-definition-uniqueness}
设 $f : \mathcal{X} \to \mathcal{Y}$ 是代数叠的态射。若对每个图
(\ref{stacks-morphisms-equation-diagram}) 及定义
\ref{stacks-morphisms-definition-fill-in-diagram} 中的 $\gamma$，虚线箭头的范畴或者为空，或者是
恰有一个同构类的集合胚，则称 $f$ 满足{\it 赋值判据的唯一性部分}。
\end{definition}

\begin{lemma}
\label{stacks-morphisms-lemma-base-change-uniqueness}
满足赋值判据唯一性部分的代数叠态射，沿任意代数叠态射作基变换后，
仍是满足赋值判据唯一性部分的代数叠态射。
\end{lemma}

\begin{proof}
由引理 \ref{stacks-morphisms-lemma-cat-dotted-arrows-base-change} 和定义得出。
\end{proof}

\begin{lemma}
\label{stacks-morphisms-lemma-composition-uniqueness}
满足赋值判据唯一性部分的代数叠态射的复合，仍是满足赋值判据唯一性部分的
代数叠态射。
\end{lemma}

\begin{proof}
由引理 \ref{stacks-morphisms-lemma-cat-dotted-arrows-composition} 和定义得出。
\end{proof}

\begin{lemma}
\label{stacks-morphisms-lemma-uniqueness-representable}
设 $f : \mathcal{X} \to \mathcal{Y}$ 是可由代数空间表示的代数叠态射。
下列条件等价：
\begin{enumerate}
\item $f$ 满足赋值判据的唯一性部分；
\item 对每个概形 $T$ 和态射 $T \to \mathcal{Y}$，态射
$\mathcal{X} \times_\mathcal{Y} T \to T$ 作为代数空间的态射，
满足赋值判据的唯一性部分。
\end{enumerate}
\end{lemma}

\begin{proof}
由引理 \ref{stacks-morphisms-lemma-cat-dotted-arrows-base-change} 和定义得出。
\end{proof}

\begin{definition}
\label{stacks-morphisms-definition-existence}
设 $f : \mathcal{X} \to \mathcal{Y}$ 是代数叠的态射。若对每个图
(\ref{stacks-morphisms-equation-diagram}) 及定义
\ref{stacks-morphisms-definition-fill-in-diagram} 中的 $\gamma$，都存在域扩张 $K'/K$
以及支配 $A$ 的赋值环 $A' \subset K'$，使得下图外部矩形
$$
\xymatrix{
\Spec(K') \ar[r] \ar@/^2em/[rr]_{x'} \ar[d]_{j'} &
\Spec(K) \ar[d]_j \ar[r]_-x &
\mathcal{X} \ar[d]^f \\
\Spec(A') \ar[r] \ar@/_2em/[rr]^{y'} &
\Spec(A) \ar[r]^-y &
\mathcal{Y}
}
$$
连同诱导的 $2$-箭头
$\gamma' : y' \circ j' \to f \circ x'$ 所对应的虚线箭头范畴非空，
则称 $f$ 满足{\it 赋值判据的存在性部分}。
\end{definition}

\noindent
我们已经在代数空间态射的情形中看到，为得到正确的赋值判据概念，必须允许
分式域的扩张。见《代数空间的态射》例
\ref{spaces-morphisms-example-finite-separable-needed}。不过，对分离代数空间之间的
态射，不需要这种扩张（《代数空间的态射》引理
\ref{spaces-morphisms-lemma-usual-enough}）。然而，对代数叠之间的态射，
即使 $\mathcal{X}$ 和 $\mathcal{Y}$ 都分离，也可能需要扩张。例如，考虑
代数叠态射
$$
[\Spec(\mathbf{C})/G] \to \Spec(\mathbf{C})
$$
，它在基概形 $\Spec(\mathbf{C})$ 上，其中 $G$ 是阶为 $2$ 的群。
源和靶都是分离代数叠，而且该态射固有。因此，它满足赋值判据的唯一性部分和
存在性部分（见引理 \ref{stacks-morphisms-lemma-criterion-proper}）。但另一方面，存在图
$$
\xymatrix{
\Spec(K) \ar[r] \ar[d] & [\Spec(\mathbf{C})/G] \ar[d] \\
\Spec(A) \ar[r] & \Spec(\mathbf{C})
}
$$
，当 $A = \mathbf{C}[[t]]$ 且 $K = \mathbf{C}((t))$ 时不存在虚线箭头。
事实上，顶部水平箭头由 $K = \mathbf{C}((t))$ 的谱上的一个 $G$-挠子给出。
由于严格 Hensel 局部环 $A = \mathbf{C}[[t]]$ 上的任意 $G$-挠子都平凡，
可见若虚线箭头总存在，则 $K$ 上的每个 $G$-挠子都平凡。这并不成立，
因为 $G = \{+1, -1\}$，且由 Kummer 理论，$K$ 上的 $G$-挠子由非平凡群
$K^*/(K^*)^2$ 分类。

\begin{lemma}
\label{stacks-morphisms-lemma-base-change-existence}
满足赋值判据存在性部分的代数叠态射，沿任意代数叠态射作基变换后，
仍是满足赋值判据存在性部分的代数叠态射。
\end{lemma}

\begin{proof}
由引理 \ref{stacks-morphisms-lemma-cat-dotted-arrows-base-change} 和定义得出。
\end{proof}

\begin{lemma}
\label{stacks-morphisms-lemma-composition-existence}
满足赋值判据存在性部分的代数叠态射的复合，仍是满足赋值判据存在性部分的
代数叠态射。
\end{lemma}

\begin{proof}
由引理 \ref{stacks-morphisms-lemma-cat-dotted-arrows-composition} 和定义得出。
\end{proof}

\begin{lemma}
\label{stacks-morphisms-lemma-existence-representable}
设 $f : \mathcal{X} \to \mathcal{Y}$ 是可由代数空间表示的代数叠态射。
下列条件等价：
\begin{enumerate}
\item $f$ 满足赋值判据的存在性部分；
\item 对每个概形 $T$ 和态射 $T \to \mathcal{Y}$，态射
$\mathcal{X} \times_\mathcal{Y} T \to T$ 作为代数空间的态射，
满足赋值判据的存在性部分。
\end{enumerate}
\end{lemma}

\begin{proof}
由引理 \ref{stacks-morphisms-lemma-cat-dotted-arrows-base-change} 和定义得出。
\end{proof}

\begin{lemma}
\label{stacks-morphisms-lemma-closed-immersion-valuative-criteria}
代数叠的闭浸入同时满足赋值判据的存在性部分和唯一性部分。
\end{lemma}

\begin{proof}
略。提示：利用引理 \ref{stacks-morphisms-lemma-uniqueness-representable} 和
\ref{stacks-morphisms-lemma-existence-representable}，归结到概形的闭浸入情形。
\end{proof}





\section{第二对角态射的赋值判据}
\label{stacks-morphisms-section-valuative-second}

\noindent
逆向陈述已经在引理 \ref{stacks-morphisms-lemma-cat-dotted-arrows} 中证明。判据本身如下。

\begin{lemma}
\label{stacks-morphisms-lemma-setoids-and-diagonal}
设 $f : \mathcal{X} \to \mathcal{Y}$ 是代数叠的态射。若
$\Delta_f$ 拟分离，并且对每个图 (\ref{stacks-morphisms-equation-diagram}) 及定义
\ref{stacks-morphisms-definition-fill-in-diagram} 中 $\gamma$ 的选取，虚线箭头范畴都是
集合胚，则 $\Delta_f$ 分离。
\end{lemma}

\begin{proof}
我们将写出详细证明，但强烈建议读者在阅读引理
\ref{stacks-morphisms-lemma-cat-dotted-arrows} 证明中的论证后，受其启发自行找出证明。

\medskip\noindent
假设 $\Delta_f$ 拟分离，并且对每个图
(\ref{stacks-morphisms-equation-diagram}) 及定义
\ref{stacks-morphisms-definition-fill-in-diagram} 中 $\gamma$ 的选取，虚线箭头范畴都是集合胚。
由引理 \ref{stacks-morphisms-lemma-diagonal-diagonal}，只需证明
$e : \mathcal{X} \to \mathcal{I}_{\mathcal{X}/\mathcal{Y}}$
是闭浸入。由引理
\ref{stacks-morphisms-lemma-first-diagonal-separated-second-diagonal-closed}，事实上只需证明
$e = \Delta_{f, 2}$ 泛闭。任一上述引理都说明，由 $\Delta_f$
拟分离的假设，$e = \Delta_{f, 2}$ 拟紧。

\medskip\noindent
本段证明 $e$ 满足赋值判据的存在性部分。考虑实线部分 $2$-交换的图
$$
\xymatrix{
\Spec(K) \ar[r]_x \ar[d]_j & \mathcal{X} \ar[d]^e \\
\Spec(A) \ar[r]^{(a, \theta)} & \mathcal{I}_{\mathcal{X}/\mathcal{Y}}
}
$$
，并任取见证该图 $2$-交换的 $2$-态射
$\alpha : (a, \theta) \circ j \to e \circ x$
（这里用 $\alpha$ 代替定义 \ref{stacks-morphisms-definition-fill-in-diagram} 中使用的字母
$\gamma$）。注意，$f \circ \theta = \text{id}$；下文将用到这一点。
又有 $e \circ x = (x, \text{id}_x)$ 以及
$(a, \theta) \circ j = (a \circ j, \theta \star \text{id}_j)$。
因此，$\alpha$ 是与 $\theta \star \text{id}_j$ 和 $\text{id}_x$
相容的 $2$-箭头 $\alpha : a \circ j \to x$。令
$y = f \circ x$ 且 $\beta = \text{id}_{f \circ a}$。反向阅读引理
\ref{stacks-morphisms-lemma-cat-dotted-arrows} 证明中的论证，可知 $\theta$
是虚线箭头 $(a, \alpha, \beta)$ 的自同构，其中
$$
\gamma : y \circ j \to f \circ x
\quad\text{equal to}\quad
\text{id}_f \star \alpha : f \circ a \circ j \to f \circ x
$$
。另一方面，$\text{id}_a$ 也是自同构，所以由关于 $f$ 的假设可知
$\theta = \text{id}_a$。于是可以取态射
$a : \Spec(A) \to \mathcal{X}$ 作为上面所示图的虚线箭头，其中
$2$-态射 $(a, \text{id}_a) \circ j \to (x, \text{id}_x)$ 由
$\alpha$ 给出，而 $(a, \theta) \to e \circ a$ 由
$\text{id}_a$ 给出。

\medskip\noindent
由引理 \ref{stacks-morphisms-lemma-base-change-existence}，$e$ 的任意基变换都满足赋值判据的
存在性部分。由于 $e$ 可由代数空间表示，只需证明沿满光滑态射
$I \to \mathcal{I}_{\mathcal{X}/\mathcal{Y}}$ 作基变换后 $e$ 泛闭，
其中 $I$ 是代数空间（见《叠的性质》一节
\ref{stacks-properties-section-properties-morphisms}）。该基变换
$e' : X' \to I'$ 是代数空间的拟紧态射，并满足赋值判据的存在性部分，
所以由《代数空间的态射》引理
\ref{spaces-morphisms-lemma-quasi-compact-existence-universally-closed} 泛闭。
\end{proof}






\section{对角态射的赋值判据}
\label{stacks-morphisms-section-valuative-diagonal}

\noindent
所需结果是引理 \ref{stacks-morphisms-lemma-uniqueness-and-diagonal}。先陈述并证明一个形式辅助引理。

\begin{lemma}
\label{stacks-morphisms-lemma-helper-diagonal}
设 $f : \mathcal{X} \to \mathcal{Y}$ 是代数叠的态射。考虑实线部分
$2$-交换的图
$$
\xymatrix{
\Spec(K) \ar[rr]_-x \ar[d]_j & &
\mathcal{X} \ar[d]^{\Delta_f} \\
\Spec(A) \ar[rr]^{(a_1, a_2, \varphi)} \ar@{..>}[rru] & &
\mathcal{X} \times_\mathcal{Y} \mathcal{X}
}
$$
，其中 $A$ 是分式域为 $K$ 的赋值环。设
$\gamma : (a_1, a_2, \varphi) \circ j \longrightarrow \Delta_f \circ x$
是见证该图 $2$-交换的 $2$-态射。则
\begin{enumerate}
\item 把 $\gamma = (\alpha_1, \alpha_2)$ 写成
$\alpha_i : a_i \circ j \to x$ 的形式，则在图
$$
\xymatrix{
\Spec(K) \ar[r]_-x \ar[d]_j & \mathcal{X} \ar[d]^f \\
\Spec(A) \ar[r]^-{f \circ a_1} \ar@{..>}[ru] & \mathcal{Y}
}
$$
中得到两个虚线箭头 $(a_1, \alpha_1, \text{id})$ 和
$(a_2, \alpha_2, \varphi)$。
\item 原图及 $\gamma$ 的虚线箭头范畴是集合胚，其对象的同构类集合等于
虚线箭头范畴中的态射集
$(a_1, \alpha_1, \text{id}) \to (a_2, \alpha_2, \varphi)$。
\end{enumerate}
\end{lemma}

\begin{proof}
由于 $\Delta_f$ 可由代数空间表示（因而 $\Delta_f$ 的对角态射分离），
由引理 \ref{stacks-morphisms-lemma-cat-dotted-arrows}，本引理第一个交换图的虚线箭头范畴
是集合胚。本引理的其余陈述都由 $2$-图式计算得出，略去这些计算。
\end{proof}

\begin{lemma}
\label{stacks-morphisms-lemma-uniqueness-and-diagonal}
设 $f : \mathcal{X} \to \mathcal{Y}$ 是代数叠的态射。假设 $f$ 拟分离。
若 $f$ 满足赋值判据的唯一性部分，则 $f$ 分离。
\end{lemma}

\begin{proof}
关于 $f$ 的假设意味着 $\Delta_f$ 拟紧且拟分离
（定义 \ref{stacks-morphisms-definition-separated}）。必须证明 $\Delta_f$ 固有。
引理 \ref{stacks-morphisms-lemma-setoids-and-diagonal} 说明 $\Delta_f$ 分离。由引理
\ref{stacks-morphisms-lemma-properties-diagonal}，$\Delta_f$ 局部有限型。为完成证明，
还需证明 $\Delta_f$ 泛闭。一个形式论证（见引理
\ref{stacks-morphisms-lemma-helper-diagonal}）说明，赋值判据的唯一性部分蕴含：任意如下
实线图中都有虚线箭头：
$$
\xymatrix{
\Spec(K) \ar[d] \ar[r] & \mathcal{X} \ar[d]^{\Delta_f} \\
\Spec(A) \ar[r] \ar@{..>}[ru] & \mathcal{X} \times_\mathcal{Y} \mathcal{X}
}
$$
利用该性质在任意基变换下保持这一点，可知 $\Delta_f$ 沿任意映入
$\mathcal{X} \times_\mathcal{Y} \mathcal{X}$ 的代数空间所作的基变换
都满足赋值判据的存在性部分；再由代数空间态射的赋值判据，可知它泛闭，
见《代数空间的态射》引理
\ref{spaces-morphisms-lemma-quasi-compact-existence-universally-closed}。
\end{proof}

\noindent
下面是逆命题。

\begin{lemma}
\label{stacks-morphisms-lemma-converse-uniqueness-and-diagonal}
设 $f : \mathcal{X} \to \mathcal{Y}$ 是代数叠的态射。若 $f$ 分离，
则 $f$ 满足赋值判据的唯一性部分。
\end{lemma}

\begin{proof}
由于 $f$ 分离，由引理 \ref{stacks-morphisms-lemma-cat-dotted-arrows}，所有虚线箭头范畴
都是集合胚。考虑图
$$
\xymatrix{
\Spec(K) \ar[r]_-x \ar[d]_j & \mathcal{X} \ar[d]^f \\
\Spec(A) \ar[r]^-y \ar@{..>}[ru] & \mathcal{Y}
}
$$
以及定义 \ref{stacks-morphisms-definition-fill-in-diagram} 中的 $2$-态射
$\gamma : y \circ j \to f \circ x$。考虑虚线箭头范畴的两个对象
$(a, \alpha, \beta)$ 和 $(a', \beta', \alpha')$。为完成证明，
必须证明这两个对象同构。同构
$$
f \circ a \xrightarrow{\beta^{-1}} y \xrightarrow{\beta'} f \circ a'
$$
意味着 $(a, a', \beta' \circ \beta^{-1})$ 是态射
$\Spec(A) \to \mathcal{X} \times_\mathcal{Y} \mathcal{X}$。
另一方面，$\alpha$ 和 $\alpha'$ 定义 $2$-箭头
$$
(a, a', \beta' \circ \beta^{-1}) \circ j =
(a \circ j, a' \circ j,
(\beta' \star \text{id}_j) \circ (\beta \star \text{id}_j)^{-1})
\xrightarrow{(\alpha, \alpha')} (x, x, \text{id}) = \Delta_f \circ x
$$
这里使用了 $(a, \alpha, \beta)$ 和 $(a', \alpha', \beta')$
都是关于 $\gamma$ 的虚线箭头这一事实。由此得到交换图
$$
\xymatrix{
\Spec(K) \ar[d]_j \ar[rr]_x & & \mathcal{X} \ar[d]^{\Delta_f} \\
\Spec(A) \ar[rr]^{(a, a', \beta' \circ \beta^{-1})} & &
\mathcal{X} \times_\mathcal{Y} \mathcal{X}
}
$$
，其 $2$-交换性由 $(\alpha, \alpha')$ 见证。现在，$\Delta_f$
可由代数空间表示（引理 \ref{stacks-morphisms-lemma-properties-diagonal}），并且因为 $f$
分离而固有。因此，由引理 \ref{stacks-morphisms-lemma-existence-representable}
和代数空间固有性的赋值判据（《代数空间的态射》引理
\ref{spaces-morphisms-lemma-characterize-proper}），可知存在虚线箭头。
展开该构造可见，这意味着 $(a, \alpha, \beta)$ 与
$(a', \alpha', \beta')$ 在虚线箭头范畴中同构，正合所需。
\end{proof}






\section{泛闭性的赋值判据}
\label{stacks-morphisms-section-valutive-criterion}

\noindent
先给出一个陈述。

\begin{lemma}
\label{stacks-morphisms-lemma-quasi-compact-existence-universally-closed}
设 $f : \mathcal{X} \to \mathcal{Y}$ 是代数叠的态射。假设
\begin{enumerate}
\item $f$ 拟紧；
\item $f$ 满足赋值判据的存在性部分。
\end{enumerate}
则 $f$ 泛闭。
\end{lemma}

\begin{proof}
由引理 \ref{stacks-morphisms-lemma-base-change-quasi-compact} 和
\ref{stacks-morphisms-lemma-base-change-existence}，性质 (1) 和 (2) 在任意基变换下保持。
由引理 \ref{stacks-morphisms-lemma-universally-closed-local}，只需证明：每当 $T$
是映入 $\mathcal{Y}$ 的仿射概形时，
$|T \times_\mathcal{Y} \mathcal{X}| \to |T|$ 是闭映射。因此，只需证明：
若 $f : \mathcal{X} \to Y$ 是从代数叠到仿射概形的拟紧态射，并满足
赋值判据的存在性部分，则 $|f|$ 是闭映射。设
$T \subset |\mathcal{X}|$ 是闭子集。为完成证明，必须证明 $f(T)$ 闭。

\medskip\noindent
设 $\mathcal{Z} \subset \mathcal{X}$ 是 $T$ 上的约化诱导代数叠结构
（《叠的性质》定义
\ref{stacks-properties-definition-reduced-induced-stack}）。则
$i : \mathcal{Z} \to \mathcal{X}$ 是闭浸入，而必须证明
$|\mathcal{Z}| \to |Y|$ 的像闭。闭浸入拟紧
（引理 \ref{stacks-morphisms-lemma-closed-immersion-quasi-compact}）且满足赋值判据的存在性部分
（引理 \ref{stacks-morphisms-lemma-closed-immersion-valuative-criteria}）；拟紧态射的复合拟紧
（引理 \ref{stacks-morphisms-lemma-composition-quasi-compact}）；复合还保持满足赋值判据存在性
部分这一性质（引理 \ref{stacks-morphisms-lemma-composition-existence}）。所以只需证明：
若 $f : \mathcal{X} \to Y$ 是从代数叠到仿射概形的拟紧态射，并满足
赋值判据的存在性部分，则 $|f|(|\mathcal{X}|)$ 闭。

\medskip\noindent
由于 $\mathcal{X}$ 拟紧（它在仿射概形 $Y$ 上拟紧），可以选取仿射概形
$U$ 以及满光滑态射 $U \to \mathcal{X}$
（《叠的性质》引理
\ref{stacks-properties-lemma-quasi-compact-stack}）。假设 $y \in Y$
位于 $U \to Y$ 的像的闭包中（换言之，位于 $|f|$ 的像的闭包中）。
由《态射》引理
\ref{morphisms-lemma-reach-points-scheme-theoretic-image}，可以找到分式域为
$K$ 的赋值环 $A$ 以及交换图
$$
\xymatrix{
\Spec(K) \ar[r] \ar[d] & U \ar[d] \\
\Spec(A) \ar[r] & Y
}
$$
，使得 $\Spec(A)$ 的闭点映到 $y$。由假设，得到扩张 $K'/K$、
支配 $A$ 的赋值环 $A' \subset K'$，以及下图中的虚线箭头：
$$
\xymatrix{
\Spec(K') \ar[r] \ar[d] &
\Spec(K) \ar[r] \ar[d] &
U \ar[d] \ar[r] &
\mathcal{X} \ar[d]^f \\
\Spec(A') \ar[r] \ar@{..>}[rrru] &
\Spec(A) \ar[r] &
Y  \ar@{=}[r] &
Y
}
$$
因此 $y$ 位于 $|f|$ 的像中，得证。
\end{proof}

\noindent
下面是逆命题。

\begin{lemma}
\label{stacks-morphisms-lemma-converse-existence-universally-closed}
设 $f : \mathcal{X} \to \mathcal{Y}$ 是代数叠的态射。假设
\begin{enumerate}
\item $f$ 拟分离；
\item $f$ 泛闭。
\end{enumerate}
则 $f$ 满足赋值判据的存在性部分。
\end{lemma}

\begin{proof}
考虑实线图
$$
\xymatrix{
\Spec(K) \ar[r]_-x \ar[d]_j & \mathcal{X} \ar[d]^f \\
\Spec(A) \ar[r]^-y \ar@{..>}[ru] & \mathcal{Y}
}
$$
，其中 $A$ 是分式域为 $K$ 的赋值环，并且
$\gamma : y \circ j \longrightarrow f \circ x$ 如定义
\ref{stacks-morphisms-definition-fill-in-diagram} 所述。由引理
\ref{stacks-morphisms-lemma-cat-dotted-arrows-base-change}，为找到虚线箭头
（可能先把 $K$ 换成一个扩张，并把 $A$ 换成支配它的赋值环），
可以用 $\Spec(A)$ 替换 $\mathcal{Y}$，并用
$\Spec(A) \times_\mathcal{Y} \mathcal{X}$ 替换 $\mathcal{X}$。
这里当然用到了拟分离和泛闭在基变换下保持。因此归结到下一段所述的情形。

\medskip\noindent
考虑实线图
$$
\xymatrix{
\Spec(K) \ar[r]_-x \ar[d]_j & \mathcal{X} \ar[d]^f \\
\Spec(A) \ar@{=}[r] \ar@{..>}[ru] & \Spec(A)
}
$$
，其中 $A$ 是分式域为 $K$ 的赋值环，如定义
\ref{stacks-morphisms-definition-fill-in-diagram} 所述。由引理
\ref{stacks-morphisms-lemma-quasi-compact-permanence} 以及 $f$ 拟分离这一事实，
态射 $x$ 拟紧。由于 $f$ 泛闭，特别地，
$|f|(\overline{\{x\}})$ 是 $\Spec(A)$ 中的闭集。
这个像包含 $\Spec(A)$ 的一般点，所以存在位于 $x$ 的闭包中的点
$x' \in |\mathcal{X}|$，它映到 $\Spec(A)$ 的闭点。
由引理 \ref{stacks-morphisms-lemma-reach-points-scheme-theoretic-image}，可以找到交换图
$$
\xymatrix{
\Spec(K') \ar[r] \ar[d] & \Spec(K) \ar[d] \\
\Spec(A') \ar[r] & \mathcal{X}
}
$$
，使得 $\Spec(A')$ 的闭点映到 $x' \in |\mathcal{X}|$。
于是 $\Spec(A') \to \Spec(A)$ 把闭点映到闭点，即 $A'$ 支配 $A$，
证明完成。
\end{proof}





\section{固有性的赋值判据}
\label{stacks-morphisms-section-valutive-criterion-properness}

\noindent
陈述如下。

\begin{lemma}
\label{stacks-morphisms-lemma-criterion-proper}
设 $f : \mathcal{X} \to \mathcal{Y}$ 是代数叠的态射。
假设 $f$ 有限型且拟分离。则下列条件等价：
\begin{enumerate}
\item $f$ 固有；
\item $f$ 同时满足赋值判据的唯一性部分和存在性部分。
\end{enumerate}
\end{lemma}

\begin{proof}
固有态射就是分离、有限型且泛闭的态射。因此，本引理由引理
\ref{stacks-morphisms-lemma-uniqueness-and-diagonal},
\ref{stacks-morphisms-lemma-converse-uniqueness-and-diagonal},
\ref{stacks-morphisms-lemma-quasi-compact-existence-universally-closed} 以及
\ref{stacks-morphisms-lemma-converse-existence-universally-closed} 得出。
\end{proof}







\section{局部完全交态射}
\label{stacks-morphisms-section-lci}

\noindent
代数空间态射“为局部完全交态射”的性质在源与靶上光滑局部；见
《代数空间上的下降》引理 \ref{spaces-descent-lemma-smooth-local-source-target}
以及《空间态射进阶》引理
\ref{spaces-more-morphisms-lemma-descending-property-lci} 和
\ref{spaces-more-morphisms-lemma-lci-syntomic-local-source}。
由上面的引理 \ref{stacks-morphisms-lemma-local-source-target}，可以如下定义一个
代数空间的态射为局部完全交态射的含义；当源和靶都是代数空间时，
它与《空间态射进阶》第 \ref{spaces-more-morphisms-section-lci} 节中
已经定义的概念一致。

\begin{definition}
\label{stacks-morphisms-definition-lci}
设 $f : \mathcal{X} \to \mathcal{Y}$ 是代数叠的态射。若引理
\ref{stacks-morphisms-lemma-local-source-target} 中的等价条件在
$\mathcal{P} = \text{local complete intersection}$ 时成立，
则称 $f$ 是{\it 局部完全交态射}或 {\it Koszul} 的。
\end{definition}

\begin{lemma}
\label{stacks-morphisms-lemma-composition-lci}
局部完全交态射的复合仍是局部完全交态射。
\end{lemma}

\begin{proof}
结合评注 \ref{stacks-morphisms-remark-composition} 与《空间态射进阶》引理
\ref{spaces-more-morphisms-lemma-composition-lci}。
\end{proof}

\begin{lemma}
\label{stacks-morphisms-lemma-flat-base-change-lci}
局部完全交态射的平坦基变换仍是局部完全交态射。
\end{lemma}

\begin{proof}
略。提示：完全按照评注 \ref{stacks-morphisms-remark-base-change} 论证（但只考虑平坦的
$\mathcal{Y}' \to \mathcal{Y}$），并使用《空间态射进阶》引理
\ref{spaces-more-morphisms-lemma-flat-base-change-lci}。
\end{proof}

\begin{lemma}
\label{stacks-morphisms-lemma-lci-permanence}
设
$$
\xymatrix{
\mathcal{X} \ar[rr]_f \ar[rd] & & \mathcal{Y} \ar[ld] \\
& \mathcal{Z}
}
$$
是代数叠态射的交换图。假设 $\mathcal{Y} \to \mathcal{Z}$ 光滑，且
$\mathcal{X} \to \mathcal{Z}$ 是局部完全交态射。则
$f : \mathcal{X} \to \mathcal{Y}$ 是局部完全交态射。
\end{lemma}

\begin{proof}
选择概形 $W$ 和满光滑态射 $W \to \mathcal{Z}$。选择概形 $V$
和满光滑态射 $V \to W \times_\mathcal{Z} \mathcal{Y}$。
再选择概形 $U$ 和满光滑态射
$U \to V \times_\mathcal{Y} \mathcal{X}$。于是 $U \to W$
是概形的局部完全交态射，而 $V \to W$ 是概形的光滑态射。
由概形情形的结论（《态射进阶》引理
\ref{more-morphisms-lemma-lci-permanence}），可知 $U \to V$
是局部完全交态射。按定义，这意味着 $f$ 是局部完全交态射。
\end{proof}





\section{稳定子保持态射}
\label{stacks-morphisms-section-stabilizer-preserving}

\noindent
文献中，若代数叠的态射 $f : \mathcal{X} \to \mathcal{Y}$ 所诱导的态射
$\mathcal{I}_\mathcal{X} \to
\mathcal{X} \times_\mathcal{Y} \mathcal{I}_\mathcal{Y}$
是同构，则称这个态射是{\it 稳定子保持的}或{\it 不动点反映的}。
这种态射在 $\mathcal{X}$ 的每一点处都诱导自同构群之间的同构
（评注 \ref{stacks-morphisms-remark-identify-automorphism-groups}）。
本节证明关于这一概念的一些简单引理。

\begin{lemma}
\label{stacks-morphisms-lemma-stabilizer-preserving}
设 $f : \mathcal{X} \to \mathcal{Y}$ 是代数叠的态射。若
$\mathcal{I}_\mathcal{X} \to
\mathcal{X} \times_\mathcal{Y} \mathcal{I}_\mathcal{Y}$ 是同构，
则 $f$ 可由代数空间表示。
\end{lemma}

\begin{proof}
由引理 \ref{stacks-morphisms-lemma-second-diagonal} 立即得到。
\end{proof}

\begin{remark}
\label{stacks-morphisms-remark-get-property-auts-from-diagonal}
设 $f : \mathcal{X} \to \mathcal{Y}$ 是代数叠的态射，并设
$U \to \mathcal{X}$ 是以代数空间为源的态射。令 $G \to H$ 为态射
$\mathcal{I}_\mathcal{X} \to
\mathcal{X} \times_\mathcal{Y} \mathcal{I}_\mathcal{Y}$
到 $U$ 上的拉回。若 $\Delta_f$ 非分歧、\'etale 等，则
$G \to H$ 也具有相应性质。这是因为
$$
\xymatrix{
U \times_\mathcal{X} U \ar[r] \ar[d] & \mathcal{X} \ar[d]^{\Delta_f} \\
U \times_\mathcal{Y} U \ar[r] & \mathcal{X} \times_\mathcal{Y} \mathcal{X}
}
$$
是笛卡儿图，并且态射 $G \to H$ 是左侧竖直箭头沿对角态射
$U \to U \times U$ 的基变换。比较引理
\ref{stacks-morphisms-lemma-separated-implies-isom} 的证明。
\end{remark}

\begin{lemma}
\label{stacks-morphisms-lemma-aut-iso-unramified}
设 $f : \mathcal{X} \to \mathcal{Y}$ 是代数叠的非分歧态射。
下列条件等价：
\begin{enumerate}
\item $\mathcal{I}_\mathcal{X} \to
\mathcal{X} \times_\mathcal{Y} \mathcal{I}_\mathcal{Y}$
是同构；
\item 对所有 $x \in |\mathcal{X}|$，$f$ 都诱导 $x$ 和 $f(x)$ 处的
自同构群之间的同构（评注 \ref{stacks-morphisms-remark-identify-automorphism-groups}）。
\end{enumerate}
\end{lemma}

\begin{proof}
选择概形 $U$ 和满光滑态射 $U \to \mathcal{X}$。记 $G \to H$ 为态射
$\mathcal{I}_\mathcal{X} \to
\mathcal{X} \times_\mathcal{Y} \mathcal{I}_\mathcal{Y}$
到 $U$ 上的拉回。由评注 \ref{stacks-morphisms-remark-get-property-auts-from-diagonal}
和引理 \ref{stacks-morphisms-lemma-characterize-unramified}，态射 $G \to H$ 是 \'etale 的。
条件 (1) 等价于 $G \to H$ 是同构（例如应用《叠的性质》引理
\ref{stacks-properties-lemma-check-property-covering} 即得）。
条件 (2) 等价于：对每个 $u \in U$，纤维的态射 $G_u \to H_u$
都是同构。因此 (1) $\Rightarrow$ (2) 是显然的。若 (2) 成立，则
$G \to H$ 是代数空间的满、泛单且 \'etale 的态射。由《代数空间的态射》引理
\ref{spaces-morphisms-lemma-etale-universally-injective-open}，
这种态射是同构。
\end{proof}

\begin{lemma}
\label{stacks-morphisms-lemma-stabilizer-preserving-unramified}
\begin{reference}
\cite[Proposition 3.5]{rydh_quotients} 和
\cite[Proposition 2.5]{alper_quotient}
\end{reference}
设 $f : \mathcal{X} \to \mathcal{Y}$ 是代数叠的态射。假设
\begin{enumerate}
\item $f$ 可由代数空间表示且非分歧；
\item $\mathcal{I}_\mathcal{Y} \to \mathcal{Y}$ 固有。
\end{enumerate}
则使 $f$ 诱导 $x$ 和 $f(x)$ 处自同构群之间同构的那些
$x \in |\mathcal{X}|$ 所成的集合是开的
（评注 \ref{stacks-morphisms-remark-identify-automorphism-groups}）。令
$\mathcal{U} \subset \mathcal{X}$ 为相应的开子叠，则态射
$\mathcal{I}_\mathcal{U} \to
\mathcal{U} \times_\mathcal{Y} \mathcal{I}_\mathcal{Y}$
是同构。
\end{lemma}

\begin{proof}
选择概形 $U$ 和满光滑态射 $U \to \mathcal{X}$。记 $G \to H$ 为态射
$\mathcal{I}_\mathcal{X} \to
\mathcal{X} \times_\mathcal{Y} \mathcal{I}_\mathcal{Y}$
到 $U$ 上的拉回。由评注 \ref{stacks-morphisms-remark-get-property-auts-from-diagonal}
和引理 \ref{stacks-morphisms-lemma-characterize-unramified}，态射 $G \to H$ 是 \'etale 的。
由于 $f$ 可由代数空间表示，$G \to H$ 是单态射。因此 $G \to H$
是开浸入；见《代数空间的态射》引理
\ref{spaces-morphisms-lemma-etale-universally-injective-open}。
依假设，$H \to U$ 固有。

\medskip\noindent
完成这些准备后，可如下证明本引理。引理中 $|\mathcal{X}|$ 的子集的
逆像，显然就是使 $G_u \to H_u$ 为同构的那些 $u \in U$
所成的集合（毕竟 $G_u$ 是 $H_u$ 的开子群代数空间）。这是开子集，
因为其补集是闭子集 $|H| \setminus |G|$ 的像，而
$|H| \to |U|$ 是闭映射。由《叠的性质》引理
\ref{stacks-properties-lemma-open-substacks}，可以考虑
$\mathcal{X}$ 的相应开子叠 $\mathcal{U}$。把引理
\ref{stacks-morphisms-lemma-aut-iso-unramified} 应用于
$\mathcal{U} \to \mathcal{Y}$，即得引理的最后一个陈述。
\end{proof}

\begin{lemma}
\label{stacks-morphisms-lemma-base-change-stabilizer-preserving}
设
$$
\xymatrix{
\mathcal{X}' \ar[r] \ar[d]_{f'} & \mathcal{X} \ar[d]^f \\
\mathcal{Y}' \ar[r] & \mathcal{Y}
}
$$
是代数叠的笛卡儿图。
\begin{enumerate}
\item 设 $x' \in |\mathcal{X}'|$ 的像为 $x \in |\mathcal{X}|$。
若 $f$ 诱导 $x$ 和 $f(x)$ 处自同构群之间的同构
（评注 \ref{stacks-morphisms-remark-identify-automorphism-groups}），则 $f'$ 诱导
$x'$ 和 $f(x')$ 处自同构群之间的同构。
\item 若 $\mathcal{I}_\mathcal{X} \to
\mathcal{X} \times_\mathcal{Y} \mathcal{I}_\mathcal{Y}$ 是同构，则
$\mathcal{I}_{\mathcal{X}'} \to
\mathcal{X}' \times_{\mathcal{Y}'} \mathcal{I}_{\mathcal{Y}'}$
是同构。
\end{enumerate}
\end{lemma}

\begin{proof}
略。
\end{proof}

\begin{lemma}
\label{stacks-morphisms-lemma-stabilizer-preserving-points-cartesian}
设
$$
\xymatrix{
\mathcal{X}' \ar[r] \ar[d]_{f'} & \mathcal{X} \ar[d]^f \\
\mathcal{Y}' \ar[r]^g & \mathcal{Y}
}
$$
是代数叠的笛卡儿图。若 $f$ 在各点处诱导自同构群之间的同构
（评注 \ref{stacks-morphisms-remark-identify-automorphism-groups}），则
$$
\Mor(\Spec(k), \mathcal{X}')
\longrightarrow
\Mor(\Spec(k), \mathcal{Y}') \times \Mor(\Spec(k), \mathcal{X})
$$
对任意域 $k$ 都在同构类上为单射。
\end{lemma}

\begin{proof}
必须证明：给定 $(y', x)$，至多有一个映到它的 $x'$。
按照我们对 $2$-纤维积的构造，态射 $x'$ 由三元组
$(x, y', \alpha)$ 给出，其中
$\alpha : g \circ y' \to f \circ x$ 是 $2$-态射。
现在假设还有第二个这样的三元组 $(x, y', \beta)$。则
$\alpha$ 和 $\beta$ 相差自同构群代数空间 $G_{f(x)}$ 的一个
$k$-值点 $\epsilon$。依假设，$f$ 诱导同构
$G_x \to G_{f(x)}$，所以可以把 $\epsilon$ 提升为 $G_x$ 的
$k$-值点 $\gamma$。于是
$(\gamma, \text{id}) : (x, y', \alpha) \to
(x, y', \beta)$ 是所需的同构。
\end{proof}

\begin{lemma}
\label{stacks-morphisms-lemma-etale-iso}
设 $f : \mathcal{X} \to \mathcal{Y}$ 是代数叠的态射。假设 $f$ 是
\'etale 的，$f$ 在各点处诱导自同构群之间的同构
（评注 \ref{stacks-morphisms-remark-identify-automorphism-groups}），并且对每个代数闭域
$k$，函子
$$
f : \Mor(\Spec(k), \mathcal{X}) \longrightarrow \Mor(\Spec(k), \mathcal{Y})
$$
都是等价。则 $f$ 是同构。
\end{lemma}

\begin{proof}
由引理 \ref{stacks-morphisms-lemma-universally-injective}，$f$ 泛单。结合引理
\ref{stacks-morphisms-lemma-stabilizer-preserving} 和
\ref{stacks-morphisms-lemma-aut-iso-unramified}
可知 $f$ 可由代数空间表示。因此，由《代数空间的态射》引理
\ref{spaces-morphisms-lemma-etale-universally-injective-open}，
$f$ 是开浸入。最后注意，本引理的条件还保证 $f$ 满。
\end{proof}







\section{正规化}
\label{stacks-morphisms-section-normalization}

\noindent
本节是《代数空间的态射》第
\ref{spaces-morphisms-section-normalization} 节的类比版本。

\begin{lemma}
\label{stacks-morphisms-lemma-prepare-normalization}
设 $\mathcal{X}$ 是代数叠。下列条件等价：
\begin{enumerate}
\item 存在满光滑态射 $U \to \mathcal{X}$，其中 $U$ 是概形，
且 $U$ 的每个拟紧开子集都有有限多个不可约分支；
\item 对每个概形 $U$ 和每个光滑态射 $U \to \mathcal{X}$，
$U$ 的每个拟紧开子集都有有限多个不可约分支；
\item 对每个代数空间 $Y$ 和光滑态射 $Y \to \mathcal{X}$，
空间 $Y$ 满足《代数空间的态射》引理
\ref{spaces-morphisms-lemma-prepare-normalization} 中的等价条件；
\item 对每个在 $\mathcal{X}$ 上光滑的拟紧代数叠 $\mathcal{Y}$，
空间 $|\mathcal{Y}|$ 有有限多个不可约分支。
\end{enumerate}
\end{lemma}

\begin{proof}
(1)、(2) 与 (3) 的等价性由《下降》引理
\ref{descent-lemma-locally-finite-nr-irred-local-fppf}、
《叠的性质》引理 \ref{stacks-properties-lemma-type-property} 以及
《代数空间的态射》引理
\ref{spaces-morphisms-lemma-prepare-normalization} 得出。
由这些引用也显然可知条件 (4) 蕴含条件 (1)。反过来，假设等价条件
(1)、(2) 与 (3) 成立，并设 $\mathcal{Y} \to \mathcal{X}$
是代数叠的光滑态射，且 $\mathcal{Y}$ 拟紧。由《叠的性质》引理
\ref{stacks-properties-lemma-quasi-compact-stack}，可以选择仿射概形
$V$ 和满光滑态射 $V \to \mathcal{Y}$。由 (2)，$V$ 有有限多个
不可约分支；又因 $|V| \to |\mathcal{Y}|$ 满且连续，故由《拓扑》引理
\ref{topology-lemma-surjective-continuous-irreducible-components}，
$|\mathcal{Y}|$ 有有限多个不可约分支。
\end{proof}

\begin{lemma}
\label{stacks-morphisms-lemma-normalization}
设 $\mathcal{X}$ 是满足引理 \ref{stacks-morphisms-lemma-prepare-normalization}
中等价条件的代数叠。则存在代数叠的整态射
$$
\mathcal{X}^\nu \longrightarrow \mathcal{X}
$$
，使得对每个概形 $U$ 和光滑态射 $U \to \mathcal{X}$，
纤维积 $\mathcal{X}^\nu \times_\mathcal{X} U$ 是 $U$ 的正规化。
\end{lemma}

\begin{proof}
设 $U \to \mathcal{X}$ 是满光滑态射，其中 $U$ 是概形。令
$R = U \times_\mathcal{X} U$。回顾一下，我们得到代数空间中的光滑群胚
$(U, R, s, t, c)$ 和 $\mathcal{X}$ 的表示
$\mathcal{X} = [U/R]$；见《代数叠》引理
\ref{algebraic-lemma-map-space-into-stack} 和
\ref{algebraic-lemma-stack-presentation} 以及定义
\ref{algebraic-definition-presentation}。关于 $\mathcal{X}$ 的假设意味着
$U$ 的正规化 $U^\nu$ 有定义；见《态射》定义
\ref{morphisms-definition-normalization}。由《代数空间的态射》引理
\ref{spaces-morphisms-lemma-normalization}
，取正规化与代数空间的光滑态射可交换。因此，$R$ 的正规化 $R^\nu$
同时同构于 $R \times_{s, U} U^\nu$ 和 $U^\nu \times_{U, t} R$。
于是得到代数空间的两个光滑态射
$s^\nu : R^\nu \to U^\nu$ 和 $t^\nu : R^\nu \to U^\nu$
。对纤维积作形式计算可知，
$R^\nu \times_{s^\nu, U^\nu, t^\nu} R^\nu$ 是
$R \times_{s, U, t} R$ 的正规化。因此，光滑态射
$c : R \times_{s, U, t} R \to R$ 也延拓为 $c^\nu$。
类似地，逆元 $i : R \to R$（一个同构）诱导同构
$i^\nu : R^\nu \to R^\nu$。最后，单位元 $e : U \to R$
提升为 $e^\nu : U^\nu \to R^\nu$；例如，这是因为 $e$ 是 $s$
的截面且 $R^\nu = R \times_{U, s} U^\nu$。我们断言
$(U^\nu, R^\nu, s^\nu, t^\nu, c^\nu)$ 是代数空间中的光滑群胚。
为此，需要对 $(U^\nu, R^\nu, s^\nu, t^\nu, c^\nu, e^\nu, i^\nu)$
验证《群胚》第 \ref{groupoids-section-groupoids} 节中的公理
(1)、(2)(a)、(2)(b)、(3)(a) 和 (3)(b)。例如，对 (1)，必须证明两个态射
$a, b : R^\nu \times_{s^\nu, U^\nu, t^\nu} R^\nu
\times_{s^\nu, U^\nu, t^\nu} R^\nu \to R^\nu$ 相同。
相同。这是因为已知相应的一对态射
$R \times_{s, U, t} R \times_{s, U, t} R \to R$
相同，而态射 $a$ 和 $b$ 都是这个态射到正规化上的唯一延拓。
其他公理同理。

\medskip\noindent
考虑代数叠 $\mathcal{X}^\nu = [U^\nu/R^\nu]$（《代数叠》定理
\ref{algebraic-theorem-smooth-groupoid-gives-algebraic-stack}).
由于有代数空间中群胚的态射
$(U^\nu, R^\nu, s^\nu, t^\nu, c^\nu) \to (U, R, s, t, c)$
，得到代数叠的态射 $\nu : \mathcal{X}^\nu \to \mathcal{X}$。
由于 $R^\nu = R \times_{s, U} U^\nu$，由《空间中的群胚》引理
\ref{spaces-groupoids-lemma-criterion-fibre-product}.
可知 $U^\nu = \mathcal{X}^\nu \times_\mathcal{X} U$。特别地，
由于 $U^\nu \to U$ 是整态射，$\nu$ 也是整态射。
略去如下验证：对每个从概形到 $\mathcal{X}$ 的光滑态射，
本引理中陈述的基变换性质都成立。
\end{proof}

\noindent
由此得到下面的定义。

\begin{definition}
\label{stacks-morphisms-definition-normalization}
设 $\mathcal{X}$ 是满足引理 \ref{stacks-morphisms-lemma-prepare-normalization}
中等价条件的代数叠。把态射
$$
\nu : \mathcal{X}^\nu \longrightarrow \mathcal{X}
$$
（由引理 \ref{stacks-morphisms-lemma-normalization} 构造）定义为 $\mathcal{X}$ 的
{\it 正规化}。
\end{definition}










\section{点与特化}
\label{stacks-morphisms-section-specializations}

\noindent
本节是《适宜代数空间》第
\ref{decent-spaces-section-specializations} 节的类比版本。

\begin{lemma}
\label{stacks-morphisms-lemma-specialization}
设 $\mathcal{X}$ 是代数叠。设 $f : U \to \mathcal{X}$ 是光滑态射，
其中 $U$ 是代数空间。设 $x' \leadsto x$ 是 $|\mathcal{X}|$
中点的一个特化。设 $u \in |U|$ 且 $f(u) = x$。若
$(\mathcal{X}, x')$ 满足《叠的性质》引理
\ref{stacks-properties-lemma-UR-quasi-compact-above-x} 中的等价条件，
则在 $|U|$ 中存在特化 $u' \leadsto u$，且 $f(u') = x'$。
\end{lemma}

\begin{proof}
选择 \'etale 态射 $(U_1, u_1) \to (U, u)$，其中 $U_1$ 是仿射概形。
于是可以并且确实用 $U_1$ 替换 $U$。因此可假设 $U$ 是仿射概形。
考虑代数空间 $R = U \times_\mathcal{X} U$ 及其光滑投影
$t, s : R \to U$。选择映到 $x'$ 的点 $w \in U$；这是可能的，
因为 $f : |U| \to |\mathcal{X}|$ 是开映射。由关于 $x'$ 的假设，
$t : R \to U$ 在 $w$ 上的纤维
$F' = t^{-1}(w) = R \times_{t, U} w$ 是拟紧代数空间。
选择仿射概形 $T$ 和满 \'etale 态射 $T \to F'$。
$x' \leadsto x$ 这一事实意味着 $u$ 位于态射
$$
T \to F' \to R \xrightarrow{s} U
$$
的像的闭包中。事实上，这个像是 $|U| \to |\mathcal{X}'|$ 在
$x'$ 上的纤维；若某个开子集 $u \in V \subset |U|$ 与该纤维不交，
则 $f(V)$ 是 $x$ 的一个不含 $x'$ 的开邻域，矛盾。因此，由《态射》引理
\ref{morphisms-lemma-reach-points-scheme-theoretic-image}
，在 $|U| \to |\mathcal{X}|$ 位于 $x'$ 上的纤维中存在
$u' \in |U|$，并且它特化到 $u$。
\end{proof}



\section{适宜代数叠}
\label{stacks-morphisms-section-reasonable-decent}

\noindent
本节是《适宜代数空间》第
\ref{decent-spaces-section-reasonable-decent} 节的类比版本。
特别地，下面的定义与该处定义的适宜代数空间概念相容。

\begin{definition}
\label{stacks-morphisms-definition-decent}
设 $\mathcal{X}$ 是代数叠。若对每个 $x \in |\mathcal{X}|$，
《叠的性质》引理
\ref{stacks-properties-lemma-UR-quasi-compact-above-x} 中的等价条件都成立，
则称 $\mathcal{X}$ 是{\it 适宜的}。
\end{definition}

\noindent
有人会把这个定义改述为：$\mathcal{X}$ 的每个点都是拟紧的。
稍强一些的条件是，要求 $x$ 的等价类中的任意态射
$\Spec(k) \to \mathcal{X}$ 既拟分离又拟紧。

\begin{lemma}
\label{stacks-morphisms-lemma-quasi-separated-decent}
拟分离代数叠 $\mathcal{X}$ 是适宜的。更一般地，若
$\Delta : \mathcal{X} \to \mathcal{X} \times \mathcal{X}$ 拟紧，
则 $\mathcal{X}$ 是适宜的。
\end{lemma}

\begin{proof}
具体地，若 $\mathcal{X}$ 拟分离，则任意源为拟紧概形 $T$ 的态射
$f : T \to \mathcal{X}$ 都拟紧；见引理
\ref{stacks-morphisms-lemma-quasi-compact-permanence}。若只知道 $\Delta$ 拟紧，
则使用描述
$$
T \times_{f, \mathcal{X}, f'} T' =
(T \times T')
\times_{(f, f'), \mathcal{X} \times \mathcal{X}, \Delta}
\mathcal{X}
$$
来证明这一点。细节略。
\end{proof}

\begin{lemma}
\label{stacks-morphisms-lemma-representable-decent}
设 $f : \mathcal{X} \to \mathcal{Y}$ 是代数叠的态射。假设 $Y$
是适宜的，并且 $f$ 可表（由概形表示），或者 $f$ 可由代数空间表示且拟分离。
则 $\mathcal{X}$ 是适宜的。
\end{lemma}

\begin{proof}
设 $x \in |\mathcal{X}|$ 的像为 $y \in |\mathcal{Y}|$。
在定义 $y$ 的等价类中选择态射 $y : \Spec(k) \to \mathcal{Y}$。
令 $\mathcal{X}_y = \Spec(k) \times_{y, \mathcal{Y}} \mathcal{X}$。
选择映到 $x$ 的点 $x' \in |\mathcal{X}_y|$；见《叠的性质》引理
\ref{stacks-properties-lemma-points-cartesian}。在 $x'$ 的等价类中选择态射
$x' : \Spec(k') \to \mathcal{X}_y$。有图
$$
\xymatrix{
\Spec(k') \ar[r]_{x'} &
\mathcal{X}_y \ar[r] \ar[d] &
\mathcal{X} \ar[d] \\
&
\Spec(k) \ar[r]^y &
\mathcal{Y}
}
$$
若 $\mathcal{Y}$ 适宜，则态射 $y$ 拟紧。因此，作为基变换，
$\mathcal{X}_y \to \mathcal{X}$ 拟紧
（引理 \ref{stacks-morphisms-lemma-base-change-quasi-compact}）。所以为得到结论，
只需证明 $x'$ 拟紧（引理 \ref{stacks-morphisms-lemma-composition-quasi-compact}）。
若 $f$ 可表，则 $\mathcal{X}_y$ 是概形，且 $x'$ 拟紧。
若 $f$ 可由代数空间表示且拟分离，则 $\mathcal{X}_y$ 是拟分离代数空间，
因而适宜（《适宜代数空间》引理
\ref{decent-spaces-lemma-properties-trivial-implications}）。
\end{proof}

\begin{lemma}
\label{stacks-morphisms-lemma-qc-decent}
设 $f : \mathcal{X} \to \mathcal{Y}$ 是代数叠的态射。
若 $f$ 拟紧且满，并且 $\mathcal{X}$ 适宜，则 $\mathcal{Y}$ 适宜。
\end{lemma}

\begin{proof}
设 $x : \Spec(k) \to \mathcal{X}$ 是态射，其中 $k$ 是域，并记
$y = f \circ x$。由于 $f$ 满，$|\mathcal{Y}|$ 的每个点都以这种方式得到；
见《叠的性质》引理
\ref{stacks-properties-lemma-characterize-surjective}。
考虑仿射概形 $T$ 和态射 $T \to \mathcal{Y}$。只需证明
$T \times_{\mathcal{Y}, y} \Spec(k)$ 拟紧；见引理
\ref{stacks-morphisms-lemma-check-quasi-compact-covering}。我们有
$$
(T \times_{\mathcal{Y}} \mathcal{X}) \times_{\mathcal{X}, x} \Spec(k) =
T \times_{\mathcal{Y}, y} \Spec(k)
$$
态射 $T \times_{\mathcal{Y}} \mathcal{X} \to T$ 拟紧，因而
$T \times_\mathcal{Y} \mathcal{X}$ 拟紧。又因 $\mathcal{X}$ 适宜，
$x$ 是拟紧态射，所以所示纤维积拟紧。
\end{proof}

\begin{lemma}
\label{stacks-morphisms-lemma-gerbe-decent}
设 $f : \mathcal{X} \to \mathcal{Y}$ 是代数叠的态射。
若 $\mathcal{X}$ 是 $\mathcal{Y}$ 上的胚，且 $\mathcal{X}$ 适宜，
则 $\mathcal{Y}$ 适宜。
\end{lemma}

\begin{proof}
假设 $\mathcal{X}$ 是 $\mathcal{Y}$ 上的胚，且 $\mathcal{X}$ 适宜。
注意，由引理 \ref{stacks-morphisms-lemma-gerbe-bijection-points}，$f$ 是泛同胚。
因此本引理由引理 \ref{stacks-morphisms-lemma-qc-decent} 得出。
\end{proof}



\section{适宜叠上的点}
\label{stacks-morphisms-section-points-decent}

\noindent
本节是《适宜代数空间》第 \ref{decent-spaces-section-points} 节的类比版本。
我们不知道适宜代数叠所关联的拓扑空间是否总是清醒的；稍弱的结果见命题
\ref{stacks-morphisms-proposition-decent-sober}。

\begin{lemma}
\label{stacks-morphisms-lemma-kolmogorov}
设 $\mathcal{X}$ 是适宜代数叠。则 $|\mathcal{X}|$ 是柯尔莫哥洛夫空间
（见《拓扑学》定义 \ref{topology-definition-generic-point}）。
\end{lemma}

\begin{proof}
设 $x_1, x_2 \in |\mathcal{X}|$，并且 $x_1 \leadsto x_2$、
$x_2 \leadsto x_1$。必须证明 $x_1 = x_2$。令
$\mathcal{Z} \subset \mathcal{X}$ 为约化闭子叠，其 $|\mathcal{Z}|$ 等于
$\overline{\{x_1\}} = \overline{\{x_2\}}$。由引理
\ref{stacks-morphisms-lemma-representable-decent}，
$\mathcal{Z}$ 适宜。用 $\mathcal{Z}$ 替换 $\mathcal{X}$ 后，
归结到下一段所述的情形。

\medskip\noindent
假设 $|\mathcal{X}|$ 不可约，并且有一般点 $x_1$ 和 $x_2$。
选取仿射概形 $U$、点 $u_1, u_2 \in U$ 以及光滑态射
$f : U \to \mathcal{X}$，使得 $f(u_i) = x_i$。于是可以找到第三个点
$u_3 \in U$，它是 $U$ 的某个不可约分支的一般点，而其像
$x_3 \in |\mathcal{X}|$ {\it 也是} $|\mathcal{X}|$ 的一般点。
具体地，只需把 $u_3$ 取为经过 $u_1$（或任选 $u_2$）的某个不可约分支的
任意一般点。下一段将证明 $x_1 = x_3$ 且 $x_2 = x_3$，从而得到所需结论。

\medskip\noindent
由对称性，只需证明 $x_1 = x_3$。由于 $x_1$ 是 $|\mathcal{X}|$
的一般点，有特化 $x_1 \leadsto x_3$。由引理
\ref{stacks-morphisms-lemma-specialization}，可以在 $U$ 中找到映到
$x_1 \leadsto x_3$ 的特化 $u'_1 \leadsto u_3$（！）。然而，
$u_3$ 是某个不可约分支的一般点，所以如所需，$u'_1 = u_3$。
\end{proof}

\begin{lemma}
\label{stacks-morphisms-lemma-decent-locally-noetherian-sober}
设 $\mathcal{X}$ 是适宜的局部诺特代数叠。则 $|\mathcal{X}|$
是清醒的局部诺特拓扑空间。
\end{lemma}

\begin{proof}
由引理 \ref{stacks-morphisms-lemma-Noetherian-topology}，拓扑空间 $|\mathcal{X}|$
局部诺特。由引理 \ref{stacks-morphisms-lemma-kolmogorov}，拓扑空间 $|\mathcal{X}|$
是柯尔莫哥洛夫空间。由引理 \ref{stacks-morphisms-lemma-locally-Noetherian-quasi-sober}，
拓扑空间 $|\mathcal{X}|$ 拟清醒。证明完成；见《拓扑学》定义
\ref{topology-definition-generic-point}。
\end{proof}

\begin{proposition}
\label{stacks-morphisms-proposition-decent-sober}
设 $\mathcal{X}$ 是适宜代数叠，并且
$\mathcal{I}_\mathcal{X} \to \mathcal{X}$ 拟紧。则
$|\mathcal{X}|$ 清醒。
\end{proposition}

\begin{proof}
由引理 \ref{stacks-morphisms-lemma-kolmogorov}，$|\mathcal{X}|$ 是柯尔莫哥洛夫空间
（实际上下面会重新证明这一点）。设
$T \subset |\mathcal{X}|$ 是不可约闭子集。必须证明 $T$ 有一般点。
令 $\mathcal{Z} \subset \mathcal{X}$ 为对应于 $T$ 的约化诱导闭子叠；
见《叠的性质》定义
\ref{stacks-properties-definition-reduced-induced-stack}。
由于 $\mathcal{Z} \to \mathcal{X}$ 是闭浸入，由引理
\ref{stacks-morphisms-lemma-representable-decent}，$\mathcal{Z}$ 是适宜代数叠。
此外，态射 $\mathcal{I}_\mathcal{Z} \to \mathcal{Z}$ 是
$\mathcal{I}_\mathcal{X} \to \mathcal{X}$ 的基变换
（引理 \ref{stacks-morphisms-lemma-monomorphism-cartesian-square-inertia}）。因此
$\mathcal{I}_\mathcal{Z} \to \mathcal{Z}$ 拟紧
（引理 \ref{stacks-morphisms-lemma-base-change-quasi-compact}）。由此归结到下一段所述的情形。

\medskip\noindent
假设 $\mathcal{X}$ 适宜，$\mathcal{I}_\mathcal{X} \to \mathcal{X}$
拟紧，$\mathcal{X}$ 约化，并且 $|\mathcal{X}|$ 不可约。必须证明
$|\mathcal{X}|$ 有一般点。由命题 \ref{stacks-morphisms-proposition-open-stratum}，
存在作为胚的稠密开子叠 $\mathcal{U} \subset \mathcal{X}$。
换言之，$|\mathcal{U}| \subset |\mathcal{X}|$ 开且稠密。
因此，除其他所有性质外，还可假设 $\mathcal{X}$ 是胚。
设 $\mathcal{X} \to X$ 使 $\mathcal{X}$ 成为代数空间 $X$ 上的胚。
由引理 \ref{stacks-morphisms-lemma-gerbe-bijection-points}，
$|\mathcal{X}| \cong |X|$。特别地，$X$ 拟紧且 $|X|$ 不可约。
又由引理 \ref{stacks-morphisms-lemma-gerbe-decent}，$X$ 是适宜代数空间。
于是，由《适宜代数空间》命题
\ref{decent-spaces-proposition-reasonable-sober}，
$|\mathcal{X}| = |X|$ 清醒，因而有一个（唯一的）一般点。
\end{proof}






\section{整代数叠}
\label{stacks-morphisms-section-integral-stacks}

\noindent
本节是《域上的空间》第
\ref{spaces-over-fields-section-integral-spaces} 节的类比版本。
受该节中的考察以及命题 \ref{stacks-morphisms-proposition-decent-sober} 的结果启发，
如下定义整代数叠（它与已有的整概形和整代数空间定义并不冲突）。

\begin{definition}
\label{stacks-morphisms-definition-integral-algebraic-stack}
若代数叠 $\mathcal{X}$ 约化、适宜，
$\mathcal{I}_\mathcal{X} \to \mathcal{X}$ 拟紧，并且
$|\mathcal{X}|$ 不可约，则称它是{\it 整的}。
\end{definition}

\noindent
注意，若 $\mathcal{X}$ 拟分离，则要使它为整的，只需
$\mathcal{X}$ 约化且 $|\mathcal{X}|$ 不可约；见引理
\ref{stacks-morphisms-lemma-totaro-case}。

\begin{lemma}
\label{stacks-morphisms-lemma-structure-integral-stack}
设 $\mathcal{X}$ 是整代数叠。则
\begin{enumerate}
\item $|\mathcal{X}|$ 清醒、不可约，并且有唯一的一般点；
\item 存在开子叠 $\mathcal{U} \subset \mathcal{X}$，它是整概形 $U$ 上的胚。
\end{enumerate}
\end{lemma}

\begin{proof}
命题 \ref{stacks-morphisms-proposition-decent-sober} 告诉我们 $|\mathcal{X}|$ 清醒。
它当然也不可约，因而有唯一的一般点 $x$（由清醒性的定义）。
命题 \ref{stacks-morphisms-proposition-open-stratum} 给出稠密开子叠
$\mathcal{U} \subset \mathcal{X}$，它是某个代数空间 $U$ 上的胚。
由引理 \ref{stacks-morphisms-lemma-gerbe-decent}，$U$ 是适宜代数空间（而
$\mathcal{U}$ 由引理 \ref{stacks-morphisms-lemma-representable-decent} 适宜）。
由于 $|U| = |\mathcal{U}|$，$|U|$ 不可约。最后，因为
$\mathcal{U}$ 约化，态射 $\mathcal{U} \to U$ 分解经过 $U_{red}$；
见《叠的性质》引理
\ref{stacks-properties-lemma-map-into-reduction}。现在，由于
$\mathcal{U} \to U$ 平坦、局部有限表现且满
（引理 \ref{stacks-morphisms-lemma-gerbe-fppf}），这蕴含 $U = U_{red}$，即 $U$
约化（略去一个小细节）。于是 $U$ 是整代数空间；见《域上的空间》定义
\ref{spaces-over-fields-definition-integral-algebraic-space}。
最后，可以用开子空间替换 $U$（并相应替换 $\mathcal{U}$），
并假设 $U$ 是整概形；见《域上的空间》第
\ref{spaces-over-fields-section-integral-spaces} 节中的讨论。
\end{proof}

\begin{lemma}
\label{stacks-morphisms-lemma-totaro-case}
设 $\mathcal{X}$ 是约化且拟分离的代数叠，并且其关联拓扑空间
$|\mathcal{X}|$ 不可约。则 $\mathcal{X}$ 是整的。
\end{lemma}

\begin{proof}
若 $\mathcal{X}$ 拟分离，则由引理
\ref{stacks-morphisms-lemma-quasi-separated-decent}，$\mathcal{X}$ 适宜。
若 $\mathcal{X}$ 拟分离，则
$\Delta : \mathcal{X} \to \mathcal{X} \times \mathcal{X}$ 拟紧；
因此，作为 $\Delta$ 沿 $\Delta$ 的基变换，
$\mathcal{I}_\mathcal{X} \to \mathcal{X}$ 拟紧；见引理
\ref{stacks-morphisms-lemma-base-change-quasi-compact}。由此可见，定义
\ref{stacks-morphisms-definition-integral-algebraic-stack} 的所有假设都成立
（同时也看出，可以把“拟分离”替换为“$\Delta_\mathcal{X}$ 拟紧”）。
\end{proof}

\begin{lemma}
\label{stacks-morphisms-lemma-decent-irreducible-closed}
设 $\mathcal{X}$ 是适宜代数叠，并且
$\mathcal{I}_\mathcal{X} \to \mathcal{X}$ 拟紧。
下列集合之间存在典范双射：
\begin{enumerate}
\item $\mathcal{X}$ 的点集，即 $|\mathcal{X}|$；
\item $|\mathcal{X}|$ 的不可约闭子集所成的集合；
\item $\mathcal{X}$ 的整闭子叠所成的集合。
\end{enumerate}
从 (1) 到 (2) 的双射把 $x$ 映到 $\overline{\{x\}}$。
从 (3) 到 (2) 的双射把 $\mathcal{Z}$ 映到 $|\mathcal{Z}|$。
\end{lemma}

\begin{proof}
由命题 \ref{stacks-morphisms-proposition-decent-sober}，$|\mathcal{X}|$ 清醒，
所以我们的映射定义了 (1) 与 (2) 之间的双射。给定闭且不可约的
$T \subset |\mathcal{X}|$，存在唯一的约化闭子叠
$\mathcal{Z} \subset \mathcal{X}$，使得 $|\mathcal{Z}| = T$；
即 $\mathcal{Z}$ 是 $T$ 上的约化诱导子空间结构；见《叠的性质》定义
\ref{stacks-properties-definition-reduced-induced-stack}。
于是 $\mathcal{Z}$ 是整代数叠，因为它适宜
（引理 \ref{stacks-morphisms-lemma-representable-decent}），态射
$\mathcal{I}_\mathcal{Z} \to \mathcal{Z}$ 拟紧
（它是 $\mathcal{I}_\mathcal{X} \to \mathcal{X}$ 的基变换；见引理
\ref{stacks-morphisms-lemma-monomorphism-cartesian-square-inertia}），
$\mathcal{Z}$ 约化，并且 $|\mathcal{Z}|$ 不可约。
\end{proof}










\section{剩余胚}
\label{stacks-morphisms-section-residual-gerbe}

\noindent
本节是《叠的性质》第
\ref{stacks-properties-section-residual-gerbe} 节的继续。

\begin{lemma}
\label{stacks-morphisms-lemma-gerbe-residual-gerbe}
设 $\pi : \mathcal{X} \to \mathcal{Y}$ 是代数叠的态射。设
$x \in |\mathcal{X}|$ 的像为 $y \in |\mathcal{Y}|$。假设
$\mathcal{Y}$ 在 $y$ 处的剩余胚
$\mathcal{Z}_y \subset \mathcal{Y}$ 存在，并且 $\mathcal{X}$ 是
$\mathcal{Y}$ 上的胚。则
$\mathcal{Z}_x = \mathcal{Z}_y \times_\mathcal{Y} \mathcal{X}$ 是
$\mathcal{X}$ 在 $x$ 处的剩余胚。
\end{lemma}

\begin{proof}
态射 $\mathcal{Z}_x \to \mathcal{X}$ 是单态射，因为它是单态射
$\mathcal{Z}_y \to \mathcal{Y}$ 的基变换。由引理
\ref{stacks-morphisms-lemma-gerbe-bijection-points}，态射 $\pi$ 是泛同胚，因而
$|\mathcal{Z}_x| = \{x\}$。最后，态射
$\mathcal{Z}_x \to \mathcal{Z}_y$ 是光滑态射 $\pi$ 的基变换，
所以它光滑；见引理 \ref{stacks-morphisms-lemma-gerbe-smooth}。因此，由于
$\mathcal{Z}_y$ 约化且局部诺特，$\mathcal{Z}_x$ 也如此
（细节略）。于是，由《叠的性质》定义
\ref{stacks-properties-definition-residual-gerbe}，$\mathcal{Z}_x$ 是
$\mathcal{X}$ 在 $x$ 处的剩余胚。
\end{proof}

\begin{lemma}
\label{stacks-morphisms-lemma-factor-through-residual-space}
设 $f : \mathcal{Y} \to \mathcal{X}$ 是代数叠的态射。
设 $x \in |\mathcal{X}|$ 是一点。假设
\begin{enumerate}
\item $\mathcal{X}$ 适宜或局部诺特（或两者皆是）；
\item $\mathcal{I}_\mathcal{X} \to \mathcal{X}$ 拟紧；
\item $|f|(|\mathcal{Y}|)$ 包含在 $\{x\} \subset |\mathcal{X}|$ 中；
\item $\mathcal{Y}$ 约化。
\end{enumerate}
则 $f$ 分解经过 $\mathcal{X}$ 在 $x$ 处的剩余胚 $\mathcal{Z}_x$
（其存在性由引理 \ref{stacks-morphisms-lemma-every-point-residual-gerbe} 或
\ref{stacks-morphisms-lemma-every-point-residual-gerbe-locally-Noetherian}).
保证）。
\end{lemma}

\begin{proof}
令 $T = \overline{\{x\}} \subset |\mathcal{X}|$ 为 $x$ 的闭包。
由《叠的性质》引理
\ref{stacks-properties-lemma-reduced-closed-substack}
，存在约化闭子叠 $\mathcal{X}' \subset \mathcal{X}$，使得
$T = |\mathcal{X}'|$。由《叠的性质》引理
\ref{stacks-properties-lemma-map-into-reduction}，态射 $f$ 分解经过
$\mathcal{X}'$。若 $\mathcal{X}$ 适宜，则由引理
\ref{stacks-morphisms-lemma-representable-decent}，叠 $\mathcal{X}'$ 适宜。
若 $\mathcal{X}$ 局部诺特，则 $\mathcal{X}'$ 局部诺特（细节略）。
注意，$\mathcal{I}_{\mathcal{X}'} \to \mathcal{X}'$ 是
$\mathcal{I}_\mathcal{X} \to \mathcal{X}$ 的基变换；由引理
\ref{stacks-morphisms-lemma-monomorphism-cartesian-square-inertia} 和
\ref{stacks-morphisms-lemma-base-change-quasi-compact}，
$\mathcal{I}_{\mathcal{X}'} \to \mathcal{X}'$ 拟紧。
由此归结到下一段所述的情形。

\medskip\noindent
假设 $\mathcal{X}$ 约化且 $x \in |\mathcal{X}|$ 是一般点。
由命题 \ref{stacks-morphisms-proposition-open-stratum}，存在作为胚的稠密开子叠
$\mathcal{U} \subset \mathcal{X}'$。注意 $x \in |\mathcal{U}|$。
重复上述论证，归结到下一段所述的情形。

\medskip\noindent
假设 $\mathcal{X} \to X$ 是代数空间 $X$ 上的胚。若
$\mathcal{X}$ 适宜，则由引理 \ref{stacks-morphisms-lemma-gerbe-bijection-points} 和
\ref{stacks-morphisms-lemma-qc-decent}，空间 $X$ 适宜。若 $\mathcal{X}$ 局部诺特，
则由 fppf 下降，$X$ 局部诺特（细节略）。因此相应结论对 $X$ 成立；
见《适宜代数空间》引理
\ref{decent-spaces-lemma-factor-through-residual-space-decent} 或
\ref{decent-spaces-lemma-factor-through-residual-space-Noetherian}
（略去一个小细节）。应用引理 \ref{stacks-morphisms-lemma-gerbe-residual-gerbe}，
可知结论对 $\mathcal{X}$ 也成立。
\end{proof}

\begin{remark}
\label{stacks-morphisms-lemma-factor-through-residual-space-Noetherian}
我们不知道仅假设 $\mathcal{X}$ 局部诺特时，引理
\ref{stacks-morphisms-lemma-factor-through-residual-space} 是否仍成立；换言之，
这里去掉关于惯性态射拟紧的假设。在这种情形下，若 $x$ 是闭点，
则下面这个简单得多的引理表明结论当然成立。
\end{remark}

\begin{lemma}
\label{stacks-morphisms-lemma-residual-space-closed}
设 $\mathcal{X}$ 是局部诺特代数叠。设 $x \in |\mathcal{X}|$ 的剩余胚为
$\mathcal{Z}_x \subset \mathcal{X}$
（引理 \ref{stacks-morphisms-lemma-every-point-residual-gerbe-locally-Noetherian}）。
则 $x$ 是 $|\mathcal{X}|$ 的闭点，当且仅当态射
$\mathcal{Z}_x \to \mathcal{X}$ 是闭浸入。
\end{lemma}

\begin{proof}
若 $\mathcal{Z}_x \to \mathcal{X}$ 是闭浸入，则 $x$ 是
$|\mathcal{X}|$ 的闭点；例如见引理 \ref{stacks-morphisms-lemma-closed-immersion-proper}。
反过来，假设 $x$ 是 $|\mathcal{X}|$ 的闭点。令
$\mathcal{Z} \subset \mathcal{X}$ 为满足 $|Z| = \{x\}$ 的约化闭子叠
（《叠的性质》引理
\ref{stacks-properties-lemma-reduced-closed-substack}）。由引理
\ref{stacks-morphisms-lemma-immersion-locally-finite-type} 和
\ref{stacks-morphisms-lemma-locally-finite-type-locally-noetherian}，$\mathcal{Z}$
是局部诺特代数叠。又因 $\mathcal{Z}$ 约化且
$|\mathcal{Z}| = \{x\}$，按定义，$\mathcal{Z} = \mathcal{Z}_x$
是剩余胚。
\end{proof}
